% ============================================================
%  preamble.tex --- styling, environments, macros
%
%  Shared by BOTH editions. main-en.tex and main-pl.tex each define \booklang
%  before % ============================================================
%  preamble.tex --- styling, environments, macros
%
%  Shared by BOTH editions. main-en.tex and main-pl.tex each define \booklang
%  before % ============================================================
%  preamble.tex --- styling, environments, macros
%
%  Shared by BOTH editions. main-en.tex and main-pl.tex each define \booklang
%  before % ============================================================
%  preamble.tex --- styling, environments, macros
%
%  Shared by BOTH editions. main-en.tex and main-pl.tex each define \booklang
%  before \input{preamble}; nothing else differs between them at this level.
%  Every user-visible string lives in lang/en.tex or lang/pl.tex, so a label
%  that appears in one edition and not the other is a missing macro rather
%  than a silent divergence.
%
%  Lineage. The bilingual wiring (\booklang, lang/, body.tex, the generated
%  structure) is the math book's; the chapter machinery (listings, the
%  admonitions, verifybox, \pyfile and \pyregion, the Mermaid pipeline) is the
%  LangChain book's. What is new here is the exercise macro, which ties every
%  exercise to a starter file and a test under code/, because this book is
%  read with an editor open beside it.
% ============================================================

% \booklang must be set by the main file. Default to English rather than
% failing, so a stray \input{preamble} still compiles and says what it did.
\providecommand{\booklang}{en}

\usepackage{etoolbox}   % loaded first: the language switch below needs it

% ---------- Page geometry ----------
% ONE format: A4, 12pt, ONE-SIDED. This book is read on a screen with an
% editor open beside it, and never printed as a bound book, so there is no
% recto, no verso, no gutter, no mirrored margin and no blank leaf before a
% chapter -- every one of those is a property of paper that a PDF viewer at
% half a screen's width turns into wasted pixels.
%
% The margins are symmetric and the measure is about seventy-five characters
% of prose, which is where the eye keeps the line return. The number that was
% actually chosen is the MONOSPACE measure: at \footnotesize on this block a
% listing line of 79 characters fits without wrapping, and 79 is the width
% every listing under code/ is held to by tools/check_structure.py --lines.
% A wrapped listing line is not an error TeX reports -- breaklines=true wraps
% it silently and prints an arrow into the middle of what the reader is meant
% to paste into an editor -- so the width is enforced in the source instead.
\usepackage[
  a4paper,
  left=3.1cm, right=3.1cm, top=2.6cm, bottom=3.0cm,
  head=26pt, headsep=0.8cm, footskip=1.4cm
]{geometry}

% ---------- Encoding & fonts ----------
\usepackage[T1]{fontenc}
\usepackage[utf8]{inputenc}
\IfFileExists{lmodern.sty}{\usepackage{lmodern}}{}
% newtxtext takes its TS1 symbols (\textcopyright, \textdegree) from TeX
% Gyre Termes, which Debian packages SEPARATELY (tex-gyre): a machine with
% newtx and without it dies on the copyright page with
% `Font TS1/ntxtlf/m/n/10.95=ts1-qtmr ... not loadable`. There is no
% inert-safe probe for a font metric in pdfTeX -- \IfFileExists searches
% TEXINPUTS and not TFMFONTS, so it reports the metric missing on a machine
% that has it, and \suppressfontnotfounderror is LuaTeX-only -- so nothing
% here guesses. CI's full TeX Live is the reference; locally, install
% tex-gyre. Both facts are recorded in CLAUDE.md under Build traps.
\IfFileExists{newtxtext.sty}{\usepackage{newtxtext}}{}
% amsmath before newtxmath, because newtxmath patches amsmath's internals.
% amssymb only when newtxmath is absent: newtxmath supplies the AMS symbols
% itself and loading both is a fatal clash (\Bbbk already defined) that is
% invisible on a machine without newtx and fatal on a full TeX Live.
\usepackage{amsmath}
\IfFileExists{newtxmath.sty}{\usepackage{newtxmath}}{\usepackage{amssymb}}
\IfFileExists{inconsolata.sty}{\usepackage[varqu,scaled=0.95]{inconsolata}}{}
\IfFileExists{lmodern.sty}{\usepackage{microtype}}{\usepackage[expansion=false]{microtype}}

% upquote is not a refinement. In T1 the input character ' is quoteright, so
% a listings body sets f('x') with U+2019 at both ends of the literal, and
% typed back in that is a SyntaxError. inconsolata's varqu option hides the
% defect on a machine that has inconsolata; upquote makes ' the ASCII
% apostrophe whatever the typewriter font is. In a book whose every listing is
% meant to be pasted into an editor, this line is load-bearing.
\IfFileExists{upquote.sty}{\usepackage{upquote}}{}

% And upquote's twin, for the same reason one character over. In T1 the pair
% -- is a ligature for an en dash, in EVERY family including the typewriter
% one, so \code{uv sync --locked} sets as "uv sync <en dash>locked" and a
% reader who copies it off the page gets an unknown-argument error from a
% command the chapter told them to run. It is invisible in review because it
% looks like a dash, and it bites exactly the flags a packaging chapter is
% made of. A `listings` body is already safe -- listings sets characters
% singly and no ligature forms -- which is what makes the defect so easy to
% miss: the same text is right inside a listing and wrong in the sentence
% above it.
%
% Scoped to tt* on purpose, and measured rather than assumed: prose keeps its
% own dashes, so \dash{} still sets an em dash and a range like 1--2 still
% sets an en dash. Probed against this preamble in all three positions before
% it was believed.
\DisableLigatures[-]{encoding = T1, family = tt*}

% ---------- Language ----------
% Loading babel with a language whose .ldf is absent is a *fatal* error, not
% a warning, so both the package and each language file are probed before
% either is requested. Both editions load the other language too, because the
% glossary rows and the cheat sheet carry words from both.
\IfFileExists{babel.sty}{%
  \IfFileExists{polish.ldf}{%
    \IfFileExists{british.ldf}{%
      \ifdefstring{\booklang}{pl}%
        {\usepackage[british,main=polish]{babel}}%
        {\usepackage[polish,main=british]{babel}}%
    }{\usepackage[polish]{babel}}%
  }{%
    \IfFileExists{british.ldf}{\usepackage[british]{babel}}{}%
  }%
}{}

% babel-polish makes " an active shorthand. listings reads its body verbatim
% and an active character in a verbatim context is a category-code accident
% waiting to happen, so the shorthand is turned off. Polish quotation marks
% come from \enquote, which csquotes sets per language.
\ifdefstring{\booklang}{pl}{\AtBeginDocument{\shorthandoff{"}}}{}

% ---------- Core packages ----------
\usepackage{graphicx}
\usepackage{xcolor}
\usepackage{booktabs}
\usepackage{tabularx}
% Appendix A's cheat sheet runs past a page in every part, and tabularx
% cannot break. longtable can, and works with booktabs.
\usepackage{longtable}
\usepackage{array}
\usepackage{enumitem}
\usepackage{listings}
\usepackage[most]{tcolorbox}
\usepackage{fancyhdr}
\usepackage{titlesec}
\usepackage{caption}
\usepackage{imakeidx}
% csquotes with autostyle follows babel's active language, so the identical
% source \enquote{x} sets 'x' in the English edition and the Polish quotation
% marks in the Polish one.
\IfFileExists{csquotes.sty}{\usepackage[autostyle=true]{csquotes}}{%
  \providecommand{\enquote}[1]{``##1''}}
\usepackage[hidelinks]{hyperref}

% ---------- Palette ----------
% The same family as the two companion volumes, so the three read as a set.
\definecolor{mblue}{HTML}{1F4E79}
\definecolor{mteal}{HTML}{0E7C7B}
\definecolor{mamber}{HTML}{B26A00}
\definecolor{mred}{HTML}{9B2C2C}
\definecolor{mviolet}{HTML}{7B4397}
\definecolor{mgrey}{HTML}{5A5A5A}
\definecolor{mlight}{HTML}{F4F6F8}
\definecolor{mfaint}{HTML}{FBFCFD}
\definecolor{codebg}{HTML}{FAFAFA}
\definecolor{codeframe}{HTML}{D8DEE4}
\definecolor{kw}{HTML}{0B5394}
\definecolor{str}{HTML}{9B2C2C}
\definecolor{cmt}{HTML}{4F7A4F}
\definecolor{typ}{HTML}{0E7C7B}
\definecolor{dec}{HTML}{7B4397}

\hypersetup{
  colorlinks=true,
  linkcolor=mblue,
  urlcolor=mteal,
  citecolor=mblue,
  pdfauthor={Konrad Cinkusz}
}

% \ifpl{Polish text}{English text} -- for the handful of places where a
% string is structural rather than content. Anything longer than a few words
% belongs in lang/ or in the edition's own source, not here.
% \DeclareRobustCommand, not \newcommand: a fragile \ifpl inside a moving
% argument is written to the .toc one level expanded and fails on the next run
% in exactly one of the two editions.
\DeclareRobustCommand{\ifpl}[2]{\ifdefstring{\booklang}{pl}{#1}{#2}}
\makeatletter
\pdfstringdefDisableCommands{%
  \ifdefstring{\booklang}{pl}{\let\ifpl\@firstoftwo}{\let\ifpl\@secondoftwo}%
  % \api{}, \pkg{} and \code{} are allowed in a section title, and a title
  % becomes a PDF bookmark. hyperref drops the \color command from a bookmark
  % string but keeps its ARGUMENT, so a title carrying \api{model\_validate}
  % made a bookmark reading "mtealmodel_validate" and warned "Token not
  % allowed" twice. The colour name is gobbled here instead. Measured on a
  % probe against this preamble before the line was added.
  \def\color#1{}%
}
\makeatother

% ---------- Language strings ----------
\input{lang/\booklang}

% The PDF's subject, keyword list and document language read \lblPdf... from
% the language file input on the line above, so they cannot sit in the colour
% \hypersetup block, which runs before it. Two blocks is the ordering, not an
% oversight; merging them raises Undefined control sequence and still writes
% a PDF, with the three fields empty.
\hypersetup{
  pdfsubject={\lblPdfSubject},
  pdfkeywords={\lblPdfKeywords},
  pdflang={\lblPdfLang}
}

% ---------- Headings ----------
% \raggedright on the title body is load-bearing: at \Huge on this block a
% title that cannot break overflows the margin by 30 pt or more, and several
% chapter titles here run past fifty characters.
\titleformat{\chapter}[display]
  {\normalfont\huge\bfseries\color{mblue}\raggedright}
  {\normalfont\normalsize\color{mgrey}\MakeUppercase{\chaptertitlename}\ \thechapter}
  {8pt}{\Huge\raggedright}
\titlespacing*{\chapter}{0pt}{10pt}{28pt}
\titleformat{\section}{\normalfont\Large\bfseries\color{mblue}\raggedright}{\thesection}{0.8em}{}
\titleformat{\subsection}{\normalfont\large\bfseries\color{mgrey}\raggedright}{\thesubsection}{0.7em}{}

% Sections in the contents, subsections not: fourteen chapters of subsections
% is a contents nobody reads, and the reader navigates by chapter and section.
\setcounter{tocdepth}{1}
\setcounter{secnumdepth}{2}

% One-sided running head: the chapter on the left, the folio on the right.
\pagestyle{fancy}
\fancyhf{}
\fancyhead[L]{\small\nouppercase{\leftmark}}
\fancyhead[R]{\small\thepage}
\renewcommand{\headrulewidth}{0.4pt}
\fancypagestyle{plain}{\fancyhf{}\fancyfoot[R]{\small\thepage}%
  \renewcommand{\headrulewidth}{0pt}}

% Drops the word "Chapter" from the head: the number is still there and the
% longer titles no longer overflow. MUST come after \pagestyle{fancy}, which
% installs its own \chaptermark and silently overwrites anything earlier.
\renewcommand{\chaptermark}[1]{\markboth{\thechapter.\ #1}{}}

% ============================================================
%  CODE LISTINGS
%  ---------------------------------------------------------
%  Every listing in a chapter is a FILE under code/, pulled in with \pyfile or
%  \pyregion, and CI runs it against the pinned versions. An in-source
%  \begin{python} block is allowed for a fragment of three or four lines that
%  exists to show a shape rather than to run; anything longer belongs in a
%  file the reader can open in the editor beside the PDF.
% ============================================================
\lstdefinelanguage{pythonx}{
  language=Python,
  morekeywords={
    async,await,match,case,nonlocal,yield,type,
    TypedDict,Annotated,Protocol,Literal,Generic,Final,ClassVar,Self,
    TypeVar,TypeGuard,TypeIs,Callable,Iterator,Iterable,Awaitable,
    dataclass,field,BaseModel,Field,Enum,StrEnum,NamedTuple,
    override,overload,property,staticmethod,classmethod
  },
  morekeywords=[2]{
    gather,create_task,TaskGroup,to_thread,timeout,wait_for,Semaphore,Queue,
    fixture,parametrize,monkeypatch,raises,mark,
    Depends,FastAPI,APIRouter,HTTPException,
    select,Session,Mapped,mapped_column,relationship,
    model_validate,model_validate_json,model_dump,model_json_schema
  },
  sensitive=true
}

% listings' C# definition predates async, records and pattern matching.
\lstdefinelanguage{csharpx}{
  language=[Sharp]C,
  morekeywords={
    async,await,var,dynamic,record,init,required,nameof,yield,partial,
    global,scoped,when,with,and,or,not,notnull,unmanaged,file,get,set,value,
    where,select,from,orderby,let,join,into,ascending,descending,on,equals,
    Task,ValueTask,CancellationToken,IEnumerable,IAsyncEnumerable,Func,Action
  },
  sensitive=true
}

\lstdefinestyle{bookcode}{
  basicstyle=\ttfamily\footnotesize,
  keywordstyle=\color{kw}\bfseries,
  keywordstyle=[2]\color{typ},
  stringstyle=\color{str},
  % Colour only, no \itshape: inconsolata (zi4) has no italic shape, so an
  % italic comment style is substituted with upright on every machine that
  % has the font, and the log carries "Font shape `T1/zi4/m/it' undefined"
  % on every build. The colour already marks a comment.
  commentstyle=\color{cmt},
  emphstyle=\color{dec}\bfseries,
  numbers=left,
  numberstyle=\ttfamily\tiny\color{mgrey},
  numbersep=8pt,
  backgroundcolor=\color{codebg},
  frame=single,
  rulecolor=\color{codeframe},
  framesep=6pt,
  breaklines=true,
  breakatwhitespace=false,
  postbreak=\mbox{\textcolor{mgrey}{$\hookrightarrow$}\space},
  showstringspaces=false,
  tabsize=4,
  columns=fullflexible,
  keepspaces=true,
  captionpos=b,
  aboveskip=1.2em,
  belowskip=1.2em,
  % Maps the characters most likely to be pasted into a listing by accident.
  %
  % Do NOT add a mapping for U+00A0 (non-breaking space). It looks like the
  % obvious next entry and it makes `listings` abort the run with "Improper
  % alphabetic constant" -- fatal, no PDF, and the message names neither the
  % character nor this line. The ASCII-in-listings convention, checked by
  % parity's C13, is the guard against it instead.
  literate={ł}{{\l}}1 {Ł}{{\L}}1 {ą}{{\k a}}1 {ę}{{\k e}}1
           {ś}{{\'s}}1 {ć}{{\'c}}1 {ż}{{\.z}}1 {ź}{{\'z}}1 {ó}{{\'o}}1 {ń}{{\'n}}1
           {Ą}{{\k A}}1 {Ę}{{\k E}}1 {Ś}{{\'S}}1 {Ć}{{\'C}}1
           {Ż}{{\.Z}}1 {Ź}{{\'Z}}1 {Ó}{{\'O}}1 {Ń}{{\'N}}1
           {—}{{-{}-}}2 {–}{{-}}1 {‘}{{'}}1 {’}{{'}}1
           {“}{{"}}1 {”}{{"}}1 {°}{{\textdegree}}1
}

\lstset{style=bookcode}

\lstnewenvironment{python}[1][]
  {\lstset{language=pythonx,style=bookcode,#1}}
  {}

% C# comparison snippets. The whole book is written for engineers arriving
% from .NET, so listings come in pairs often enough that the C# half has its
% own environment rather than a generic one.
\lstnewenvironment{csharp}[1][]
  {\lstset{language=csharpx,style=bookcode,#1}}
  {}

\lstnewenvironment{shellcmd}[1][]
  {\lstset{language=bash,style=bookcode,numbers=none,keywordstyle=\color{mteal}\bfseries,#1}}
  {}

\lstnewenvironment{tomlcode}[1][]
  {\lstset{style=bookcode,numbers=none,keywordstyle=\color{black},#1}}
  {}

\lstnewenvironment{yamlcode}[1][]
  {\lstset{style=bookcode,numbers=none,keywordstyle=\color{black},#1}}
  {}

\lstnewenvironment{jsoncode}[1][]
  {\lstset{style=bookcode,numbers=none,keywordstyle=\color{black},#1}}
  {}

% Terminal transcripts typed into the source. Allowed for a shell session the
% reader is meant to reproduce; NOT for program output, which is \transcript
% below, because an in-source transcript nobody ran reads exactly like one
% that was, and that is where a remembered number hides.
\lstnewenvironment{console}[1][]
  {\lstset{style=bookcode,numbers=none,basicstyle=\ttfamily\footnotesize,
           keywordstyle=\color{black},backgroundcolor=\color{white},#1}}
  {}

% The path is printed above every file listing, in the repository-relative
% form the reader opens in the editor beside the PDF. That is the whole point
% of a listing being a file, and the front matter promises it. The \nobreak
% keeps the path on the same page as the listing it names.
% ---------------------------------------------------------------------------
%  Appendix B prints the trap catalogue from notes/02-traps.md. One entry is
%  its number -- which a chapter may cite, and which is never reused -- the
%  habit in the reader's own voice, and what Python does instead.
%
%  The body is set RAGGED RIGHT on purpose. An entry is mostly \code{} spans,
%  which are unbreakable runs, and sixty-seven justified paragraphs of them
%  is sixty-seven chances at an overfull hbox in the edition with the longer
%  words. Ragged right removes the class rather than fixing it one box at a
%  time, and a reference list is the one place in this book where it reads
%  better anyway.
%
%  The habit is set with \enquote{} rather than \emph{}, and that is not a
%  style choice. Every habit here carries \code{} spans, and \code{} inside
%  \emph{} asks inconsolata for an italic it does not have: the first build
%  of this appendix raised `Font shape T1/zi4/m/it undefined', which
%  checklog.py is hard on for exactly this reason. It is the trap CLAUDE.md
%  records, met sixty-seven times at once -- and fixed once, because the
%  entries go through a macro. A quotation is also what the habit IS, and
%  \enquote{} is what the Polish edition owes anyway.
% ---------------------------------------------------------------------------
%  Appendix A's cheat sheet: one table per part, C# on the left, Python on
%  the right, and the chapter that teaches the row on the end. The three
%  header strings are arguments so that one environment serves both
%  editions and the structural signature stays identical.
%
%  Columns are ragged right for the reason Appendix B's entries are: a cell
%  is mostly \code{} spans, which are unbreakable runs, and a justified
%  narrow column full of them is an overfull box waiting for the edition
%  with the longer words.
% ---------------------------------------------------------------------------
%  Appendix C's tool matrix: the .NET tool, the Python tool this book uses,
%  the version it was verified against and the chapter that introduces it.
%
%  The version column prints \pydanticver and its siblings and is never
%  typed, which is the whole point of the appendix: it is the one page where
%  a reader sees every pin at once, and a pin that had to be re-typed here
%  would be a second copy of a fact the preamble already owns.
%  tools/check_structure.py --tools fails when a pin macro exists and this
%  table does not print it.
\NewDocumentEnvironment{toolmatrix}{mmmm}{%
  \small
  \begin{longtable}{@{}%
    >{\raggedright\arraybackslash}p{0.27\linewidth}%
    >{\raggedright\arraybackslash}p{0.31\linewidth}%
    >{\raggedright\arraybackslash}p{0.20\linewidth}%
    >{\raggedright\arraybackslash}p{0.08\linewidth}@{}}
  \toprule
  \textbf{#1} & \textbf{#2} & \textbf{#3} & \textbf{#4} \\
  \midrule
  \endfirsthead
  \toprule
  \textbf{#1} & \textbf{#2} & \textbf{#3} & \textbf{#4} \\
  \midrule
  \endhead
  \midrule
  \endfoot
  \bottomrule
  \endlastfoot
}{%
  \end{longtable}%
}

\NewDocumentEnvironment{cheatsheet}{mmm}{%
  \small
  \begin{longtable}{@{}%
    >{\raggedright\arraybackslash}p{0.41\linewidth}%
    >{\raggedright\arraybackslash}p{0.41\linewidth}%
    >{\raggedright\arraybackslash}p{0.08\linewidth}@{}}
  \toprule
  \textbf{#1} & \textbf{#2} & \textbf{#3} \\
  \midrule
  \endfirsthead
  \toprule
  \textbf{#1} & \textbf{#2} & \textbf{#3} \\
  \midrule
  \endhead
  \midrule
  \endfoot
  \bottomrule
  \endlastfoot
}{%
  \end{longtable}%
}

\newcommand{\trapentry}[3]{%
  \par\medskip
  \noindent\textbf{#1.}\enspace\enquote{#2}\par
  \nopagebreak
  {\raggedright\noindent#3\par}%
}

\newcommand{\listingpath}[1]{%
  \par\medskip\noindent
  {\ttfamily\footnotesize\color{mgrey}\detokenize{#1}}%
  \par\nobreak}

% External source file -- the preferred form. The path is relative to the
% repository root, which is also the path the reader opens in the editor.
% Usage: \pyfile{code/ch05/decorators.py}{Caption}{lst:label}
\newcommand{\pyfile}[3]{%
  \listingpath{#1}%
  \lstinputlisting[language=pythonx,style=bookcode,aboveskip=0.3em,
    caption={#2},label={#3}]{#1}%
}

% A named region of an external file, delimited in the source by
% `# --8<-- [start:name]` / `# --8<-- [end:name]`, so the book and the running
% code cannot drift apart. tools/check_structure.py --listings fails when the
% file or the region is missing.
%
% The markers are matched through listings' range prefix and suffix keys,
% with linerange={name-name}. The form this was inherited with -- a linerange
% whose two ends were the whole marker strings -- matched nothing and printed
% the WHOLE FILE, with no warning: measured on a three-region probe file
% before this was rewritten, and the same definition sits unused in the
% sibling book it came from. A region name is a word, never a number, because
% a number is a line number to listings. The line numbers printed are the
% file's own, so a region's listing says where in the file it sits.
% Usage: \pyregion{code/ch05/decorators.py}{timing}{Caption}{lst:label}
\lstset{
  rangebeginprefix=\#\ --8<--\ [start:, rangebeginsuffix=],
  rangeendprefix=\#\ --8<--\ [end:, rangeendsuffix=],
  includerangemarker=false}
\newcommand{\pyregion}[4]{%
  \listingpath{#1}%
  \lstinputlisting[language=pythonx,style=bookcode,aboveskip=0.3em,
    linerange={#2-#2},caption={#3},label={#4}]{#1}%
}

% The C# twin of \pyfile, for the comparison half of a pair.
\newcommand{\csfile}[3]{%
  \listingpath{#1}%
  \lstinputlisting[language=csharpx,style=bookcode,aboveskip=0.3em,
    caption={#2},label={#3}]{#1}%
}

% A program's output, written by a script under code/measure/ and pulled in
% whole. There is no typed form of this and there must not be: \val{} does
% not expand inside a verbatim body, so a transcript typed into a chapter is
% one that cannot quote a computed value, and a transcript nobody ran is
% indistinguishable on the page from one that was.
\newcommand{\transcript}[1]{%
  \par\smallskip
  \IfFileExists{figures/transcripts/#1.txt}{%
    \lstinputlisting[style=bookcode,numbers=none,basicstyle=\ttfamily\footnotesize,
      keywordstyle=\color{black},backgroundcolor=\color{white},
      language={}]{figures/transcripts/#1.txt}%
  }{%
    \begin{tcolorbox}[enhanced, sharp corners, width=\linewidth,
      colback=mlight, colframe=mgrey, boxrule=0.6pt,
      left=8pt, right=8pt, top=6pt, bottom=6pt]
      {\bfseries\color{mgrey}\small \lblTranscriptMissing}\\[2pt]
      {\ttfamily\scriptsize\color{mgrey} figures/transcripts/#1.txt}
    \end{tcolorbox}%
  }%
  \par\smallskip
}

% Inline literals. Underscores inside these must be written \_.
\newcommand{\code}[1]{\texttt{\small #1}}
\newcommand{\pkg}[1]{\texttt{\small\color{mblue}#1}}
\newcommand{\api}[1]{\texttt{\small\color{mteal}#1}}

% ============================================================
%  ADMONITIONS
%  ---------------------------------------------------------
%  A deliberately small, fixed vocabulary. Two visual treatments carry the
%  whole grammar: `tinted` is a block the reader may read in passing -- a
%  tint and a thick left bar, no outline; `outlined` is a block the reader
%  must stop at, with a full rule on white. A third treatment would mean the
%  first two are not carrying their meaning.
% ============================================================
\tcbset{
  mfa tinted/.style={
    enhanced, breakable, sharp corners,
    colback=mlight, boxrule=0pt, leftrule=3pt,
    left=8pt, right=8pt, top=6pt, bottom=6pt,
    before skip=13pt, after skip=13pt,
    fonttitle=\bfseries\small,
    coltitle=black, colbacktitle=mlight,
    attach title to upper={\par\smallskip},
  },
  mfa outlined/.style={
    enhanced, breakable, sharp corners,
    colback=white, boxrule=0.6pt,
    left=8pt, right=8pt, top=6pt, bottom=6pt,
    before skip=13pt, after skip=13pt,
    fonttitle=\bfseries\small,
    coltitle=black, colbacktitle=white,
    attach title to upper={\par\smallskip},
  },
}

% Read in passing.
\newtcolorbox{note}[1][]{mfa tinted, colframe=mblue, title={\lblNote}, #1}
\newtcolorbox{warning}[1][]{mfa tinted, colframe=mred, title={\lblWarning}, #1}
% The translation. The .NET reader is the book's only assumed reader, so this
% box carries the C# side of a comparison and nothing else. Its use is
% counted: a chapter with none is a chapter that has stopped translating.
\newtcolorbox{csbox}[1][]{mfa tinted, colframe=mteal, title={\lblCsharp}, #1}
% Anything that is preview, experimental or likely to move. Python moves more
% slowly than the agent frameworks do, and this box should be rare here.
\newtcolorbox{versionbox}[1][]{mfa tinted, colframe=mamber, title={\lblVersion}, #1}

% Stop and read.
% The misconception, named AFTER the reader has walked into it. Appendix B is
% the catalogue; a trap box in a chapter is one entry of it, in place.
\newtcolorbox{trapbox}[1][]{mfa outlined, colframe=mred, title={\lblTrap}, #1}
% Marks a listing that has not been executed against the pinned versions.
% `make debt` counts these, CI prints the count, and nothing ships to a reader
% with one attached. Removing it means you ran the code.
\newtcolorbox{verifybox}[1][]{mfa outlined, colframe=mgrey, title={\lblVerify}, #1}
\newtcolorbox{exercisebox}[1][]{mfa outlined, colframe=mteal, title={\lblExercise}, #1}
% The guiding project. Points at the stage of trace-assert this section
% builds, under code/src/trace_assert/.
\newtcolorbox{projectbox}[1][]{mfa outlined, colframe=mamber, title={\lblProject}, #1}

% ============================================================
%  EXERCISES
%  ---------------------------------------------------------
%  \begin{exercise}{key}{title} ... \end{exercise}
%
%  This is the mechanism the whole book is read through: the PDF on one half
%  of the screen, an editor on the other. An exercise is therefore not a
%  paragraph but a FILE the reader opens -- a starter under
%  code/exercises/chNN/<key>.py -- and a TEST that decides whether it is done,
%  test_<key>.py beside it. The box prints both paths and the one command
%  that runs the test, and the manifest at the back lists every exercise.
%
%  An ENVIRONMENT, not a macro with a body argument: a macro argument cannot
%  contain verbatim material, so an exercise written as \exercise{key}{title}
%  {body} could never show the reader a code fragment, and the first chapter
%  that tried would have died on `\begin{python}` inside the braces.
%
%  The key is the file stem, with underscores, and it is written with
%  \detokenize so the underscore reaches the page as a character rather than
%  as a subscript. tools/check_structure.py --exercises fails when the starter
%  or the test is missing, or when the key's chapter prefix disagrees with the
%  chapter it sits in.
% ============================================================
\newcounter{exercise}[chapter]
\renewcommand{\theexercise}{\thechapter.\arabic{exercise}}
\makeatletter
\newcommand{\exercisedir}{ch\two@digits{\value{chapter}}}
\newcommand{\listofexercises}{%
  \section*{\lblExerciseManifest}%
  \phantomsection\addcontentsline{toc}{section}{\lblExerciseManifest}%
  % The entries are subsection-level and tocdepth is 1 for the contents,
  % so the depth is raised for this list alone or it prints empty --
  % which it did, on the scaffold's first clean build.
  \begingroup\c@tocdepth=2\relax\@starttoc{exr}\endgroup%
}
\makeatother
\NewDocumentEnvironment{exercise}{mm}{%
  \refstepcounter{exercise}%
  \addcontentsline{exr}{subsection}{\protect\numberline{\theexercise}\texttt{\protect\detokenize{#1}} --- #2}%
  \begin{exercisebox}[title={\lblExercise~\theexercise\ \dash{} #2}]%
}{%
  \par\medskip
  {\small\setlength{\parindent}{0pt}%
  \lblExerciseStarter: \texttt{code/exercises/\exercisedir/\detokenize{#1}.py}\\
  \lblExerciseTest: \texttt{code/exercises/\exercisedir/test\_\detokenize{#1}.py}\\
  \lblExerciseRun: \texttt{cd code \&\& uv run pytest -k \detokenize{#1}}\par}%
  \end{exercisebox}%
}

% ============================================================
%  DIAGRAMS --- Mermaid pipeline
%  ---------------------------------------------------------
%  \mermaidfig{key}{caption}{one-line manifest description}
%
%  Source of truth is figures/mermaid/<lang>/<key>.mmd, committed. `make
%  diagrams` renders each to figures/diagrams/<lang>/<key>.pdf, which is build
%  output and is not committed, so a diagram change reviews as a text diff.
%  The source is per-language because a diagram with English labels in the
%  Polish edition is exactly the half-translation that makes a bilingual book
%  feel machine-made; CI checks that both languages carry the same keys.
%
%  If the rendered PDF is absent the macro draws a placeholder AND typesets
%  the Mermaid source, in the flow rather than as a float, because a source
%  typeset verbatim is often taller than a page and a float cannot break.
% ============================================================
\makeatletter
\newcommand{\listofdiagrams}{%
  \section*{\lblDiagramManifest}%
  \phantomsection\addcontentsline{toc}{section}{\lblDiagramManifest}%
  % The entries are subsection-level and tocdepth is 1 for the contents,
  % so the depth is raised for this list alone or it prints empty --
  % which it did, on the scaffold's first clean build.
  \begingroup\c@tocdepth=2\relax\@starttoc{dgm}\endgroup%
}
\makeatother

\newcommand{\mermaidfig}[3]{%
  % \phantomsection is load-bearing and not decoration. \addcontentsline
  % records hyperref's CURRENT anchor, and here that is whatever preceded
  % the figure -- for a figure placed after a listing it is the listing's
  % line-number anchor, and listings has already advanced the counter past
  % the last line it printed, so the manifest entry linked to a
  % destination that is never created: "name{lstnumber.10.4.46} has been
  % referenced but does not exist, replaced by a fixed one". The two
  % front-matter figures had the same wrong target and were silent about
  % it, because the lines they named happen to be printed. The exercise
  % manifest never had the bug, because \begin{exercise} steps a counter
  % and that is what sets the anchor.
  \phantomsection%
  \addcontentsline{dgm}{subsection}{\protect\numberline{}\texttt{#1.mmd} --- #3}%
  \IfFileExists{figures/diagrams/\booklang/#1.pdf}{%
    \begin{figure}[htbp]
    \centering
    \includegraphics[width=0.95\linewidth,height=0.42\textheight,keepaspectratio]{figures/diagrams/\booklang/#1.pdf}%
    \caption{#2}\label{fig:#1}
    \end{figure}%
  }{%
    \par\medskip
    \begin{tcolorbox}[
      enhanced, breakable, width=\linewidth, sharp corners,
      colback=white, colframe=mgrey, boxrule=0.8pt,
      left=8pt, right=8pt, top=8pt, bottom=8pt]
      {\bfseries\color{mgrey}\small \lblNotRendered}\\[4pt]
      {\ttfamily\scriptsize\color{mgrey} figures/mermaid/\booklang/#1.mmd}\\[6pt]
      \IfFileExists{figures/mermaid/\booklang/#1.mmd}{%
        \lstinputlisting[style=bookcode,numbers=none,basicstyle=\ttfamily\scriptsize,
                         backgroundcolor=\color{mlight},frame=none,
                         keywordstyle=\color{black},language={}]{figures/mermaid/\booklang/#1.mmd}%
      }{%
        {\small\color{mred} \lblSourceMissing: \texttt{figures/mermaid/\booklang/#1.mmd}}%
      }%
    \end{tcolorbox}%
    \captionof{figure}{#2}\label{fig:#1}
    \par\medskip
  }%
}

% ============================================================
%  CHAPTER STUBS
%  ---------------------------------------------------------
%  Prints a visible marker so an unwritten chapter cannot be mistaken for a
%  written one and the page count never flatters the draft. `make debt`
%  counts them, and CI publishes the count on every build.
% ============================================================
\newcommand{\chapterstub}[1]{%
  \begin{tcolorbox}[
    enhanced, breakable, width=\linewidth, sharp corners,
    colback=mlight, colframe=mred, boxrule=1pt,
    borderline={0.6pt}{3pt}{mred, dashed},
    left=12pt, right=12pt, top=12pt, bottom=12pt]
    {\bfseries\color{mred}\large \lblNotWritten}\\[8pt]
    \small\raggedright #1
  \end{tcolorbox}%
}

% ============================================================
%  COMPUTED VALUES
%  ---------------------------------------------------------
%  A number that came out of a measurement is written by a script under
%  code/measure/ into figures/values/<name>.tex as
%
%      \pyval{e01.threads.speedup}{3.71}
%
%  and pulled into the text with \val{e01.threads.speedup}. The script is the
%  source of truth and the book cannot disagree with it, because the book does
%  not contain the digits. `make verify` fails when a committed value no
%  longer matches its script. A value the scripts do not produce prints a
%  visible marker.
%
%  \val also applies the edition's decimal separator, which is how the Polish
%  edition gets its decimal comma without a translator having to remember.
% ============================================================
\newcommand{\bookdecimalmarker}{\ifpl{,}{.}}
\IfFileExists{siunitx.sty}{%
  \usepackage{siunitx}%
  \sisetup{
    group-digits = integer,
    group-minimum-digits = 5,
    group-separator = {\,},
    output-decimal-marker = {\bookdecimalmarker}
  }%
  \newcommand{\booknum}[1]{\num{##1}}%
}{%
  \newcommand{\booknum}[1]{##1}%
}

\newcommand{\pyvalues}{}
\makeatletter
\newcommand{\pyval}[2]{\csgdef{pyval@#1}{#2}\listxadd{\pyvalues}{#1}}
% \num on a non-numeric body is a FATAL siunitx error. The generator knows
% whether it emitted a number or a word, so it writes \pyval for a number and
% \pyvaltext for a word, and this end only enforces which macro reads which.
\newcommand{\val}[1]{%
  \ifcsdef{pyval@#1}{%
    \ifcsdef{pyvaltext@#1}{%
      \PackageError{pybook}{Value `#1' is not a number}%
        {\string\val\space passes its body to siunitx, which refuses anything
         that is not a number. Use \string\valtext{#1} instead.}%
    }{}%
    \booknum{\csuse{pyval@#1}}%
  }{\textcolor{mred}{\textbf{[\lblValueMissing: #1]}}}%
}
\newcommand{\valtext}[1]{%
  \ifcsdef{pyval@#1}{\csuse{pyval@#1}}%
    {\textcolor{mred}{\textbf{[\lblValueMissing: #1]}}}%
}
\newcommand{\pyvaltext}[2]{%
  \csgdef{pyval@#1}{#2}\csgdef{pyvaltext@#1}{}\listxadd{\pyvalues}{#1}%
}
\makeatother

% figures/values/all.tex is generated by `make numbers` and does nothing but
% \input each value file. A build with no values yet still compiles; every
% \val{} in it prints a visible marker instead of a number.
\IfFileExists{figures/values/all.tex}{\input{figures/values/all}}{%
  \AtBeginDocument{\PackageWarning{pybook}{No computed values found. Run `make numbers`.}}}

% ============================================================
%  PINNED VERSIONS --- single source of truth
%  Change these here; they propagate through both editions and Appendix C.
%  Verified against pypi.org on \pinnedon (lang/<lang>.tex).
%  tools/check_versions.py compares every one of them with PyPI on each build.
%  Python itself is not on PyPI and is checked by hand against python.org.
% ============================================================
\newcommand{\bookedition}{0.1-dev}
\newcommand{\bookyear}{2026}
\newcommand{\pyver}{3.14}              % the minor the book targets
\newcommand{\pypatch}{3.14.7}          % the exact interpreter the code ran on
\newcommand{\uvver}{0.12.13}           % uv
\newcommand{\dotnetver}{10.0.100}      % dotnet
\newcommand{\ruffver}{0.16.7}          % ruff
\newcommand{\pyrightver}{1.1.414}      % pyright
\newcommand{\mypyver}{2.3.1}           % mypy
\newcommand{\pydanticver}{2.13.5}      % pydantic
\newcommand{\pysettingsver}{2.15.0}    % pydantic-settings
\newcommand{\fastapiver}{0.141.1}      % fastapi
\newcommand{\uvicornver}{0.53.0}       % uvicorn
\newcommand{\sqlalchemyver}{2.0.52}    % sqlalchemy
\newcommand{\alembicver}{1.20.0}       % alembic
\newcommand{\aiosqlitever}{0.22.1}     % aiosqlite
\newcommand{\pytestver}{9.1.1}         % pytest
\newcommand{\pytestasyncver}{1.4.0}    % pytest-asyncio
\newcommand{\hypothesisver}{6.168.0}   % hypothesis
\newcommand{\coveragever}{7.16.1}      % coverage
\newcommand{\httpxver}{0.28.1}         % httpx
\newcommand{\structlogver}{26.1.0}     % structlog
\newcommand{\otelver}{1.44.0}          % opentelemetry-sdk
\newcommand{\polarsver}{1.44.2}        % polars
\newcommand{\anthropicver}{1.5.0}      % anthropic
\newcommand{\openaiver}{3.13.0}        % openai
\newcommand{\pipauditver}{2.10.1}      % pip-audit

% ============================================================
%  INDEX
%  Two-column, and its entries include \texttt{} names that do not hyphenate.
%  Ragged right keeps multi-word entries off the margin; it cannot help a
%  single token wider than the column, so keep verbatim index entries under
%  about 29 characters and index longer names by concept instead.
%  \raggedcolumns is the switch imakeidx's multicols actually reads.
% ============================================================
\apptocmd{\theindex}{\raggedcolumns\raggedright}{}{%
  \PackageWarning{pybook}{Could not patch theindex for ragged setting}}
\providecommand*\see[2]{}
\renewcommand*\see[2]{\emph{\lblIndexSee} #1}
\providecommand*\seealso[2]{}
\renewcommand*\seealso[2]{\emph{\lblIndexSeeAlso} #1}

\makeindex
; nothing else differs between them at this level.
%  Every user-visible string lives in lang/en.tex or lang/pl.tex, so a label
%  that appears in one edition and not the other is a missing macro rather
%  than a silent divergence.
%
%  Lineage. The bilingual wiring (\booklang, lang/, body.tex, the generated
%  structure) is the math book's; the chapter machinery (listings, the
%  admonitions, verifybox, \pyfile and \pyregion, the Mermaid pipeline) is the
%  LangChain book's. What is new here is the exercise macro, which ties every
%  exercise to a starter file and a test under code/, because this book is
%  read with an editor open beside it.
% ============================================================

% \booklang must be set by the main file. Default to English rather than
% failing, so a stray % ============================================================
%  preamble.tex --- styling, environments, macros
%
%  Shared by BOTH editions. main-en.tex and main-pl.tex each define \booklang
%  before \input{preamble}; nothing else differs between them at this level.
%  Every user-visible string lives in lang/en.tex or lang/pl.tex, so a label
%  that appears in one edition and not the other is a missing macro rather
%  than a silent divergence.
%
%  Lineage. The bilingual wiring (\booklang, lang/, body.tex, the generated
%  structure) is the math book's; the chapter machinery (listings, the
%  admonitions, verifybox, \pyfile and \pyregion, the Mermaid pipeline) is the
%  LangChain book's. What is new here is the exercise macro, which ties every
%  exercise to a starter file and a test under code/, because this book is
%  read with an editor open beside it.
% ============================================================

% \booklang must be set by the main file. Default to English rather than
% failing, so a stray \input{preamble} still compiles and says what it did.
\providecommand{\booklang}{en}

\usepackage{etoolbox}   % loaded first: the language switch below needs it

% ---------- Page geometry ----------
% ONE format: A4, 12pt, ONE-SIDED. This book is read on a screen with an
% editor open beside it, and never printed as a bound book, so there is no
% recto, no verso, no gutter, no mirrored margin and no blank leaf before a
% chapter -- every one of those is a property of paper that a PDF viewer at
% half a screen's width turns into wasted pixels.
%
% The margins are symmetric and the measure is about seventy-five characters
% of prose, which is where the eye keeps the line return. The number that was
% actually chosen is the MONOSPACE measure: at \footnotesize on this block a
% listing line of 79 characters fits without wrapping, and 79 is the width
% every listing under code/ is held to by tools/check_structure.py --lines.
% A wrapped listing line is not an error TeX reports -- breaklines=true wraps
% it silently and prints an arrow into the middle of what the reader is meant
% to paste into an editor -- so the width is enforced in the source instead.
\usepackage[
  a4paper,
  left=3.1cm, right=3.1cm, top=2.6cm, bottom=3.0cm,
  head=26pt, headsep=0.8cm, footskip=1.4cm
]{geometry}

% ---------- Encoding & fonts ----------
\usepackage[T1]{fontenc}
\usepackage[utf8]{inputenc}
\IfFileExists{lmodern.sty}{\usepackage{lmodern}}{}
% newtxtext takes its TS1 symbols (\textcopyright, \textdegree) from TeX
% Gyre Termes, which Debian packages SEPARATELY (tex-gyre): a machine with
% newtx and without it dies on the copyright page with
% `Font TS1/ntxtlf/m/n/10.95=ts1-qtmr ... not loadable`. There is no
% inert-safe probe for a font metric in pdfTeX -- \IfFileExists searches
% TEXINPUTS and not TFMFONTS, so it reports the metric missing on a machine
% that has it, and \suppressfontnotfounderror is LuaTeX-only -- so nothing
% here guesses. CI's full TeX Live is the reference; locally, install
% tex-gyre. Both facts are recorded in CLAUDE.md under Build traps.
\IfFileExists{newtxtext.sty}{\usepackage{newtxtext}}{}
% amsmath before newtxmath, because newtxmath patches amsmath's internals.
% amssymb only when newtxmath is absent: newtxmath supplies the AMS symbols
% itself and loading both is a fatal clash (\Bbbk already defined) that is
% invisible on a machine without newtx and fatal on a full TeX Live.
\usepackage{amsmath}
\IfFileExists{newtxmath.sty}{\usepackage{newtxmath}}{\usepackage{amssymb}}
\IfFileExists{inconsolata.sty}{\usepackage[varqu,scaled=0.95]{inconsolata}}{}
\IfFileExists{lmodern.sty}{\usepackage{microtype}}{\usepackage[expansion=false]{microtype}}

% upquote is not a refinement. In T1 the input character ' is quoteright, so
% a listings body sets f('x') with U+2019 at both ends of the literal, and
% typed back in that is a SyntaxError. inconsolata's varqu option hides the
% defect on a machine that has inconsolata; upquote makes ' the ASCII
% apostrophe whatever the typewriter font is. In a book whose every listing is
% meant to be pasted into an editor, this line is load-bearing.
\IfFileExists{upquote.sty}{\usepackage{upquote}}{}

% And upquote's twin, for the same reason one character over. In T1 the pair
% -- is a ligature for an en dash, in EVERY family including the typewriter
% one, so \code{uv sync --locked} sets as "uv sync <en dash>locked" and a
% reader who copies it off the page gets an unknown-argument error from a
% command the chapter told them to run. It is invisible in review because it
% looks like a dash, and it bites exactly the flags a packaging chapter is
% made of. A `listings` body is already safe -- listings sets characters
% singly and no ligature forms -- which is what makes the defect so easy to
% miss: the same text is right inside a listing and wrong in the sentence
% above it.
%
% Scoped to tt* on purpose, and measured rather than assumed: prose keeps its
% own dashes, so \dash{} still sets an em dash and a range like 1--2 still
% sets an en dash. Probed against this preamble in all three positions before
% it was believed.
\DisableLigatures[-]{encoding = T1, family = tt*}

% ---------- Language ----------
% Loading babel with a language whose .ldf is absent is a *fatal* error, not
% a warning, so both the package and each language file are probed before
% either is requested. Both editions load the other language too, because the
% glossary rows and the cheat sheet carry words from both.
\IfFileExists{babel.sty}{%
  \IfFileExists{polish.ldf}{%
    \IfFileExists{british.ldf}{%
      \ifdefstring{\booklang}{pl}%
        {\usepackage[british,main=polish]{babel}}%
        {\usepackage[polish,main=british]{babel}}%
    }{\usepackage[polish]{babel}}%
  }{%
    \IfFileExists{british.ldf}{\usepackage[british]{babel}}{}%
  }%
}{}

% babel-polish makes " an active shorthand. listings reads its body verbatim
% and an active character in a verbatim context is a category-code accident
% waiting to happen, so the shorthand is turned off. Polish quotation marks
% come from \enquote, which csquotes sets per language.
\ifdefstring{\booklang}{pl}{\AtBeginDocument{\shorthandoff{"}}}{}

% ---------- Core packages ----------
\usepackage{graphicx}
\usepackage{xcolor}
\usepackage{booktabs}
\usepackage{tabularx}
% Appendix A's cheat sheet runs past a page in every part, and tabularx
% cannot break. longtable can, and works with booktabs.
\usepackage{longtable}
\usepackage{array}
\usepackage{enumitem}
\usepackage{listings}
\usepackage[most]{tcolorbox}
\usepackage{fancyhdr}
\usepackage{titlesec}
\usepackage{caption}
\usepackage{imakeidx}
% csquotes with autostyle follows babel's active language, so the identical
% source \enquote{x} sets 'x' in the English edition and the Polish quotation
% marks in the Polish one.
\IfFileExists{csquotes.sty}{\usepackage[autostyle=true]{csquotes}}{%
  \providecommand{\enquote}[1]{``##1''}}
\usepackage[hidelinks]{hyperref}

% ---------- Palette ----------
% The same family as the two companion volumes, so the three read as a set.
\definecolor{mblue}{HTML}{1F4E79}
\definecolor{mteal}{HTML}{0E7C7B}
\definecolor{mamber}{HTML}{B26A00}
\definecolor{mred}{HTML}{9B2C2C}
\definecolor{mviolet}{HTML}{7B4397}
\definecolor{mgrey}{HTML}{5A5A5A}
\definecolor{mlight}{HTML}{F4F6F8}
\definecolor{mfaint}{HTML}{FBFCFD}
\definecolor{codebg}{HTML}{FAFAFA}
\definecolor{codeframe}{HTML}{D8DEE4}
\definecolor{kw}{HTML}{0B5394}
\definecolor{str}{HTML}{9B2C2C}
\definecolor{cmt}{HTML}{4F7A4F}
\definecolor{typ}{HTML}{0E7C7B}
\definecolor{dec}{HTML}{7B4397}

\hypersetup{
  colorlinks=true,
  linkcolor=mblue,
  urlcolor=mteal,
  citecolor=mblue,
  pdfauthor={Konrad Cinkusz}
}

% \ifpl{Polish text}{English text} -- for the handful of places where a
% string is structural rather than content. Anything longer than a few words
% belongs in lang/ or in the edition's own source, not here.
% \DeclareRobustCommand, not \newcommand: a fragile \ifpl inside a moving
% argument is written to the .toc one level expanded and fails on the next run
% in exactly one of the two editions.
\DeclareRobustCommand{\ifpl}[2]{\ifdefstring{\booklang}{pl}{#1}{#2}}
\makeatletter
\pdfstringdefDisableCommands{%
  \ifdefstring{\booklang}{pl}{\let\ifpl\@firstoftwo}{\let\ifpl\@secondoftwo}%
  % \api{}, \pkg{} and \code{} are allowed in a section title, and a title
  % becomes a PDF bookmark. hyperref drops the \color command from a bookmark
  % string but keeps its ARGUMENT, so a title carrying \api{model\_validate}
  % made a bookmark reading "mtealmodel_validate" and warned "Token not
  % allowed" twice. The colour name is gobbled here instead. Measured on a
  % probe against this preamble before the line was added.
  \def\color#1{}%
}
\makeatother

% ---------- Language strings ----------
\input{lang/\booklang}

% The PDF's subject, keyword list and document language read \lblPdf... from
% the language file input on the line above, so they cannot sit in the colour
% \hypersetup block, which runs before it. Two blocks is the ordering, not an
% oversight; merging them raises Undefined control sequence and still writes
% a PDF, with the three fields empty.
\hypersetup{
  pdfsubject={\lblPdfSubject},
  pdfkeywords={\lblPdfKeywords},
  pdflang={\lblPdfLang}
}

% ---------- Headings ----------
% \raggedright on the title body is load-bearing: at \Huge on this block a
% title that cannot break overflows the margin by 30 pt or more, and several
% chapter titles here run past fifty characters.
\titleformat{\chapter}[display]
  {\normalfont\huge\bfseries\color{mblue}\raggedright}
  {\normalfont\normalsize\color{mgrey}\MakeUppercase{\chaptertitlename}\ \thechapter}
  {8pt}{\Huge\raggedright}
\titlespacing*{\chapter}{0pt}{10pt}{28pt}
\titleformat{\section}{\normalfont\Large\bfseries\color{mblue}\raggedright}{\thesection}{0.8em}{}
\titleformat{\subsection}{\normalfont\large\bfseries\color{mgrey}\raggedright}{\thesubsection}{0.7em}{}

% Sections in the contents, subsections not: fourteen chapters of subsections
% is a contents nobody reads, and the reader navigates by chapter and section.
\setcounter{tocdepth}{1}
\setcounter{secnumdepth}{2}

% One-sided running head: the chapter on the left, the folio on the right.
\pagestyle{fancy}
\fancyhf{}
\fancyhead[L]{\small\nouppercase{\leftmark}}
\fancyhead[R]{\small\thepage}
\renewcommand{\headrulewidth}{0.4pt}
\fancypagestyle{plain}{\fancyhf{}\fancyfoot[R]{\small\thepage}%
  \renewcommand{\headrulewidth}{0pt}}

% Drops the word "Chapter" from the head: the number is still there and the
% longer titles no longer overflow. MUST come after \pagestyle{fancy}, which
% installs its own \chaptermark and silently overwrites anything earlier.
\renewcommand{\chaptermark}[1]{\markboth{\thechapter.\ #1}{}}

% ============================================================
%  CODE LISTINGS
%  ---------------------------------------------------------
%  Every listing in a chapter is a FILE under code/, pulled in with \pyfile or
%  \pyregion, and CI runs it against the pinned versions. An in-source
%  \begin{python} block is allowed for a fragment of three or four lines that
%  exists to show a shape rather than to run; anything longer belongs in a
%  file the reader can open in the editor beside the PDF.
% ============================================================
\lstdefinelanguage{pythonx}{
  language=Python,
  morekeywords={
    async,await,match,case,nonlocal,yield,type,
    TypedDict,Annotated,Protocol,Literal,Generic,Final,ClassVar,Self,
    TypeVar,TypeGuard,TypeIs,Callable,Iterator,Iterable,Awaitable,
    dataclass,field,BaseModel,Field,Enum,StrEnum,NamedTuple,
    override,overload,property,staticmethod,classmethod
  },
  morekeywords=[2]{
    gather,create_task,TaskGroup,to_thread,timeout,wait_for,Semaphore,Queue,
    fixture,parametrize,monkeypatch,raises,mark,
    Depends,FastAPI,APIRouter,HTTPException,
    select,Session,Mapped,mapped_column,relationship,
    model_validate,model_validate_json,model_dump,model_json_schema
  },
  sensitive=true
}

% listings' C# definition predates async, records and pattern matching.
\lstdefinelanguage{csharpx}{
  language=[Sharp]C,
  morekeywords={
    async,await,var,dynamic,record,init,required,nameof,yield,partial,
    global,scoped,when,with,and,or,not,notnull,unmanaged,file,get,set,value,
    where,select,from,orderby,let,join,into,ascending,descending,on,equals,
    Task,ValueTask,CancellationToken,IEnumerable,IAsyncEnumerable,Func,Action
  },
  sensitive=true
}

\lstdefinestyle{bookcode}{
  basicstyle=\ttfamily\footnotesize,
  keywordstyle=\color{kw}\bfseries,
  keywordstyle=[2]\color{typ},
  stringstyle=\color{str},
  % Colour only, no \itshape: inconsolata (zi4) has no italic shape, so an
  % italic comment style is substituted with upright on every machine that
  % has the font, and the log carries "Font shape `T1/zi4/m/it' undefined"
  % on every build. The colour already marks a comment.
  commentstyle=\color{cmt},
  emphstyle=\color{dec}\bfseries,
  numbers=left,
  numberstyle=\ttfamily\tiny\color{mgrey},
  numbersep=8pt,
  backgroundcolor=\color{codebg},
  frame=single,
  rulecolor=\color{codeframe},
  framesep=6pt,
  breaklines=true,
  breakatwhitespace=false,
  postbreak=\mbox{\textcolor{mgrey}{$\hookrightarrow$}\space},
  showstringspaces=false,
  tabsize=4,
  columns=fullflexible,
  keepspaces=true,
  captionpos=b,
  aboveskip=1.2em,
  belowskip=1.2em,
  % Maps the characters most likely to be pasted into a listing by accident.
  %
  % Do NOT add a mapping for U+00A0 (non-breaking space). It looks like the
  % obvious next entry and it makes `listings` abort the run with "Improper
  % alphabetic constant" -- fatal, no PDF, and the message names neither the
  % character nor this line. The ASCII-in-listings convention, checked by
  % parity's C13, is the guard against it instead.
  literate={ł}{{\l}}1 {Ł}{{\L}}1 {ą}{{\k a}}1 {ę}{{\k e}}1
           {ś}{{\'s}}1 {ć}{{\'c}}1 {ż}{{\.z}}1 {ź}{{\'z}}1 {ó}{{\'o}}1 {ń}{{\'n}}1
           {Ą}{{\k A}}1 {Ę}{{\k E}}1 {Ś}{{\'S}}1 {Ć}{{\'C}}1
           {Ż}{{\.Z}}1 {Ź}{{\'Z}}1 {Ó}{{\'O}}1 {Ń}{{\'N}}1
           {—}{{-{}-}}2 {–}{{-}}1 {‘}{{'}}1 {’}{{'}}1
           {“}{{"}}1 {”}{{"}}1 {°}{{\textdegree}}1
}

\lstset{style=bookcode}

\lstnewenvironment{python}[1][]
  {\lstset{language=pythonx,style=bookcode,#1}}
  {}

% C# comparison snippets. The whole book is written for engineers arriving
% from .NET, so listings come in pairs often enough that the C# half has its
% own environment rather than a generic one.
\lstnewenvironment{csharp}[1][]
  {\lstset{language=csharpx,style=bookcode,#1}}
  {}

\lstnewenvironment{shellcmd}[1][]
  {\lstset{language=bash,style=bookcode,numbers=none,keywordstyle=\color{mteal}\bfseries,#1}}
  {}

\lstnewenvironment{tomlcode}[1][]
  {\lstset{style=bookcode,numbers=none,keywordstyle=\color{black},#1}}
  {}

\lstnewenvironment{yamlcode}[1][]
  {\lstset{style=bookcode,numbers=none,keywordstyle=\color{black},#1}}
  {}

\lstnewenvironment{jsoncode}[1][]
  {\lstset{style=bookcode,numbers=none,keywordstyle=\color{black},#1}}
  {}

% Terminal transcripts typed into the source. Allowed for a shell session the
% reader is meant to reproduce; NOT for program output, which is \transcript
% below, because an in-source transcript nobody ran reads exactly like one
% that was, and that is where a remembered number hides.
\lstnewenvironment{console}[1][]
  {\lstset{style=bookcode,numbers=none,basicstyle=\ttfamily\footnotesize,
           keywordstyle=\color{black},backgroundcolor=\color{white},#1}}
  {}

% The path is printed above every file listing, in the repository-relative
% form the reader opens in the editor beside the PDF. That is the whole point
% of a listing being a file, and the front matter promises it. The \nobreak
% keeps the path on the same page as the listing it names.
% ---------------------------------------------------------------------------
%  Appendix B prints the trap catalogue from notes/02-traps.md. One entry is
%  its number -- which a chapter may cite, and which is never reused -- the
%  habit in the reader's own voice, and what Python does instead.
%
%  The body is set RAGGED RIGHT on purpose. An entry is mostly \code{} spans,
%  which are unbreakable runs, and sixty-seven justified paragraphs of them
%  is sixty-seven chances at an overfull hbox in the edition with the longer
%  words. Ragged right removes the class rather than fixing it one box at a
%  time, and a reference list is the one place in this book where it reads
%  better anyway.
%
%  The habit is set with \enquote{} rather than \emph{}, and that is not a
%  style choice. Every habit here carries \code{} spans, and \code{} inside
%  \emph{} asks inconsolata for an italic it does not have: the first build
%  of this appendix raised `Font shape T1/zi4/m/it undefined', which
%  checklog.py is hard on for exactly this reason. It is the trap CLAUDE.md
%  records, met sixty-seven times at once -- and fixed once, because the
%  entries go through a macro. A quotation is also what the habit IS, and
%  \enquote{} is what the Polish edition owes anyway.
% ---------------------------------------------------------------------------
%  Appendix A's cheat sheet: one table per part, C# on the left, Python on
%  the right, and the chapter that teaches the row on the end. The three
%  header strings are arguments so that one environment serves both
%  editions and the structural signature stays identical.
%
%  Columns are ragged right for the reason Appendix B's entries are: a cell
%  is mostly \code{} spans, which are unbreakable runs, and a justified
%  narrow column full of them is an overfull box waiting for the edition
%  with the longer words.
% ---------------------------------------------------------------------------
%  Appendix C's tool matrix: the .NET tool, the Python tool this book uses,
%  the version it was verified against and the chapter that introduces it.
%
%  The version column prints \pydanticver and its siblings and is never
%  typed, which is the whole point of the appendix: it is the one page where
%  a reader sees every pin at once, and a pin that had to be re-typed here
%  would be a second copy of a fact the preamble already owns.
%  tools/check_structure.py --tools fails when a pin macro exists and this
%  table does not print it.
\NewDocumentEnvironment{toolmatrix}{mmmm}{%
  \small
  \begin{longtable}{@{}%
    >{\raggedright\arraybackslash}p{0.27\linewidth}%
    >{\raggedright\arraybackslash}p{0.31\linewidth}%
    >{\raggedright\arraybackslash}p{0.20\linewidth}%
    >{\raggedright\arraybackslash}p{0.08\linewidth}@{}}
  \toprule
  \textbf{#1} & \textbf{#2} & \textbf{#3} & \textbf{#4} \\
  \midrule
  \endfirsthead
  \toprule
  \textbf{#1} & \textbf{#2} & \textbf{#3} & \textbf{#4} \\
  \midrule
  \endhead
  \midrule
  \endfoot
  \bottomrule
  \endlastfoot
}{%
  \end{longtable}%
}

\NewDocumentEnvironment{cheatsheet}{mmm}{%
  \small
  \begin{longtable}{@{}%
    >{\raggedright\arraybackslash}p{0.41\linewidth}%
    >{\raggedright\arraybackslash}p{0.41\linewidth}%
    >{\raggedright\arraybackslash}p{0.08\linewidth}@{}}
  \toprule
  \textbf{#1} & \textbf{#2} & \textbf{#3} \\
  \midrule
  \endfirsthead
  \toprule
  \textbf{#1} & \textbf{#2} & \textbf{#3} \\
  \midrule
  \endhead
  \midrule
  \endfoot
  \bottomrule
  \endlastfoot
}{%
  \end{longtable}%
}

\newcommand{\trapentry}[3]{%
  \par\medskip
  \noindent\textbf{#1.}\enspace\enquote{#2}\par
  \nopagebreak
  {\raggedright\noindent#3\par}%
}

\newcommand{\listingpath}[1]{%
  \par\medskip\noindent
  {\ttfamily\footnotesize\color{mgrey}\detokenize{#1}}%
  \par\nobreak}

% External source file -- the preferred form. The path is relative to the
% repository root, which is also the path the reader opens in the editor.
% Usage: \pyfile{code/ch05/decorators.py}{Caption}{lst:label}
\newcommand{\pyfile}[3]{%
  \listingpath{#1}%
  \lstinputlisting[language=pythonx,style=bookcode,aboveskip=0.3em,
    caption={#2},label={#3}]{#1}%
}

% A named region of an external file, delimited in the source by
% `# --8<-- [start:name]` / `# --8<-- [end:name]`, so the book and the running
% code cannot drift apart. tools/check_structure.py --listings fails when the
% file or the region is missing.
%
% The markers are matched through listings' range prefix and suffix keys,
% with linerange={name-name}. The form this was inherited with -- a linerange
% whose two ends were the whole marker strings -- matched nothing and printed
% the WHOLE FILE, with no warning: measured on a three-region probe file
% before this was rewritten, and the same definition sits unused in the
% sibling book it came from. A region name is a word, never a number, because
% a number is a line number to listings. The line numbers printed are the
% file's own, so a region's listing says where in the file it sits.
% Usage: \pyregion{code/ch05/decorators.py}{timing}{Caption}{lst:label}
\lstset{
  rangebeginprefix=\#\ --8<--\ [start:, rangebeginsuffix=],
  rangeendprefix=\#\ --8<--\ [end:, rangeendsuffix=],
  includerangemarker=false}
\newcommand{\pyregion}[4]{%
  \listingpath{#1}%
  \lstinputlisting[language=pythonx,style=bookcode,aboveskip=0.3em,
    linerange={#2-#2},caption={#3},label={#4}]{#1}%
}

% The C# twin of \pyfile, for the comparison half of a pair.
\newcommand{\csfile}[3]{%
  \listingpath{#1}%
  \lstinputlisting[language=csharpx,style=bookcode,aboveskip=0.3em,
    caption={#2},label={#3}]{#1}%
}

% A program's output, written by a script under code/measure/ and pulled in
% whole. There is no typed form of this and there must not be: \val{} does
% not expand inside a verbatim body, so a transcript typed into a chapter is
% one that cannot quote a computed value, and a transcript nobody ran is
% indistinguishable on the page from one that was.
\newcommand{\transcript}[1]{%
  \par\smallskip
  \IfFileExists{figures/transcripts/#1.txt}{%
    \lstinputlisting[style=bookcode,numbers=none,basicstyle=\ttfamily\footnotesize,
      keywordstyle=\color{black},backgroundcolor=\color{white},
      language={}]{figures/transcripts/#1.txt}%
  }{%
    \begin{tcolorbox}[enhanced, sharp corners, width=\linewidth,
      colback=mlight, colframe=mgrey, boxrule=0.6pt,
      left=8pt, right=8pt, top=6pt, bottom=6pt]
      {\bfseries\color{mgrey}\small \lblTranscriptMissing}\\[2pt]
      {\ttfamily\scriptsize\color{mgrey} figures/transcripts/#1.txt}
    \end{tcolorbox}%
  }%
  \par\smallskip
}

% Inline literals. Underscores inside these must be written \_.
\newcommand{\code}[1]{\texttt{\small #1}}
\newcommand{\pkg}[1]{\texttt{\small\color{mblue}#1}}
\newcommand{\api}[1]{\texttt{\small\color{mteal}#1}}

% ============================================================
%  ADMONITIONS
%  ---------------------------------------------------------
%  A deliberately small, fixed vocabulary. Two visual treatments carry the
%  whole grammar: `tinted` is a block the reader may read in passing -- a
%  tint and a thick left bar, no outline; `outlined` is a block the reader
%  must stop at, with a full rule on white. A third treatment would mean the
%  first two are not carrying their meaning.
% ============================================================
\tcbset{
  mfa tinted/.style={
    enhanced, breakable, sharp corners,
    colback=mlight, boxrule=0pt, leftrule=3pt,
    left=8pt, right=8pt, top=6pt, bottom=6pt,
    before skip=13pt, after skip=13pt,
    fonttitle=\bfseries\small,
    coltitle=black, colbacktitle=mlight,
    attach title to upper={\par\smallskip},
  },
  mfa outlined/.style={
    enhanced, breakable, sharp corners,
    colback=white, boxrule=0.6pt,
    left=8pt, right=8pt, top=6pt, bottom=6pt,
    before skip=13pt, after skip=13pt,
    fonttitle=\bfseries\small,
    coltitle=black, colbacktitle=white,
    attach title to upper={\par\smallskip},
  },
}

% Read in passing.
\newtcolorbox{note}[1][]{mfa tinted, colframe=mblue, title={\lblNote}, #1}
\newtcolorbox{warning}[1][]{mfa tinted, colframe=mred, title={\lblWarning}, #1}
% The translation. The .NET reader is the book's only assumed reader, so this
% box carries the C# side of a comparison and nothing else. Its use is
% counted: a chapter with none is a chapter that has stopped translating.
\newtcolorbox{csbox}[1][]{mfa tinted, colframe=mteal, title={\lblCsharp}, #1}
% Anything that is preview, experimental or likely to move. Python moves more
% slowly than the agent frameworks do, and this box should be rare here.
\newtcolorbox{versionbox}[1][]{mfa tinted, colframe=mamber, title={\lblVersion}, #1}

% Stop and read.
% The misconception, named AFTER the reader has walked into it. Appendix B is
% the catalogue; a trap box in a chapter is one entry of it, in place.
\newtcolorbox{trapbox}[1][]{mfa outlined, colframe=mred, title={\lblTrap}, #1}
% Marks a listing that has not been executed against the pinned versions.
% `make debt` counts these, CI prints the count, and nothing ships to a reader
% with one attached. Removing it means you ran the code.
\newtcolorbox{verifybox}[1][]{mfa outlined, colframe=mgrey, title={\lblVerify}, #1}
\newtcolorbox{exercisebox}[1][]{mfa outlined, colframe=mteal, title={\lblExercise}, #1}
% The guiding project. Points at the stage of trace-assert this section
% builds, under code/src/trace_assert/.
\newtcolorbox{projectbox}[1][]{mfa outlined, colframe=mamber, title={\lblProject}, #1}

% ============================================================
%  EXERCISES
%  ---------------------------------------------------------
%  \begin{exercise}{key}{title} ... \end{exercise}
%
%  This is the mechanism the whole book is read through: the PDF on one half
%  of the screen, an editor on the other. An exercise is therefore not a
%  paragraph but a FILE the reader opens -- a starter under
%  code/exercises/chNN/<key>.py -- and a TEST that decides whether it is done,
%  test_<key>.py beside it. The box prints both paths and the one command
%  that runs the test, and the manifest at the back lists every exercise.
%
%  An ENVIRONMENT, not a macro with a body argument: a macro argument cannot
%  contain verbatim material, so an exercise written as \exercise{key}{title}
%  {body} could never show the reader a code fragment, and the first chapter
%  that tried would have died on `\begin{python}` inside the braces.
%
%  The key is the file stem, with underscores, and it is written with
%  \detokenize so the underscore reaches the page as a character rather than
%  as a subscript. tools/check_structure.py --exercises fails when the starter
%  or the test is missing, or when the key's chapter prefix disagrees with the
%  chapter it sits in.
% ============================================================
\newcounter{exercise}[chapter]
\renewcommand{\theexercise}{\thechapter.\arabic{exercise}}
\makeatletter
\newcommand{\exercisedir}{ch\two@digits{\value{chapter}}}
\newcommand{\listofexercises}{%
  \section*{\lblExerciseManifest}%
  \phantomsection\addcontentsline{toc}{section}{\lblExerciseManifest}%
  % The entries are subsection-level and tocdepth is 1 for the contents,
  % so the depth is raised for this list alone or it prints empty --
  % which it did, on the scaffold's first clean build.
  \begingroup\c@tocdepth=2\relax\@starttoc{exr}\endgroup%
}
\makeatother
\NewDocumentEnvironment{exercise}{mm}{%
  \refstepcounter{exercise}%
  \addcontentsline{exr}{subsection}{\protect\numberline{\theexercise}\texttt{\protect\detokenize{#1}} --- #2}%
  \begin{exercisebox}[title={\lblExercise~\theexercise\ \dash{} #2}]%
}{%
  \par\medskip
  {\small\setlength{\parindent}{0pt}%
  \lblExerciseStarter: \texttt{code/exercises/\exercisedir/\detokenize{#1}.py}\\
  \lblExerciseTest: \texttt{code/exercises/\exercisedir/test\_\detokenize{#1}.py}\\
  \lblExerciseRun: \texttt{cd code \&\& uv run pytest -k \detokenize{#1}}\par}%
  \end{exercisebox}%
}

% ============================================================
%  DIAGRAMS --- Mermaid pipeline
%  ---------------------------------------------------------
%  \mermaidfig{key}{caption}{one-line manifest description}
%
%  Source of truth is figures/mermaid/<lang>/<key>.mmd, committed. `make
%  diagrams` renders each to figures/diagrams/<lang>/<key>.pdf, which is build
%  output and is not committed, so a diagram change reviews as a text diff.
%  The source is per-language because a diagram with English labels in the
%  Polish edition is exactly the half-translation that makes a bilingual book
%  feel machine-made; CI checks that both languages carry the same keys.
%
%  If the rendered PDF is absent the macro draws a placeholder AND typesets
%  the Mermaid source, in the flow rather than as a float, because a source
%  typeset verbatim is often taller than a page and a float cannot break.
% ============================================================
\makeatletter
\newcommand{\listofdiagrams}{%
  \section*{\lblDiagramManifest}%
  \phantomsection\addcontentsline{toc}{section}{\lblDiagramManifest}%
  % The entries are subsection-level and tocdepth is 1 for the contents,
  % so the depth is raised for this list alone or it prints empty --
  % which it did, on the scaffold's first clean build.
  \begingroup\c@tocdepth=2\relax\@starttoc{dgm}\endgroup%
}
\makeatother

\newcommand{\mermaidfig}[3]{%
  % \phantomsection is load-bearing and not decoration. \addcontentsline
  % records hyperref's CURRENT anchor, and here that is whatever preceded
  % the figure -- for a figure placed after a listing it is the listing's
  % line-number anchor, and listings has already advanced the counter past
  % the last line it printed, so the manifest entry linked to a
  % destination that is never created: "name{lstnumber.10.4.46} has been
  % referenced but does not exist, replaced by a fixed one". The two
  % front-matter figures had the same wrong target and were silent about
  % it, because the lines they named happen to be printed. The exercise
  % manifest never had the bug, because \begin{exercise} steps a counter
  % and that is what sets the anchor.
  \phantomsection%
  \addcontentsline{dgm}{subsection}{\protect\numberline{}\texttt{#1.mmd} --- #3}%
  \IfFileExists{figures/diagrams/\booklang/#1.pdf}{%
    \begin{figure}[htbp]
    \centering
    \includegraphics[width=0.95\linewidth,height=0.42\textheight,keepaspectratio]{figures/diagrams/\booklang/#1.pdf}%
    \caption{#2}\label{fig:#1}
    \end{figure}%
  }{%
    \par\medskip
    \begin{tcolorbox}[
      enhanced, breakable, width=\linewidth, sharp corners,
      colback=white, colframe=mgrey, boxrule=0.8pt,
      left=8pt, right=8pt, top=8pt, bottom=8pt]
      {\bfseries\color{mgrey}\small \lblNotRendered}\\[4pt]
      {\ttfamily\scriptsize\color{mgrey} figures/mermaid/\booklang/#1.mmd}\\[6pt]
      \IfFileExists{figures/mermaid/\booklang/#1.mmd}{%
        \lstinputlisting[style=bookcode,numbers=none,basicstyle=\ttfamily\scriptsize,
                         backgroundcolor=\color{mlight},frame=none,
                         keywordstyle=\color{black},language={}]{figures/mermaid/\booklang/#1.mmd}%
      }{%
        {\small\color{mred} \lblSourceMissing: \texttt{figures/mermaid/\booklang/#1.mmd}}%
      }%
    \end{tcolorbox}%
    \captionof{figure}{#2}\label{fig:#1}
    \par\medskip
  }%
}

% ============================================================
%  CHAPTER STUBS
%  ---------------------------------------------------------
%  Prints a visible marker so an unwritten chapter cannot be mistaken for a
%  written one and the page count never flatters the draft. `make debt`
%  counts them, and CI publishes the count on every build.
% ============================================================
\newcommand{\chapterstub}[1]{%
  \begin{tcolorbox}[
    enhanced, breakable, width=\linewidth, sharp corners,
    colback=mlight, colframe=mred, boxrule=1pt,
    borderline={0.6pt}{3pt}{mred, dashed},
    left=12pt, right=12pt, top=12pt, bottom=12pt]
    {\bfseries\color{mred}\large \lblNotWritten}\\[8pt]
    \small\raggedright #1
  \end{tcolorbox}%
}

% ============================================================
%  COMPUTED VALUES
%  ---------------------------------------------------------
%  A number that came out of a measurement is written by a script under
%  code/measure/ into figures/values/<name>.tex as
%
%      \pyval{e01.threads.speedup}{3.71}
%
%  and pulled into the text with \val{e01.threads.speedup}. The script is the
%  source of truth and the book cannot disagree with it, because the book does
%  not contain the digits. `make verify` fails when a committed value no
%  longer matches its script. A value the scripts do not produce prints a
%  visible marker.
%
%  \val also applies the edition's decimal separator, which is how the Polish
%  edition gets its decimal comma without a translator having to remember.
% ============================================================
\newcommand{\bookdecimalmarker}{\ifpl{,}{.}}
\IfFileExists{siunitx.sty}{%
  \usepackage{siunitx}%
  \sisetup{
    group-digits = integer,
    group-minimum-digits = 5,
    group-separator = {\,},
    output-decimal-marker = {\bookdecimalmarker}
  }%
  \newcommand{\booknum}[1]{\num{##1}}%
}{%
  \newcommand{\booknum}[1]{##1}%
}

\newcommand{\pyvalues}{}
\makeatletter
\newcommand{\pyval}[2]{\csgdef{pyval@#1}{#2}\listxadd{\pyvalues}{#1}}
% \num on a non-numeric body is a FATAL siunitx error. The generator knows
% whether it emitted a number or a word, so it writes \pyval for a number and
% \pyvaltext for a word, and this end only enforces which macro reads which.
\newcommand{\val}[1]{%
  \ifcsdef{pyval@#1}{%
    \ifcsdef{pyvaltext@#1}{%
      \PackageError{pybook}{Value `#1' is not a number}%
        {\string\val\space passes its body to siunitx, which refuses anything
         that is not a number. Use \string\valtext{#1} instead.}%
    }{}%
    \booknum{\csuse{pyval@#1}}%
  }{\textcolor{mred}{\textbf{[\lblValueMissing: #1]}}}%
}
\newcommand{\valtext}[1]{%
  \ifcsdef{pyval@#1}{\csuse{pyval@#1}}%
    {\textcolor{mred}{\textbf{[\lblValueMissing: #1]}}}%
}
\newcommand{\pyvaltext}[2]{%
  \csgdef{pyval@#1}{#2}\csgdef{pyvaltext@#1}{}\listxadd{\pyvalues}{#1}%
}
\makeatother

% figures/values/all.tex is generated by `make numbers` and does nothing but
% \input each value file. A build with no values yet still compiles; every
% \val{} in it prints a visible marker instead of a number.
\IfFileExists{figures/values/all.tex}{\input{figures/values/all}}{%
  \AtBeginDocument{\PackageWarning{pybook}{No computed values found. Run `make numbers`.}}}

% ============================================================
%  PINNED VERSIONS --- single source of truth
%  Change these here; they propagate through both editions and Appendix C.
%  Verified against pypi.org on \pinnedon (lang/<lang>.tex).
%  tools/check_versions.py compares every one of them with PyPI on each build.
%  Python itself is not on PyPI and is checked by hand against python.org.
% ============================================================
\newcommand{\bookedition}{0.1-dev}
\newcommand{\bookyear}{2026}
\newcommand{\pyver}{3.14}              % the minor the book targets
\newcommand{\pypatch}{3.14.7}          % the exact interpreter the code ran on
\newcommand{\uvver}{0.12.13}           % uv
\newcommand{\dotnetver}{10.0.100}      % dotnet
\newcommand{\ruffver}{0.16.7}          % ruff
\newcommand{\pyrightver}{1.1.414}      % pyright
\newcommand{\mypyver}{2.3.1}           % mypy
\newcommand{\pydanticver}{2.13.5}      % pydantic
\newcommand{\pysettingsver}{2.15.0}    % pydantic-settings
\newcommand{\fastapiver}{0.141.1}      % fastapi
\newcommand{\uvicornver}{0.53.0}       % uvicorn
\newcommand{\sqlalchemyver}{2.0.52}    % sqlalchemy
\newcommand{\alembicver}{1.20.0}       % alembic
\newcommand{\aiosqlitever}{0.22.1}     % aiosqlite
\newcommand{\pytestver}{9.1.1}         % pytest
\newcommand{\pytestasyncver}{1.4.0}    % pytest-asyncio
\newcommand{\hypothesisver}{6.168.0}   % hypothesis
\newcommand{\coveragever}{7.16.1}      % coverage
\newcommand{\httpxver}{0.28.1}         % httpx
\newcommand{\structlogver}{26.1.0}     % structlog
\newcommand{\otelver}{1.44.0}          % opentelemetry-sdk
\newcommand{\polarsver}{1.44.2}        % polars
\newcommand{\anthropicver}{1.5.0}      % anthropic
\newcommand{\openaiver}{3.13.0}        % openai
\newcommand{\pipauditver}{2.10.1}      % pip-audit

% ============================================================
%  INDEX
%  Two-column, and its entries include \texttt{} names that do not hyphenate.
%  Ragged right keeps multi-word entries off the margin; it cannot help a
%  single token wider than the column, so keep verbatim index entries under
%  about 29 characters and index longer names by concept instead.
%  \raggedcolumns is the switch imakeidx's multicols actually reads.
% ============================================================
\apptocmd{\theindex}{\raggedcolumns\raggedright}{}{%
  \PackageWarning{pybook}{Could not patch theindex for ragged setting}}
\providecommand*\see[2]{}
\renewcommand*\see[2]{\emph{\lblIndexSee} #1}
\providecommand*\seealso[2]{}
\renewcommand*\seealso[2]{\emph{\lblIndexSeeAlso} #1}

\makeindex
 still compiles and says what it did.
\providecommand{\booklang}{en}

\usepackage{etoolbox}   % loaded first: the language switch below needs it

% ---------- Page geometry ----------
% ONE format: A4, 12pt, ONE-SIDED. This book is read on a screen with an
% editor open beside it, and never printed as a bound book, so there is no
% recto, no verso, no gutter, no mirrored margin and no blank leaf before a
% chapter -- every one of those is a property of paper that a PDF viewer at
% half a screen's width turns into wasted pixels.
%
% The margins are symmetric and the measure is about seventy-five characters
% of prose, which is where the eye keeps the line return. The number that was
% actually chosen is the MONOSPACE measure: at \footnotesize on this block a
% listing line of 79 characters fits without wrapping, and 79 is the width
% every listing under code/ is held to by tools/check_structure.py --lines.
% A wrapped listing line is not an error TeX reports -- breaklines=true wraps
% it silently and prints an arrow into the middle of what the reader is meant
% to paste into an editor -- so the width is enforced in the source instead.
\usepackage[
  a4paper,
  left=3.1cm, right=3.1cm, top=2.6cm, bottom=3.0cm,
  head=26pt, headsep=0.8cm, footskip=1.4cm
]{geometry}

% ---------- Encoding & fonts ----------
\usepackage[T1]{fontenc}
\usepackage[utf8]{inputenc}
\IfFileExists{lmodern.sty}{\usepackage{lmodern}}{}
% newtxtext takes its TS1 symbols (\textcopyright, \textdegree) from TeX
% Gyre Termes, which Debian packages SEPARATELY (tex-gyre): a machine with
% newtx and without it dies on the copyright page with
% `Font TS1/ntxtlf/m/n/10.95=ts1-qtmr ... not loadable`. There is no
% inert-safe probe for a font metric in pdfTeX -- \IfFileExists searches
% TEXINPUTS and not TFMFONTS, so it reports the metric missing on a machine
% that has it, and \suppressfontnotfounderror is LuaTeX-only -- so nothing
% here guesses. CI's full TeX Live is the reference; locally, install
% tex-gyre. Both facts are recorded in CLAUDE.md under Build traps.
\IfFileExists{newtxtext.sty}{\usepackage{newtxtext}}{}
% amsmath before newtxmath, because newtxmath patches amsmath's internals.
% amssymb only when newtxmath is absent: newtxmath supplies the AMS symbols
% itself and loading both is a fatal clash (\Bbbk already defined) that is
% invisible on a machine without newtx and fatal on a full TeX Live.
\usepackage{amsmath}
\IfFileExists{newtxmath.sty}{\usepackage{newtxmath}}{\usepackage{amssymb}}
\IfFileExists{inconsolata.sty}{\usepackage[varqu,scaled=0.95]{inconsolata}}{}
\IfFileExists{lmodern.sty}{\usepackage{microtype}}{\usepackage[expansion=false]{microtype}}

% upquote is not a refinement. In T1 the input character ' is quoteright, so
% a listings body sets f('x') with U+2019 at both ends of the literal, and
% typed back in that is a SyntaxError. inconsolata's varqu option hides the
% defect on a machine that has inconsolata; upquote makes ' the ASCII
% apostrophe whatever the typewriter font is. In a book whose every listing is
% meant to be pasted into an editor, this line is load-bearing.
\IfFileExists{upquote.sty}{\usepackage{upquote}}{}

% And upquote's twin, for the same reason one character over. In T1 the pair
% -- is a ligature for an en dash, in EVERY family including the typewriter
% one, so \code{uv sync --locked} sets as "uv sync <en dash>locked" and a
% reader who copies it off the page gets an unknown-argument error from a
% command the chapter told them to run. It is invisible in review because it
% looks like a dash, and it bites exactly the flags a packaging chapter is
% made of. A `listings` body is already safe -- listings sets characters
% singly and no ligature forms -- which is what makes the defect so easy to
% miss: the same text is right inside a listing and wrong in the sentence
% above it.
%
% Scoped to tt* on purpose, and measured rather than assumed: prose keeps its
% own dashes, so \dash{} still sets an em dash and a range like 1--2 still
% sets an en dash. Probed against this preamble in all three positions before
% it was believed.
\DisableLigatures[-]{encoding = T1, family = tt*}

% ---------- Language ----------
% Loading babel with a language whose .ldf is absent is a *fatal* error, not
% a warning, so both the package and each language file are probed before
% either is requested. Both editions load the other language too, because the
% glossary rows and the cheat sheet carry words from both.
\IfFileExists{babel.sty}{%
  \IfFileExists{polish.ldf}{%
    \IfFileExists{british.ldf}{%
      \ifdefstring{\booklang}{pl}%
        {\usepackage[british,main=polish]{babel}}%
        {\usepackage[polish,main=british]{babel}}%
    }{\usepackage[polish]{babel}}%
  }{%
    \IfFileExists{british.ldf}{\usepackage[british]{babel}}{}%
  }%
}{}

% babel-polish makes " an active shorthand. listings reads its body verbatim
% and an active character in a verbatim context is a category-code accident
% waiting to happen, so the shorthand is turned off. Polish quotation marks
% come from \enquote, which csquotes sets per language.
\ifdefstring{\booklang}{pl}{\AtBeginDocument{\shorthandoff{"}}}{}

% ---------- Core packages ----------
\usepackage{graphicx}
\usepackage{xcolor}
\usepackage{booktabs}
\usepackage{tabularx}
% Appendix A's cheat sheet runs past a page in every part, and tabularx
% cannot break. longtable can, and works with booktabs.
\usepackage{longtable}
\usepackage{array}
\usepackage{enumitem}
\usepackage{listings}
\usepackage[most]{tcolorbox}
\usepackage{fancyhdr}
\usepackage{titlesec}
\usepackage{caption}
\usepackage{imakeidx}
% csquotes with autostyle follows babel's active language, so the identical
% source \enquote{x} sets 'x' in the English edition and the Polish quotation
% marks in the Polish one.
\IfFileExists{csquotes.sty}{\usepackage[autostyle=true]{csquotes}}{%
  \providecommand{\enquote}[1]{``##1''}}
\usepackage[hidelinks]{hyperref}

% ---------- Palette ----------
% The same family as the two companion volumes, so the three read as a set.
\definecolor{mblue}{HTML}{1F4E79}
\definecolor{mteal}{HTML}{0E7C7B}
\definecolor{mamber}{HTML}{B26A00}
\definecolor{mred}{HTML}{9B2C2C}
\definecolor{mviolet}{HTML}{7B4397}
\definecolor{mgrey}{HTML}{5A5A5A}
\definecolor{mlight}{HTML}{F4F6F8}
\definecolor{mfaint}{HTML}{FBFCFD}
\definecolor{codebg}{HTML}{FAFAFA}
\definecolor{codeframe}{HTML}{D8DEE4}
\definecolor{kw}{HTML}{0B5394}
\definecolor{str}{HTML}{9B2C2C}
\definecolor{cmt}{HTML}{4F7A4F}
\definecolor{typ}{HTML}{0E7C7B}
\definecolor{dec}{HTML}{7B4397}

\hypersetup{
  colorlinks=true,
  linkcolor=mblue,
  urlcolor=mteal,
  citecolor=mblue,
  pdfauthor={Konrad Cinkusz}
}

% \ifpl{Polish text}{English text} -- for the handful of places where a
% string is structural rather than content. Anything longer than a few words
% belongs in lang/ or in the edition's own source, not here.
% \DeclareRobustCommand, not \newcommand: a fragile \ifpl inside a moving
% argument is written to the .toc one level expanded and fails on the next run
% in exactly one of the two editions.
\DeclareRobustCommand{\ifpl}[2]{\ifdefstring{\booklang}{pl}{#1}{#2}}
\makeatletter
\pdfstringdefDisableCommands{%
  \ifdefstring{\booklang}{pl}{\let\ifpl\@firstoftwo}{\let\ifpl\@secondoftwo}%
  % \api{}, \pkg{} and \code{} are allowed in a section title, and a title
  % becomes a PDF bookmark. hyperref drops the \color command from a bookmark
  % string but keeps its ARGUMENT, so a title carrying \api{model\_validate}
  % made a bookmark reading "mtealmodel_validate" and warned "Token not
  % allowed" twice. The colour name is gobbled here instead. Measured on a
  % probe against this preamble before the line was added.
  \def\color#1{}%
}
\makeatother

% ---------- Language strings ----------
\input{lang/\booklang}

% The PDF's subject, keyword list and document language read \lblPdf... from
% the language file input on the line above, so they cannot sit in the colour
% \hypersetup block, which runs before it. Two blocks is the ordering, not an
% oversight; merging them raises Undefined control sequence and still writes
% a PDF, with the three fields empty.
\hypersetup{
  pdfsubject={\lblPdfSubject},
  pdfkeywords={\lblPdfKeywords},
  pdflang={\lblPdfLang}
}

% ---------- Headings ----------
% \raggedright on the title body is load-bearing: at \Huge on this block a
% title that cannot break overflows the margin by 30 pt or more, and several
% chapter titles here run past fifty characters.
\titleformat{\chapter}[display]
  {\normalfont\huge\bfseries\color{mblue}\raggedright}
  {\normalfont\normalsize\color{mgrey}\MakeUppercase{\chaptertitlename}\ \thechapter}
  {8pt}{\Huge\raggedright}
\titlespacing*{\chapter}{0pt}{10pt}{28pt}
\titleformat{\section}{\normalfont\Large\bfseries\color{mblue}\raggedright}{\thesection}{0.8em}{}
\titleformat{\subsection}{\normalfont\large\bfseries\color{mgrey}\raggedright}{\thesubsection}{0.7em}{}

% Sections in the contents, subsections not: fourteen chapters of subsections
% is a contents nobody reads, and the reader navigates by chapter and section.
\setcounter{tocdepth}{1}
\setcounter{secnumdepth}{2}

% One-sided running head: the chapter on the left, the folio on the right.
\pagestyle{fancy}
\fancyhf{}
\fancyhead[L]{\small\nouppercase{\leftmark}}
\fancyhead[R]{\small\thepage}
\renewcommand{\headrulewidth}{0.4pt}
\fancypagestyle{plain}{\fancyhf{}\fancyfoot[R]{\small\thepage}%
  \renewcommand{\headrulewidth}{0pt}}

% Drops the word "Chapter" from the head: the number is still there and the
% longer titles no longer overflow. MUST come after \pagestyle{fancy}, which
% installs its own \chaptermark and silently overwrites anything earlier.
\renewcommand{\chaptermark}[1]{\markboth{\thechapter.\ #1}{}}

% ============================================================
%  CODE LISTINGS
%  ---------------------------------------------------------
%  Every listing in a chapter is a FILE under code/, pulled in with \pyfile or
%  \pyregion, and CI runs it against the pinned versions. An in-source
%  \begin{python} block is allowed for a fragment of three or four lines that
%  exists to show a shape rather than to run; anything longer belongs in a
%  file the reader can open in the editor beside the PDF.
% ============================================================
\lstdefinelanguage{pythonx}{
  language=Python,
  morekeywords={
    async,await,match,case,nonlocal,yield,type,
    TypedDict,Annotated,Protocol,Literal,Generic,Final,ClassVar,Self,
    TypeVar,TypeGuard,TypeIs,Callable,Iterator,Iterable,Awaitable,
    dataclass,field,BaseModel,Field,Enum,StrEnum,NamedTuple,
    override,overload,property,staticmethod,classmethod
  },
  morekeywords=[2]{
    gather,create_task,TaskGroup,to_thread,timeout,wait_for,Semaphore,Queue,
    fixture,parametrize,monkeypatch,raises,mark,
    Depends,FastAPI,APIRouter,HTTPException,
    select,Session,Mapped,mapped_column,relationship,
    model_validate,model_validate_json,model_dump,model_json_schema
  },
  sensitive=true
}

% listings' C# definition predates async, records and pattern matching.
\lstdefinelanguage{csharpx}{
  language=[Sharp]C,
  morekeywords={
    async,await,var,dynamic,record,init,required,nameof,yield,partial,
    global,scoped,when,with,and,or,not,notnull,unmanaged,file,get,set,value,
    where,select,from,orderby,let,join,into,ascending,descending,on,equals,
    Task,ValueTask,CancellationToken,IEnumerable,IAsyncEnumerable,Func,Action
  },
  sensitive=true
}

\lstdefinestyle{bookcode}{
  basicstyle=\ttfamily\footnotesize,
  keywordstyle=\color{kw}\bfseries,
  keywordstyle=[2]\color{typ},
  stringstyle=\color{str},
  % Colour only, no \itshape: inconsolata (zi4) has no italic shape, so an
  % italic comment style is substituted with upright on every machine that
  % has the font, and the log carries "Font shape `T1/zi4/m/it' undefined"
  % on every build. The colour already marks a comment.
  commentstyle=\color{cmt},
  emphstyle=\color{dec}\bfseries,
  numbers=left,
  numberstyle=\ttfamily\tiny\color{mgrey},
  numbersep=8pt,
  backgroundcolor=\color{codebg},
  frame=single,
  rulecolor=\color{codeframe},
  framesep=6pt,
  breaklines=true,
  breakatwhitespace=false,
  postbreak=\mbox{\textcolor{mgrey}{$\hookrightarrow$}\space},
  showstringspaces=false,
  tabsize=4,
  columns=fullflexible,
  keepspaces=true,
  captionpos=b,
  aboveskip=1.2em,
  belowskip=1.2em,
  % Maps the characters most likely to be pasted into a listing by accident.
  %
  % Do NOT add a mapping for U+00A0 (non-breaking space). It looks like the
  % obvious next entry and it makes `listings` abort the run with "Improper
  % alphabetic constant" -- fatal, no PDF, and the message names neither the
  % character nor this line. The ASCII-in-listings convention, checked by
  % parity's C13, is the guard against it instead.
  literate={ł}{{\l}}1 {Ł}{{\L}}1 {ą}{{\k a}}1 {ę}{{\k e}}1
           {ś}{{\'s}}1 {ć}{{\'c}}1 {ż}{{\.z}}1 {ź}{{\'z}}1 {ó}{{\'o}}1 {ń}{{\'n}}1
           {Ą}{{\k A}}1 {Ę}{{\k E}}1 {Ś}{{\'S}}1 {Ć}{{\'C}}1
           {Ż}{{\.Z}}1 {Ź}{{\'Z}}1 {Ó}{{\'O}}1 {Ń}{{\'N}}1
           {—}{{-{}-}}2 {–}{{-}}1 {‘}{{'}}1 {’}{{'}}1
           {“}{{"}}1 {”}{{"}}1 {°}{{\textdegree}}1
}

\lstset{style=bookcode}

\lstnewenvironment{python}[1][]
  {\lstset{language=pythonx,style=bookcode,#1}}
  {}

% C# comparison snippets. The whole book is written for engineers arriving
% from .NET, so listings come in pairs often enough that the C# half has its
% own environment rather than a generic one.
\lstnewenvironment{csharp}[1][]
  {\lstset{language=csharpx,style=bookcode,#1}}
  {}

\lstnewenvironment{shellcmd}[1][]
  {\lstset{language=bash,style=bookcode,numbers=none,keywordstyle=\color{mteal}\bfseries,#1}}
  {}

\lstnewenvironment{tomlcode}[1][]
  {\lstset{style=bookcode,numbers=none,keywordstyle=\color{black},#1}}
  {}

\lstnewenvironment{yamlcode}[1][]
  {\lstset{style=bookcode,numbers=none,keywordstyle=\color{black},#1}}
  {}

\lstnewenvironment{jsoncode}[1][]
  {\lstset{style=bookcode,numbers=none,keywordstyle=\color{black},#1}}
  {}

% Terminal transcripts typed into the source. Allowed for a shell session the
% reader is meant to reproduce; NOT for program output, which is \transcript
% below, because an in-source transcript nobody ran reads exactly like one
% that was, and that is where a remembered number hides.
\lstnewenvironment{console}[1][]
  {\lstset{style=bookcode,numbers=none,basicstyle=\ttfamily\footnotesize,
           keywordstyle=\color{black},backgroundcolor=\color{white},#1}}
  {}

% The path is printed above every file listing, in the repository-relative
% form the reader opens in the editor beside the PDF. That is the whole point
% of a listing being a file, and the front matter promises it. The \nobreak
% keeps the path on the same page as the listing it names.
% ---------------------------------------------------------------------------
%  Appendix B prints the trap catalogue from notes/02-traps.md. One entry is
%  its number -- which a chapter may cite, and which is never reused -- the
%  habit in the reader's own voice, and what Python does instead.
%
%  The body is set RAGGED RIGHT on purpose. An entry is mostly \code{} spans,
%  which are unbreakable runs, and sixty-seven justified paragraphs of them
%  is sixty-seven chances at an overfull hbox in the edition with the longer
%  words. Ragged right removes the class rather than fixing it one box at a
%  time, and a reference list is the one place in this book where it reads
%  better anyway.
%
%  The habit is set with \enquote{} rather than \emph{}, and that is not a
%  style choice. Every habit here carries \code{} spans, and \code{} inside
%  \emph{} asks inconsolata for an italic it does not have: the first build
%  of this appendix raised `Font shape T1/zi4/m/it undefined', which
%  checklog.py is hard on for exactly this reason. It is the trap CLAUDE.md
%  records, met sixty-seven times at once -- and fixed once, because the
%  entries go through a macro. A quotation is also what the habit IS, and
%  \enquote{} is what the Polish edition owes anyway.
% ---------------------------------------------------------------------------
%  Appendix A's cheat sheet: one table per part, C# on the left, Python on
%  the right, and the chapter that teaches the row on the end. The three
%  header strings are arguments so that one environment serves both
%  editions and the structural signature stays identical.
%
%  Columns are ragged right for the reason Appendix B's entries are: a cell
%  is mostly \code{} spans, which are unbreakable runs, and a justified
%  narrow column full of them is an overfull box waiting for the edition
%  with the longer words.
% ---------------------------------------------------------------------------
%  Appendix C's tool matrix: the .NET tool, the Python tool this book uses,
%  the version it was verified against and the chapter that introduces it.
%
%  The version column prints \pydanticver and its siblings and is never
%  typed, which is the whole point of the appendix: it is the one page where
%  a reader sees every pin at once, and a pin that had to be re-typed here
%  would be a second copy of a fact the preamble already owns.
%  tools/check_structure.py --tools fails when a pin macro exists and this
%  table does not print it.
\NewDocumentEnvironment{toolmatrix}{mmmm}{%
  \small
  \begin{longtable}{@{}%
    >{\raggedright\arraybackslash}p{0.27\linewidth}%
    >{\raggedright\arraybackslash}p{0.31\linewidth}%
    >{\raggedright\arraybackslash}p{0.20\linewidth}%
    >{\raggedright\arraybackslash}p{0.08\linewidth}@{}}
  \toprule
  \textbf{#1} & \textbf{#2} & \textbf{#3} & \textbf{#4} \\
  \midrule
  \endfirsthead
  \toprule
  \textbf{#1} & \textbf{#2} & \textbf{#3} & \textbf{#4} \\
  \midrule
  \endhead
  \midrule
  \endfoot
  \bottomrule
  \endlastfoot
}{%
  \end{longtable}%
}

\NewDocumentEnvironment{cheatsheet}{mmm}{%
  \small
  \begin{longtable}{@{}%
    >{\raggedright\arraybackslash}p{0.41\linewidth}%
    >{\raggedright\arraybackslash}p{0.41\linewidth}%
    >{\raggedright\arraybackslash}p{0.08\linewidth}@{}}
  \toprule
  \textbf{#1} & \textbf{#2} & \textbf{#3} \\
  \midrule
  \endfirsthead
  \toprule
  \textbf{#1} & \textbf{#2} & \textbf{#3} \\
  \midrule
  \endhead
  \midrule
  \endfoot
  \bottomrule
  \endlastfoot
}{%
  \end{longtable}%
}

\newcommand{\trapentry}[3]{%
  \par\medskip
  \noindent\textbf{#1.}\enspace\enquote{#2}\par
  \nopagebreak
  {\raggedright\noindent#3\par}%
}

\newcommand{\listingpath}[1]{%
  \par\medskip\noindent
  {\ttfamily\footnotesize\color{mgrey}\detokenize{#1}}%
  \par\nobreak}

% External source file -- the preferred form. The path is relative to the
% repository root, which is also the path the reader opens in the editor.
% Usage: \pyfile{code/ch05/decorators.py}{Caption}{lst:label}
\newcommand{\pyfile}[3]{%
  \listingpath{#1}%
  \lstinputlisting[language=pythonx,style=bookcode,aboveskip=0.3em,
    caption={#2},label={#3}]{#1}%
}

% A named region of an external file, delimited in the source by
% `# --8<-- [start:name]` / `# --8<-- [end:name]`, so the book and the running
% code cannot drift apart. tools/check_structure.py --listings fails when the
% file or the region is missing.
%
% The markers are matched through listings' range prefix and suffix keys,
% with linerange={name-name}. The form this was inherited with -- a linerange
% whose two ends were the whole marker strings -- matched nothing and printed
% the WHOLE FILE, with no warning: measured on a three-region probe file
% before this was rewritten, and the same definition sits unused in the
% sibling book it came from. A region name is a word, never a number, because
% a number is a line number to listings. The line numbers printed are the
% file's own, so a region's listing says where in the file it sits.
% Usage: \pyregion{code/ch05/decorators.py}{timing}{Caption}{lst:label}
\lstset{
  rangebeginprefix=\#\ --8<--\ [start:, rangebeginsuffix=],
  rangeendprefix=\#\ --8<--\ [end:, rangeendsuffix=],
  includerangemarker=false}
\newcommand{\pyregion}[4]{%
  \listingpath{#1}%
  \lstinputlisting[language=pythonx,style=bookcode,aboveskip=0.3em,
    linerange={#2-#2},caption={#3},label={#4}]{#1}%
}

% The C# twin of \pyfile, for the comparison half of a pair.
\newcommand{\csfile}[3]{%
  \listingpath{#1}%
  \lstinputlisting[language=csharpx,style=bookcode,aboveskip=0.3em,
    caption={#2},label={#3}]{#1}%
}

% A program's output, written by a script under code/measure/ and pulled in
% whole. There is no typed form of this and there must not be: \val{} does
% not expand inside a verbatim body, so a transcript typed into a chapter is
% one that cannot quote a computed value, and a transcript nobody ran is
% indistinguishable on the page from one that was.
\newcommand{\transcript}[1]{%
  \par\smallskip
  \IfFileExists{figures/transcripts/#1.txt}{%
    \lstinputlisting[style=bookcode,numbers=none,basicstyle=\ttfamily\footnotesize,
      keywordstyle=\color{black},backgroundcolor=\color{white},
      language={}]{figures/transcripts/#1.txt}%
  }{%
    \begin{tcolorbox}[enhanced, sharp corners, width=\linewidth,
      colback=mlight, colframe=mgrey, boxrule=0.6pt,
      left=8pt, right=8pt, top=6pt, bottom=6pt]
      {\bfseries\color{mgrey}\small \lblTranscriptMissing}\\[2pt]
      {\ttfamily\scriptsize\color{mgrey} figures/transcripts/#1.txt}
    \end{tcolorbox}%
  }%
  \par\smallskip
}

% Inline literals. Underscores inside these must be written \_.
\newcommand{\code}[1]{\texttt{\small #1}}
\newcommand{\pkg}[1]{\texttt{\small\color{mblue}#1}}
\newcommand{\api}[1]{\texttt{\small\color{mteal}#1}}

% ============================================================
%  ADMONITIONS
%  ---------------------------------------------------------
%  A deliberately small, fixed vocabulary. Two visual treatments carry the
%  whole grammar: `tinted` is a block the reader may read in passing -- a
%  tint and a thick left bar, no outline; `outlined` is a block the reader
%  must stop at, with a full rule on white. A third treatment would mean the
%  first two are not carrying their meaning.
% ============================================================
\tcbset{
  mfa tinted/.style={
    enhanced, breakable, sharp corners,
    colback=mlight, boxrule=0pt, leftrule=3pt,
    left=8pt, right=8pt, top=6pt, bottom=6pt,
    before skip=13pt, after skip=13pt,
    fonttitle=\bfseries\small,
    coltitle=black, colbacktitle=mlight,
    attach title to upper={\par\smallskip},
  },
  mfa outlined/.style={
    enhanced, breakable, sharp corners,
    colback=white, boxrule=0.6pt,
    left=8pt, right=8pt, top=6pt, bottom=6pt,
    before skip=13pt, after skip=13pt,
    fonttitle=\bfseries\small,
    coltitle=black, colbacktitle=white,
    attach title to upper={\par\smallskip},
  },
}

% Read in passing.
\newtcolorbox{note}[1][]{mfa tinted, colframe=mblue, title={\lblNote}, #1}
\newtcolorbox{warning}[1][]{mfa tinted, colframe=mred, title={\lblWarning}, #1}
% The translation. The .NET reader is the book's only assumed reader, so this
% box carries the C# side of a comparison and nothing else. Its use is
% counted: a chapter with none is a chapter that has stopped translating.
\newtcolorbox{csbox}[1][]{mfa tinted, colframe=mteal, title={\lblCsharp}, #1}
% Anything that is preview, experimental or likely to move. Python moves more
% slowly than the agent frameworks do, and this box should be rare here.
\newtcolorbox{versionbox}[1][]{mfa tinted, colframe=mamber, title={\lblVersion}, #1}

% Stop and read.
% The misconception, named AFTER the reader has walked into it. Appendix B is
% the catalogue; a trap box in a chapter is one entry of it, in place.
\newtcolorbox{trapbox}[1][]{mfa outlined, colframe=mred, title={\lblTrap}, #1}
% Marks a listing that has not been executed against the pinned versions.
% `make debt` counts these, CI prints the count, and nothing ships to a reader
% with one attached. Removing it means you ran the code.
\newtcolorbox{verifybox}[1][]{mfa outlined, colframe=mgrey, title={\lblVerify}, #1}
\newtcolorbox{exercisebox}[1][]{mfa outlined, colframe=mteal, title={\lblExercise}, #1}
% The guiding project. Points at the stage of trace-assert this section
% builds, under code/src/trace_assert/.
\newtcolorbox{projectbox}[1][]{mfa outlined, colframe=mamber, title={\lblProject}, #1}

% ============================================================
%  EXERCISES
%  ---------------------------------------------------------
%  \begin{exercise}{key}{title} ... \end{exercise}
%
%  This is the mechanism the whole book is read through: the PDF on one half
%  of the screen, an editor on the other. An exercise is therefore not a
%  paragraph but a FILE the reader opens -- a starter under
%  code/exercises/chNN/<key>.py -- and a TEST that decides whether it is done,
%  test_<key>.py beside it. The box prints both paths and the one command
%  that runs the test, and the manifest at the back lists every exercise.
%
%  An ENVIRONMENT, not a macro with a body argument: a macro argument cannot
%  contain verbatim material, so an exercise written as \exercise{key}{title}
%  {body} could never show the reader a code fragment, and the first chapter
%  that tried would have died on `\begin{python}` inside the braces.
%
%  The key is the file stem, with underscores, and it is written with
%  \detokenize so the underscore reaches the page as a character rather than
%  as a subscript. tools/check_structure.py --exercises fails when the starter
%  or the test is missing, or when the key's chapter prefix disagrees with the
%  chapter it sits in.
% ============================================================
\newcounter{exercise}[chapter]
\renewcommand{\theexercise}{\thechapter.\arabic{exercise}}
\makeatletter
\newcommand{\exercisedir}{ch\two@digits{\value{chapter}}}
\newcommand{\listofexercises}{%
  \section*{\lblExerciseManifest}%
  \phantomsection\addcontentsline{toc}{section}{\lblExerciseManifest}%
  % The entries are subsection-level and tocdepth is 1 for the contents,
  % so the depth is raised for this list alone or it prints empty --
  % which it did, on the scaffold's first clean build.
  \begingroup\c@tocdepth=2\relax\@starttoc{exr}\endgroup%
}
\makeatother
\NewDocumentEnvironment{exercise}{mm}{%
  \refstepcounter{exercise}%
  \addcontentsline{exr}{subsection}{\protect\numberline{\theexercise}\texttt{\protect\detokenize{#1}} --- #2}%
  \begin{exercisebox}[title={\lblExercise~\theexercise\ \dash{} #2}]%
}{%
  \par\medskip
  {\small\setlength{\parindent}{0pt}%
  \lblExerciseStarter: \texttt{code/exercises/\exercisedir/\detokenize{#1}.py}\\
  \lblExerciseTest: \texttt{code/exercises/\exercisedir/test\_\detokenize{#1}.py}\\
  \lblExerciseRun: \texttt{cd code \&\& uv run pytest -k \detokenize{#1}}\par}%
  \end{exercisebox}%
}

% ============================================================
%  DIAGRAMS --- Mermaid pipeline
%  ---------------------------------------------------------
%  \mermaidfig{key}{caption}{one-line manifest description}
%
%  Source of truth is figures/mermaid/<lang>/<key>.mmd, committed. `make
%  diagrams` renders each to figures/diagrams/<lang>/<key>.pdf, which is build
%  output and is not committed, so a diagram change reviews as a text diff.
%  The source is per-language because a diagram with English labels in the
%  Polish edition is exactly the half-translation that makes a bilingual book
%  feel machine-made; CI checks that both languages carry the same keys.
%
%  If the rendered PDF is absent the macro draws a placeholder AND typesets
%  the Mermaid source, in the flow rather than as a float, because a source
%  typeset verbatim is often taller than a page and a float cannot break.
% ============================================================
\makeatletter
\newcommand{\listofdiagrams}{%
  \section*{\lblDiagramManifest}%
  \phantomsection\addcontentsline{toc}{section}{\lblDiagramManifest}%
  % The entries are subsection-level and tocdepth is 1 for the contents,
  % so the depth is raised for this list alone or it prints empty --
  % which it did, on the scaffold's first clean build.
  \begingroup\c@tocdepth=2\relax\@starttoc{dgm}\endgroup%
}
\makeatother

\newcommand{\mermaidfig}[3]{%
  % \phantomsection is load-bearing and not decoration. \addcontentsline
  % records hyperref's CURRENT anchor, and here that is whatever preceded
  % the figure -- for a figure placed after a listing it is the listing's
  % line-number anchor, and listings has already advanced the counter past
  % the last line it printed, so the manifest entry linked to a
  % destination that is never created: "name{lstnumber.10.4.46} has been
  % referenced but does not exist, replaced by a fixed one". The two
  % front-matter figures had the same wrong target and were silent about
  % it, because the lines they named happen to be printed. The exercise
  % manifest never had the bug, because \begin{exercise} steps a counter
  % and that is what sets the anchor.
  \phantomsection%
  \addcontentsline{dgm}{subsection}{\protect\numberline{}\texttt{#1.mmd} --- #3}%
  \IfFileExists{figures/diagrams/\booklang/#1.pdf}{%
    \begin{figure}[htbp]
    \centering
    \includegraphics[width=0.95\linewidth,height=0.42\textheight,keepaspectratio]{figures/diagrams/\booklang/#1.pdf}%
    \caption{#2}\label{fig:#1}
    \end{figure}%
  }{%
    \par\medskip
    \begin{tcolorbox}[
      enhanced, breakable, width=\linewidth, sharp corners,
      colback=white, colframe=mgrey, boxrule=0.8pt,
      left=8pt, right=8pt, top=8pt, bottom=8pt]
      {\bfseries\color{mgrey}\small \lblNotRendered}\\[4pt]
      {\ttfamily\scriptsize\color{mgrey} figures/mermaid/\booklang/#1.mmd}\\[6pt]
      \IfFileExists{figures/mermaid/\booklang/#1.mmd}{%
        \lstinputlisting[style=bookcode,numbers=none,basicstyle=\ttfamily\scriptsize,
                         backgroundcolor=\color{mlight},frame=none,
                         keywordstyle=\color{black},language={}]{figures/mermaid/\booklang/#1.mmd}%
      }{%
        {\small\color{mred} \lblSourceMissing: \texttt{figures/mermaid/\booklang/#1.mmd}}%
      }%
    \end{tcolorbox}%
    \captionof{figure}{#2}\label{fig:#1}
    \par\medskip
  }%
}

% ============================================================
%  CHAPTER STUBS
%  ---------------------------------------------------------
%  Prints a visible marker so an unwritten chapter cannot be mistaken for a
%  written one and the page count never flatters the draft. `make debt`
%  counts them, and CI publishes the count on every build.
% ============================================================
\newcommand{\chapterstub}[1]{%
  \begin{tcolorbox}[
    enhanced, breakable, width=\linewidth, sharp corners,
    colback=mlight, colframe=mred, boxrule=1pt,
    borderline={0.6pt}{3pt}{mred, dashed},
    left=12pt, right=12pt, top=12pt, bottom=12pt]
    {\bfseries\color{mred}\large \lblNotWritten}\\[8pt]
    \small\raggedright #1
  \end{tcolorbox}%
}

% ============================================================
%  COMPUTED VALUES
%  ---------------------------------------------------------
%  A number that came out of a measurement is written by a script under
%  code/measure/ into figures/values/<name>.tex as
%
%      \pyval{e01.threads.speedup}{3.71}
%
%  and pulled into the text with \val{e01.threads.speedup}. The script is the
%  source of truth and the book cannot disagree with it, because the book does
%  not contain the digits. `make verify` fails when a committed value no
%  longer matches its script. A value the scripts do not produce prints a
%  visible marker.
%
%  \val also applies the edition's decimal separator, which is how the Polish
%  edition gets its decimal comma without a translator having to remember.
% ============================================================
\newcommand{\bookdecimalmarker}{\ifpl{,}{.}}
\IfFileExists{siunitx.sty}{%
  \usepackage{siunitx}%
  \sisetup{
    group-digits = integer,
    group-minimum-digits = 5,
    group-separator = {\,},
    output-decimal-marker = {\bookdecimalmarker}
  }%
  \newcommand{\booknum}[1]{\num{##1}}%
}{%
  \newcommand{\booknum}[1]{##1}%
}

\newcommand{\pyvalues}{}
\makeatletter
\newcommand{\pyval}[2]{\csgdef{pyval@#1}{#2}\listxadd{\pyvalues}{#1}}
% \num on a non-numeric body is a FATAL siunitx error. The generator knows
% whether it emitted a number or a word, so it writes \pyval for a number and
% \pyvaltext for a word, and this end only enforces which macro reads which.
\newcommand{\val}[1]{%
  \ifcsdef{pyval@#1}{%
    \ifcsdef{pyvaltext@#1}{%
      \PackageError{pybook}{Value `#1' is not a number}%
        {\string\val\space passes its body to siunitx, which refuses anything
         that is not a number. Use \string\valtext{#1} instead.}%
    }{}%
    \booknum{\csuse{pyval@#1}}%
  }{\textcolor{mred}{\textbf{[\lblValueMissing: #1]}}}%
}
\newcommand{\valtext}[1]{%
  \ifcsdef{pyval@#1}{\csuse{pyval@#1}}%
    {\textcolor{mred}{\textbf{[\lblValueMissing: #1]}}}%
}
\newcommand{\pyvaltext}[2]{%
  \csgdef{pyval@#1}{#2}\csgdef{pyvaltext@#1}{}\listxadd{\pyvalues}{#1}%
}
\makeatother

% figures/values/all.tex is generated by `make numbers` and does nothing but
% \input each value file. A build with no values yet still compiles; every
% \val{} in it prints a visible marker instead of a number.
\IfFileExists{figures/values/all.tex}{\input{figures/values/all}}{%
  \AtBeginDocument{\PackageWarning{pybook}{No computed values found. Run `make numbers`.}}}

% ============================================================
%  PINNED VERSIONS --- single source of truth
%  Change these here; they propagate through both editions and Appendix C.
%  Verified against pypi.org on \pinnedon (lang/<lang>.tex).
%  tools/check_versions.py compares every one of them with PyPI on each build.
%  Python itself is not on PyPI and is checked by hand against python.org.
% ============================================================
\newcommand{\bookedition}{0.1-dev}
\newcommand{\bookyear}{2026}
\newcommand{\pyver}{3.14}              % the minor the book targets
\newcommand{\pypatch}{3.14.7}          % the exact interpreter the code ran on
\newcommand{\uvver}{0.12.13}           % uv
\newcommand{\dotnetver}{10.0.100}      % dotnet
\newcommand{\ruffver}{0.16.7}          % ruff
\newcommand{\pyrightver}{1.1.414}      % pyright
\newcommand{\mypyver}{2.3.1}           % mypy
\newcommand{\pydanticver}{2.13.5}      % pydantic
\newcommand{\pysettingsver}{2.15.0}    % pydantic-settings
\newcommand{\fastapiver}{0.141.1}      % fastapi
\newcommand{\uvicornver}{0.53.0}       % uvicorn
\newcommand{\sqlalchemyver}{2.0.52}    % sqlalchemy
\newcommand{\alembicver}{1.20.0}       % alembic
\newcommand{\aiosqlitever}{0.22.1}     % aiosqlite
\newcommand{\pytestver}{9.1.1}         % pytest
\newcommand{\pytestasyncver}{1.4.0}    % pytest-asyncio
\newcommand{\hypothesisver}{6.168.0}   % hypothesis
\newcommand{\coveragever}{7.16.1}      % coverage
\newcommand{\httpxver}{0.28.1}         % httpx
\newcommand{\structlogver}{26.1.0}     % structlog
\newcommand{\otelver}{1.44.0}          % opentelemetry-sdk
\newcommand{\polarsver}{1.44.2}        % polars
\newcommand{\anthropicver}{1.5.0}      % anthropic
\newcommand{\openaiver}{3.13.0}        % openai
\newcommand{\pipauditver}{2.10.1}      % pip-audit

% ============================================================
%  INDEX
%  Two-column, and its entries include \texttt{} names that do not hyphenate.
%  Ragged right keeps multi-word entries off the margin; it cannot help a
%  single token wider than the column, so keep verbatim index entries under
%  about 29 characters and index longer names by concept instead.
%  \raggedcolumns is the switch imakeidx's multicols actually reads.
% ============================================================
\apptocmd{\theindex}{\raggedcolumns\raggedright}{}{%
  \PackageWarning{pybook}{Could not patch theindex for ragged setting}}
\providecommand*\see[2]{}
\renewcommand*\see[2]{\emph{\lblIndexSee} #1}
\providecommand*\seealso[2]{}
\renewcommand*\seealso[2]{\emph{\lblIndexSeeAlso} #1}

\makeindex
; nothing else differs between them at this level.
%  Every user-visible string lives in lang/en.tex or lang/pl.tex, so a label
%  that appears in one edition and not the other is a missing macro rather
%  than a silent divergence.
%
%  Lineage. The bilingual wiring (\booklang, lang/, body.tex, the generated
%  structure) is the math book's; the chapter machinery (listings, the
%  admonitions, verifybox, \pyfile and \pyregion, the Mermaid pipeline) is the
%  LangChain book's. What is new here is the exercise macro, which ties every
%  exercise to a starter file and a test under code/, because this book is
%  read with an editor open beside it.
% ============================================================

% \booklang must be set by the main file. Default to English rather than
% failing, so a stray % ============================================================
%  preamble.tex --- styling, environments, macros
%
%  Shared by BOTH editions. main-en.tex and main-pl.tex each define \booklang
%  before % ============================================================
%  preamble.tex --- styling, environments, macros
%
%  Shared by BOTH editions. main-en.tex and main-pl.tex each define \booklang
%  before \input{preamble}; nothing else differs between them at this level.
%  Every user-visible string lives in lang/en.tex or lang/pl.tex, so a label
%  that appears in one edition and not the other is a missing macro rather
%  than a silent divergence.
%
%  Lineage. The bilingual wiring (\booklang, lang/, body.tex, the generated
%  structure) is the math book's; the chapter machinery (listings, the
%  admonitions, verifybox, \pyfile and \pyregion, the Mermaid pipeline) is the
%  LangChain book's. What is new here is the exercise macro, which ties every
%  exercise to a starter file and a test under code/, because this book is
%  read with an editor open beside it.
% ============================================================

% \booklang must be set by the main file. Default to English rather than
% failing, so a stray \input{preamble} still compiles and says what it did.
\providecommand{\booklang}{en}

\usepackage{etoolbox}   % loaded first: the language switch below needs it

% ---------- Page geometry ----------
% ONE format: A4, 12pt, ONE-SIDED. This book is read on a screen with an
% editor open beside it, and never printed as a bound book, so there is no
% recto, no verso, no gutter, no mirrored margin and no blank leaf before a
% chapter -- every one of those is a property of paper that a PDF viewer at
% half a screen's width turns into wasted pixels.
%
% The margins are symmetric and the measure is about seventy-five characters
% of prose, which is where the eye keeps the line return. The number that was
% actually chosen is the MONOSPACE measure: at \footnotesize on this block a
% listing line of 79 characters fits without wrapping, and 79 is the width
% every listing under code/ is held to by tools/check_structure.py --lines.
% A wrapped listing line is not an error TeX reports -- breaklines=true wraps
% it silently and prints an arrow into the middle of what the reader is meant
% to paste into an editor -- so the width is enforced in the source instead.
\usepackage[
  a4paper,
  left=3.1cm, right=3.1cm, top=2.6cm, bottom=3.0cm,
  head=26pt, headsep=0.8cm, footskip=1.4cm
]{geometry}

% ---------- Encoding & fonts ----------
\usepackage[T1]{fontenc}
\usepackage[utf8]{inputenc}
\IfFileExists{lmodern.sty}{\usepackage{lmodern}}{}
% newtxtext takes its TS1 symbols (\textcopyright, \textdegree) from TeX
% Gyre Termes, which Debian packages SEPARATELY (tex-gyre): a machine with
% newtx and without it dies on the copyright page with
% `Font TS1/ntxtlf/m/n/10.95=ts1-qtmr ... not loadable`. There is no
% inert-safe probe for a font metric in pdfTeX -- \IfFileExists searches
% TEXINPUTS and not TFMFONTS, so it reports the metric missing on a machine
% that has it, and \suppressfontnotfounderror is LuaTeX-only -- so nothing
% here guesses. CI's full TeX Live is the reference; locally, install
% tex-gyre. Both facts are recorded in CLAUDE.md under Build traps.
\IfFileExists{newtxtext.sty}{\usepackage{newtxtext}}{}
% amsmath before newtxmath, because newtxmath patches amsmath's internals.
% amssymb only when newtxmath is absent: newtxmath supplies the AMS symbols
% itself and loading both is a fatal clash (\Bbbk already defined) that is
% invisible on a machine without newtx and fatal on a full TeX Live.
\usepackage{amsmath}
\IfFileExists{newtxmath.sty}{\usepackage{newtxmath}}{\usepackage{amssymb}}
\IfFileExists{inconsolata.sty}{\usepackage[varqu,scaled=0.95]{inconsolata}}{}
\IfFileExists{lmodern.sty}{\usepackage{microtype}}{\usepackage[expansion=false]{microtype}}

% upquote is not a refinement. In T1 the input character ' is quoteright, so
% a listings body sets f('x') with U+2019 at both ends of the literal, and
% typed back in that is a SyntaxError. inconsolata's varqu option hides the
% defect on a machine that has inconsolata; upquote makes ' the ASCII
% apostrophe whatever the typewriter font is. In a book whose every listing is
% meant to be pasted into an editor, this line is load-bearing.
\IfFileExists{upquote.sty}{\usepackage{upquote}}{}

% And upquote's twin, for the same reason one character over. In T1 the pair
% -- is a ligature for an en dash, in EVERY family including the typewriter
% one, so \code{uv sync --locked} sets as "uv sync <en dash>locked" and a
% reader who copies it off the page gets an unknown-argument error from a
% command the chapter told them to run. It is invisible in review because it
% looks like a dash, and it bites exactly the flags a packaging chapter is
% made of. A `listings` body is already safe -- listings sets characters
% singly and no ligature forms -- which is what makes the defect so easy to
% miss: the same text is right inside a listing and wrong in the sentence
% above it.
%
% Scoped to tt* on purpose, and measured rather than assumed: prose keeps its
% own dashes, so \dash{} still sets an em dash and a range like 1--2 still
% sets an en dash. Probed against this preamble in all three positions before
% it was believed.
\DisableLigatures[-]{encoding = T1, family = tt*}

% ---------- Language ----------
% Loading babel with a language whose .ldf is absent is a *fatal* error, not
% a warning, so both the package and each language file are probed before
% either is requested. Both editions load the other language too, because the
% glossary rows and the cheat sheet carry words from both.
\IfFileExists{babel.sty}{%
  \IfFileExists{polish.ldf}{%
    \IfFileExists{british.ldf}{%
      \ifdefstring{\booklang}{pl}%
        {\usepackage[british,main=polish]{babel}}%
        {\usepackage[polish,main=british]{babel}}%
    }{\usepackage[polish]{babel}}%
  }{%
    \IfFileExists{british.ldf}{\usepackage[british]{babel}}{}%
  }%
}{}

% babel-polish makes " an active shorthand. listings reads its body verbatim
% and an active character in a verbatim context is a category-code accident
% waiting to happen, so the shorthand is turned off. Polish quotation marks
% come from \enquote, which csquotes sets per language.
\ifdefstring{\booklang}{pl}{\AtBeginDocument{\shorthandoff{"}}}{}

% ---------- Core packages ----------
\usepackage{graphicx}
\usepackage{xcolor}
\usepackage{booktabs}
\usepackage{tabularx}
% Appendix A's cheat sheet runs past a page in every part, and tabularx
% cannot break. longtable can, and works with booktabs.
\usepackage{longtable}
\usepackage{array}
\usepackage{enumitem}
\usepackage{listings}
\usepackage[most]{tcolorbox}
\usepackage{fancyhdr}
\usepackage{titlesec}
\usepackage{caption}
\usepackage{imakeidx}
% csquotes with autostyle follows babel's active language, so the identical
% source \enquote{x} sets 'x' in the English edition and the Polish quotation
% marks in the Polish one.
\IfFileExists{csquotes.sty}{\usepackage[autostyle=true]{csquotes}}{%
  \providecommand{\enquote}[1]{``##1''}}
\usepackage[hidelinks]{hyperref}

% ---------- Palette ----------
% The same family as the two companion volumes, so the three read as a set.
\definecolor{mblue}{HTML}{1F4E79}
\definecolor{mteal}{HTML}{0E7C7B}
\definecolor{mamber}{HTML}{B26A00}
\definecolor{mred}{HTML}{9B2C2C}
\definecolor{mviolet}{HTML}{7B4397}
\definecolor{mgrey}{HTML}{5A5A5A}
\definecolor{mlight}{HTML}{F4F6F8}
\definecolor{mfaint}{HTML}{FBFCFD}
\definecolor{codebg}{HTML}{FAFAFA}
\definecolor{codeframe}{HTML}{D8DEE4}
\definecolor{kw}{HTML}{0B5394}
\definecolor{str}{HTML}{9B2C2C}
\definecolor{cmt}{HTML}{4F7A4F}
\definecolor{typ}{HTML}{0E7C7B}
\definecolor{dec}{HTML}{7B4397}

\hypersetup{
  colorlinks=true,
  linkcolor=mblue,
  urlcolor=mteal,
  citecolor=mblue,
  pdfauthor={Konrad Cinkusz}
}

% \ifpl{Polish text}{English text} -- for the handful of places where a
% string is structural rather than content. Anything longer than a few words
% belongs in lang/ or in the edition's own source, not here.
% \DeclareRobustCommand, not \newcommand: a fragile \ifpl inside a moving
% argument is written to the .toc one level expanded and fails on the next run
% in exactly one of the two editions.
\DeclareRobustCommand{\ifpl}[2]{\ifdefstring{\booklang}{pl}{#1}{#2}}
\makeatletter
\pdfstringdefDisableCommands{%
  \ifdefstring{\booklang}{pl}{\let\ifpl\@firstoftwo}{\let\ifpl\@secondoftwo}%
  % \api{}, \pkg{} and \code{} are allowed in a section title, and a title
  % becomes a PDF bookmark. hyperref drops the \color command from a bookmark
  % string but keeps its ARGUMENT, so a title carrying \api{model\_validate}
  % made a bookmark reading "mtealmodel_validate" and warned "Token not
  % allowed" twice. The colour name is gobbled here instead. Measured on a
  % probe against this preamble before the line was added.
  \def\color#1{}%
}
\makeatother

% ---------- Language strings ----------
\input{lang/\booklang}

% The PDF's subject, keyword list and document language read \lblPdf... from
% the language file input on the line above, so they cannot sit in the colour
% \hypersetup block, which runs before it. Two blocks is the ordering, not an
% oversight; merging them raises Undefined control sequence and still writes
% a PDF, with the three fields empty.
\hypersetup{
  pdfsubject={\lblPdfSubject},
  pdfkeywords={\lblPdfKeywords},
  pdflang={\lblPdfLang}
}

% ---------- Headings ----------
% \raggedright on the title body is load-bearing: at \Huge on this block a
% title that cannot break overflows the margin by 30 pt or more, and several
% chapter titles here run past fifty characters.
\titleformat{\chapter}[display]
  {\normalfont\huge\bfseries\color{mblue}\raggedright}
  {\normalfont\normalsize\color{mgrey}\MakeUppercase{\chaptertitlename}\ \thechapter}
  {8pt}{\Huge\raggedright}
\titlespacing*{\chapter}{0pt}{10pt}{28pt}
\titleformat{\section}{\normalfont\Large\bfseries\color{mblue}\raggedright}{\thesection}{0.8em}{}
\titleformat{\subsection}{\normalfont\large\bfseries\color{mgrey}\raggedright}{\thesubsection}{0.7em}{}

% Sections in the contents, subsections not: fourteen chapters of subsections
% is a contents nobody reads, and the reader navigates by chapter and section.
\setcounter{tocdepth}{1}
\setcounter{secnumdepth}{2}

% One-sided running head: the chapter on the left, the folio on the right.
\pagestyle{fancy}
\fancyhf{}
\fancyhead[L]{\small\nouppercase{\leftmark}}
\fancyhead[R]{\small\thepage}
\renewcommand{\headrulewidth}{0.4pt}
\fancypagestyle{plain}{\fancyhf{}\fancyfoot[R]{\small\thepage}%
  \renewcommand{\headrulewidth}{0pt}}

% Drops the word "Chapter" from the head: the number is still there and the
% longer titles no longer overflow. MUST come after \pagestyle{fancy}, which
% installs its own \chaptermark and silently overwrites anything earlier.
\renewcommand{\chaptermark}[1]{\markboth{\thechapter.\ #1}{}}

% ============================================================
%  CODE LISTINGS
%  ---------------------------------------------------------
%  Every listing in a chapter is a FILE under code/, pulled in with \pyfile or
%  \pyregion, and CI runs it against the pinned versions. An in-source
%  \begin{python} block is allowed for a fragment of three or four lines that
%  exists to show a shape rather than to run; anything longer belongs in a
%  file the reader can open in the editor beside the PDF.
% ============================================================
\lstdefinelanguage{pythonx}{
  language=Python,
  morekeywords={
    async,await,match,case,nonlocal,yield,type,
    TypedDict,Annotated,Protocol,Literal,Generic,Final,ClassVar,Self,
    TypeVar,TypeGuard,TypeIs,Callable,Iterator,Iterable,Awaitable,
    dataclass,field,BaseModel,Field,Enum,StrEnum,NamedTuple,
    override,overload,property,staticmethod,classmethod
  },
  morekeywords=[2]{
    gather,create_task,TaskGroup,to_thread,timeout,wait_for,Semaphore,Queue,
    fixture,parametrize,monkeypatch,raises,mark,
    Depends,FastAPI,APIRouter,HTTPException,
    select,Session,Mapped,mapped_column,relationship,
    model_validate,model_validate_json,model_dump,model_json_schema
  },
  sensitive=true
}

% listings' C# definition predates async, records and pattern matching.
\lstdefinelanguage{csharpx}{
  language=[Sharp]C,
  morekeywords={
    async,await,var,dynamic,record,init,required,nameof,yield,partial,
    global,scoped,when,with,and,or,not,notnull,unmanaged,file,get,set,value,
    where,select,from,orderby,let,join,into,ascending,descending,on,equals,
    Task,ValueTask,CancellationToken,IEnumerable,IAsyncEnumerable,Func,Action
  },
  sensitive=true
}

\lstdefinestyle{bookcode}{
  basicstyle=\ttfamily\footnotesize,
  keywordstyle=\color{kw}\bfseries,
  keywordstyle=[2]\color{typ},
  stringstyle=\color{str},
  % Colour only, no \itshape: inconsolata (zi4) has no italic shape, so an
  % italic comment style is substituted with upright on every machine that
  % has the font, and the log carries "Font shape `T1/zi4/m/it' undefined"
  % on every build. The colour already marks a comment.
  commentstyle=\color{cmt},
  emphstyle=\color{dec}\bfseries,
  numbers=left,
  numberstyle=\ttfamily\tiny\color{mgrey},
  numbersep=8pt,
  backgroundcolor=\color{codebg},
  frame=single,
  rulecolor=\color{codeframe},
  framesep=6pt,
  breaklines=true,
  breakatwhitespace=false,
  postbreak=\mbox{\textcolor{mgrey}{$\hookrightarrow$}\space},
  showstringspaces=false,
  tabsize=4,
  columns=fullflexible,
  keepspaces=true,
  captionpos=b,
  aboveskip=1.2em,
  belowskip=1.2em,
  % Maps the characters most likely to be pasted into a listing by accident.
  %
  % Do NOT add a mapping for U+00A0 (non-breaking space). It looks like the
  % obvious next entry and it makes `listings` abort the run with "Improper
  % alphabetic constant" -- fatal, no PDF, and the message names neither the
  % character nor this line. The ASCII-in-listings convention, checked by
  % parity's C13, is the guard against it instead.
  literate={ł}{{\l}}1 {Ł}{{\L}}1 {ą}{{\k a}}1 {ę}{{\k e}}1
           {ś}{{\'s}}1 {ć}{{\'c}}1 {ż}{{\.z}}1 {ź}{{\'z}}1 {ó}{{\'o}}1 {ń}{{\'n}}1
           {Ą}{{\k A}}1 {Ę}{{\k E}}1 {Ś}{{\'S}}1 {Ć}{{\'C}}1
           {Ż}{{\.Z}}1 {Ź}{{\'Z}}1 {Ó}{{\'O}}1 {Ń}{{\'N}}1
           {—}{{-{}-}}2 {–}{{-}}1 {‘}{{'}}1 {’}{{'}}1
           {“}{{"}}1 {”}{{"}}1 {°}{{\textdegree}}1
}

\lstset{style=bookcode}

\lstnewenvironment{python}[1][]
  {\lstset{language=pythonx,style=bookcode,#1}}
  {}

% C# comparison snippets. The whole book is written for engineers arriving
% from .NET, so listings come in pairs often enough that the C# half has its
% own environment rather than a generic one.
\lstnewenvironment{csharp}[1][]
  {\lstset{language=csharpx,style=bookcode,#1}}
  {}

\lstnewenvironment{shellcmd}[1][]
  {\lstset{language=bash,style=bookcode,numbers=none,keywordstyle=\color{mteal}\bfseries,#1}}
  {}

\lstnewenvironment{tomlcode}[1][]
  {\lstset{style=bookcode,numbers=none,keywordstyle=\color{black},#1}}
  {}

\lstnewenvironment{yamlcode}[1][]
  {\lstset{style=bookcode,numbers=none,keywordstyle=\color{black},#1}}
  {}

\lstnewenvironment{jsoncode}[1][]
  {\lstset{style=bookcode,numbers=none,keywordstyle=\color{black},#1}}
  {}

% Terminal transcripts typed into the source. Allowed for a shell session the
% reader is meant to reproduce; NOT for program output, which is \transcript
% below, because an in-source transcript nobody ran reads exactly like one
% that was, and that is where a remembered number hides.
\lstnewenvironment{console}[1][]
  {\lstset{style=bookcode,numbers=none,basicstyle=\ttfamily\footnotesize,
           keywordstyle=\color{black},backgroundcolor=\color{white},#1}}
  {}

% The path is printed above every file listing, in the repository-relative
% form the reader opens in the editor beside the PDF. That is the whole point
% of a listing being a file, and the front matter promises it. The \nobreak
% keeps the path on the same page as the listing it names.
% ---------------------------------------------------------------------------
%  Appendix B prints the trap catalogue from notes/02-traps.md. One entry is
%  its number -- which a chapter may cite, and which is never reused -- the
%  habit in the reader's own voice, and what Python does instead.
%
%  The body is set RAGGED RIGHT on purpose. An entry is mostly \code{} spans,
%  which are unbreakable runs, and sixty-seven justified paragraphs of them
%  is sixty-seven chances at an overfull hbox in the edition with the longer
%  words. Ragged right removes the class rather than fixing it one box at a
%  time, and a reference list is the one place in this book where it reads
%  better anyway.
%
%  The habit is set with \enquote{} rather than \emph{}, and that is not a
%  style choice. Every habit here carries \code{} spans, and \code{} inside
%  \emph{} asks inconsolata for an italic it does not have: the first build
%  of this appendix raised `Font shape T1/zi4/m/it undefined', which
%  checklog.py is hard on for exactly this reason. It is the trap CLAUDE.md
%  records, met sixty-seven times at once -- and fixed once, because the
%  entries go through a macro. A quotation is also what the habit IS, and
%  \enquote{} is what the Polish edition owes anyway.
% ---------------------------------------------------------------------------
%  Appendix A's cheat sheet: one table per part, C# on the left, Python on
%  the right, and the chapter that teaches the row on the end. The three
%  header strings are arguments so that one environment serves both
%  editions and the structural signature stays identical.
%
%  Columns are ragged right for the reason Appendix B's entries are: a cell
%  is mostly \code{} spans, which are unbreakable runs, and a justified
%  narrow column full of them is an overfull box waiting for the edition
%  with the longer words.
% ---------------------------------------------------------------------------
%  Appendix C's tool matrix: the .NET tool, the Python tool this book uses,
%  the version it was verified against and the chapter that introduces it.
%
%  The version column prints \pydanticver and its siblings and is never
%  typed, which is the whole point of the appendix: it is the one page where
%  a reader sees every pin at once, and a pin that had to be re-typed here
%  would be a second copy of a fact the preamble already owns.
%  tools/check_structure.py --tools fails when a pin macro exists and this
%  table does not print it.
\NewDocumentEnvironment{toolmatrix}{mmmm}{%
  \small
  \begin{longtable}{@{}%
    >{\raggedright\arraybackslash}p{0.27\linewidth}%
    >{\raggedright\arraybackslash}p{0.31\linewidth}%
    >{\raggedright\arraybackslash}p{0.20\linewidth}%
    >{\raggedright\arraybackslash}p{0.08\linewidth}@{}}
  \toprule
  \textbf{#1} & \textbf{#2} & \textbf{#3} & \textbf{#4} \\
  \midrule
  \endfirsthead
  \toprule
  \textbf{#1} & \textbf{#2} & \textbf{#3} & \textbf{#4} \\
  \midrule
  \endhead
  \midrule
  \endfoot
  \bottomrule
  \endlastfoot
}{%
  \end{longtable}%
}

\NewDocumentEnvironment{cheatsheet}{mmm}{%
  \small
  \begin{longtable}{@{}%
    >{\raggedright\arraybackslash}p{0.41\linewidth}%
    >{\raggedright\arraybackslash}p{0.41\linewidth}%
    >{\raggedright\arraybackslash}p{0.08\linewidth}@{}}
  \toprule
  \textbf{#1} & \textbf{#2} & \textbf{#3} \\
  \midrule
  \endfirsthead
  \toprule
  \textbf{#1} & \textbf{#2} & \textbf{#3} \\
  \midrule
  \endhead
  \midrule
  \endfoot
  \bottomrule
  \endlastfoot
}{%
  \end{longtable}%
}

\newcommand{\trapentry}[3]{%
  \par\medskip
  \noindent\textbf{#1.}\enspace\enquote{#2}\par
  \nopagebreak
  {\raggedright\noindent#3\par}%
}

\newcommand{\listingpath}[1]{%
  \par\medskip\noindent
  {\ttfamily\footnotesize\color{mgrey}\detokenize{#1}}%
  \par\nobreak}

% External source file -- the preferred form. The path is relative to the
% repository root, which is also the path the reader opens in the editor.
% Usage: \pyfile{code/ch05/decorators.py}{Caption}{lst:label}
\newcommand{\pyfile}[3]{%
  \listingpath{#1}%
  \lstinputlisting[language=pythonx,style=bookcode,aboveskip=0.3em,
    caption={#2},label={#3}]{#1}%
}

% A named region of an external file, delimited in the source by
% `# --8<-- [start:name]` / `# --8<-- [end:name]`, so the book and the running
% code cannot drift apart. tools/check_structure.py --listings fails when the
% file or the region is missing.
%
% The markers are matched through listings' range prefix and suffix keys,
% with linerange={name-name}. The form this was inherited with -- a linerange
% whose two ends were the whole marker strings -- matched nothing and printed
% the WHOLE FILE, with no warning: measured on a three-region probe file
% before this was rewritten, and the same definition sits unused in the
% sibling book it came from. A region name is a word, never a number, because
% a number is a line number to listings. The line numbers printed are the
% file's own, so a region's listing says where in the file it sits.
% Usage: \pyregion{code/ch05/decorators.py}{timing}{Caption}{lst:label}
\lstset{
  rangebeginprefix=\#\ --8<--\ [start:, rangebeginsuffix=],
  rangeendprefix=\#\ --8<--\ [end:, rangeendsuffix=],
  includerangemarker=false}
\newcommand{\pyregion}[4]{%
  \listingpath{#1}%
  \lstinputlisting[language=pythonx,style=bookcode,aboveskip=0.3em,
    linerange={#2-#2},caption={#3},label={#4}]{#1}%
}

% The C# twin of \pyfile, for the comparison half of a pair.
\newcommand{\csfile}[3]{%
  \listingpath{#1}%
  \lstinputlisting[language=csharpx,style=bookcode,aboveskip=0.3em,
    caption={#2},label={#3}]{#1}%
}

% A program's output, written by a script under code/measure/ and pulled in
% whole. There is no typed form of this and there must not be: \val{} does
% not expand inside a verbatim body, so a transcript typed into a chapter is
% one that cannot quote a computed value, and a transcript nobody ran is
% indistinguishable on the page from one that was.
\newcommand{\transcript}[1]{%
  \par\smallskip
  \IfFileExists{figures/transcripts/#1.txt}{%
    \lstinputlisting[style=bookcode,numbers=none,basicstyle=\ttfamily\footnotesize,
      keywordstyle=\color{black},backgroundcolor=\color{white},
      language={}]{figures/transcripts/#1.txt}%
  }{%
    \begin{tcolorbox}[enhanced, sharp corners, width=\linewidth,
      colback=mlight, colframe=mgrey, boxrule=0.6pt,
      left=8pt, right=8pt, top=6pt, bottom=6pt]
      {\bfseries\color{mgrey}\small \lblTranscriptMissing}\\[2pt]
      {\ttfamily\scriptsize\color{mgrey} figures/transcripts/#1.txt}
    \end{tcolorbox}%
  }%
  \par\smallskip
}

% Inline literals. Underscores inside these must be written \_.
\newcommand{\code}[1]{\texttt{\small #1}}
\newcommand{\pkg}[1]{\texttt{\small\color{mblue}#1}}
\newcommand{\api}[1]{\texttt{\small\color{mteal}#1}}

% ============================================================
%  ADMONITIONS
%  ---------------------------------------------------------
%  A deliberately small, fixed vocabulary. Two visual treatments carry the
%  whole grammar: `tinted` is a block the reader may read in passing -- a
%  tint and a thick left bar, no outline; `outlined` is a block the reader
%  must stop at, with a full rule on white. A third treatment would mean the
%  first two are not carrying their meaning.
% ============================================================
\tcbset{
  mfa tinted/.style={
    enhanced, breakable, sharp corners,
    colback=mlight, boxrule=0pt, leftrule=3pt,
    left=8pt, right=8pt, top=6pt, bottom=6pt,
    before skip=13pt, after skip=13pt,
    fonttitle=\bfseries\small,
    coltitle=black, colbacktitle=mlight,
    attach title to upper={\par\smallskip},
  },
  mfa outlined/.style={
    enhanced, breakable, sharp corners,
    colback=white, boxrule=0.6pt,
    left=8pt, right=8pt, top=6pt, bottom=6pt,
    before skip=13pt, after skip=13pt,
    fonttitle=\bfseries\small,
    coltitle=black, colbacktitle=white,
    attach title to upper={\par\smallskip},
  },
}

% Read in passing.
\newtcolorbox{note}[1][]{mfa tinted, colframe=mblue, title={\lblNote}, #1}
\newtcolorbox{warning}[1][]{mfa tinted, colframe=mred, title={\lblWarning}, #1}
% The translation. The .NET reader is the book's only assumed reader, so this
% box carries the C# side of a comparison and nothing else. Its use is
% counted: a chapter with none is a chapter that has stopped translating.
\newtcolorbox{csbox}[1][]{mfa tinted, colframe=mteal, title={\lblCsharp}, #1}
% Anything that is preview, experimental or likely to move. Python moves more
% slowly than the agent frameworks do, and this box should be rare here.
\newtcolorbox{versionbox}[1][]{mfa tinted, colframe=mamber, title={\lblVersion}, #1}

% Stop and read.
% The misconception, named AFTER the reader has walked into it. Appendix B is
% the catalogue; a trap box in a chapter is one entry of it, in place.
\newtcolorbox{trapbox}[1][]{mfa outlined, colframe=mred, title={\lblTrap}, #1}
% Marks a listing that has not been executed against the pinned versions.
% `make debt` counts these, CI prints the count, and nothing ships to a reader
% with one attached. Removing it means you ran the code.
\newtcolorbox{verifybox}[1][]{mfa outlined, colframe=mgrey, title={\lblVerify}, #1}
\newtcolorbox{exercisebox}[1][]{mfa outlined, colframe=mteal, title={\lblExercise}, #1}
% The guiding project. Points at the stage of trace-assert this section
% builds, under code/src/trace_assert/.
\newtcolorbox{projectbox}[1][]{mfa outlined, colframe=mamber, title={\lblProject}, #1}

% ============================================================
%  EXERCISES
%  ---------------------------------------------------------
%  \begin{exercise}{key}{title} ... \end{exercise}
%
%  This is the mechanism the whole book is read through: the PDF on one half
%  of the screen, an editor on the other. An exercise is therefore not a
%  paragraph but a FILE the reader opens -- a starter under
%  code/exercises/chNN/<key>.py -- and a TEST that decides whether it is done,
%  test_<key>.py beside it. The box prints both paths and the one command
%  that runs the test, and the manifest at the back lists every exercise.
%
%  An ENVIRONMENT, not a macro with a body argument: a macro argument cannot
%  contain verbatim material, so an exercise written as \exercise{key}{title}
%  {body} could never show the reader a code fragment, and the first chapter
%  that tried would have died on `\begin{python}` inside the braces.
%
%  The key is the file stem, with underscores, and it is written with
%  \detokenize so the underscore reaches the page as a character rather than
%  as a subscript. tools/check_structure.py --exercises fails when the starter
%  or the test is missing, or when the key's chapter prefix disagrees with the
%  chapter it sits in.
% ============================================================
\newcounter{exercise}[chapter]
\renewcommand{\theexercise}{\thechapter.\arabic{exercise}}
\makeatletter
\newcommand{\exercisedir}{ch\two@digits{\value{chapter}}}
\newcommand{\listofexercises}{%
  \section*{\lblExerciseManifest}%
  \phantomsection\addcontentsline{toc}{section}{\lblExerciseManifest}%
  % The entries are subsection-level and tocdepth is 1 for the contents,
  % so the depth is raised for this list alone or it prints empty --
  % which it did, on the scaffold's first clean build.
  \begingroup\c@tocdepth=2\relax\@starttoc{exr}\endgroup%
}
\makeatother
\NewDocumentEnvironment{exercise}{mm}{%
  \refstepcounter{exercise}%
  \addcontentsline{exr}{subsection}{\protect\numberline{\theexercise}\texttt{\protect\detokenize{#1}} --- #2}%
  \begin{exercisebox}[title={\lblExercise~\theexercise\ \dash{} #2}]%
}{%
  \par\medskip
  {\small\setlength{\parindent}{0pt}%
  \lblExerciseStarter: \texttt{code/exercises/\exercisedir/\detokenize{#1}.py}\\
  \lblExerciseTest: \texttt{code/exercises/\exercisedir/test\_\detokenize{#1}.py}\\
  \lblExerciseRun: \texttt{cd code \&\& uv run pytest -k \detokenize{#1}}\par}%
  \end{exercisebox}%
}

% ============================================================
%  DIAGRAMS --- Mermaid pipeline
%  ---------------------------------------------------------
%  \mermaidfig{key}{caption}{one-line manifest description}
%
%  Source of truth is figures/mermaid/<lang>/<key>.mmd, committed. `make
%  diagrams` renders each to figures/diagrams/<lang>/<key>.pdf, which is build
%  output and is not committed, so a diagram change reviews as a text diff.
%  The source is per-language because a diagram with English labels in the
%  Polish edition is exactly the half-translation that makes a bilingual book
%  feel machine-made; CI checks that both languages carry the same keys.
%
%  If the rendered PDF is absent the macro draws a placeholder AND typesets
%  the Mermaid source, in the flow rather than as a float, because a source
%  typeset verbatim is often taller than a page and a float cannot break.
% ============================================================
\makeatletter
\newcommand{\listofdiagrams}{%
  \section*{\lblDiagramManifest}%
  \phantomsection\addcontentsline{toc}{section}{\lblDiagramManifest}%
  % The entries are subsection-level and tocdepth is 1 for the contents,
  % so the depth is raised for this list alone or it prints empty --
  % which it did, on the scaffold's first clean build.
  \begingroup\c@tocdepth=2\relax\@starttoc{dgm}\endgroup%
}
\makeatother

\newcommand{\mermaidfig}[3]{%
  % \phantomsection is load-bearing and not decoration. \addcontentsline
  % records hyperref's CURRENT anchor, and here that is whatever preceded
  % the figure -- for a figure placed after a listing it is the listing's
  % line-number anchor, and listings has already advanced the counter past
  % the last line it printed, so the manifest entry linked to a
  % destination that is never created: "name{lstnumber.10.4.46} has been
  % referenced but does not exist, replaced by a fixed one". The two
  % front-matter figures had the same wrong target and were silent about
  % it, because the lines they named happen to be printed. The exercise
  % manifest never had the bug, because \begin{exercise} steps a counter
  % and that is what sets the anchor.
  \phantomsection%
  \addcontentsline{dgm}{subsection}{\protect\numberline{}\texttt{#1.mmd} --- #3}%
  \IfFileExists{figures/diagrams/\booklang/#1.pdf}{%
    \begin{figure}[htbp]
    \centering
    \includegraphics[width=0.95\linewidth,height=0.42\textheight,keepaspectratio]{figures/diagrams/\booklang/#1.pdf}%
    \caption{#2}\label{fig:#1}
    \end{figure}%
  }{%
    \par\medskip
    \begin{tcolorbox}[
      enhanced, breakable, width=\linewidth, sharp corners,
      colback=white, colframe=mgrey, boxrule=0.8pt,
      left=8pt, right=8pt, top=8pt, bottom=8pt]
      {\bfseries\color{mgrey}\small \lblNotRendered}\\[4pt]
      {\ttfamily\scriptsize\color{mgrey} figures/mermaid/\booklang/#1.mmd}\\[6pt]
      \IfFileExists{figures/mermaid/\booklang/#1.mmd}{%
        \lstinputlisting[style=bookcode,numbers=none,basicstyle=\ttfamily\scriptsize,
                         backgroundcolor=\color{mlight},frame=none,
                         keywordstyle=\color{black},language={}]{figures/mermaid/\booklang/#1.mmd}%
      }{%
        {\small\color{mred} \lblSourceMissing: \texttt{figures/mermaid/\booklang/#1.mmd}}%
      }%
    \end{tcolorbox}%
    \captionof{figure}{#2}\label{fig:#1}
    \par\medskip
  }%
}

% ============================================================
%  CHAPTER STUBS
%  ---------------------------------------------------------
%  Prints a visible marker so an unwritten chapter cannot be mistaken for a
%  written one and the page count never flatters the draft. `make debt`
%  counts them, and CI publishes the count on every build.
% ============================================================
\newcommand{\chapterstub}[1]{%
  \begin{tcolorbox}[
    enhanced, breakable, width=\linewidth, sharp corners,
    colback=mlight, colframe=mred, boxrule=1pt,
    borderline={0.6pt}{3pt}{mred, dashed},
    left=12pt, right=12pt, top=12pt, bottom=12pt]
    {\bfseries\color{mred}\large \lblNotWritten}\\[8pt]
    \small\raggedright #1
  \end{tcolorbox}%
}

% ============================================================
%  COMPUTED VALUES
%  ---------------------------------------------------------
%  A number that came out of a measurement is written by a script under
%  code/measure/ into figures/values/<name>.tex as
%
%      \pyval{e01.threads.speedup}{3.71}
%
%  and pulled into the text with \val{e01.threads.speedup}. The script is the
%  source of truth and the book cannot disagree with it, because the book does
%  not contain the digits. `make verify` fails when a committed value no
%  longer matches its script. A value the scripts do not produce prints a
%  visible marker.
%
%  \val also applies the edition's decimal separator, which is how the Polish
%  edition gets its decimal comma without a translator having to remember.
% ============================================================
\newcommand{\bookdecimalmarker}{\ifpl{,}{.}}
\IfFileExists{siunitx.sty}{%
  \usepackage{siunitx}%
  \sisetup{
    group-digits = integer,
    group-minimum-digits = 5,
    group-separator = {\,},
    output-decimal-marker = {\bookdecimalmarker}
  }%
  \newcommand{\booknum}[1]{\num{##1}}%
}{%
  \newcommand{\booknum}[1]{##1}%
}

\newcommand{\pyvalues}{}
\makeatletter
\newcommand{\pyval}[2]{\csgdef{pyval@#1}{#2}\listxadd{\pyvalues}{#1}}
% \num on a non-numeric body is a FATAL siunitx error. The generator knows
% whether it emitted a number or a word, so it writes \pyval for a number and
% \pyvaltext for a word, and this end only enforces which macro reads which.
\newcommand{\val}[1]{%
  \ifcsdef{pyval@#1}{%
    \ifcsdef{pyvaltext@#1}{%
      \PackageError{pybook}{Value `#1' is not a number}%
        {\string\val\space passes its body to siunitx, which refuses anything
         that is not a number. Use \string\valtext{#1} instead.}%
    }{}%
    \booknum{\csuse{pyval@#1}}%
  }{\textcolor{mred}{\textbf{[\lblValueMissing: #1]}}}%
}
\newcommand{\valtext}[1]{%
  \ifcsdef{pyval@#1}{\csuse{pyval@#1}}%
    {\textcolor{mred}{\textbf{[\lblValueMissing: #1]}}}%
}
\newcommand{\pyvaltext}[2]{%
  \csgdef{pyval@#1}{#2}\csgdef{pyvaltext@#1}{}\listxadd{\pyvalues}{#1}%
}
\makeatother

% figures/values/all.tex is generated by `make numbers` and does nothing but
% \input each value file. A build with no values yet still compiles; every
% \val{} in it prints a visible marker instead of a number.
\IfFileExists{figures/values/all.tex}{\input{figures/values/all}}{%
  \AtBeginDocument{\PackageWarning{pybook}{No computed values found. Run `make numbers`.}}}

% ============================================================
%  PINNED VERSIONS --- single source of truth
%  Change these here; they propagate through both editions and Appendix C.
%  Verified against pypi.org on \pinnedon (lang/<lang>.tex).
%  tools/check_versions.py compares every one of them with PyPI on each build.
%  Python itself is not on PyPI and is checked by hand against python.org.
% ============================================================
\newcommand{\bookedition}{0.1-dev}
\newcommand{\bookyear}{2026}
\newcommand{\pyver}{3.14}              % the minor the book targets
\newcommand{\pypatch}{3.14.7}          % the exact interpreter the code ran on
\newcommand{\uvver}{0.12.13}           % uv
\newcommand{\dotnetver}{10.0.100}      % dotnet
\newcommand{\ruffver}{0.16.7}          % ruff
\newcommand{\pyrightver}{1.1.414}      % pyright
\newcommand{\mypyver}{2.3.1}           % mypy
\newcommand{\pydanticver}{2.13.5}      % pydantic
\newcommand{\pysettingsver}{2.15.0}    % pydantic-settings
\newcommand{\fastapiver}{0.141.1}      % fastapi
\newcommand{\uvicornver}{0.53.0}       % uvicorn
\newcommand{\sqlalchemyver}{2.0.52}    % sqlalchemy
\newcommand{\alembicver}{1.20.0}       % alembic
\newcommand{\aiosqlitever}{0.22.1}     % aiosqlite
\newcommand{\pytestver}{9.1.1}         % pytest
\newcommand{\pytestasyncver}{1.4.0}    % pytest-asyncio
\newcommand{\hypothesisver}{6.168.0}   % hypothesis
\newcommand{\coveragever}{7.16.1}      % coverage
\newcommand{\httpxver}{0.28.1}         % httpx
\newcommand{\structlogver}{26.1.0}     % structlog
\newcommand{\otelver}{1.44.0}          % opentelemetry-sdk
\newcommand{\polarsver}{1.44.2}        % polars
\newcommand{\anthropicver}{1.5.0}      % anthropic
\newcommand{\openaiver}{3.13.0}        % openai
\newcommand{\pipauditver}{2.10.1}      % pip-audit

% ============================================================
%  INDEX
%  Two-column, and its entries include \texttt{} names that do not hyphenate.
%  Ragged right keeps multi-word entries off the margin; it cannot help a
%  single token wider than the column, so keep verbatim index entries under
%  about 29 characters and index longer names by concept instead.
%  \raggedcolumns is the switch imakeidx's multicols actually reads.
% ============================================================
\apptocmd{\theindex}{\raggedcolumns\raggedright}{}{%
  \PackageWarning{pybook}{Could not patch theindex for ragged setting}}
\providecommand*\see[2]{}
\renewcommand*\see[2]{\emph{\lblIndexSee} #1}
\providecommand*\seealso[2]{}
\renewcommand*\seealso[2]{\emph{\lblIndexSeeAlso} #1}

\makeindex
; nothing else differs between them at this level.
%  Every user-visible string lives in lang/en.tex or lang/pl.tex, so a label
%  that appears in one edition and not the other is a missing macro rather
%  than a silent divergence.
%
%  Lineage. The bilingual wiring (\booklang, lang/, body.tex, the generated
%  structure) is the math book's; the chapter machinery (listings, the
%  admonitions, verifybox, \pyfile and \pyregion, the Mermaid pipeline) is the
%  LangChain book's. What is new here is the exercise macro, which ties every
%  exercise to a starter file and a test under code/, because this book is
%  read with an editor open beside it.
% ============================================================

% \booklang must be set by the main file. Default to English rather than
% failing, so a stray % ============================================================
%  preamble.tex --- styling, environments, macros
%
%  Shared by BOTH editions. main-en.tex and main-pl.tex each define \booklang
%  before \input{preamble}; nothing else differs between them at this level.
%  Every user-visible string lives in lang/en.tex or lang/pl.tex, so a label
%  that appears in one edition and not the other is a missing macro rather
%  than a silent divergence.
%
%  Lineage. The bilingual wiring (\booklang, lang/, body.tex, the generated
%  structure) is the math book's; the chapter machinery (listings, the
%  admonitions, verifybox, \pyfile and \pyregion, the Mermaid pipeline) is the
%  LangChain book's. What is new here is the exercise macro, which ties every
%  exercise to a starter file and a test under code/, because this book is
%  read with an editor open beside it.
% ============================================================

% \booklang must be set by the main file. Default to English rather than
% failing, so a stray \input{preamble} still compiles and says what it did.
\providecommand{\booklang}{en}

\usepackage{etoolbox}   % loaded first: the language switch below needs it

% ---------- Page geometry ----------
% ONE format: A4, 12pt, ONE-SIDED. This book is read on a screen with an
% editor open beside it, and never printed as a bound book, so there is no
% recto, no verso, no gutter, no mirrored margin and no blank leaf before a
% chapter -- every one of those is a property of paper that a PDF viewer at
% half a screen's width turns into wasted pixels.
%
% The margins are symmetric and the measure is about seventy-five characters
% of prose, which is where the eye keeps the line return. The number that was
% actually chosen is the MONOSPACE measure: at \footnotesize on this block a
% listing line of 79 characters fits without wrapping, and 79 is the width
% every listing under code/ is held to by tools/check_structure.py --lines.
% A wrapped listing line is not an error TeX reports -- breaklines=true wraps
% it silently and prints an arrow into the middle of what the reader is meant
% to paste into an editor -- so the width is enforced in the source instead.
\usepackage[
  a4paper,
  left=3.1cm, right=3.1cm, top=2.6cm, bottom=3.0cm,
  head=26pt, headsep=0.8cm, footskip=1.4cm
]{geometry}

% ---------- Encoding & fonts ----------
\usepackage[T1]{fontenc}
\usepackage[utf8]{inputenc}
\IfFileExists{lmodern.sty}{\usepackage{lmodern}}{}
% newtxtext takes its TS1 symbols (\textcopyright, \textdegree) from TeX
% Gyre Termes, which Debian packages SEPARATELY (tex-gyre): a machine with
% newtx and without it dies on the copyright page with
% `Font TS1/ntxtlf/m/n/10.95=ts1-qtmr ... not loadable`. There is no
% inert-safe probe for a font metric in pdfTeX -- \IfFileExists searches
% TEXINPUTS and not TFMFONTS, so it reports the metric missing on a machine
% that has it, and \suppressfontnotfounderror is LuaTeX-only -- so nothing
% here guesses. CI's full TeX Live is the reference; locally, install
% tex-gyre. Both facts are recorded in CLAUDE.md under Build traps.
\IfFileExists{newtxtext.sty}{\usepackage{newtxtext}}{}
% amsmath before newtxmath, because newtxmath patches amsmath's internals.
% amssymb only when newtxmath is absent: newtxmath supplies the AMS symbols
% itself and loading both is a fatal clash (\Bbbk already defined) that is
% invisible on a machine without newtx and fatal on a full TeX Live.
\usepackage{amsmath}
\IfFileExists{newtxmath.sty}{\usepackage{newtxmath}}{\usepackage{amssymb}}
\IfFileExists{inconsolata.sty}{\usepackage[varqu,scaled=0.95]{inconsolata}}{}
\IfFileExists{lmodern.sty}{\usepackage{microtype}}{\usepackage[expansion=false]{microtype}}

% upquote is not a refinement. In T1 the input character ' is quoteright, so
% a listings body sets f('x') with U+2019 at both ends of the literal, and
% typed back in that is a SyntaxError. inconsolata's varqu option hides the
% defect on a machine that has inconsolata; upquote makes ' the ASCII
% apostrophe whatever the typewriter font is. In a book whose every listing is
% meant to be pasted into an editor, this line is load-bearing.
\IfFileExists{upquote.sty}{\usepackage{upquote}}{}

% And upquote's twin, for the same reason one character over. In T1 the pair
% -- is a ligature for an en dash, in EVERY family including the typewriter
% one, so \code{uv sync --locked} sets as "uv sync <en dash>locked" and a
% reader who copies it off the page gets an unknown-argument error from a
% command the chapter told them to run. It is invisible in review because it
% looks like a dash, and it bites exactly the flags a packaging chapter is
% made of. A `listings` body is already safe -- listings sets characters
% singly and no ligature forms -- which is what makes the defect so easy to
% miss: the same text is right inside a listing and wrong in the sentence
% above it.
%
% Scoped to tt* on purpose, and measured rather than assumed: prose keeps its
% own dashes, so \dash{} still sets an em dash and a range like 1--2 still
% sets an en dash. Probed against this preamble in all three positions before
% it was believed.
\DisableLigatures[-]{encoding = T1, family = tt*}

% ---------- Language ----------
% Loading babel with a language whose .ldf is absent is a *fatal* error, not
% a warning, so both the package and each language file are probed before
% either is requested. Both editions load the other language too, because the
% glossary rows and the cheat sheet carry words from both.
\IfFileExists{babel.sty}{%
  \IfFileExists{polish.ldf}{%
    \IfFileExists{british.ldf}{%
      \ifdefstring{\booklang}{pl}%
        {\usepackage[british,main=polish]{babel}}%
        {\usepackage[polish,main=british]{babel}}%
    }{\usepackage[polish]{babel}}%
  }{%
    \IfFileExists{british.ldf}{\usepackage[british]{babel}}{}%
  }%
}{}

% babel-polish makes " an active shorthand. listings reads its body verbatim
% and an active character in a verbatim context is a category-code accident
% waiting to happen, so the shorthand is turned off. Polish quotation marks
% come from \enquote, which csquotes sets per language.
\ifdefstring{\booklang}{pl}{\AtBeginDocument{\shorthandoff{"}}}{}

% ---------- Core packages ----------
\usepackage{graphicx}
\usepackage{xcolor}
\usepackage{booktabs}
\usepackage{tabularx}
% Appendix A's cheat sheet runs past a page in every part, and tabularx
% cannot break. longtable can, and works with booktabs.
\usepackage{longtable}
\usepackage{array}
\usepackage{enumitem}
\usepackage{listings}
\usepackage[most]{tcolorbox}
\usepackage{fancyhdr}
\usepackage{titlesec}
\usepackage{caption}
\usepackage{imakeidx}
% csquotes with autostyle follows babel's active language, so the identical
% source \enquote{x} sets 'x' in the English edition and the Polish quotation
% marks in the Polish one.
\IfFileExists{csquotes.sty}{\usepackage[autostyle=true]{csquotes}}{%
  \providecommand{\enquote}[1]{``##1''}}
\usepackage[hidelinks]{hyperref}

% ---------- Palette ----------
% The same family as the two companion volumes, so the three read as a set.
\definecolor{mblue}{HTML}{1F4E79}
\definecolor{mteal}{HTML}{0E7C7B}
\definecolor{mamber}{HTML}{B26A00}
\definecolor{mred}{HTML}{9B2C2C}
\definecolor{mviolet}{HTML}{7B4397}
\definecolor{mgrey}{HTML}{5A5A5A}
\definecolor{mlight}{HTML}{F4F6F8}
\definecolor{mfaint}{HTML}{FBFCFD}
\definecolor{codebg}{HTML}{FAFAFA}
\definecolor{codeframe}{HTML}{D8DEE4}
\definecolor{kw}{HTML}{0B5394}
\definecolor{str}{HTML}{9B2C2C}
\definecolor{cmt}{HTML}{4F7A4F}
\definecolor{typ}{HTML}{0E7C7B}
\definecolor{dec}{HTML}{7B4397}

\hypersetup{
  colorlinks=true,
  linkcolor=mblue,
  urlcolor=mteal,
  citecolor=mblue,
  pdfauthor={Konrad Cinkusz}
}

% \ifpl{Polish text}{English text} -- for the handful of places where a
% string is structural rather than content. Anything longer than a few words
% belongs in lang/ or in the edition's own source, not here.
% \DeclareRobustCommand, not \newcommand: a fragile \ifpl inside a moving
% argument is written to the .toc one level expanded and fails on the next run
% in exactly one of the two editions.
\DeclareRobustCommand{\ifpl}[2]{\ifdefstring{\booklang}{pl}{#1}{#2}}
\makeatletter
\pdfstringdefDisableCommands{%
  \ifdefstring{\booklang}{pl}{\let\ifpl\@firstoftwo}{\let\ifpl\@secondoftwo}%
  % \api{}, \pkg{} and \code{} are allowed in a section title, and a title
  % becomes a PDF bookmark. hyperref drops the \color command from a bookmark
  % string but keeps its ARGUMENT, so a title carrying \api{model\_validate}
  % made a bookmark reading "mtealmodel_validate" and warned "Token not
  % allowed" twice. The colour name is gobbled here instead. Measured on a
  % probe against this preamble before the line was added.
  \def\color#1{}%
}
\makeatother

% ---------- Language strings ----------
\input{lang/\booklang}

% The PDF's subject, keyword list and document language read \lblPdf... from
% the language file input on the line above, so they cannot sit in the colour
% \hypersetup block, which runs before it. Two blocks is the ordering, not an
% oversight; merging them raises Undefined control sequence and still writes
% a PDF, with the three fields empty.
\hypersetup{
  pdfsubject={\lblPdfSubject},
  pdfkeywords={\lblPdfKeywords},
  pdflang={\lblPdfLang}
}

% ---------- Headings ----------
% \raggedright on the title body is load-bearing: at \Huge on this block a
% title that cannot break overflows the margin by 30 pt or more, and several
% chapter titles here run past fifty characters.
\titleformat{\chapter}[display]
  {\normalfont\huge\bfseries\color{mblue}\raggedright}
  {\normalfont\normalsize\color{mgrey}\MakeUppercase{\chaptertitlename}\ \thechapter}
  {8pt}{\Huge\raggedright}
\titlespacing*{\chapter}{0pt}{10pt}{28pt}
\titleformat{\section}{\normalfont\Large\bfseries\color{mblue}\raggedright}{\thesection}{0.8em}{}
\titleformat{\subsection}{\normalfont\large\bfseries\color{mgrey}\raggedright}{\thesubsection}{0.7em}{}

% Sections in the contents, subsections not: fourteen chapters of subsections
% is a contents nobody reads, and the reader navigates by chapter and section.
\setcounter{tocdepth}{1}
\setcounter{secnumdepth}{2}

% One-sided running head: the chapter on the left, the folio on the right.
\pagestyle{fancy}
\fancyhf{}
\fancyhead[L]{\small\nouppercase{\leftmark}}
\fancyhead[R]{\small\thepage}
\renewcommand{\headrulewidth}{0.4pt}
\fancypagestyle{plain}{\fancyhf{}\fancyfoot[R]{\small\thepage}%
  \renewcommand{\headrulewidth}{0pt}}

% Drops the word "Chapter" from the head: the number is still there and the
% longer titles no longer overflow. MUST come after \pagestyle{fancy}, which
% installs its own \chaptermark and silently overwrites anything earlier.
\renewcommand{\chaptermark}[1]{\markboth{\thechapter.\ #1}{}}

% ============================================================
%  CODE LISTINGS
%  ---------------------------------------------------------
%  Every listing in a chapter is a FILE under code/, pulled in with \pyfile or
%  \pyregion, and CI runs it against the pinned versions. An in-source
%  \begin{python} block is allowed for a fragment of three or four lines that
%  exists to show a shape rather than to run; anything longer belongs in a
%  file the reader can open in the editor beside the PDF.
% ============================================================
\lstdefinelanguage{pythonx}{
  language=Python,
  morekeywords={
    async,await,match,case,nonlocal,yield,type,
    TypedDict,Annotated,Protocol,Literal,Generic,Final,ClassVar,Self,
    TypeVar,TypeGuard,TypeIs,Callable,Iterator,Iterable,Awaitable,
    dataclass,field,BaseModel,Field,Enum,StrEnum,NamedTuple,
    override,overload,property,staticmethod,classmethod
  },
  morekeywords=[2]{
    gather,create_task,TaskGroup,to_thread,timeout,wait_for,Semaphore,Queue,
    fixture,parametrize,monkeypatch,raises,mark,
    Depends,FastAPI,APIRouter,HTTPException,
    select,Session,Mapped,mapped_column,relationship,
    model_validate,model_validate_json,model_dump,model_json_schema
  },
  sensitive=true
}

% listings' C# definition predates async, records and pattern matching.
\lstdefinelanguage{csharpx}{
  language=[Sharp]C,
  morekeywords={
    async,await,var,dynamic,record,init,required,nameof,yield,partial,
    global,scoped,when,with,and,or,not,notnull,unmanaged,file,get,set,value,
    where,select,from,orderby,let,join,into,ascending,descending,on,equals,
    Task,ValueTask,CancellationToken,IEnumerable,IAsyncEnumerable,Func,Action
  },
  sensitive=true
}

\lstdefinestyle{bookcode}{
  basicstyle=\ttfamily\footnotesize,
  keywordstyle=\color{kw}\bfseries,
  keywordstyle=[2]\color{typ},
  stringstyle=\color{str},
  % Colour only, no \itshape: inconsolata (zi4) has no italic shape, so an
  % italic comment style is substituted with upright on every machine that
  % has the font, and the log carries "Font shape `T1/zi4/m/it' undefined"
  % on every build. The colour already marks a comment.
  commentstyle=\color{cmt},
  emphstyle=\color{dec}\bfseries,
  numbers=left,
  numberstyle=\ttfamily\tiny\color{mgrey},
  numbersep=8pt,
  backgroundcolor=\color{codebg},
  frame=single,
  rulecolor=\color{codeframe},
  framesep=6pt,
  breaklines=true,
  breakatwhitespace=false,
  postbreak=\mbox{\textcolor{mgrey}{$\hookrightarrow$}\space},
  showstringspaces=false,
  tabsize=4,
  columns=fullflexible,
  keepspaces=true,
  captionpos=b,
  aboveskip=1.2em,
  belowskip=1.2em,
  % Maps the characters most likely to be pasted into a listing by accident.
  %
  % Do NOT add a mapping for U+00A0 (non-breaking space). It looks like the
  % obvious next entry and it makes `listings` abort the run with "Improper
  % alphabetic constant" -- fatal, no PDF, and the message names neither the
  % character nor this line. The ASCII-in-listings convention, checked by
  % parity's C13, is the guard against it instead.
  literate={ł}{{\l}}1 {Ł}{{\L}}1 {ą}{{\k a}}1 {ę}{{\k e}}1
           {ś}{{\'s}}1 {ć}{{\'c}}1 {ż}{{\.z}}1 {ź}{{\'z}}1 {ó}{{\'o}}1 {ń}{{\'n}}1
           {Ą}{{\k A}}1 {Ę}{{\k E}}1 {Ś}{{\'S}}1 {Ć}{{\'C}}1
           {Ż}{{\.Z}}1 {Ź}{{\'Z}}1 {Ó}{{\'O}}1 {Ń}{{\'N}}1
           {—}{{-{}-}}2 {–}{{-}}1 {‘}{{'}}1 {’}{{'}}1
           {“}{{"}}1 {”}{{"}}1 {°}{{\textdegree}}1
}

\lstset{style=bookcode}

\lstnewenvironment{python}[1][]
  {\lstset{language=pythonx,style=bookcode,#1}}
  {}

% C# comparison snippets. The whole book is written for engineers arriving
% from .NET, so listings come in pairs often enough that the C# half has its
% own environment rather than a generic one.
\lstnewenvironment{csharp}[1][]
  {\lstset{language=csharpx,style=bookcode,#1}}
  {}

\lstnewenvironment{shellcmd}[1][]
  {\lstset{language=bash,style=bookcode,numbers=none,keywordstyle=\color{mteal}\bfseries,#1}}
  {}

\lstnewenvironment{tomlcode}[1][]
  {\lstset{style=bookcode,numbers=none,keywordstyle=\color{black},#1}}
  {}

\lstnewenvironment{yamlcode}[1][]
  {\lstset{style=bookcode,numbers=none,keywordstyle=\color{black},#1}}
  {}

\lstnewenvironment{jsoncode}[1][]
  {\lstset{style=bookcode,numbers=none,keywordstyle=\color{black},#1}}
  {}

% Terminal transcripts typed into the source. Allowed for a shell session the
% reader is meant to reproduce; NOT for program output, which is \transcript
% below, because an in-source transcript nobody ran reads exactly like one
% that was, and that is where a remembered number hides.
\lstnewenvironment{console}[1][]
  {\lstset{style=bookcode,numbers=none,basicstyle=\ttfamily\footnotesize,
           keywordstyle=\color{black},backgroundcolor=\color{white},#1}}
  {}

% The path is printed above every file listing, in the repository-relative
% form the reader opens in the editor beside the PDF. That is the whole point
% of a listing being a file, and the front matter promises it. The \nobreak
% keeps the path on the same page as the listing it names.
% ---------------------------------------------------------------------------
%  Appendix B prints the trap catalogue from notes/02-traps.md. One entry is
%  its number -- which a chapter may cite, and which is never reused -- the
%  habit in the reader's own voice, and what Python does instead.
%
%  The body is set RAGGED RIGHT on purpose. An entry is mostly \code{} spans,
%  which are unbreakable runs, and sixty-seven justified paragraphs of them
%  is sixty-seven chances at an overfull hbox in the edition with the longer
%  words. Ragged right removes the class rather than fixing it one box at a
%  time, and a reference list is the one place in this book where it reads
%  better anyway.
%
%  The habit is set with \enquote{} rather than \emph{}, and that is not a
%  style choice. Every habit here carries \code{} spans, and \code{} inside
%  \emph{} asks inconsolata for an italic it does not have: the first build
%  of this appendix raised `Font shape T1/zi4/m/it undefined', which
%  checklog.py is hard on for exactly this reason. It is the trap CLAUDE.md
%  records, met sixty-seven times at once -- and fixed once, because the
%  entries go through a macro. A quotation is also what the habit IS, and
%  \enquote{} is what the Polish edition owes anyway.
% ---------------------------------------------------------------------------
%  Appendix A's cheat sheet: one table per part, C# on the left, Python on
%  the right, and the chapter that teaches the row on the end. The three
%  header strings are arguments so that one environment serves both
%  editions and the structural signature stays identical.
%
%  Columns are ragged right for the reason Appendix B's entries are: a cell
%  is mostly \code{} spans, which are unbreakable runs, and a justified
%  narrow column full of them is an overfull box waiting for the edition
%  with the longer words.
% ---------------------------------------------------------------------------
%  Appendix C's tool matrix: the .NET tool, the Python tool this book uses,
%  the version it was verified against and the chapter that introduces it.
%
%  The version column prints \pydanticver and its siblings and is never
%  typed, which is the whole point of the appendix: it is the one page where
%  a reader sees every pin at once, and a pin that had to be re-typed here
%  would be a second copy of a fact the preamble already owns.
%  tools/check_structure.py --tools fails when a pin macro exists and this
%  table does not print it.
\NewDocumentEnvironment{toolmatrix}{mmmm}{%
  \small
  \begin{longtable}{@{}%
    >{\raggedright\arraybackslash}p{0.27\linewidth}%
    >{\raggedright\arraybackslash}p{0.31\linewidth}%
    >{\raggedright\arraybackslash}p{0.20\linewidth}%
    >{\raggedright\arraybackslash}p{0.08\linewidth}@{}}
  \toprule
  \textbf{#1} & \textbf{#2} & \textbf{#3} & \textbf{#4} \\
  \midrule
  \endfirsthead
  \toprule
  \textbf{#1} & \textbf{#2} & \textbf{#3} & \textbf{#4} \\
  \midrule
  \endhead
  \midrule
  \endfoot
  \bottomrule
  \endlastfoot
}{%
  \end{longtable}%
}

\NewDocumentEnvironment{cheatsheet}{mmm}{%
  \small
  \begin{longtable}{@{}%
    >{\raggedright\arraybackslash}p{0.41\linewidth}%
    >{\raggedright\arraybackslash}p{0.41\linewidth}%
    >{\raggedright\arraybackslash}p{0.08\linewidth}@{}}
  \toprule
  \textbf{#1} & \textbf{#2} & \textbf{#3} \\
  \midrule
  \endfirsthead
  \toprule
  \textbf{#1} & \textbf{#2} & \textbf{#3} \\
  \midrule
  \endhead
  \midrule
  \endfoot
  \bottomrule
  \endlastfoot
}{%
  \end{longtable}%
}

\newcommand{\trapentry}[3]{%
  \par\medskip
  \noindent\textbf{#1.}\enspace\enquote{#2}\par
  \nopagebreak
  {\raggedright\noindent#3\par}%
}

\newcommand{\listingpath}[1]{%
  \par\medskip\noindent
  {\ttfamily\footnotesize\color{mgrey}\detokenize{#1}}%
  \par\nobreak}

% External source file -- the preferred form. The path is relative to the
% repository root, which is also the path the reader opens in the editor.
% Usage: \pyfile{code/ch05/decorators.py}{Caption}{lst:label}
\newcommand{\pyfile}[3]{%
  \listingpath{#1}%
  \lstinputlisting[language=pythonx,style=bookcode,aboveskip=0.3em,
    caption={#2},label={#3}]{#1}%
}

% A named region of an external file, delimited in the source by
% `# --8<-- [start:name]` / `# --8<-- [end:name]`, so the book and the running
% code cannot drift apart. tools/check_structure.py --listings fails when the
% file or the region is missing.
%
% The markers are matched through listings' range prefix and suffix keys,
% with linerange={name-name}. The form this was inherited with -- a linerange
% whose two ends were the whole marker strings -- matched nothing and printed
% the WHOLE FILE, with no warning: measured on a three-region probe file
% before this was rewritten, and the same definition sits unused in the
% sibling book it came from. A region name is a word, never a number, because
% a number is a line number to listings. The line numbers printed are the
% file's own, so a region's listing says where in the file it sits.
% Usage: \pyregion{code/ch05/decorators.py}{timing}{Caption}{lst:label}
\lstset{
  rangebeginprefix=\#\ --8<--\ [start:, rangebeginsuffix=],
  rangeendprefix=\#\ --8<--\ [end:, rangeendsuffix=],
  includerangemarker=false}
\newcommand{\pyregion}[4]{%
  \listingpath{#1}%
  \lstinputlisting[language=pythonx,style=bookcode,aboveskip=0.3em,
    linerange={#2-#2},caption={#3},label={#4}]{#1}%
}

% The C# twin of \pyfile, for the comparison half of a pair.
\newcommand{\csfile}[3]{%
  \listingpath{#1}%
  \lstinputlisting[language=csharpx,style=bookcode,aboveskip=0.3em,
    caption={#2},label={#3}]{#1}%
}

% A program's output, written by a script under code/measure/ and pulled in
% whole. There is no typed form of this and there must not be: \val{} does
% not expand inside a verbatim body, so a transcript typed into a chapter is
% one that cannot quote a computed value, and a transcript nobody ran is
% indistinguishable on the page from one that was.
\newcommand{\transcript}[1]{%
  \par\smallskip
  \IfFileExists{figures/transcripts/#1.txt}{%
    \lstinputlisting[style=bookcode,numbers=none,basicstyle=\ttfamily\footnotesize,
      keywordstyle=\color{black},backgroundcolor=\color{white},
      language={}]{figures/transcripts/#1.txt}%
  }{%
    \begin{tcolorbox}[enhanced, sharp corners, width=\linewidth,
      colback=mlight, colframe=mgrey, boxrule=0.6pt,
      left=8pt, right=8pt, top=6pt, bottom=6pt]
      {\bfseries\color{mgrey}\small \lblTranscriptMissing}\\[2pt]
      {\ttfamily\scriptsize\color{mgrey} figures/transcripts/#1.txt}
    \end{tcolorbox}%
  }%
  \par\smallskip
}

% Inline literals. Underscores inside these must be written \_.
\newcommand{\code}[1]{\texttt{\small #1}}
\newcommand{\pkg}[1]{\texttt{\small\color{mblue}#1}}
\newcommand{\api}[1]{\texttt{\small\color{mteal}#1}}

% ============================================================
%  ADMONITIONS
%  ---------------------------------------------------------
%  A deliberately small, fixed vocabulary. Two visual treatments carry the
%  whole grammar: `tinted` is a block the reader may read in passing -- a
%  tint and a thick left bar, no outline; `outlined` is a block the reader
%  must stop at, with a full rule on white. A third treatment would mean the
%  first two are not carrying their meaning.
% ============================================================
\tcbset{
  mfa tinted/.style={
    enhanced, breakable, sharp corners,
    colback=mlight, boxrule=0pt, leftrule=3pt,
    left=8pt, right=8pt, top=6pt, bottom=6pt,
    before skip=13pt, after skip=13pt,
    fonttitle=\bfseries\small,
    coltitle=black, colbacktitle=mlight,
    attach title to upper={\par\smallskip},
  },
  mfa outlined/.style={
    enhanced, breakable, sharp corners,
    colback=white, boxrule=0.6pt,
    left=8pt, right=8pt, top=6pt, bottom=6pt,
    before skip=13pt, after skip=13pt,
    fonttitle=\bfseries\small,
    coltitle=black, colbacktitle=white,
    attach title to upper={\par\smallskip},
  },
}

% Read in passing.
\newtcolorbox{note}[1][]{mfa tinted, colframe=mblue, title={\lblNote}, #1}
\newtcolorbox{warning}[1][]{mfa tinted, colframe=mred, title={\lblWarning}, #1}
% The translation. The .NET reader is the book's only assumed reader, so this
% box carries the C# side of a comparison and nothing else. Its use is
% counted: a chapter with none is a chapter that has stopped translating.
\newtcolorbox{csbox}[1][]{mfa tinted, colframe=mteal, title={\lblCsharp}, #1}
% Anything that is preview, experimental or likely to move. Python moves more
% slowly than the agent frameworks do, and this box should be rare here.
\newtcolorbox{versionbox}[1][]{mfa tinted, colframe=mamber, title={\lblVersion}, #1}

% Stop and read.
% The misconception, named AFTER the reader has walked into it. Appendix B is
% the catalogue; a trap box in a chapter is one entry of it, in place.
\newtcolorbox{trapbox}[1][]{mfa outlined, colframe=mred, title={\lblTrap}, #1}
% Marks a listing that has not been executed against the pinned versions.
% `make debt` counts these, CI prints the count, and nothing ships to a reader
% with one attached. Removing it means you ran the code.
\newtcolorbox{verifybox}[1][]{mfa outlined, colframe=mgrey, title={\lblVerify}, #1}
\newtcolorbox{exercisebox}[1][]{mfa outlined, colframe=mteal, title={\lblExercise}, #1}
% The guiding project. Points at the stage of trace-assert this section
% builds, under code/src/trace_assert/.
\newtcolorbox{projectbox}[1][]{mfa outlined, colframe=mamber, title={\lblProject}, #1}

% ============================================================
%  EXERCISES
%  ---------------------------------------------------------
%  \begin{exercise}{key}{title} ... \end{exercise}
%
%  This is the mechanism the whole book is read through: the PDF on one half
%  of the screen, an editor on the other. An exercise is therefore not a
%  paragraph but a FILE the reader opens -- a starter under
%  code/exercises/chNN/<key>.py -- and a TEST that decides whether it is done,
%  test_<key>.py beside it. The box prints both paths and the one command
%  that runs the test, and the manifest at the back lists every exercise.
%
%  An ENVIRONMENT, not a macro with a body argument: a macro argument cannot
%  contain verbatim material, so an exercise written as \exercise{key}{title}
%  {body} could never show the reader a code fragment, and the first chapter
%  that tried would have died on `\begin{python}` inside the braces.
%
%  The key is the file stem, with underscores, and it is written with
%  \detokenize so the underscore reaches the page as a character rather than
%  as a subscript. tools/check_structure.py --exercises fails when the starter
%  or the test is missing, or when the key's chapter prefix disagrees with the
%  chapter it sits in.
% ============================================================
\newcounter{exercise}[chapter]
\renewcommand{\theexercise}{\thechapter.\arabic{exercise}}
\makeatletter
\newcommand{\exercisedir}{ch\two@digits{\value{chapter}}}
\newcommand{\listofexercises}{%
  \section*{\lblExerciseManifest}%
  \phantomsection\addcontentsline{toc}{section}{\lblExerciseManifest}%
  % The entries are subsection-level and tocdepth is 1 for the contents,
  % so the depth is raised for this list alone or it prints empty --
  % which it did, on the scaffold's first clean build.
  \begingroup\c@tocdepth=2\relax\@starttoc{exr}\endgroup%
}
\makeatother
\NewDocumentEnvironment{exercise}{mm}{%
  \refstepcounter{exercise}%
  \addcontentsline{exr}{subsection}{\protect\numberline{\theexercise}\texttt{\protect\detokenize{#1}} --- #2}%
  \begin{exercisebox}[title={\lblExercise~\theexercise\ \dash{} #2}]%
}{%
  \par\medskip
  {\small\setlength{\parindent}{0pt}%
  \lblExerciseStarter: \texttt{code/exercises/\exercisedir/\detokenize{#1}.py}\\
  \lblExerciseTest: \texttt{code/exercises/\exercisedir/test\_\detokenize{#1}.py}\\
  \lblExerciseRun: \texttt{cd code \&\& uv run pytest -k \detokenize{#1}}\par}%
  \end{exercisebox}%
}

% ============================================================
%  DIAGRAMS --- Mermaid pipeline
%  ---------------------------------------------------------
%  \mermaidfig{key}{caption}{one-line manifest description}
%
%  Source of truth is figures/mermaid/<lang>/<key>.mmd, committed. `make
%  diagrams` renders each to figures/diagrams/<lang>/<key>.pdf, which is build
%  output and is not committed, so a diagram change reviews as a text diff.
%  The source is per-language because a diagram with English labels in the
%  Polish edition is exactly the half-translation that makes a bilingual book
%  feel machine-made; CI checks that both languages carry the same keys.
%
%  If the rendered PDF is absent the macro draws a placeholder AND typesets
%  the Mermaid source, in the flow rather than as a float, because a source
%  typeset verbatim is often taller than a page and a float cannot break.
% ============================================================
\makeatletter
\newcommand{\listofdiagrams}{%
  \section*{\lblDiagramManifest}%
  \phantomsection\addcontentsline{toc}{section}{\lblDiagramManifest}%
  % The entries are subsection-level and tocdepth is 1 for the contents,
  % so the depth is raised for this list alone or it prints empty --
  % which it did, on the scaffold's first clean build.
  \begingroup\c@tocdepth=2\relax\@starttoc{dgm}\endgroup%
}
\makeatother

\newcommand{\mermaidfig}[3]{%
  % \phantomsection is load-bearing and not decoration. \addcontentsline
  % records hyperref's CURRENT anchor, and here that is whatever preceded
  % the figure -- for a figure placed after a listing it is the listing's
  % line-number anchor, and listings has already advanced the counter past
  % the last line it printed, so the manifest entry linked to a
  % destination that is never created: "name{lstnumber.10.4.46} has been
  % referenced but does not exist, replaced by a fixed one". The two
  % front-matter figures had the same wrong target and were silent about
  % it, because the lines they named happen to be printed. The exercise
  % manifest never had the bug, because \begin{exercise} steps a counter
  % and that is what sets the anchor.
  \phantomsection%
  \addcontentsline{dgm}{subsection}{\protect\numberline{}\texttt{#1.mmd} --- #3}%
  \IfFileExists{figures/diagrams/\booklang/#1.pdf}{%
    \begin{figure}[htbp]
    \centering
    \includegraphics[width=0.95\linewidth,height=0.42\textheight,keepaspectratio]{figures/diagrams/\booklang/#1.pdf}%
    \caption{#2}\label{fig:#1}
    \end{figure}%
  }{%
    \par\medskip
    \begin{tcolorbox}[
      enhanced, breakable, width=\linewidth, sharp corners,
      colback=white, colframe=mgrey, boxrule=0.8pt,
      left=8pt, right=8pt, top=8pt, bottom=8pt]
      {\bfseries\color{mgrey}\small \lblNotRendered}\\[4pt]
      {\ttfamily\scriptsize\color{mgrey} figures/mermaid/\booklang/#1.mmd}\\[6pt]
      \IfFileExists{figures/mermaid/\booklang/#1.mmd}{%
        \lstinputlisting[style=bookcode,numbers=none,basicstyle=\ttfamily\scriptsize,
                         backgroundcolor=\color{mlight},frame=none,
                         keywordstyle=\color{black},language={}]{figures/mermaid/\booklang/#1.mmd}%
      }{%
        {\small\color{mred} \lblSourceMissing: \texttt{figures/mermaid/\booklang/#1.mmd}}%
      }%
    \end{tcolorbox}%
    \captionof{figure}{#2}\label{fig:#1}
    \par\medskip
  }%
}

% ============================================================
%  CHAPTER STUBS
%  ---------------------------------------------------------
%  Prints a visible marker so an unwritten chapter cannot be mistaken for a
%  written one and the page count never flatters the draft. `make debt`
%  counts them, and CI publishes the count on every build.
% ============================================================
\newcommand{\chapterstub}[1]{%
  \begin{tcolorbox}[
    enhanced, breakable, width=\linewidth, sharp corners,
    colback=mlight, colframe=mred, boxrule=1pt,
    borderline={0.6pt}{3pt}{mred, dashed},
    left=12pt, right=12pt, top=12pt, bottom=12pt]
    {\bfseries\color{mred}\large \lblNotWritten}\\[8pt]
    \small\raggedright #1
  \end{tcolorbox}%
}

% ============================================================
%  COMPUTED VALUES
%  ---------------------------------------------------------
%  A number that came out of a measurement is written by a script under
%  code/measure/ into figures/values/<name>.tex as
%
%      \pyval{e01.threads.speedup}{3.71}
%
%  and pulled into the text with \val{e01.threads.speedup}. The script is the
%  source of truth and the book cannot disagree with it, because the book does
%  not contain the digits. `make verify` fails when a committed value no
%  longer matches its script. A value the scripts do not produce prints a
%  visible marker.
%
%  \val also applies the edition's decimal separator, which is how the Polish
%  edition gets its decimal comma without a translator having to remember.
% ============================================================
\newcommand{\bookdecimalmarker}{\ifpl{,}{.}}
\IfFileExists{siunitx.sty}{%
  \usepackage{siunitx}%
  \sisetup{
    group-digits = integer,
    group-minimum-digits = 5,
    group-separator = {\,},
    output-decimal-marker = {\bookdecimalmarker}
  }%
  \newcommand{\booknum}[1]{\num{##1}}%
}{%
  \newcommand{\booknum}[1]{##1}%
}

\newcommand{\pyvalues}{}
\makeatletter
\newcommand{\pyval}[2]{\csgdef{pyval@#1}{#2}\listxadd{\pyvalues}{#1}}
% \num on a non-numeric body is a FATAL siunitx error. The generator knows
% whether it emitted a number or a word, so it writes \pyval for a number and
% \pyvaltext for a word, and this end only enforces which macro reads which.
\newcommand{\val}[1]{%
  \ifcsdef{pyval@#1}{%
    \ifcsdef{pyvaltext@#1}{%
      \PackageError{pybook}{Value `#1' is not a number}%
        {\string\val\space passes its body to siunitx, which refuses anything
         that is not a number. Use \string\valtext{#1} instead.}%
    }{}%
    \booknum{\csuse{pyval@#1}}%
  }{\textcolor{mred}{\textbf{[\lblValueMissing: #1]}}}%
}
\newcommand{\valtext}[1]{%
  \ifcsdef{pyval@#1}{\csuse{pyval@#1}}%
    {\textcolor{mred}{\textbf{[\lblValueMissing: #1]}}}%
}
\newcommand{\pyvaltext}[2]{%
  \csgdef{pyval@#1}{#2}\csgdef{pyvaltext@#1}{}\listxadd{\pyvalues}{#1}%
}
\makeatother

% figures/values/all.tex is generated by `make numbers` and does nothing but
% \input each value file. A build with no values yet still compiles; every
% \val{} in it prints a visible marker instead of a number.
\IfFileExists{figures/values/all.tex}{\input{figures/values/all}}{%
  \AtBeginDocument{\PackageWarning{pybook}{No computed values found. Run `make numbers`.}}}

% ============================================================
%  PINNED VERSIONS --- single source of truth
%  Change these here; they propagate through both editions and Appendix C.
%  Verified against pypi.org on \pinnedon (lang/<lang>.tex).
%  tools/check_versions.py compares every one of them with PyPI on each build.
%  Python itself is not on PyPI and is checked by hand against python.org.
% ============================================================
\newcommand{\bookedition}{0.1-dev}
\newcommand{\bookyear}{2026}
\newcommand{\pyver}{3.14}              % the minor the book targets
\newcommand{\pypatch}{3.14.7}          % the exact interpreter the code ran on
\newcommand{\uvver}{0.12.13}           % uv
\newcommand{\dotnetver}{10.0.100}      % dotnet
\newcommand{\ruffver}{0.16.7}          % ruff
\newcommand{\pyrightver}{1.1.414}      % pyright
\newcommand{\mypyver}{2.3.1}           % mypy
\newcommand{\pydanticver}{2.13.5}      % pydantic
\newcommand{\pysettingsver}{2.15.0}    % pydantic-settings
\newcommand{\fastapiver}{0.141.1}      % fastapi
\newcommand{\uvicornver}{0.53.0}       % uvicorn
\newcommand{\sqlalchemyver}{2.0.52}    % sqlalchemy
\newcommand{\alembicver}{1.20.0}       % alembic
\newcommand{\aiosqlitever}{0.22.1}     % aiosqlite
\newcommand{\pytestver}{9.1.1}         % pytest
\newcommand{\pytestasyncver}{1.4.0}    % pytest-asyncio
\newcommand{\hypothesisver}{6.168.0}   % hypothesis
\newcommand{\coveragever}{7.16.1}      % coverage
\newcommand{\httpxver}{0.28.1}         % httpx
\newcommand{\structlogver}{26.1.0}     % structlog
\newcommand{\otelver}{1.44.0}          % opentelemetry-sdk
\newcommand{\polarsver}{1.44.2}        % polars
\newcommand{\anthropicver}{1.5.0}      % anthropic
\newcommand{\openaiver}{3.13.0}        % openai
\newcommand{\pipauditver}{2.10.1}      % pip-audit

% ============================================================
%  INDEX
%  Two-column, and its entries include \texttt{} names that do not hyphenate.
%  Ragged right keeps multi-word entries off the margin; it cannot help a
%  single token wider than the column, so keep verbatim index entries under
%  about 29 characters and index longer names by concept instead.
%  \raggedcolumns is the switch imakeidx's multicols actually reads.
% ============================================================
\apptocmd{\theindex}{\raggedcolumns\raggedright}{}{%
  \PackageWarning{pybook}{Could not patch theindex for ragged setting}}
\providecommand*\see[2]{}
\renewcommand*\see[2]{\emph{\lblIndexSee} #1}
\providecommand*\seealso[2]{}
\renewcommand*\seealso[2]{\emph{\lblIndexSeeAlso} #1}

\makeindex
 still compiles and says what it did.
\providecommand{\booklang}{en}

\usepackage{etoolbox}   % loaded first: the language switch below needs it

% ---------- Page geometry ----------
% ONE format: A4, 12pt, ONE-SIDED. This book is read on a screen with an
% editor open beside it, and never printed as a bound book, so there is no
% recto, no verso, no gutter, no mirrored margin and no blank leaf before a
% chapter -- every one of those is a property of paper that a PDF viewer at
% half a screen's width turns into wasted pixels.
%
% The margins are symmetric and the measure is about seventy-five characters
% of prose, which is where the eye keeps the line return. The number that was
% actually chosen is the MONOSPACE measure: at \footnotesize on this block a
% listing line of 79 characters fits without wrapping, and 79 is the width
% every listing under code/ is held to by tools/check_structure.py --lines.
% A wrapped listing line is not an error TeX reports -- breaklines=true wraps
% it silently and prints an arrow into the middle of what the reader is meant
% to paste into an editor -- so the width is enforced in the source instead.
\usepackage[
  a4paper,
  left=3.1cm, right=3.1cm, top=2.6cm, bottom=3.0cm,
  head=26pt, headsep=0.8cm, footskip=1.4cm
]{geometry}

% ---------- Encoding & fonts ----------
\usepackage[T1]{fontenc}
\usepackage[utf8]{inputenc}
\IfFileExists{lmodern.sty}{\usepackage{lmodern}}{}
% newtxtext takes its TS1 symbols (\textcopyright, \textdegree) from TeX
% Gyre Termes, which Debian packages SEPARATELY (tex-gyre): a machine with
% newtx and without it dies on the copyright page with
% `Font TS1/ntxtlf/m/n/10.95=ts1-qtmr ... not loadable`. There is no
% inert-safe probe for a font metric in pdfTeX -- \IfFileExists searches
% TEXINPUTS and not TFMFONTS, so it reports the metric missing on a machine
% that has it, and \suppressfontnotfounderror is LuaTeX-only -- so nothing
% here guesses. CI's full TeX Live is the reference; locally, install
% tex-gyre. Both facts are recorded in CLAUDE.md under Build traps.
\IfFileExists{newtxtext.sty}{\usepackage{newtxtext}}{}
% amsmath before newtxmath, because newtxmath patches amsmath's internals.
% amssymb only when newtxmath is absent: newtxmath supplies the AMS symbols
% itself and loading both is a fatal clash (\Bbbk already defined) that is
% invisible on a machine without newtx and fatal on a full TeX Live.
\usepackage{amsmath}
\IfFileExists{newtxmath.sty}{\usepackage{newtxmath}}{\usepackage{amssymb}}
\IfFileExists{inconsolata.sty}{\usepackage[varqu,scaled=0.95]{inconsolata}}{}
\IfFileExists{lmodern.sty}{\usepackage{microtype}}{\usepackage[expansion=false]{microtype}}

% upquote is not a refinement. In T1 the input character ' is quoteright, so
% a listings body sets f('x') with U+2019 at both ends of the literal, and
% typed back in that is a SyntaxError. inconsolata's varqu option hides the
% defect on a machine that has inconsolata; upquote makes ' the ASCII
% apostrophe whatever the typewriter font is. In a book whose every listing is
% meant to be pasted into an editor, this line is load-bearing.
\IfFileExists{upquote.sty}{\usepackage{upquote}}{}

% And upquote's twin, for the same reason one character over. In T1 the pair
% -- is a ligature for an en dash, in EVERY family including the typewriter
% one, so \code{uv sync --locked} sets as "uv sync <en dash>locked" and a
% reader who copies it off the page gets an unknown-argument error from a
% command the chapter told them to run. It is invisible in review because it
% looks like a dash, and it bites exactly the flags a packaging chapter is
% made of. A `listings` body is already safe -- listings sets characters
% singly and no ligature forms -- which is what makes the defect so easy to
% miss: the same text is right inside a listing and wrong in the sentence
% above it.
%
% Scoped to tt* on purpose, and measured rather than assumed: prose keeps its
% own dashes, so \dash{} still sets an em dash and a range like 1--2 still
% sets an en dash. Probed against this preamble in all three positions before
% it was believed.
\DisableLigatures[-]{encoding = T1, family = tt*}

% ---------- Language ----------
% Loading babel with a language whose .ldf is absent is a *fatal* error, not
% a warning, so both the package and each language file are probed before
% either is requested. Both editions load the other language too, because the
% glossary rows and the cheat sheet carry words from both.
\IfFileExists{babel.sty}{%
  \IfFileExists{polish.ldf}{%
    \IfFileExists{british.ldf}{%
      \ifdefstring{\booklang}{pl}%
        {\usepackage[british,main=polish]{babel}}%
        {\usepackage[polish,main=british]{babel}}%
    }{\usepackage[polish]{babel}}%
  }{%
    \IfFileExists{british.ldf}{\usepackage[british]{babel}}{}%
  }%
}{}

% babel-polish makes " an active shorthand. listings reads its body verbatim
% and an active character in a verbatim context is a category-code accident
% waiting to happen, so the shorthand is turned off. Polish quotation marks
% come from \enquote, which csquotes sets per language.
\ifdefstring{\booklang}{pl}{\AtBeginDocument{\shorthandoff{"}}}{}

% ---------- Core packages ----------
\usepackage{graphicx}
\usepackage{xcolor}
\usepackage{booktabs}
\usepackage{tabularx}
% Appendix A's cheat sheet runs past a page in every part, and tabularx
% cannot break. longtable can, and works with booktabs.
\usepackage{longtable}
\usepackage{array}
\usepackage{enumitem}
\usepackage{listings}
\usepackage[most]{tcolorbox}
\usepackage{fancyhdr}
\usepackage{titlesec}
\usepackage{caption}
\usepackage{imakeidx}
% csquotes with autostyle follows babel's active language, so the identical
% source \enquote{x} sets 'x' in the English edition and the Polish quotation
% marks in the Polish one.
\IfFileExists{csquotes.sty}{\usepackage[autostyle=true]{csquotes}}{%
  \providecommand{\enquote}[1]{``##1''}}
\usepackage[hidelinks]{hyperref}

% ---------- Palette ----------
% The same family as the two companion volumes, so the three read as a set.
\definecolor{mblue}{HTML}{1F4E79}
\definecolor{mteal}{HTML}{0E7C7B}
\definecolor{mamber}{HTML}{B26A00}
\definecolor{mred}{HTML}{9B2C2C}
\definecolor{mviolet}{HTML}{7B4397}
\definecolor{mgrey}{HTML}{5A5A5A}
\definecolor{mlight}{HTML}{F4F6F8}
\definecolor{mfaint}{HTML}{FBFCFD}
\definecolor{codebg}{HTML}{FAFAFA}
\definecolor{codeframe}{HTML}{D8DEE4}
\definecolor{kw}{HTML}{0B5394}
\definecolor{str}{HTML}{9B2C2C}
\definecolor{cmt}{HTML}{4F7A4F}
\definecolor{typ}{HTML}{0E7C7B}
\definecolor{dec}{HTML}{7B4397}

\hypersetup{
  colorlinks=true,
  linkcolor=mblue,
  urlcolor=mteal,
  citecolor=mblue,
  pdfauthor={Konrad Cinkusz}
}

% \ifpl{Polish text}{English text} -- for the handful of places where a
% string is structural rather than content. Anything longer than a few words
% belongs in lang/ or in the edition's own source, not here.
% \DeclareRobustCommand, not \newcommand: a fragile \ifpl inside a moving
% argument is written to the .toc one level expanded and fails on the next run
% in exactly one of the two editions.
\DeclareRobustCommand{\ifpl}[2]{\ifdefstring{\booklang}{pl}{#1}{#2}}
\makeatletter
\pdfstringdefDisableCommands{%
  \ifdefstring{\booklang}{pl}{\let\ifpl\@firstoftwo}{\let\ifpl\@secondoftwo}%
  % \api{}, \pkg{} and \code{} are allowed in a section title, and a title
  % becomes a PDF bookmark. hyperref drops the \color command from a bookmark
  % string but keeps its ARGUMENT, so a title carrying \api{model\_validate}
  % made a bookmark reading "mtealmodel_validate" and warned "Token not
  % allowed" twice. The colour name is gobbled here instead. Measured on a
  % probe against this preamble before the line was added.
  \def\color#1{}%
}
\makeatother

% ---------- Language strings ----------
\input{lang/\booklang}

% The PDF's subject, keyword list and document language read \lblPdf... from
% the language file input on the line above, so they cannot sit in the colour
% \hypersetup block, which runs before it. Two blocks is the ordering, not an
% oversight; merging them raises Undefined control sequence and still writes
% a PDF, with the three fields empty.
\hypersetup{
  pdfsubject={\lblPdfSubject},
  pdfkeywords={\lblPdfKeywords},
  pdflang={\lblPdfLang}
}

% ---------- Headings ----------
% \raggedright on the title body is load-bearing: at \Huge on this block a
% title that cannot break overflows the margin by 30 pt or more, and several
% chapter titles here run past fifty characters.
\titleformat{\chapter}[display]
  {\normalfont\huge\bfseries\color{mblue}\raggedright}
  {\normalfont\normalsize\color{mgrey}\MakeUppercase{\chaptertitlename}\ \thechapter}
  {8pt}{\Huge\raggedright}
\titlespacing*{\chapter}{0pt}{10pt}{28pt}
\titleformat{\section}{\normalfont\Large\bfseries\color{mblue}\raggedright}{\thesection}{0.8em}{}
\titleformat{\subsection}{\normalfont\large\bfseries\color{mgrey}\raggedright}{\thesubsection}{0.7em}{}

% Sections in the contents, subsections not: fourteen chapters of subsections
% is a contents nobody reads, and the reader navigates by chapter and section.
\setcounter{tocdepth}{1}
\setcounter{secnumdepth}{2}

% One-sided running head: the chapter on the left, the folio on the right.
\pagestyle{fancy}
\fancyhf{}
\fancyhead[L]{\small\nouppercase{\leftmark}}
\fancyhead[R]{\small\thepage}
\renewcommand{\headrulewidth}{0.4pt}
\fancypagestyle{plain}{\fancyhf{}\fancyfoot[R]{\small\thepage}%
  \renewcommand{\headrulewidth}{0pt}}

% Drops the word "Chapter" from the head: the number is still there and the
% longer titles no longer overflow. MUST come after \pagestyle{fancy}, which
% installs its own \chaptermark and silently overwrites anything earlier.
\renewcommand{\chaptermark}[1]{\markboth{\thechapter.\ #1}{}}

% ============================================================
%  CODE LISTINGS
%  ---------------------------------------------------------
%  Every listing in a chapter is a FILE under code/, pulled in with \pyfile or
%  \pyregion, and CI runs it against the pinned versions. An in-source
%  \begin{python} block is allowed for a fragment of three or four lines that
%  exists to show a shape rather than to run; anything longer belongs in a
%  file the reader can open in the editor beside the PDF.
% ============================================================
\lstdefinelanguage{pythonx}{
  language=Python,
  morekeywords={
    async,await,match,case,nonlocal,yield,type,
    TypedDict,Annotated,Protocol,Literal,Generic,Final,ClassVar,Self,
    TypeVar,TypeGuard,TypeIs,Callable,Iterator,Iterable,Awaitable,
    dataclass,field,BaseModel,Field,Enum,StrEnum,NamedTuple,
    override,overload,property,staticmethod,classmethod
  },
  morekeywords=[2]{
    gather,create_task,TaskGroup,to_thread,timeout,wait_for,Semaphore,Queue,
    fixture,parametrize,monkeypatch,raises,mark,
    Depends,FastAPI,APIRouter,HTTPException,
    select,Session,Mapped,mapped_column,relationship,
    model_validate,model_validate_json,model_dump,model_json_schema
  },
  sensitive=true
}

% listings' C# definition predates async, records and pattern matching.
\lstdefinelanguage{csharpx}{
  language=[Sharp]C,
  morekeywords={
    async,await,var,dynamic,record,init,required,nameof,yield,partial,
    global,scoped,when,with,and,or,not,notnull,unmanaged,file,get,set,value,
    where,select,from,orderby,let,join,into,ascending,descending,on,equals,
    Task,ValueTask,CancellationToken,IEnumerable,IAsyncEnumerable,Func,Action
  },
  sensitive=true
}

\lstdefinestyle{bookcode}{
  basicstyle=\ttfamily\footnotesize,
  keywordstyle=\color{kw}\bfseries,
  keywordstyle=[2]\color{typ},
  stringstyle=\color{str},
  % Colour only, no \itshape: inconsolata (zi4) has no italic shape, so an
  % italic comment style is substituted with upright on every machine that
  % has the font, and the log carries "Font shape `T1/zi4/m/it' undefined"
  % on every build. The colour already marks a comment.
  commentstyle=\color{cmt},
  emphstyle=\color{dec}\bfseries,
  numbers=left,
  numberstyle=\ttfamily\tiny\color{mgrey},
  numbersep=8pt,
  backgroundcolor=\color{codebg},
  frame=single,
  rulecolor=\color{codeframe},
  framesep=6pt,
  breaklines=true,
  breakatwhitespace=false,
  postbreak=\mbox{\textcolor{mgrey}{$\hookrightarrow$}\space},
  showstringspaces=false,
  tabsize=4,
  columns=fullflexible,
  keepspaces=true,
  captionpos=b,
  aboveskip=1.2em,
  belowskip=1.2em,
  % Maps the characters most likely to be pasted into a listing by accident.
  %
  % Do NOT add a mapping for U+00A0 (non-breaking space). It looks like the
  % obvious next entry and it makes `listings` abort the run with "Improper
  % alphabetic constant" -- fatal, no PDF, and the message names neither the
  % character nor this line. The ASCII-in-listings convention, checked by
  % parity's C13, is the guard against it instead.
  literate={ł}{{\l}}1 {Ł}{{\L}}1 {ą}{{\k a}}1 {ę}{{\k e}}1
           {ś}{{\'s}}1 {ć}{{\'c}}1 {ż}{{\.z}}1 {ź}{{\'z}}1 {ó}{{\'o}}1 {ń}{{\'n}}1
           {Ą}{{\k A}}1 {Ę}{{\k E}}1 {Ś}{{\'S}}1 {Ć}{{\'C}}1
           {Ż}{{\.Z}}1 {Ź}{{\'Z}}1 {Ó}{{\'O}}1 {Ń}{{\'N}}1
           {—}{{-{}-}}2 {–}{{-}}1 {‘}{{'}}1 {’}{{'}}1
           {“}{{"}}1 {”}{{"}}1 {°}{{\textdegree}}1
}

\lstset{style=bookcode}

\lstnewenvironment{python}[1][]
  {\lstset{language=pythonx,style=bookcode,#1}}
  {}

% C# comparison snippets. The whole book is written for engineers arriving
% from .NET, so listings come in pairs often enough that the C# half has its
% own environment rather than a generic one.
\lstnewenvironment{csharp}[1][]
  {\lstset{language=csharpx,style=bookcode,#1}}
  {}

\lstnewenvironment{shellcmd}[1][]
  {\lstset{language=bash,style=bookcode,numbers=none,keywordstyle=\color{mteal}\bfseries,#1}}
  {}

\lstnewenvironment{tomlcode}[1][]
  {\lstset{style=bookcode,numbers=none,keywordstyle=\color{black},#1}}
  {}

\lstnewenvironment{yamlcode}[1][]
  {\lstset{style=bookcode,numbers=none,keywordstyle=\color{black},#1}}
  {}

\lstnewenvironment{jsoncode}[1][]
  {\lstset{style=bookcode,numbers=none,keywordstyle=\color{black},#1}}
  {}

% Terminal transcripts typed into the source. Allowed for a shell session the
% reader is meant to reproduce; NOT for program output, which is \transcript
% below, because an in-source transcript nobody ran reads exactly like one
% that was, and that is where a remembered number hides.
\lstnewenvironment{console}[1][]
  {\lstset{style=bookcode,numbers=none,basicstyle=\ttfamily\footnotesize,
           keywordstyle=\color{black},backgroundcolor=\color{white},#1}}
  {}

% The path is printed above every file listing, in the repository-relative
% form the reader opens in the editor beside the PDF. That is the whole point
% of a listing being a file, and the front matter promises it. The \nobreak
% keeps the path on the same page as the listing it names.
% ---------------------------------------------------------------------------
%  Appendix B prints the trap catalogue from notes/02-traps.md. One entry is
%  its number -- which a chapter may cite, and which is never reused -- the
%  habit in the reader's own voice, and what Python does instead.
%
%  The body is set RAGGED RIGHT on purpose. An entry is mostly \code{} spans,
%  which are unbreakable runs, and sixty-seven justified paragraphs of them
%  is sixty-seven chances at an overfull hbox in the edition with the longer
%  words. Ragged right removes the class rather than fixing it one box at a
%  time, and a reference list is the one place in this book where it reads
%  better anyway.
%
%  The habit is set with \enquote{} rather than \emph{}, and that is not a
%  style choice. Every habit here carries \code{} spans, and \code{} inside
%  \emph{} asks inconsolata for an italic it does not have: the first build
%  of this appendix raised `Font shape T1/zi4/m/it undefined', which
%  checklog.py is hard on for exactly this reason. It is the trap CLAUDE.md
%  records, met sixty-seven times at once -- and fixed once, because the
%  entries go through a macro. A quotation is also what the habit IS, and
%  \enquote{} is what the Polish edition owes anyway.
% ---------------------------------------------------------------------------
%  Appendix A's cheat sheet: one table per part, C# on the left, Python on
%  the right, and the chapter that teaches the row on the end. The three
%  header strings are arguments so that one environment serves both
%  editions and the structural signature stays identical.
%
%  Columns are ragged right for the reason Appendix B's entries are: a cell
%  is mostly \code{} spans, which are unbreakable runs, and a justified
%  narrow column full of them is an overfull box waiting for the edition
%  with the longer words.
% ---------------------------------------------------------------------------
%  Appendix C's tool matrix: the .NET tool, the Python tool this book uses,
%  the version it was verified against and the chapter that introduces it.
%
%  The version column prints \pydanticver and its siblings and is never
%  typed, which is the whole point of the appendix: it is the one page where
%  a reader sees every pin at once, and a pin that had to be re-typed here
%  would be a second copy of a fact the preamble already owns.
%  tools/check_structure.py --tools fails when a pin macro exists and this
%  table does not print it.
\NewDocumentEnvironment{toolmatrix}{mmmm}{%
  \small
  \begin{longtable}{@{}%
    >{\raggedright\arraybackslash}p{0.27\linewidth}%
    >{\raggedright\arraybackslash}p{0.31\linewidth}%
    >{\raggedright\arraybackslash}p{0.20\linewidth}%
    >{\raggedright\arraybackslash}p{0.08\linewidth}@{}}
  \toprule
  \textbf{#1} & \textbf{#2} & \textbf{#3} & \textbf{#4} \\
  \midrule
  \endfirsthead
  \toprule
  \textbf{#1} & \textbf{#2} & \textbf{#3} & \textbf{#4} \\
  \midrule
  \endhead
  \midrule
  \endfoot
  \bottomrule
  \endlastfoot
}{%
  \end{longtable}%
}

\NewDocumentEnvironment{cheatsheet}{mmm}{%
  \small
  \begin{longtable}{@{}%
    >{\raggedright\arraybackslash}p{0.41\linewidth}%
    >{\raggedright\arraybackslash}p{0.41\linewidth}%
    >{\raggedright\arraybackslash}p{0.08\linewidth}@{}}
  \toprule
  \textbf{#1} & \textbf{#2} & \textbf{#3} \\
  \midrule
  \endfirsthead
  \toprule
  \textbf{#1} & \textbf{#2} & \textbf{#3} \\
  \midrule
  \endhead
  \midrule
  \endfoot
  \bottomrule
  \endlastfoot
}{%
  \end{longtable}%
}

\newcommand{\trapentry}[3]{%
  \par\medskip
  \noindent\textbf{#1.}\enspace\enquote{#2}\par
  \nopagebreak
  {\raggedright\noindent#3\par}%
}

\newcommand{\listingpath}[1]{%
  \par\medskip\noindent
  {\ttfamily\footnotesize\color{mgrey}\detokenize{#1}}%
  \par\nobreak}

% External source file -- the preferred form. The path is relative to the
% repository root, which is also the path the reader opens in the editor.
% Usage: \pyfile{code/ch05/decorators.py}{Caption}{lst:label}
\newcommand{\pyfile}[3]{%
  \listingpath{#1}%
  \lstinputlisting[language=pythonx,style=bookcode,aboveskip=0.3em,
    caption={#2},label={#3}]{#1}%
}

% A named region of an external file, delimited in the source by
% `# --8<-- [start:name]` / `# --8<-- [end:name]`, so the book and the running
% code cannot drift apart. tools/check_structure.py --listings fails when the
% file or the region is missing.
%
% The markers are matched through listings' range prefix and suffix keys,
% with linerange={name-name}. The form this was inherited with -- a linerange
% whose two ends were the whole marker strings -- matched nothing and printed
% the WHOLE FILE, with no warning: measured on a three-region probe file
% before this was rewritten, and the same definition sits unused in the
% sibling book it came from. A region name is a word, never a number, because
% a number is a line number to listings. The line numbers printed are the
% file's own, so a region's listing says where in the file it sits.
% Usage: \pyregion{code/ch05/decorators.py}{timing}{Caption}{lst:label}
\lstset{
  rangebeginprefix=\#\ --8<--\ [start:, rangebeginsuffix=],
  rangeendprefix=\#\ --8<--\ [end:, rangeendsuffix=],
  includerangemarker=false}
\newcommand{\pyregion}[4]{%
  \listingpath{#1}%
  \lstinputlisting[language=pythonx,style=bookcode,aboveskip=0.3em,
    linerange={#2-#2},caption={#3},label={#4}]{#1}%
}

% The C# twin of \pyfile, for the comparison half of a pair.
\newcommand{\csfile}[3]{%
  \listingpath{#1}%
  \lstinputlisting[language=csharpx,style=bookcode,aboveskip=0.3em,
    caption={#2},label={#3}]{#1}%
}

% A program's output, written by a script under code/measure/ and pulled in
% whole. There is no typed form of this and there must not be: \val{} does
% not expand inside a verbatim body, so a transcript typed into a chapter is
% one that cannot quote a computed value, and a transcript nobody ran is
% indistinguishable on the page from one that was.
\newcommand{\transcript}[1]{%
  \par\smallskip
  \IfFileExists{figures/transcripts/#1.txt}{%
    \lstinputlisting[style=bookcode,numbers=none,basicstyle=\ttfamily\footnotesize,
      keywordstyle=\color{black},backgroundcolor=\color{white},
      language={}]{figures/transcripts/#1.txt}%
  }{%
    \begin{tcolorbox}[enhanced, sharp corners, width=\linewidth,
      colback=mlight, colframe=mgrey, boxrule=0.6pt,
      left=8pt, right=8pt, top=6pt, bottom=6pt]
      {\bfseries\color{mgrey}\small \lblTranscriptMissing}\\[2pt]
      {\ttfamily\scriptsize\color{mgrey} figures/transcripts/#1.txt}
    \end{tcolorbox}%
  }%
  \par\smallskip
}

% Inline literals. Underscores inside these must be written \_.
\newcommand{\code}[1]{\texttt{\small #1}}
\newcommand{\pkg}[1]{\texttt{\small\color{mblue}#1}}
\newcommand{\api}[1]{\texttt{\small\color{mteal}#1}}

% ============================================================
%  ADMONITIONS
%  ---------------------------------------------------------
%  A deliberately small, fixed vocabulary. Two visual treatments carry the
%  whole grammar: `tinted` is a block the reader may read in passing -- a
%  tint and a thick left bar, no outline; `outlined` is a block the reader
%  must stop at, with a full rule on white. A third treatment would mean the
%  first two are not carrying their meaning.
% ============================================================
\tcbset{
  mfa tinted/.style={
    enhanced, breakable, sharp corners,
    colback=mlight, boxrule=0pt, leftrule=3pt,
    left=8pt, right=8pt, top=6pt, bottom=6pt,
    before skip=13pt, after skip=13pt,
    fonttitle=\bfseries\small,
    coltitle=black, colbacktitle=mlight,
    attach title to upper={\par\smallskip},
  },
  mfa outlined/.style={
    enhanced, breakable, sharp corners,
    colback=white, boxrule=0.6pt,
    left=8pt, right=8pt, top=6pt, bottom=6pt,
    before skip=13pt, after skip=13pt,
    fonttitle=\bfseries\small,
    coltitle=black, colbacktitle=white,
    attach title to upper={\par\smallskip},
  },
}

% Read in passing.
\newtcolorbox{note}[1][]{mfa tinted, colframe=mblue, title={\lblNote}, #1}
\newtcolorbox{warning}[1][]{mfa tinted, colframe=mred, title={\lblWarning}, #1}
% The translation. The .NET reader is the book's only assumed reader, so this
% box carries the C# side of a comparison and nothing else. Its use is
% counted: a chapter with none is a chapter that has stopped translating.
\newtcolorbox{csbox}[1][]{mfa tinted, colframe=mteal, title={\lblCsharp}, #1}
% Anything that is preview, experimental or likely to move. Python moves more
% slowly than the agent frameworks do, and this box should be rare here.
\newtcolorbox{versionbox}[1][]{mfa tinted, colframe=mamber, title={\lblVersion}, #1}

% Stop and read.
% The misconception, named AFTER the reader has walked into it. Appendix B is
% the catalogue; a trap box in a chapter is one entry of it, in place.
\newtcolorbox{trapbox}[1][]{mfa outlined, colframe=mred, title={\lblTrap}, #1}
% Marks a listing that has not been executed against the pinned versions.
% `make debt` counts these, CI prints the count, and nothing ships to a reader
% with one attached. Removing it means you ran the code.
\newtcolorbox{verifybox}[1][]{mfa outlined, colframe=mgrey, title={\lblVerify}, #1}
\newtcolorbox{exercisebox}[1][]{mfa outlined, colframe=mteal, title={\lblExercise}, #1}
% The guiding project. Points at the stage of trace-assert this section
% builds, under code/src/trace_assert/.
\newtcolorbox{projectbox}[1][]{mfa outlined, colframe=mamber, title={\lblProject}, #1}

% ============================================================
%  EXERCISES
%  ---------------------------------------------------------
%  \begin{exercise}{key}{title} ... \end{exercise}
%
%  This is the mechanism the whole book is read through: the PDF on one half
%  of the screen, an editor on the other. An exercise is therefore not a
%  paragraph but a FILE the reader opens -- a starter under
%  code/exercises/chNN/<key>.py -- and a TEST that decides whether it is done,
%  test_<key>.py beside it. The box prints both paths and the one command
%  that runs the test, and the manifest at the back lists every exercise.
%
%  An ENVIRONMENT, not a macro with a body argument: a macro argument cannot
%  contain verbatim material, so an exercise written as \exercise{key}{title}
%  {body} could never show the reader a code fragment, and the first chapter
%  that tried would have died on `\begin{python}` inside the braces.
%
%  The key is the file stem, with underscores, and it is written with
%  \detokenize so the underscore reaches the page as a character rather than
%  as a subscript. tools/check_structure.py --exercises fails when the starter
%  or the test is missing, or when the key's chapter prefix disagrees with the
%  chapter it sits in.
% ============================================================
\newcounter{exercise}[chapter]
\renewcommand{\theexercise}{\thechapter.\arabic{exercise}}
\makeatletter
\newcommand{\exercisedir}{ch\two@digits{\value{chapter}}}
\newcommand{\listofexercises}{%
  \section*{\lblExerciseManifest}%
  \phantomsection\addcontentsline{toc}{section}{\lblExerciseManifest}%
  % The entries are subsection-level and tocdepth is 1 for the contents,
  % so the depth is raised for this list alone or it prints empty --
  % which it did, on the scaffold's first clean build.
  \begingroup\c@tocdepth=2\relax\@starttoc{exr}\endgroup%
}
\makeatother
\NewDocumentEnvironment{exercise}{mm}{%
  \refstepcounter{exercise}%
  \addcontentsline{exr}{subsection}{\protect\numberline{\theexercise}\texttt{\protect\detokenize{#1}} --- #2}%
  \begin{exercisebox}[title={\lblExercise~\theexercise\ \dash{} #2}]%
}{%
  \par\medskip
  {\small\setlength{\parindent}{0pt}%
  \lblExerciseStarter: \texttt{code/exercises/\exercisedir/\detokenize{#1}.py}\\
  \lblExerciseTest: \texttt{code/exercises/\exercisedir/test\_\detokenize{#1}.py}\\
  \lblExerciseRun: \texttt{cd code \&\& uv run pytest -k \detokenize{#1}}\par}%
  \end{exercisebox}%
}

% ============================================================
%  DIAGRAMS --- Mermaid pipeline
%  ---------------------------------------------------------
%  \mermaidfig{key}{caption}{one-line manifest description}
%
%  Source of truth is figures/mermaid/<lang>/<key>.mmd, committed. `make
%  diagrams` renders each to figures/diagrams/<lang>/<key>.pdf, which is build
%  output and is not committed, so a diagram change reviews as a text diff.
%  The source is per-language because a diagram with English labels in the
%  Polish edition is exactly the half-translation that makes a bilingual book
%  feel machine-made; CI checks that both languages carry the same keys.
%
%  If the rendered PDF is absent the macro draws a placeholder AND typesets
%  the Mermaid source, in the flow rather than as a float, because a source
%  typeset verbatim is often taller than a page and a float cannot break.
% ============================================================
\makeatletter
\newcommand{\listofdiagrams}{%
  \section*{\lblDiagramManifest}%
  \phantomsection\addcontentsline{toc}{section}{\lblDiagramManifest}%
  % The entries are subsection-level and tocdepth is 1 for the contents,
  % so the depth is raised for this list alone or it prints empty --
  % which it did, on the scaffold's first clean build.
  \begingroup\c@tocdepth=2\relax\@starttoc{dgm}\endgroup%
}
\makeatother

\newcommand{\mermaidfig}[3]{%
  % \phantomsection is load-bearing and not decoration. \addcontentsline
  % records hyperref's CURRENT anchor, and here that is whatever preceded
  % the figure -- for a figure placed after a listing it is the listing's
  % line-number anchor, and listings has already advanced the counter past
  % the last line it printed, so the manifest entry linked to a
  % destination that is never created: "name{lstnumber.10.4.46} has been
  % referenced but does not exist, replaced by a fixed one". The two
  % front-matter figures had the same wrong target and were silent about
  % it, because the lines they named happen to be printed. The exercise
  % manifest never had the bug, because \begin{exercise} steps a counter
  % and that is what sets the anchor.
  \phantomsection%
  \addcontentsline{dgm}{subsection}{\protect\numberline{}\texttt{#1.mmd} --- #3}%
  \IfFileExists{figures/diagrams/\booklang/#1.pdf}{%
    \begin{figure}[htbp]
    \centering
    \includegraphics[width=0.95\linewidth,height=0.42\textheight,keepaspectratio]{figures/diagrams/\booklang/#1.pdf}%
    \caption{#2}\label{fig:#1}
    \end{figure}%
  }{%
    \par\medskip
    \begin{tcolorbox}[
      enhanced, breakable, width=\linewidth, sharp corners,
      colback=white, colframe=mgrey, boxrule=0.8pt,
      left=8pt, right=8pt, top=8pt, bottom=8pt]
      {\bfseries\color{mgrey}\small \lblNotRendered}\\[4pt]
      {\ttfamily\scriptsize\color{mgrey} figures/mermaid/\booklang/#1.mmd}\\[6pt]
      \IfFileExists{figures/mermaid/\booklang/#1.mmd}{%
        \lstinputlisting[style=bookcode,numbers=none,basicstyle=\ttfamily\scriptsize,
                         backgroundcolor=\color{mlight},frame=none,
                         keywordstyle=\color{black},language={}]{figures/mermaid/\booklang/#1.mmd}%
      }{%
        {\small\color{mred} \lblSourceMissing: \texttt{figures/mermaid/\booklang/#1.mmd}}%
      }%
    \end{tcolorbox}%
    \captionof{figure}{#2}\label{fig:#1}
    \par\medskip
  }%
}

% ============================================================
%  CHAPTER STUBS
%  ---------------------------------------------------------
%  Prints a visible marker so an unwritten chapter cannot be mistaken for a
%  written one and the page count never flatters the draft. `make debt`
%  counts them, and CI publishes the count on every build.
% ============================================================
\newcommand{\chapterstub}[1]{%
  \begin{tcolorbox}[
    enhanced, breakable, width=\linewidth, sharp corners,
    colback=mlight, colframe=mred, boxrule=1pt,
    borderline={0.6pt}{3pt}{mred, dashed},
    left=12pt, right=12pt, top=12pt, bottom=12pt]
    {\bfseries\color{mred}\large \lblNotWritten}\\[8pt]
    \small\raggedright #1
  \end{tcolorbox}%
}

% ============================================================
%  COMPUTED VALUES
%  ---------------------------------------------------------
%  A number that came out of a measurement is written by a script under
%  code/measure/ into figures/values/<name>.tex as
%
%      \pyval{e01.threads.speedup}{3.71}
%
%  and pulled into the text with \val{e01.threads.speedup}. The script is the
%  source of truth and the book cannot disagree with it, because the book does
%  not contain the digits. `make verify` fails when a committed value no
%  longer matches its script. A value the scripts do not produce prints a
%  visible marker.
%
%  \val also applies the edition's decimal separator, which is how the Polish
%  edition gets its decimal comma without a translator having to remember.
% ============================================================
\newcommand{\bookdecimalmarker}{\ifpl{,}{.}}
\IfFileExists{siunitx.sty}{%
  \usepackage{siunitx}%
  \sisetup{
    group-digits = integer,
    group-minimum-digits = 5,
    group-separator = {\,},
    output-decimal-marker = {\bookdecimalmarker}
  }%
  \newcommand{\booknum}[1]{\num{##1}}%
}{%
  \newcommand{\booknum}[1]{##1}%
}

\newcommand{\pyvalues}{}
\makeatletter
\newcommand{\pyval}[2]{\csgdef{pyval@#1}{#2}\listxadd{\pyvalues}{#1}}
% \num on a non-numeric body is a FATAL siunitx error. The generator knows
% whether it emitted a number or a word, so it writes \pyval for a number and
% \pyvaltext for a word, and this end only enforces which macro reads which.
\newcommand{\val}[1]{%
  \ifcsdef{pyval@#1}{%
    \ifcsdef{pyvaltext@#1}{%
      \PackageError{pybook}{Value `#1' is not a number}%
        {\string\val\space passes its body to siunitx, which refuses anything
         that is not a number. Use \string\valtext{#1} instead.}%
    }{}%
    \booknum{\csuse{pyval@#1}}%
  }{\textcolor{mred}{\textbf{[\lblValueMissing: #1]}}}%
}
\newcommand{\valtext}[1]{%
  \ifcsdef{pyval@#1}{\csuse{pyval@#1}}%
    {\textcolor{mred}{\textbf{[\lblValueMissing: #1]}}}%
}
\newcommand{\pyvaltext}[2]{%
  \csgdef{pyval@#1}{#2}\csgdef{pyvaltext@#1}{}\listxadd{\pyvalues}{#1}%
}
\makeatother

% figures/values/all.tex is generated by `make numbers` and does nothing but
% \input each value file. A build with no values yet still compiles; every
% \val{} in it prints a visible marker instead of a number.
\IfFileExists{figures/values/all.tex}{\input{figures/values/all}}{%
  \AtBeginDocument{\PackageWarning{pybook}{No computed values found. Run `make numbers`.}}}

% ============================================================
%  PINNED VERSIONS --- single source of truth
%  Change these here; they propagate through both editions and Appendix C.
%  Verified against pypi.org on \pinnedon (lang/<lang>.tex).
%  tools/check_versions.py compares every one of them with PyPI on each build.
%  Python itself is not on PyPI and is checked by hand against python.org.
% ============================================================
\newcommand{\bookedition}{0.1-dev}
\newcommand{\bookyear}{2026}
\newcommand{\pyver}{3.14}              % the minor the book targets
\newcommand{\pypatch}{3.14.7}          % the exact interpreter the code ran on
\newcommand{\uvver}{0.12.13}           % uv
\newcommand{\dotnetver}{10.0.100}      % dotnet
\newcommand{\ruffver}{0.16.7}          % ruff
\newcommand{\pyrightver}{1.1.414}      % pyright
\newcommand{\mypyver}{2.3.1}           % mypy
\newcommand{\pydanticver}{2.13.5}      % pydantic
\newcommand{\pysettingsver}{2.15.0}    % pydantic-settings
\newcommand{\fastapiver}{0.141.1}      % fastapi
\newcommand{\uvicornver}{0.53.0}       % uvicorn
\newcommand{\sqlalchemyver}{2.0.52}    % sqlalchemy
\newcommand{\alembicver}{1.20.0}       % alembic
\newcommand{\aiosqlitever}{0.22.1}     % aiosqlite
\newcommand{\pytestver}{9.1.1}         % pytest
\newcommand{\pytestasyncver}{1.4.0}    % pytest-asyncio
\newcommand{\hypothesisver}{6.168.0}   % hypothesis
\newcommand{\coveragever}{7.16.1}      % coverage
\newcommand{\httpxver}{0.28.1}         % httpx
\newcommand{\structlogver}{26.1.0}     % structlog
\newcommand{\otelver}{1.44.0}          % opentelemetry-sdk
\newcommand{\polarsver}{1.44.2}        % polars
\newcommand{\anthropicver}{1.5.0}      % anthropic
\newcommand{\openaiver}{3.13.0}        % openai
\newcommand{\pipauditver}{2.10.1}      % pip-audit

% ============================================================
%  INDEX
%  Two-column, and its entries include \texttt{} names that do not hyphenate.
%  Ragged right keeps multi-word entries off the margin; it cannot help a
%  single token wider than the column, so keep verbatim index entries under
%  about 29 characters and index longer names by concept instead.
%  \raggedcolumns is the switch imakeidx's multicols actually reads.
% ============================================================
\apptocmd{\theindex}{\raggedcolumns\raggedright}{}{%
  \PackageWarning{pybook}{Could not patch theindex for ragged setting}}
\providecommand*\see[2]{}
\renewcommand*\see[2]{\emph{\lblIndexSee} #1}
\providecommand*\seealso[2]{}
\renewcommand*\seealso[2]{\emph{\lblIndexSeeAlso} #1}

\makeindex
 still compiles and says what it did.
\providecommand{\booklang}{en}

\usepackage{etoolbox}   % loaded first: the language switch below needs it

% ---------- Page geometry ----------
% ONE format: A4, 12pt, ONE-SIDED. This book is read on a screen with an
% editor open beside it, and never printed as a bound book, so there is no
% recto, no verso, no gutter, no mirrored margin and no blank leaf before a
% chapter -- every one of those is a property of paper that a PDF viewer at
% half a screen's width turns into wasted pixels.
%
% The margins are symmetric and the measure is about seventy-five characters
% of prose, which is where the eye keeps the line return. The number that was
% actually chosen is the MONOSPACE measure: at \footnotesize on this block a
% listing line of 79 characters fits without wrapping, and 79 is the width
% every listing under code/ is held to by tools/check_structure.py --lines.
% A wrapped listing line is not an error TeX reports -- breaklines=true wraps
% it silently and prints an arrow into the middle of what the reader is meant
% to paste into an editor -- so the width is enforced in the source instead.
\usepackage[
  a4paper,
  left=3.1cm, right=3.1cm, top=2.6cm, bottom=3.0cm,
  head=26pt, headsep=0.8cm, footskip=1.4cm
]{geometry}

% ---------- Encoding & fonts ----------
\usepackage[T1]{fontenc}
\usepackage[utf8]{inputenc}
\IfFileExists{lmodern.sty}{\usepackage{lmodern}}{}
% newtxtext takes its TS1 symbols (\textcopyright, \textdegree) from TeX
% Gyre Termes, which Debian packages SEPARATELY (tex-gyre): a machine with
% newtx and without it dies on the copyright page with
% `Font TS1/ntxtlf/m/n/10.95=ts1-qtmr ... not loadable`. There is no
% inert-safe probe for a font metric in pdfTeX -- \IfFileExists searches
% TEXINPUTS and not TFMFONTS, so it reports the metric missing on a machine
% that has it, and \suppressfontnotfounderror is LuaTeX-only -- so nothing
% here guesses. CI's full TeX Live is the reference; locally, install
% tex-gyre. Both facts are recorded in CLAUDE.md under Build traps.
\IfFileExists{newtxtext.sty}{\usepackage{newtxtext}}{}
% amsmath before newtxmath, because newtxmath patches amsmath's internals.
% amssymb only when newtxmath is absent: newtxmath supplies the AMS symbols
% itself and loading both is a fatal clash (\Bbbk already defined) that is
% invisible on a machine without newtx and fatal on a full TeX Live.
\usepackage{amsmath}
\IfFileExists{newtxmath.sty}{\usepackage{newtxmath}}{\usepackage{amssymb}}
\IfFileExists{inconsolata.sty}{\usepackage[varqu,scaled=0.95]{inconsolata}}{}
\IfFileExists{lmodern.sty}{\usepackage{microtype}}{\usepackage[expansion=false]{microtype}}

% upquote is not a refinement. In T1 the input character ' is quoteright, so
% a listings body sets f('x') with U+2019 at both ends of the literal, and
% typed back in that is a SyntaxError. inconsolata's varqu option hides the
% defect on a machine that has inconsolata; upquote makes ' the ASCII
% apostrophe whatever the typewriter font is. In a book whose every listing is
% meant to be pasted into an editor, this line is load-bearing.
\IfFileExists{upquote.sty}{\usepackage{upquote}}{}

% And upquote's twin, for the same reason one character over. In T1 the pair
% -- is a ligature for an en dash, in EVERY family including the typewriter
% one, so \code{uv sync --locked} sets as "uv sync <en dash>locked" and a
% reader who copies it off the page gets an unknown-argument error from a
% command the chapter told them to run. It is invisible in review because it
% looks like a dash, and it bites exactly the flags a packaging chapter is
% made of. A `listings` body is already safe -- listings sets characters
% singly and no ligature forms -- which is what makes the defect so easy to
% miss: the same text is right inside a listing and wrong in the sentence
% above it.
%
% Scoped to tt* on purpose, and measured rather than assumed: prose keeps its
% own dashes, so \dash{} still sets an em dash and a range like 1--2 still
% sets an en dash. Probed against this preamble in all three positions before
% it was believed.
\DisableLigatures[-]{encoding = T1, family = tt*}

% ---------- Language ----------
% Loading babel with a language whose .ldf is absent is a *fatal* error, not
% a warning, so both the package and each language file are probed before
% either is requested. Both editions load the other language too, because the
% glossary rows and the cheat sheet carry words from both.
\IfFileExists{babel.sty}{%
  \IfFileExists{polish.ldf}{%
    \IfFileExists{british.ldf}{%
      \ifdefstring{\booklang}{pl}%
        {\usepackage[british,main=polish]{babel}}%
        {\usepackage[polish,main=british]{babel}}%
    }{\usepackage[polish]{babel}}%
  }{%
    \IfFileExists{british.ldf}{\usepackage[british]{babel}}{}%
  }%
}{}

% babel-polish makes " an active shorthand. listings reads its body verbatim
% and an active character in a verbatim context is a category-code accident
% waiting to happen, so the shorthand is turned off. Polish quotation marks
% come from \enquote, which csquotes sets per language.
\ifdefstring{\booklang}{pl}{\AtBeginDocument{\shorthandoff{"}}}{}

% ---------- Core packages ----------
\usepackage{graphicx}
\usepackage{xcolor}
\usepackage{booktabs}
\usepackage{tabularx}
% Appendix A's cheat sheet runs past a page in every part, and tabularx
% cannot break. longtable can, and works with booktabs.
\usepackage{longtable}
\usepackage{array}
\usepackage{enumitem}
\usepackage{listings}
\usepackage[most]{tcolorbox}
\usepackage{fancyhdr}
\usepackage{titlesec}
\usepackage{caption}
\usepackage{imakeidx}
% csquotes with autostyle follows babel's active language, so the identical
% source \enquote{x} sets 'x' in the English edition and the Polish quotation
% marks in the Polish one.
\IfFileExists{csquotes.sty}{\usepackage[autostyle=true]{csquotes}}{%
  \providecommand{\enquote}[1]{``##1''}}
\usepackage[hidelinks]{hyperref}

% ---------- Palette ----------
% The same family as the two companion volumes, so the three read as a set.
\definecolor{mblue}{HTML}{1F4E79}
\definecolor{mteal}{HTML}{0E7C7B}
\definecolor{mamber}{HTML}{B26A00}
\definecolor{mred}{HTML}{9B2C2C}
\definecolor{mviolet}{HTML}{7B4397}
\definecolor{mgrey}{HTML}{5A5A5A}
\definecolor{mlight}{HTML}{F4F6F8}
\definecolor{mfaint}{HTML}{FBFCFD}
\definecolor{codebg}{HTML}{FAFAFA}
\definecolor{codeframe}{HTML}{D8DEE4}
\definecolor{kw}{HTML}{0B5394}
\definecolor{str}{HTML}{9B2C2C}
\definecolor{cmt}{HTML}{4F7A4F}
\definecolor{typ}{HTML}{0E7C7B}
\definecolor{dec}{HTML}{7B4397}

\hypersetup{
  colorlinks=true,
  linkcolor=mblue,
  urlcolor=mteal,
  citecolor=mblue,
  pdfauthor={Konrad Cinkusz}
}

% \ifpl{Polish text}{English text} -- for the handful of places where a
% string is structural rather than content. Anything longer than a few words
% belongs in lang/ or in the edition's own source, not here.
% \DeclareRobustCommand, not \newcommand: a fragile \ifpl inside a moving
% argument is written to the .toc one level expanded and fails on the next run
% in exactly one of the two editions.
\DeclareRobustCommand{\ifpl}[2]{\ifdefstring{\booklang}{pl}{#1}{#2}}
\makeatletter
\pdfstringdefDisableCommands{%
  \ifdefstring{\booklang}{pl}{\let\ifpl\@firstoftwo}{\let\ifpl\@secondoftwo}%
  % \api{}, \pkg{} and \code{} are allowed in a section title, and a title
  % becomes a PDF bookmark. hyperref drops the \color command from a bookmark
  % string but keeps its ARGUMENT, so a title carrying \api{model\_validate}
  % made a bookmark reading "mtealmodel_validate" and warned "Token not
  % allowed" twice. The colour name is gobbled here instead. Measured on a
  % probe against this preamble before the line was added.
  \def\color#1{}%
}
\makeatother

% ---------- Language strings ----------
\input{lang/\booklang}

% The PDF's subject, keyword list and document language read \lblPdf... from
% the language file input on the line above, so they cannot sit in the colour
% \hypersetup block, which runs before it. Two blocks is the ordering, not an
% oversight; merging them raises Undefined control sequence and still writes
% a PDF, with the three fields empty.
\hypersetup{
  pdfsubject={\lblPdfSubject},
  pdfkeywords={\lblPdfKeywords},
  pdflang={\lblPdfLang}
}

% ---------- Headings ----------
% \raggedright on the title body is load-bearing: at \Huge on this block a
% title that cannot break overflows the margin by 30 pt or more, and several
% chapter titles here run past fifty characters.
\titleformat{\chapter}[display]
  {\normalfont\huge\bfseries\color{mblue}\raggedright}
  {\normalfont\normalsize\color{mgrey}\MakeUppercase{\chaptertitlename}\ \thechapter}
  {8pt}{\Huge\raggedright}
\titlespacing*{\chapter}{0pt}{10pt}{28pt}
\titleformat{\section}{\normalfont\Large\bfseries\color{mblue}\raggedright}{\thesection}{0.8em}{}
\titleformat{\subsection}{\normalfont\large\bfseries\color{mgrey}\raggedright}{\thesubsection}{0.7em}{}

% Sections in the contents, subsections not: fourteen chapters of subsections
% is a contents nobody reads, and the reader navigates by chapter and section.
\setcounter{tocdepth}{1}
\setcounter{secnumdepth}{2}

% One-sided running head: the chapter on the left, the folio on the right.
\pagestyle{fancy}
\fancyhf{}
\fancyhead[L]{\small\nouppercase{\leftmark}}
\fancyhead[R]{\small\thepage}
\renewcommand{\headrulewidth}{0.4pt}
\fancypagestyle{plain}{\fancyhf{}\fancyfoot[R]{\small\thepage}%
  \renewcommand{\headrulewidth}{0pt}}

% Drops the word "Chapter" from the head: the number is still there and the
% longer titles no longer overflow. MUST come after \pagestyle{fancy}, which
% installs its own \chaptermark and silently overwrites anything earlier.
\renewcommand{\chaptermark}[1]{\markboth{\thechapter.\ #1}{}}

% ============================================================
%  CODE LISTINGS
%  ---------------------------------------------------------
%  Every listing in a chapter is a FILE under code/, pulled in with \pyfile or
%  \pyregion, and CI runs it against the pinned versions. An in-source
%  \begin{python} block is allowed for a fragment of three or four lines that
%  exists to show a shape rather than to run; anything longer belongs in a
%  file the reader can open in the editor beside the PDF.
% ============================================================
\lstdefinelanguage{pythonx}{
  language=Python,
  morekeywords={
    async,await,match,case,nonlocal,yield,type,
    TypedDict,Annotated,Protocol,Literal,Generic,Final,ClassVar,Self,
    TypeVar,TypeGuard,TypeIs,Callable,Iterator,Iterable,Awaitable,
    dataclass,field,BaseModel,Field,Enum,StrEnum,NamedTuple,
    override,overload,property,staticmethod,classmethod
  },
  morekeywords=[2]{
    gather,create_task,TaskGroup,to_thread,timeout,wait_for,Semaphore,Queue,
    fixture,parametrize,monkeypatch,raises,mark,
    Depends,FastAPI,APIRouter,HTTPException,
    select,Session,Mapped,mapped_column,relationship,
    model_validate,model_validate_json,model_dump,model_json_schema
  },
  sensitive=true
}

% listings' C# definition predates async, records and pattern matching.
\lstdefinelanguage{csharpx}{
  language=[Sharp]C,
  morekeywords={
    async,await,var,dynamic,record,init,required,nameof,yield,partial,
    global,scoped,when,with,and,or,not,notnull,unmanaged,file,get,set,value,
    where,select,from,orderby,let,join,into,ascending,descending,on,equals,
    Task,ValueTask,CancellationToken,IEnumerable,IAsyncEnumerable,Func,Action
  },
  sensitive=true
}

\lstdefinestyle{bookcode}{
  basicstyle=\ttfamily\footnotesize,
  keywordstyle=\color{kw}\bfseries,
  keywordstyle=[2]\color{typ},
  stringstyle=\color{str},
  % Colour only, no \itshape: inconsolata (zi4) has no italic shape, so an
  % italic comment style is substituted with upright on every machine that
  % has the font, and the log carries "Font shape `T1/zi4/m/it' undefined"
  % on every build. The colour already marks a comment.
  commentstyle=\color{cmt},
  emphstyle=\color{dec}\bfseries,
  numbers=left,
  numberstyle=\ttfamily\tiny\color{mgrey},
  numbersep=8pt,
  backgroundcolor=\color{codebg},
  frame=single,
  rulecolor=\color{codeframe},
  framesep=6pt,
  breaklines=true,
  breakatwhitespace=false,
  postbreak=\mbox{\textcolor{mgrey}{$\hookrightarrow$}\space},
  showstringspaces=false,
  tabsize=4,
  columns=fullflexible,
  keepspaces=true,
  captionpos=b,
  aboveskip=1.2em,
  belowskip=1.2em,
  % Maps the characters most likely to be pasted into a listing by accident.
  %
  % Do NOT add a mapping for U+00A0 (non-breaking space). It looks like the
  % obvious next entry and it makes `listings` abort the run with "Improper
  % alphabetic constant" -- fatal, no PDF, and the message names neither the
  % character nor this line. The ASCII-in-listings convention, checked by
  % parity's C13, is the guard against it instead.
  literate={ł}{{\l}}1 {Ł}{{\L}}1 {ą}{{\k a}}1 {ę}{{\k e}}1
           {ś}{{\'s}}1 {ć}{{\'c}}1 {ż}{{\.z}}1 {ź}{{\'z}}1 {ó}{{\'o}}1 {ń}{{\'n}}1
           {Ą}{{\k A}}1 {Ę}{{\k E}}1 {Ś}{{\'S}}1 {Ć}{{\'C}}1
           {Ż}{{\.Z}}1 {Ź}{{\'Z}}1 {Ó}{{\'O}}1 {Ń}{{\'N}}1
           {—}{{-{}-}}2 {–}{{-}}1 {‘}{{'}}1 {’}{{'}}1
           {“}{{"}}1 {”}{{"}}1 {°}{{\textdegree}}1
}

\lstset{style=bookcode}

\lstnewenvironment{python}[1][]
  {\lstset{language=pythonx,style=bookcode,#1}}
  {}

% C# comparison snippets. The whole book is written for engineers arriving
% from .NET, so listings come in pairs often enough that the C# half has its
% own environment rather than a generic one.
\lstnewenvironment{csharp}[1][]
  {\lstset{language=csharpx,style=bookcode,#1}}
  {}

\lstnewenvironment{shellcmd}[1][]
  {\lstset{language=bash,style=bookcode,numbers=none,keywordstyle=\color{mteal}\bfseries,#1}}
  {}

\lstnewenvironment{tomlcode}[1][]
  {\lstset{style=bookcode,numbers=none,keywordstyle=\color{black},#1}}
  {}

\lstnewenvironment{yamlcode}[1][]
  {\lstset{style=bookcode,numbers=none,keywordstyle=\color{black},#1}}
  {}

\lstnewenvironment{jsoncode}[1][]
  {\lstset{style=bookcode,numbers=none,keywordstyle=\color{black},#1}}
  {}

% Terminal transcripts typed into the source. Allowed for a shell session the
% reader is meant to reproduce; NOT for program output, which is \transcript
% below, because an in-source transcript nobody ran reads exactly like one
% that was, and that is where a remembered number hides.
\lstnewenvironment{console}[1][]
  {\lstset{style=bookcode,numbers=none,basicstyle=\ttfamily\footnotesize,
           keywordstyle=\color{black},backgroundcolor=\color{white},#1}}
  {}

% The path is printed above every file listing, in the repository-relative
% form the reader opens in the editor beside the PDF. That is the whole point
% of a listing being a file, and the front matter promises it. The \nobreak
% keeps the path on the same page as the listing it names.
% ---------------------------------------------------------------------------
%  Appendix B prints the trap catalogue from notes/02-traps.md. One entry is
%  its number -- which a chapter may cite, and which is never reused -- the
%  habit in the reader's own voice, and what Python does instead.
%
%  The body is set RAGGED RIGHT on purpose. An entry is mostly \code{} spans,
%  which are unbreakable runs, and sixty-seven justified paragraphs of them
%  is sixty-seven chances at an overfull hbox in the edition with the longer
%  words. Ragged right removes the class rather than fixing it one box at a
%  time, and a reference list is the one place in this book where it reads
%  better anyway.
%
%  The habit is set with \enquote{} rather than \emph{}, and that is not a
%  style choice. Every habit here carries \code{} spans, and \code{} inside
%  \emph{} asks inconsolata for an italic it does not have: the first build
%  of this appendix raised `Font shape T1/zi4/m/it undefined', which
%  checklog.py is hard on for exactly this reason. It is the trap CLAUDE.md
%  records, met sixty-seven times at once -- and fixed once, because the
%  entries go through a macro. A quotation is also what the habit IS, and
%  \enquote{} is what the Polish edition owes anyway.
% ---------------------------------------------------------------------------
%  Appendix A's cheat sheet: one table per part, C# on the left, Python on
%  the right, and the chapter that teaches the row on the end. The three
%  header strings are arguments so that one environment serves both
%  editions and the structural signature stays identical.
%
%  Columns are ragged right for the reason Appendix B's entries are: a cell
%  is mostly \code{} spans, which are unbreakable runs, and a justified
%  narrow column full of them is an overfull box waiting for the edition
%  with the longer words.
% ---------------------------------------------------------------------------
%  Appendix C's tool matrix: the .NET tool, the Python tool this book uses,
%  the version it was verified against and the chapter that introduces it.
%
%  The version column prints \pydanticver and its siblings and is never
%  typed, which is the whole point of the appendix: it is the one page where
%  a reader sees every pin at once, and a pin that had to be re-typed here
%  would be a second copy of a fact the preamble already owns.
%  tools/check_structure.py --tools fails when a pin macro exists and this
%  table does not print it.
\NewDocumentEnvironment{toolmatrix}{mmmm}{%
  \small
  \begin{longtable}{@{}%
    >{\raggedright\arraybackslash}p{0.27\linewidth}%
    >{\raggedright\arraybackslash}p{0.31\linewidth}%
    >{\raggedright\arraybackslash}p{0.20\linewidth}%
    >{\raggedright\arraybackslash}p{0.08\linewidth}@{}}
  \toprule
  \textbf{#1} & \textbf{#2} & \textbf{#3} & \textbf{#4} \\
  \midrule
  \endfirsthead
  \toprule
  \textbf{#1} & \textbf{#2} & \textbf{#3} & \textbf{#4} \\
  \midrule
  \endhead
  \midrule
  \endfoot
  \bottomrule
  \endlastfoot
}{%
  \end{longtable}%
}

\NewDocumentEnvironment{cheatsheet}{mmm}{%
  \small
  \begin{longtable}{@{}%
    >{\raggedright\arraybackslash}p{0.41\linewidth}%
    >{\raggedright\arraybackslash}p{0.41\linewidth}%
    >{\raggedright\arraybackslash}p{0.08\linewidth}@{}}
  \toprule
  \textbf{#1} & \textbf{#2} & \textbf{#3} \\
  \midrule
  \endfirsthead
  \toprule
  \textbf{#1} & \textbf{#2} & \textbf{#3} \\
  \midrule
  \endhead
  \midrule
  \endfoot
  \bottomrule
  \endlastfoot
}{%
  \end{longtable}%
}

\newcommand{\trapentry}[3]{%
  \par\medskip
  \noindent\textbf{#1.}\enspace\enquote{#2}\par
  \nopagebreak
  {\raggedright\noindent#3\par}%
}

\newcommand{\listingpath}[1]{%
  \par\medskip\noindent
  {\ttfamily\footnotesize\color{mgrey}\detokenize{#1}}%
  \par\nobreak}

% External source file -- the preferred form. The path is relative to the
% repository root, which is also the path the reader opens in the editor.
% Usage: \pyfile{code/ch05/decorators.py}{Caption}{lst:label}
\newcommand{\pyfile}[3]{%
  \listingpath{#1}%
  \lstinputlisting[language=pythonx,style=bookcode,aboveskip=0.3em,
    caption={#2},label={#3}]{#1}%
}

% A named region of an external file, delimited in the source by
% `# --8<-- [start:name]` / `# --8<-- [end:name]`, so the book and the running
% code cannot drift apart. tools/check_structure.py --listings fails when the
% file or the region is missing.
%
% The markers are matched through listings' range prefix and suffix keys,
% with linerange={name-name}. The form this was inherited with -- a linerange
% whose two ends were the whole marker strings -- matched nothing and printed
% the WHOLE FILE, with no warning: measured on a three-region probe file
% before this was rewritten, and the same definition sits unused in the
% sibling book it came from. A region name is a word, never a number, because
% a number is a line number to listings. The line numbers printed are the
% file's own, so a region's listing says where in the file it sits.
% Usage: \pyregion{code/ch05/decorators.py}{timing}{Caption}{lst:label}
\lstset{
  rangebeginprefix=\#\ --8<--\ [start:, rangebeginsuffix=],
  rangeendprefix=\#\ --8<--\ [end:, rangeendsuffix=],
  includerangemarker=false}
\newcommand{\pyregion}[4]{%
  \listingpath{#1}%
  \lstinputlisting[language=pythonx,style=bookcode,aboveskip=0.3em,
    linerange={#2-#2},caption={#3},label={#4}]{#1}%
}

% The C# twin of \pyfile, for the comparison half of a pair.
\newcommand{\csfile}[3]{%
  \listingpath{#1}%
  \lstinputlisting[language=csharpx,style=bookcode,aboveskip=0.3em,
    caption={#2},label={#3}]{#1}%
}

% A program's output, written by a script under code/measure/ and pulled in
% whole. There is no typed form of this and there must not be: \val{} does
% not expand inside a verbatim body, so a transcript typed into a chapter is
% one that cannot quote a computed value, and a transcript nobody ran is
% indistinguishable on the page from one that was.
\newcommand{\transcript}[1]{%
  \par\smallskip
  \IfFileExists{figures/transcripts/#1.txt}{%
    \lstinputlisting[style=bookcode,numbers=none,basicstyle=\ttfamily\footnotesize,
      keywordstyle=\color{black},backgroundcolor=\color{white},
      language={}]{figures/transcripts/#1.txt}%
  }{%
    \begin{tcolorbox}[enhanced, sharp corners, width=\linewidth,
      colback=mlight, colframe=mgrey, boxrule=0.6pt,
      left=8pt, right=8pt, top=6pt, bottom=6pt]
      {\bfseries\color{mgrey}\small \lblTranscriptMissing}\\[2pt]
      {\ttfamily\scriptsize\color{mgrey} figures/transcripts/#1.txt}
    \end{tcolorbox}%
  }%
  \par\smallskip
}

% Inline literals. Underscores inside these must be written \_.
\newcommand{\code}[1]{\texttt{\small #1}}
\newcommand{\pkg}[1]{\texttt{\small\color{mblue}#1}}
\newcommand{\api}[1]{\texttt{\small\color{mteal}#1}}

% ============================================================
%  ADMONITIONS
%  ---------------------------------------------------------
%  A deliberately small, fixed vocabulary. Two visual treatments carry the
%  whole grammar: `tinted` is a block the reader may read in passing -- a
%  tint and a thick left bar, no outline; `outlined` is a block the reader
%  must stop at, with a full rule on white. A third treatment would mean the
%  first two are not carrying their meaning.
% ============================================================
\tcbset{
  mfa tinted/.style={
    enhanced, breakable, sharp corners,
    colback=mlight, boxrule=0pt, leftrule=3pt,
    left=8pt, right=8pt, top=6pt, bottom=6pt,
    before skip=13pt, after skip=13pt,
    fonttitle=\bfseries\small,
    coltitle=black, colbacktitle=mlight,
    attach title to upper={\par\smallskip},
  },
  mfa outlined/.style={
    enhanced, breakable, sharp corners,
    colback=white, boxrule=0.6pt,
    left=8pt, right=8pt, top=6pt, bottom=6pt,
    before skip=13pt, after skip=13pt,
    fonttitle=\bfseries\small,
    coltitle=black, colbacktitle=white,
    attach title to upper={\par\smallskip},
  },
}

% Read in passing.
\newtcolorbox{note}[1][]{mfa tinted, colframe=mblue, title={\lblNote}, #1}
\newtcolorbox{warning}[1][]{mfa tinted, colframe=mred, title={\lblWarning}, #1}
% The translation. The .NET reader is the book's only assumed reader, so this
% box carries the C# side of a comparison and nothing else. Its use is
% counted: a chapter with none is a chapter that has stopped translating.
\newtcolorbox{csbox}[1][]{mfa tinted, colframe=mteal, title={\lblCsharp}, #1}
% Anything that is preview, experimental or likely to move. Python moves more
% slowly than the agent frameworks do, and this box should be rare here.
\newtcolorbox{versionbox}[1][]{mfa tinted, colframe=mamber, title={\lblVersion}, #1}

% Stop and read.
% The misconception, named AFTER the reader has walked into it. Appendix B is
% the catalogue; a trap box in a chapter is one entry of it, in place.
\newtcolorbox{trapbox}[1][]{mfa outlined, colframe=mred, title={\lblTrap}, #1}
% Marks a listing that has not been executed against the pinned versions.
% `make debt` counts these, CI prints the count, and nothing ships to a reader
% with one attached. Removing it means you ran the code.
\newtcolorbox{verifybox}[1][]{mfa outlined, colframe=mgrey, title={\lblVerify}, #1}
\newtcolorbox{exercisebox}[1][]{mfa outlined, colframe=mteal, title={\lblExercise}, #1}
% The guiding project. Points at the stage of trace-assert this section
% builds, under code/src/trace_assert/.
\newtcolorbox{projectbox}[1][]{mfa outlined, colframe=mamber, title={\lblProject}, #1}

% ============================================================
%  EXERCISES
%  ---------------------------------------------------------
%  \begin{exercise}{key}{title} ... \end{exercise}
%
%  This is the mechanism the whole book is read through: the PDF on one half
%  of the screen, an editor on the other. An exercise is therefore not a
%  paragraph but a FILE the reader opens -- a starter under
%  code/exercises/chNN/<key>.py -- and a TEST that decides whether it is done,
%  test_<key>.py beside it. The box prints both paths and the one command
%  that runs the test, and the manifest at the back lists every exercise.
%
%  An ENVIRONMENT, not a macro with a body argument: a macro argument cannot
%  contain verbatim material, so an exercise written as \exercise{key}{title}
%  {body} could never show the reader a code fragment, and the first chapter
%  that tried would have died on `\begin{python}` inside the braces.
%
%  The key is the file stem, with underscores, and it is written with
%  \detokenize so the underscore reaches the page as a character rather than
%  as a subscript. tools/check_structure.py --exercises fails when the starter
%  or the test is missing, or when the key's chapter prefix disagrees with the
%  chapter it sits in.
% ============================================================
\newcounter{exercise}[chapter]
\renewcommand{\theexercise}{\thechapter.\arabic{exercise}}
\makeatletter
\newcommand{\exercisedir}{ch\two@digits{\value{chapter}}}
\newcommand{\listofexercises}{%
  \section*{\lblExerciseManifest}%
  \phantomsection\addcontentsline{toc}{section}{\lblExerciseManifest}%
  % The entries are subsection-level and tocdepth is 1 for the contents,
  % so the depth is raised for this list alone or it prints empty --
  % which it did, on the scaffold's first clean build.
  \begingroup\c@tocdepth=2\relax\@starttoc{exr}\endgroup%
}
\makeatother
\NewDocumentEnvironment{exercise}{mm}{%
  \refstepcounter{exercise}%
  \addcontentsline{exr}{subsection}{\protect\numberline{\theexercise}\texttt{\protect\detokenize{#1}} --- #2}%
  \begin{exercisebox}[title={\lblExercise~\theexercise\ \dash{} #2}]%
}{%
  \par\medskip
  {\small\setlength{\parindent}{0pt}%
  \lblExerciseStarter: \texttt{code/exercises/\exercisedir/\detokenize{#1}.py}\\
  \lblExerciseTest: \texttt{code/exercises/\exercisedir/test\_\detokenize{#1}.py}\\
  \lblExerciseRun: \texttt{cd code \&\& uv run pytest -k \detokenize{#1}}\par}%
  \end{exercisebox}%
}

% ============================================================
%  DIAGRAMS --- Mermaid pipeline
%  ---------------------------------------------------------
%  \mermaidfig{key}{caption}{one-line manifest description}
%
%  Source of truth is figures/mermaid/<lang>/<key>.mmd, committed. `make
%  diagrams` renders each to figures/diagrams/<lang>/<key>.pdf, which is build
%  output and is not committed, so a diagram change reviews as a text diff.
%  The source is per-language because a diagram with English labels in the
%  Polish edition is exactly the half-translation that makes a bilingual book
%  feel machine-made; CI checks that both languages carry the same keys.
%
%  If the rendered PDF is absent the macro draws a placeholder AND typesets
%  the Mermaid source, in the flow rather than as a float, because a source
%  typeset verbatim is often taller than a page and a float cannot break.
% ============================================================
\makeatletter
\newcommand{\listofdiagrams}{%
  \section*{\lblDiagramManifest}%
  \phantomsection\addcontentsline{toc}{section}{\lblDiagramManifest}%
  % The entries are subsection-level and tocdepth is 1 for the contents,
  % so the depth is raised for this list alone or it prints empty --
  % which it did, on the scaffold's first clean build.
  \begingroup\c@tocdepth=2\relax\@starttoc{dgm}\endgroup%
}
\makeatother

\newcommand{\mermaidfig}[3]{%
  % \phantomsection is load-bearing and not decoration. \addcontentsline
  % records hyperref's CURRENT anchor, and here that is whatever preceded
  % the figure -- for a figure placed after a listing it is the listing's
  % line-number anchor, and listings has already advanced the counter past
  % the last line it printed, so the manifest entry linked to a
  % destination that is never created: "name{lstnumber.10.4.46} has been
  % referenced but does not exist, replaced by a fixed one". The two
  % front-matter figures had the same wrong target and were silent about
  % it, because the lines they named happen to be printed. The exercise
  % manifest never had the bug, because \begin{exercise} steps a counter
  % and that is what sets the anchor.
  \phantomsection%
  \addcontentsline{dgm}{subsection}{\protect\numberline{}\texttt{#1.mmd} --- #3}%
  \IfFileExists{figures/diagrams/\booklang/#1.pdf}{%
    \begin{figure}[htbp]
    \centering
    \includegraphics[width=0.95\linewidth,height=0.42\textheight,keepaspectratio]{figures/diagrams/\booklang/#1.pdf}%
    \caption{#2}\label{fig:#1}
    \end{figure}%
  }{%
    \par\medskip
    \begin{tcolorbox}[
      enhanced, breakable, width=\linewidth, sharp corners,
      colback=white, colframe=mgrey, boxrule=0.8pt,
      left=8pt, right=8pt, top=8pt, bottom=8pt]
      {\bfseries\color{mgrey}\small \lblNotRendered}\\[4pt]
      {\ttfamily\scriptsize\color{mgrey} figures/mermaid/\booklang/#1.mmd}\\[6pt]
      \IfFileExists{figures/mermaid/\booklang/#1.mmd}{%
        \lstinputlisting[style=bookcode,numbers=none,basicstyle=\ttfamily\scriptsize,
                         backgroundcolor=\color{mlight},frame=none,
                         keywordstyle=\color{black},language={}]{figures/mermaid/\booklang/#1.mmd}%
      }{%
        {\small\color{mred} \lblSourceMissing: \texttt{figures/mermaid/\booklang/#1.mmd}}%
      }%
    \end{tcolorbox}%
    \captionof{figure}{#2}\label{fig:#1}
    \par\medskip
  }%
}

% ============================================================
%  CHAPTER STUBS
%  ---------------------------------------------------------
%  Prints a visible marker so an unwritten chapter cannot be mistaken for a
%  written one and the page count never flatters the draft. `make debt`
%  counts them, and CI publishes the count on every build.
% ============================================================
\newcommand{\chapterstub}[1]{%
  \begin{tcolorbox}[
    enhanced, breakable, width=\linewidth, sharp corners,
    colback=mlight, colframe=mred, boxrule=1pt,
    borderline={0.6pt}{3pt}{mred, dashed},
    left=12pt, right=12pt, top=12pt, bottom=12pt]
    {\bfseries\color{mred}\large \lblNotWritten}\\[8pt]
    \small\raggedright #1
  \end{tcolorbox}%
}

% ============================================================
%  COMPUTED VALUES
%  ---------------------------------------------------------
%  A number that came out of a measurement is written by a script under
%  code/measure/ into figures/values/<name>.tex as
%
%      \pyval{e01.threads.speedup}{3.71}
%
%  and pulled into the text with \val{e01.threads.speedup}. The script is the
%  source of truth and the book cannot disagree with it, because the book does
%  not contain the digits. `make verify` fails when a committed value no
%  longer matches its script. A value the scripts do not produce prints a
%  visible marker.
%
%  \val also applies the edition's decimal separator, which is how the Polish
%  edition gets its decimal comma without a translator having to remember.
% ============================================================
\newcommand{\bookdecimalmarker}{\ifpl{,}{.}}
\IfFileExists{siunitx.sty}{%
  \usepackage{siunitx}%
  \sisetup{
    group-digits = integer,
    group-minimum-digits = 5,
    group-separator = {\,},
    output-decimal-marker = {\bookdecimalmarker}
  }%
  \newcommand{\booknum}[1]{\num{##1}}%
}{%
  \newcommand{\booknum}[1]{##1}%
}

\newcommand{\pyvalues}{}
\makeatletter
\newcommand{\pyval}[2]{\csgdef{pyval@#1}{#2}\listxadd{\pyvalues}{#1}}
% \num on a non-numeric body is a FATAL siunitx error. The generator knows
% whether it emitted a number or a word, so it writes \pyval for a number and
% \pyvaltext for a word, and this end only enforces which macro reads which.
\newcommand{\val}[1]{%
  \ifcsdef{pyval@#1}{%
    \ifcsdef{pyvaltext@#1}{%
      \PackageError{pybook}{Value `#1' is not a number}%
        {\string\val\space passes its body to siunitx, which refuses anything
         that is not a number. Use \string\valtext{#1} instead.}%
    }{}%
    \booknum{\csuse{pyval@#1}}%
  }{\textcolor{mred}{\textbf{[\lblValueMissing: #1]}}}%
}
\newcommand{\valtext}[1]{%
  \ifcsdef{pyval@#1}{\csuse{pyval@#1}}%
    {\textcolor{mred}{\textbf{[\lblValueMissing: #1]}}}%
}
\newcommand{\pyvaltext}[2]{%
  \csgdef{pyval@#1}{#2}\csgdef{pyvaltext@#1}{}\listxadd{\pyvalues}{#1}%
}
\makeatother

% figures/values/all.tex is generated by `make numbers` and does nothing but
% \input each value file. A build with no values yet still compiles; every
% \val{} in it prints a visible marker instead of a number.
\IfFileExists{figures/values/all.tex}{\input{figures/values/all}}{%
  \AtBeginDocument{\PackageWarning{pybook}{No computed values found. Run `make numbers`.}}}

% ============================================================
%  PINNED VERSIONS --- single source of truth
%  Change these here; they propagate through both editions and Appendix C.
%  Verified against pypi.org on \pinnedon (lang/<lang>.tex).
%  tools/check_versions.py compares every one of them with PyPI on each build.
%  Python itself is not on PyPI and is checked by hand against python.org.
% ============================================================
\newcommand{\bookedition}{0.1-dev}
\newcommand{\bookyear}{2026}
\newcommand{\pyver}{3.14}              % the minor the book targets
\newcommand{\pypatch}{3.14.7}          % the exact interpreter the code ran on
\newcommand{\uvver}{0.12.13}           % uv
\newcommand{\dotnetver}{10.0.100}      % dotnet
\newcommand{\ruffver}{0.16.7}          % ruff
\newcommand{\pyrightver}{1.1.414}      % pyright
\newcommand{\mypyver}{2.3.1}           % mypy
\newcommand{\pydanticver}{2.13.5}      % pydantic
\newcommand{\pysettingsver}{2.15.0}    % pydantic-settings
\newcommand{\fastapiver}{0.141.1}      % fastapi
\newcommand{\uvicornver}{0.53.0}       % uvicorn
\newcommand{\sqlalchemyver}{2.0.52}    % sqlalchemy
\newcommand{\alembicver}{1.20.0}       % alembic
\newcommand{\aiosqlitever}{0.22.1}     % aiosqlite
\newcommand{\pytestver}{9.1.1}         % pytest
\newcommand{\pytestasyncver}{1.4.0}    % pytest-asyncio
\newcommand{\hypothesisver}{6.168.0}   % hypothesis
\newcommand{\coveragever}{7.16.1}      % coverage
\newcommand{\httpxver}{0.28.1}         % httpx
\newcommand{\structlogver}{26.1.0}     % structlog
\newcommand{\otelver}{1.44.0}          % opentelemetry-sdk
\newcommand{\polarsver}{1.44.2}        % polars
\newcommand{\anthropicver}{1.5.0}      % anthropic
\newcommand{\openaiver}{3.13.0}        % openai
\newcommand{\pipauditver}{2.10.1}      % pip-audit

% ============================================================
%  INDEX
%  Two-column, and its entries include \texttt{} names that do not hyphenate.
%  Ragged right keeps multi-word entries off the margin; it cannot help a
%  single token wider than the column, so keep verbatim index entries under
%  about 29 characters and index longer names by concept instead.
%  \raggedcolumns is the switch imakeidx's multicols actually reads.
% ============================================================
\apptocmd{\theindex}{\raggedcolumns\raggedright}{}{%
  \PackageWarning{pybook}{Could not patch theindex for ragged setting}}
\providecommand*\see[2]{}
\renewcommand*\see[2]{\emph{\lblIndexSee} #1}
\providecommand*\seealso[2]{}
\renewcommand*\seealso[2]{\emph{\lblIndexSeeAlso} #1}

\makeindex
; nothing else differs between them at this level.
%  Every user-visible string lives in lang/en.tex or lang/pl.tex, so a label
%  that appears in one edition and not the other is a missing macro rather
%  than a silent divergence.
%
%  Lineage. The bilingual wiring (\booklang, lang/, body.tex, the generated
%  structure) is the math book's; the chapter machinery (listings, the
%  admonitions, verifybox, \pyfile and \pyregion, the Mermaid pipeline) is the
%  LangChain book's. What is new here is the exercise macro, which ties every
%  exercise to a starter file and a test under code/, because this book is
%  read with an editor open beside it.
% ============================================================

% \booklang must be set by the main file. Default to English rather than
% failing, so a stray % ============================================================
%  preamble.tex --- styling, environments, macros
%
%  Shared by BOTH editions. main-en.tex and main-pl.tex each define \booklang
%  before % ============================================================
%  preamble.tex --- styling, environments, macros
%
%  Shared by BOTH editions. main-en.tex and main-pl.tex each define \booklang
%  before % ============================================================
%  preamble.tex --- styling, environments, macros
%
%  Shared by BOTH editions. main-en.tex and main-pl.tex each define \booklang
%  before \input{preamble}; nothing else differs between them at this level.
%  Every user-visible string lives in lang/en.tex or lang/pl.tex, so a label
%  that appears in one edition and not the other is a missing macro rather
%  than a silent divergence.
%
%  Lineage. The bilingual wiring (\booklang, lang/, body.tex, the generated
%  structure) is the math book's; the chapter machinery (listings, the
%  admonitions, verifybox, \pyfile and \pyregion, the Mermaid pipeline) is the
%  LangChain book's. What is new here is the exercise macro, which ties every
%  exercise to a starter file and a test under code/, because this book is
%  read with an editor open beside it.
% ============================================================

% \booklang must be set by the main file. Default to English rather than
% failing, so a stray \input{preamble} still compiles and says what it did.
\providecommand{\booklang}{en}

\usepackage{etoolbox}   % loaded first: the language switch below needs it

% ---------- Page geometry ----------
% ONE format: A4, 12pt, ONE-SIDED. This book is read on a screen with an
% editor open beside it, and never printed as a bound book, so there is no
% recto, no verso, no gutter, no mirrored margin and no blank leaf before a
% chapter -- every one of those is a property of paper that a PDF viewer at
% half a screen's width turns into wasted pixels.
%
% The margins are symmetric and the measure is about seventy-five characters
% of prose, which is where the eye keeps the line return. The number that was
% actually chosen is the MONOSPACE measure: at \footnotesize on this block a
% listing line of 79 characters fits without wrapping, and 79 is the width
% every listing under code/ is held to by tools/check_structure.py --lines.
% A wrapped listing line is not an error TeX reports -- breaklines=true wraps
% it silently and prints an arrow into the middle of what the reader is meant
% to paste into an editor -- so the width is enforced in the source instead.
\usepackage[
  a4paper,
  left=3.1cm, right=3.1cm, top=2.6cm, bottom=3.0cm,
  head=26pt, headsep=0.8cm, footskip=1.4cm
]{geometry}

% ---------- Encoding & fonts ----------
\usepackage[T1]{fontenc}
\usepackage[utf8]{inputenc}
\IfFileExists{lmodern.sty}{\usepackage{lmodern}}{}
% newtxtext takes its TS1 symbols (\textcopyright, \textdegree) from TeX
% Gyre Termes, which Debian packages SEPARATELY (tex-gyre): a machine with
% newtx and without it dies on the copyright page with
% `Font TS1/ntxtlf/m/n/10.95=ts1-qtmr ... not loadable`. There is no
% inert-safe probe for a font metric in pdfTeX -- \IfFileExists searches
% TEXINPUTS and not TFMFONTS, so it reports the metric missing on a machine
% that has it, and \suppressfontnotfounderror is LuaTeX-only -- so nothing
% here guesses. CI's full TeX Live is the reference; locally, install
% tex-gyre. Both facts are recorded in CLAUDE.md under Build traps.
\IfFileExists{newtxtext.sty}{\usepackage{newtxtext}}{}
% amsmath before newtxmath, because newtxmath patches amsmath's internals.
% amssymb only when newtxmath is absent: newtxmath supplies the AMS symbols
% itself and loading both is a fatal clash (\Bbbk already defined) that is
% invisible on a machine without newtx and fatal on a full TeX Live.
\usepackage{amsmath}
\IfFileExists{newtxmath.sty}{\usepackage{newtxmath}}{\usepackage{amssymb}}
\IfFileExists{inconsolata.sty}{\usepackage[varqu,scaled=0.95]{inconsolata}}{}
\IfFileExists{lmodern.sty}{\usepackage{microtype}}{\usepackage[expansion=false]{microtype}}

% upquote is not a refinement. In T1 the input character ' is quoteright, so
% a listings body sets f('x') with U+2019 at both ends of the literal, and
% typed back in that is a SyntaxError. inconsolata's varqu option hides the
% defect on a machine that has inconsolata; upquote makes ' the ASCII
% apostrophe whatever the typewriter font is. In a book whose every listing is
% meant to be pasted into an editor, this line is load-bearing.
\IfFileExists{upquote.sty}{\usepackage{upquote}}{}

% And upquote's twin, for the same reason one character over. In T1 the pair
% -- is a ligature for an en dash, in EVERY family including the typewriter
% one, so \code{uv sync --locked} sets as "uv sync <en dash>locked" and a
% reader who copies it off the page gets an unknown-argument error from a
% command the chapter told them to run. It is invisible in review because it
% looks like a dash, and it bites exactly the flags a packaging chapter is
% made of. A `listings` body is already safe -- listings sets characters
% singly and no ligature forms -- which is what makes the defect so easy to
% miss: the same text is right inside a listing and wrong in the sentence
% above it.
%
% Scoped to tt* on purpose, and measured rather than assumed: prose keeps its
% own dashes, so \dash{} still sets an em dash and a range like 1--2 still
% sets an en dash. Probed against this preamble in all three positions before
% it was believed.
\DisableLigatures[-]{encoding = T1, family = tt*}

% ---------- Language ----------
% Loading babel with a language whose .ldf is absent is a *fatal* error, not
% a warning, so both the package and each language file are probed before
% either is requested. Both editions load the other language too, because the
% glossary rows and the cheat sheet carry words from both.
\IfFileExists{babel.sty}{%
  \IfFileExists{polish.ldf}{%
    \IfFileExists{british.ldf}{%
      \ifdefstring{\booklang}{pl}%
        {\usepackage[british,main=polish]{babel}}%
        {\usepackage[polish,main=british]{babel}}%
    }{\usepackage[polish]{babel}}%
  }{%
    \IfFileExists{british.ldf}{\usepackage[british]{babel}}{}%
  }%
}{}

% babel-polish makes " an active shorthand. listings reads its body verbatim
% and an active character in a verbatim context is a category-code accident
% waiting to happen, so the shorthand is turned off. Polish quotation marks
% come from \enquote, which csquotes sets per language.
\ifdefstring{\booklang}{pl}{\AtBeginDocument{\shorthandoff{"}}}{}

% ---------- Core packages ----------
\usepackage{graphicx}
\usepackage{xcolor}
\usepackage{booktabs}
\usepackage{tabularx}
% Appendix A's cheat sheet runs past a page in every part, and tabularx
% cannot break. longtable can, and works with booktabs.
\usepackage{longtable}
\usepackage{array}
\usepackage{enumitem}
\usepackage{listings}
\usepackage[most]{tcolorbox}
\usepackage{fancyhdr}
\usepackage{titlesec}
\usepackage{caption}
\usepackage{imakeidx}
% csquotes with autostyle follows babel's active language, so the identical
% source \enquote{x} sets 'x' in the English edition and the Polish quotation
% marks in the Polish one.
\IfFileExists{csquotes.sty}{\usepackage[autostyle=true]{csquotes}}{%
  \providecommand{\enquote}[1]{``##1''}}
\usepackage[hidelinks]{hyperref}

% ---------- Palette ----------
% The same family as the two companion volumes, so the three read as a set.
\definecolor{mblue}{HTML}{1F4E79}
\definecolor{mteal}{HTML}{0E7C7B}
\definecolor{mamber}{HTML}{B26A00}
\definecolor{mred}{HTML}{9B2C2C}
\definecolor{mviolet}{HTML}{7B4397}
\definecolor{mgrey}{HTML}{5A5A5A}
\definecolor{mlight}{HTML}{F4F6F8}
\definecolor{mfaint}{HTML}{FBFCFD}
\definecolor{codebg}{HTML}{FAFAFA}
\definecolor{codeframe}{HTML}{D8DEE4}
\definecolor{kw}{HTML}{0B5394}
\definecolor{str}{HTML}{9B2C2C}
\definecolor{cmt}{HTML}{4F7A4F}
\definecolor{typ}{HTML}{0E7C7B}
\definecolor{dec}{HTML}{7B4397}

\hypersetup{
  colorlinks=true,
  linkcolor=mblue,
  urlcolor=mteal,
  citecolor=mblue,
  pdfauthor={Konrad Cinkusz}
}

% \ifpl{Polish text}{English text} -- for the handful of places where a
% string is structural rather than content. Anything longer than a few words
% belongs in lang/ or in the edition's own source, not here.
% \DeclareRobustCommand, not \newcommand: a fragile \ifpl inside a moving
% argument is written to the .toc one level expanded and fails on the next run
% in exactly one of the two editions.
\DeclareRobustCommand{\ifpl}[2]{\ifdefstring{\booklang}{pl}{#1}{#2}}
\makeatletter
\pdfstringdefDisableCommands{%
  \ifdefstring{\booklang}{pl}{\let\ifpl\@firstoftwo}{\let\ifpl\@secondoftwo}%
  % \api{}, \pkg{} and \code{} are allowed in a section title, and a title
  % becomes a PDF bookmark. hyperref drops the \color command from a bookmark
  % string but keeps its ARGUMENT, so a title carrying \api{model\_validate}
  % made a bookmark reading "mtealmodel_validate" and warned "Token not
  % allowed" twice. The colour name is gobbled here instead. Measured on a
  % probe against this preamble before the line was added.
  \def\color#1{}%
}
\makeatother

% ---------- Language strings ----------
\input{lang/\booklang}

% The PDF's subject, keyword list and document language read \lblPdf... from
% the language file input on the line above, so they cannot sit in the colour
% \hypersetup block, which runs before it. Two blocks is the ordering, not an
% oversight; merging them raises Undefined control sequence and still writes
% a PDF, with the three fields empty.
\hypersetup{
  pdfsubject={\lblPdfSubject},
  pdfkeywords={\lblPdfKeywords},
  pdflang={\lblPdfLang}
}

% ---------- Headings ----------
% \raggedright on the title body is load-bearing: at \Huge on this block a
% title that cannot break overflows the margin by 30 pt or more, and several
% chapter titles here run past fifty characters.
\titleformat{\chapter}[display]
  {\normalfont\huge\bfseries\color{mblue}\raggedright}
  {\normalfont\normalsize\color{mgrey}\MakeUppercase{\chaptertitlename}\ \thechapter}
  {8pt}{\Huge\raggedright}
\titlespacing*{\chapter}{0pt}{10pt}{28pt}
\titleformat{\section}{\normalfont\Large\bfseries\color{mblue}\raggedright}{\thesection}{0.8em}{}
\titleformat{\subsection}{\normalfont\large\bfseries\color{mgrey}\raggedright}{\thesubsection}{0.7em}{}

% Sections in the contents, subsections not: fourteen chapters of subsections
% is a contents nobody reads, and the reader navigates by chapter and section.
\setcounter{tocdepth}{1}
\setcounter{secnumdepth}{2}

% One-sided running head: the chapter on the left, the folio on the right.
\pagestyle{fancy}
\fancyhf{}
\fancyhead[L]{\small\nouppercase{\leftmark}}
\fancyhead[R]{\small\thepage}
\renewcommand{\headrulewidth}{0.4pt}
\fancypagestyle{plain}{\fancyhf{}\fancyfoot[R]{\small\thepage}%
  \renewcommand{\headrulewidth}{0pt}}

% Drops the word "Chapter" from the head: the number is still there and the
% longer titles no longer overflow. MUST come after \pagestyle{fancy}, which
% installs its own \chaptermark and silently overwrites anything earlier.
\renewcommand{\chaptermark}[1]{\markboth{\thechapter.\ #1}{}}

% ============================================================
%  CODE LISTINGS
%  ---------------------------------------------------------
%  Every listing in a chapter is a FILE under code/, pulled in with \pyfile or
%  \pyregion, and CI runs it against the pinned versions. An in-source
%  \begin{python} block is allowed for a fragment of three or four lines that
%  exists to show a shape rather than to run; anything longer belongs in a
%  file the reader can open in the editor beside the PDF.
% ============================================================
\lstdefinelanguage{pythonx}{
  language=Python,
  morekeywords={
    async,await,match,case,nonlocal,yield,type,
    TypedDict,Annotated,Protocol,Literal,Generic,Final,ClassVar,Self,
    TypeVar,TypeGuard,TypeIs,Callable,Iterator,Iterable,Awaitable,
    dataclass,field,BaseModel,Field,Enum,StrEnum,NamedTuple,
    override,overload,property,staticmethod,classmethod
  },
  morekeywords=[2]{
    gather,create_task,TaskGroup,to_thread,timeout,wait_for,Semaphore,Queue,
    fixture,parametrize,monkeypatch,raises,mark,
    Depends,FastAPI,APIRouter,HTTPException,
    select,Session,Mapped,mapped_column,relationship,
    model_validate,model_validate_json,model_dump,model_json_schema
  },
  sensitive=true
}

% listings' C# definition predates async, records and pattern matching.
\lstdefinelanguage{csharpx}{
  language=[Sharp]C,
  morekeywords={
    async,await,var,dynamic,record,init,required,nameof,yield,partial,
    global,scoped,when,with,and,or,not,notnull,unmanaged,file,get,set,value,
    where,select,from,orderby,let,join,into,ascending,descending,on,equals,
    Task,ValueTask,CancellationToken,IEnumerable,IAsyncEnumerable,Func,Action
  },
  sensitive=true
}

\lstdefinestyle{bookcode}{
  basicstyle=\ttfamily\footnotesize,
  keywordstyle=\color{kw}\bfseries,
  keywordstyle=[2]\color{typ},
  stringstyle=\color{str},
  % Colour only, no \itshape: inconsolata (zi4) has no italic shape, so an
  % italic comment style is substituted with upright on every machine that
  % has the font, and the log carries "Font shape `T1/zi4/m/it' undefined"
  % on every build. The colour already marks a comment.
  commentstyle=\color{cmt},
  emphstyle=\color{dec}\bfseries,
  numbers=left,
  numberstyle=\ttfamily\tiny\color{mgrey},
  numbersep=8pt,
  backgroundcolor=\color{codebg},
  frame=single,
  rulecolor=\color{codeframe},
  framesep=6pt,
  breaklines=true,
  breakatwhitespace=false,
  postbreak=\mbox{\textcolor{mgrey}{$\hookrightarrow$}\space},
  showstringspaces=false,
  tabsize=4,
  columns=fullflexible,
  keepspaces=true,
  captionpos=b,
  aboveskip=1.2em,
  belowskip=1.2em,
  % Maps the characters most likely to be pasted into a listing by accident.
  %
  % Do NOT add a mapping for U+00A0 (non-breaking space). It looks like the
  % obvious next entry and it makes `listings` abort the run with "Improper
  % alphabetic constant" -- fatal, no PDF, and the message names neither the
  % character nor this line. The ASCII-in-listings convention, checked by
  % parity's C13, is the guard against it instead.
  literate={ł}{{\l}}1 {Ł}{{\L}}1 {ą}{{\k a}}1 {ę}{{\k e}}1
           {ś}{{\'s}}1 {ć}{{\'c}}1 {ż}{{\.z}}1 {ź}{{\'z}}1 {ó}{{\'o}}1 {ń}{{\'n}}1
           {Ą}{{\k A}}1 {Ę}{{\k E}}1 {Ś}{{\'S}}1 {Ć}{{\'C}}1
           {Ż}{{\.Z}}1 {Ź}{{\'Z}}1 {Ó}{{\'O}}1 {Ń}{{\'N}}1
           {—}{{-{}-}}2 {–}{{-}}1 {‘}{{'}}1 {’}{{'}}1
           {“}{{"}}1 {”}{{"}}1 {°}{{\textdegree}}1
}

\lstset{style=bookcode}

\lstnewenvironment{python}[1][]
  {\lstset{language=pythonx,style=bookcode,#1}}
  {}

% C# comparison snippets. The whole book is written for engineers arriving
% from .NET, so listings come in pairs often enough that the C# half has its
% own environment rather than a generic one.
\lstnewenvironment{csharp}[1][]
  {\lstset{language=csharpx,style=bookcode,#1}}
  {}

\lstnewenvironment{shellcmd}[1][]
  {\lstset{language=bash,style=bookcode,numbers=none,keywordstyle=\color{mteal}\bfseries,#1}}
  {}

\lstnewenvironment{tomlcode}[1][]
  {\lstset{style=bookcode,numbers=none,keywordstyle=\color{black},#1}}
  {}

\lstnewenvironment{yamlcode}[1][]
  {\lstset{style=bookcode,numbers=none,keywordstyle=\color{black},#1}}
  {}

\lstnewenvironment{jsoncode}[1][]
  {\lstset{style=bookcode,numbers=none,keywordstyle=\color{black},#1}}
  {}

% Terminal transcripts typed into the source. Allowed for a shell session the
% reader is meant to reproduce; NOT for program output, which is \transcript
% below, because an in-source transcript nobody ran reads exactly like one
% that was, and that is where a remembered number hides.
\lstnewenvironment{console}[1][]
  {\lstset{style=bookcode,numbers=none,basicstyle=\ttfamily\footnotesize,
           keywordstyle=\color{black},backgroundcolor=\color{white},#1}}
  {}

% The path is printed above every file listing, in the repository-relative
% form the reader opens in the editor beside the PDF. That is the whole point
% of a listing being a file, and the front matter promises it. The \nobreak
% keeps the path on the same page as the listing it names.
% ---------------------------------------------------------------------------
%  Appendix B prints the trap catalogue from notes/02-traps.md. One entry is
%  its number -- which a chapter may cite, and which is never reused -- the
%  habit in the reader's own voice, and what Python does instead.
%
%  The body is set RAGGED RIGHT on purpose. An entry is mostly \code{} spans,
%  which are unbreakable runs, and sixty-seven justified paragraphs of them
%  is sixty-seven chances at an overfull hbox in the edition with the longer
%  words. Ragged right removes the class rather than fixing it one box at a
%  time, and a reference list is the one place in this book where it reads
%  better anyway.
%
%  The habit is set with \enquote{} rather than \emph{}, and that is not a
%  style choice. Every habit here carries \code{} spans, and \code{} inside
%  \emph{} asks inconsolata for an italic it does not have: the first build
%  of this appendix raised `Font shape T1/zi4/m/it undefined', which
%  checklog.py is hard on for exactly this reason. It is the trap CLAUDE.md
%  records, met sixty-seven times at once -- and fixed once, because the
%  entries go through a macro. A quotation is also what the habit IS, and
%  \enquote{} is what the Polish edition owes anyway.
% ---------------------------------------------------------------------------
%  Appendix A's cheat sheet: one table per part, C# on the left, Python on
%  the right, and the chapter that teaches the row on the end. The three
%  header strings are arguments so that one environment serves both
%  editions and the structural signature stays identical.
%
%  Columns are ragged right for the reason Appendix B's entries are: a cell
%  is mostly \code{} spans, which are unbreakable runs, and a justified
%  narrow column full of them is an overfull box waiting for the edition
%  with the longer words.
% ---------------------------------------------------------------------------
%  Appendix C's tool matrix: the .NET tool, the Python tool this book uses,
%  the version it was verified against and the chapter that introduces it.
%
%  The version column prints \pydanticver and its siblings and is never
%  typed, which is the whole point of the appendix: it is the one page where
%  a reader sees every pin at once, and a pin that had to be re-typed here
%  would be a second copy of a fact the preamble already owns.
%  tools/check_structure.py --tools fails when a pin macro exists and this
%  table does not print it.
\NewDocumentEnvironment{toolmatrix}{mmmm}{%
  \small
  \begin{longtable}{@{}%
    >{\raggedright\arraybackslash}p{0.27\linewidth}%
    >{\raggedright\arraybackslash}p{0.31\linewidth}%
    >{\raggedright\arraybackslash}p{0.20\linewidth}%
    >{\raggedright\arraybackslash}p{0.08\linewidth}@{}}
  \toprule
  \textbf{#1} & \textbf{#2} & \textbf{#3} & \textbf{#4} \\
  \midrule
  \endfirsthead
  \toprule
  \textbf{#1} & \textbf{#2} & \textbf{#3} & \textbf{#4} \\
  \midrule
  \endhead
  \midrule
  \endfoot
  \bottomrule
  \endlastfoot
}{%
  \end{longtable}%
}

\NewDocumentEnvironment{cheatsheet}{mmm}{%
  \small
  \begin{longtable}{@{}%
    >{\raggedright\arraybackslash}p{0.41\linewidth}%
    >{\raggedright\arraybackslash}p{0.41\linewidth}%
    >{\raggedright\arraybackslash}p{0.08\linewidth}@{}}
  \toprule
  \textbf{#1} & \textbf{#2} & \textbf{#3} \\
  \midrule
  \endfirsthead
  \toprule
  \textbf{#1} & \textbf{#2} & \textbf{#3} \\
  \midrule
  \endhead
  \midrule
  \endfoot
  \bottomrule
  \endlastfoot
}{%
  \end{longtable}%
}

\newcommand{\trapentry}[3]{%
  \par\medskip
  \noindent\textbf{#1.}\enspace\enquote{#2}\par
  \nopagebreak
  {\raggedright\noindent#3\par}%
}

\newcommand{\listingpath}[1]{%
  \par\medskip\noindent
  {\ttfamily\footnotesize\color{mgrey}\detokenize{#1}}%
  \par\nobreak}

% External source file -- the preferred form. The path is relative to the
% repository root, which is also the path the reader opens in the editor.
% Usage: \pyfile{code/ch05/decorators.py}{Caption}{lst:label}
\newcommand{\pyfile}[3]{%
  \listingpath{#1}%
  \lstinputlisting[language=pythonx,style=bookcode,aboveskip=0.3em,
    caption={#2},label={#3}]{#1}%
}

% A named region of an external file, delimited in the source by
% `# --8<-- [start:name]` / `# --8<-- [end:name]`, so the book and the running
% code cannot drift apart. tools/check_structure.py --listings fails when the
% file or the region is missing.
%
% The markers are matched through listings' range prefix and suffix keys,
% with linerange={name-name}. The form this was inherited with -- a linerange
% whose two ends were the whole marker strings -- matched nothing and printed
% the WHOLE FILE, with no warning: measured on a three-region probe file
% before this was rewritten, and the same definition sits unused in the
% sibling book it came from. A region name is a word, never a number, because
% a number is a line number to listings. The line numbers printed are the
% file's own, so a region's listing says where in the file it sits.
% Usage: \pyregion{code/ch05/decorators.py}{timing}{Caption}{lst:label}
\lstset{
  rangebeginprefix=\#\ --8<--\ [start:, rangebeginsuffix=],
  rangeendprefix=\#\ --8<--\ [end:, rangeendsuffix=],
  includerangemarker=false}
\newcommand{\pyregion}[4]{%
  \listingpath{#1}%
  \lstinputlisting[language=pythonx,style=bookcode,aboveskip=0.3em,
    linerange={#2-#2},caption={#3},label={#4}]{#1}%
}

% The C# twin of \pyfile, for the comparison half of a pair.
\newcommand{\csfile}[3]{%
  \listingpath{#1}%
  \lstinputlisting[language=csharpx,style=bookcode,aboveskip=0.3em,
    caption={#2},label={#3}]{#1}%
}

% A program's output, written by a script under code/measure/ and pulled in
% whole. There is no typed form of this and there must not be: \val{} does
% not expand inside a verbatim body, so a transcript typed into a chapter is
% one that cannot quote a computed value, and a transcript nobody ran is
% indistinguishable on the page from one that was.
\newcommand{\transcript}[1]{%
  \par\smallskip
  \IfFileExists{figures/transcripts/#1.txt}{%
    \lstinputlisting[style=bookcode,numbers=none,basicstyle=\ttfamily\footnotesize,
      keywordstyle=\color{black},backgroundcolor=\color{white},
      language={}]{figures/transcripts/#1.txt}%
  }{%
    \begin{tcolorbox}[enhanced, sharp corners, width=\linewidth,
      colback=mlight, colframe=mgrey, boxrule=0.6pt,
      left=8pt, right=8pt, top=6pt, bottom=6pt]
      {\bfseries\color{mgrey}\small \lblTranscriptMissing}\\[2pt]
      {\ttfamily\scriptsize\color{mgrey} figures/transcripts/#1.txt}
    \end{tcolorbox}%
  }%
  \par\smallskip
}

% Inline literals. Underscores inside these must be written \_.
\newcommand{\code}[1]{\texttt{\small #1}}
\newcommand{\pkg}[1]{\texttt{\small\color{mblue}#1}}
\newcommand{\api}[1]{\texttt{\small\color{mteal}#1}}

% ============================================================
%  ADMONITIONS
%  ---------------------------------------------------------
%  A deliberately small, fixed vocabulary. Two visual treatments carry the
%  whole grammar: `tinted` is a block the reader may read in passing -- a
%  tint and a thick left bar, no outline; `outlined` is a block the reader
%  must stop at, with a full rule on white. A third treatment would mean the
%  first two are not carrying their meaning.
% ============================================================
\tcbset{
  mfa tinted/.style={
    enhanced, breakable, sharp corners,
    colback=mlight, boxrule=0pt, leftrule=3pt,
    left=8pt, right=8pt, top=6pt, bottom=6pt,
    before skip=13pt, after skip=13pt,
    fonttitle=\bfseries\small,
    coltitle=black, colbacktitle=mlight,
    attach title to upper={\par\smallskip},
  },
  mfa outlined/.style={
    enhanced, breakable, sharp corners,
    colback=white, boxrule=0.6pt,
    left=8pt, right=8pt, top=6pt, bottom=6pt,
    before skip=13pt, after skip=13pt,
    fonttitle=\bfseries\small,
    coltitle=black, colbacktitle=white,
    attach title to upper={\par\smallskip},
  },
}

% Read in passing.
\newtcolorbox{note}[1][]{mfa tinted, colframe=mblue, title={\lblNote}, #1}
\newtcolorbox{warning}[1][]{mfa tinted, colframe=mred, title={\lblWarning}, #1}
% The translation. The .NET reader is the book's only assumed reader, so this
% box carries the C# side of a comparison and nothing else. Its use is
% counted: a chapter with none is a chapter that has stopped translating.
\newtcolorbox{csbox}[1][]{mfa tinted, colframe=mteal, title={\lblCsharp}, #1}
% Anything that is preview, experimental or likely to move. Python moves more
% slowly than the agent frameworks do, and this box should be rare here.
\newtcolorbox{versionbox}[1][]{mfa tinted, colframe=mamber, title={\lblVersion}, #1}

% Stop and read.
% The misconception, named AFTER the reader has walked into it. Appendix B is
% the catalogue; a trap box in a chapter is one entry of it, in place.
\newtcolorbox{trapbox}[1][]{mfa outlined, colframe=mred, title={\lblTrap}, #1}
% Marks a listing that has not been executed against the pinned versions.
% `make debt` counts these, CI prints the count, and nothing ships to a reader
% with one attached. Removing it means you ran the code.
\newtcolorbox{verifybox}[1][]{mfa outlined, colframe=mgrey, title={\lblVerify}, #1}
\newtcolorbox{exercisebox}[1][]{mfa outlined, colframe=mteal, title={\lblExercise}, #1}
% The guiding project. Points at the stage of trace-assert this section
% builds, under code/src/trace_assert/.
\newtcolorbox{projectbox}[1][]{mfa outlined, colframe=mamber, title={\lblProject}, #1}

% ============================================================
%  EXERCISES
%  ---------------------------------------------------------
%  \begin{exercise}{key}{title} ... \end{exercise}
%
%  This is the mechanism the whole book is read through: the PDF on one half
%  of the screen, an editor on the other. An exercise is therefore not a
%  paragraph but a FILE the reader opens -- a starter under
%  code/exercises/chNN/<key>.py -- and a TEST that decides whether it is done,
%  test_<key>.py beside it. The box prints both paths and the one command
%  that runs the test, and the manifest at the back lists every exercise.
%
%  An ENVIRONMENT, not a macro with a body argument: a macro argument cannot
%  contain verbatim material, so an exercise written as \exercise{key}{title}
%  {body} could never show the reader a code fragment, and the first chapter
%  that tried would have died on `\begin{python}` inside the braces.
%
%  The key is the file stem, with underscores, and it is written with
%  \detokenize so the underscore reaches the page as a character rather than
%  as a subscript. tools/check_structure.py --exercises fails when the starter
%  or the test is missing, or when the key's chapter prefix disagrees with the
%  chapter it sits in.
% ============================================================
\newcounter{exercise}[chapter]
\renewcommand{\theexercise}{\thechapter.\arabic{exercise}}
\makeatletter
\newcommand{\exercisedir}{ch\two@digits{\value{chapter}}}
\newcommand{\listofexercises}{%
  \section*{\lblExerciseManifest}%
  \phantomsection\addcontentsline{toc}{section}{\lblExerciseManifest}%
  % The entries are subsection-level and tocdepth is 1 for the contents,
  % so the depth is raised for this list alone or it prints empty --
  % which it did, on the scaffold's first clean build.
  \begingroup\c@tocdepth=2\relax\@starttoc{exr}\endgroup%
}
\makeatother
\NewDocumentEnvironment{exercise}{mm}{%
  \refstepcounter{exercise}%
  \addcontentsline{exr}{subsection}{\protect\numberline{\theexercise}\texttt{\protect\detokenize{#1}} --- #2}%
  \begin{exercisebox}[title={\lblExercise~\theexercise\ \dash{} #2}]%
}{%
  \par\medskip
  {\small\setlength{\parindent}{0pt}%
  \lblExerciseStarter: \texttt{code/exercises/\exercisedir/\detokenize{#1}.py}\\
  \lblExerciseTest: \texttt{code/exercises/\exercisedir/test\_\detokenize{#1}.py}\\
  \lblExerciseRun: \texttt{cd code \&\& uv run pytest -k \detokenize{#1}}\par}%
  \end{exercisebox}%
}

% ============================================================
%  DIAGRAMS --- Mermaid pipeline
%  ---------------------------------------------------------
%  \mermaidfig{key}{caption}{one-line manifest description}
%
%  Source of truth is figures/mermaid/<lang>/<key>.mmd, committed. `make
%  diagrams` renders each to figures/diagrams/<lang>/<key>.pdf, which is build
%  output and is not committed, so a diagram change reviews as a text diff.
%  The source is per-language because a diagram with English labels in the
%  Polish edition is exactly the half-translation that makes a bilingual book
%  feel machine-made; CI checks that both languages carry the same keys.
%
%  If the rendered PDF is absent the macro draws a placeholder AND typesets
%  the Mermaid source, in the flow rather than as a float, because a source
%  typeset verbatim is often taller than a page and a float cannot break.
% ============================================================
\makeatletter
\newcommand{\listofdiagrams}{%
  \section*{\lblDiagramManifest}%
  \phantomsection\addcontentsline{toc}{section}{\lblDiagramManifest}%
  % The entries are subsection-level and tocdepth is 1 for the contents,
  % so the depth is raised for this list alone or it prints empty --
  % which it did, on the scaffold's first clean build.
  \begingroup\c@tocdepth=2\relax\@starttoc{dgm}\endgroup%
}
\makeatother

\newcommand{\mermaidfig}[3]{%
  % \phantomsection is load-bearing and not decoration. \addcontentsline
  % records hyperref's CURRENT anchor, and here that is whatever preceded
  % the figure -- for a figure placed after a listing it is the listing's
  % line-number anchor, and listings has already advanced the counter past
  % the last line it printed, so the manifest entry linked to a
  % destination that is never created: "name{lstnumber.10.4.46} has been
  % referenced but does not exist, replaced by a fixed one". The two
  % front-matter figures had the same wrong target and were silent about
  % it, because the lines they named happen to be printed. The exercise
  % manifest never had the bug, because \begin{exercise} steps a counter
  % and that is what sets the anchor.
  \phantomsection%
  \addcontentsline{dgm}{subsection}{\protect\numberline{}\texttt{#1.mmd} --- #3}%
  \IfFileExists{figures/diagrams/\booklang/#1.pdf}{%
    \begin{figure}[htbp]
    \centering
    \includegraphics[width=0.95\linewidth,height=0.42\textheight,keepaspectratio]{figures/diagrams/\booklang/#1.pdf}%
    \caption{#2}\label{fig:#1}
    \end{figure}%
  }{%
    \par\medskip
    \begin{tcolorbox}[
      enhanced, breakable, width=\linewidth, sharp corners,
      colback=white, colframe=mgrey, boxrule=0.8pt,
      left=8pt, right=8pt, top=8pt, bottom=8pt]
      {\bfseries\color{mgrey}\small \lblNotRendered}\\[4pt]
      {\ttfamily\scriptsize\color{mgrey} figures/mermaid/\booklang/#1.mmd}\\[6pt]
      \IfFileExists{figures/mermaid/\booklang/#1.mmd}{%
        \lstinputlisting[style=bookcode,numbers=none,basicstyle=\ttfamily\scriptsize,
                         backgroundcolor=\color{mlight},frame=none,
                         keywordstyle=\color{black},language={}]{figures/mermaid/\booklang/#1.mmd}%
      }{%
        {\small\color{mred} \lblSourceMissing: \texttt{figures/mermaid/\booklang/#1.mmd}}%
      }%
    \end{tcolorbox}%
    \captionof{figure}{#2}\label{fig:#1}
    \par\medskip
  }%
}

% ============================================================
%  CHAPTER STUBS
%  ---------------------------------------------------------
%  Prints a visible marker so an unwritten chapter cannot be mistaken for a
%  written one and the page count never flatters the draft. `make debt`
%  counts them, and CI publishes the count on every build.
% ============================================================
\newcommand{\chapterstub}[1]{%
  \begin{tcolorbox}[
    enhanced, breakable, width=\linewidth, sharp corners,
    colback=mlight, colframe=mred, boxrule=1pt,
    borderline={0.6pt}{3pt}{mred, dashed},
    left=12pt, right=12pt, top=12pt, bottom=12pt]
    {\bfseries\color{mred}\large \lblNotWritten}\\[8pt]
    \small\raggedright #1
  \end{tcolorbox}%
}

% ============================================================
%  COMPUTED VALUES
%  ---------------------------------------------------------
%  A number that came out of a measurement is written by a script under
%  code/measure/ into figures/values/<name>.tex as
%
%      \pyval{e01.threads.speedup}{3.71}
%
%  and pulled into the text with \val{e01.threads.speedup}. The script is the
%  source of truth and the book cannot disagree with it, because the book does
%  not contain the digits. `make verify` fails when a committed value no
%  longer matches its script. A value the scripts do not produce prints a
%  visible marker.
%
%  \val also applies the edition's decimal separator, which is how the Polish
%  edition gets its decimal comma without a translator having to remember.
% ============================================================
\newcommand{\bookdecimalmarker}{\ifpl{,}{.}}
\IfFileExists{siunitx.sty}{%
  \usepackage{siunitx}%
  \sisetup{
    group-digits = integer,
    group-minimum-digits = 5,
    group-separator = {\,},
    output-decimal-marker = {\bookdecimalmarker}
  }%
  \newcommand{\booknum}[1]{\num{##1}}%
}{%
  \newcommand{\booknum}[1]{##1}%
}

\newcommand{\pyvalues}{}
\makeatletter
\newcommand{\pyval}[2]{\csgdef{pyval@#1}{#2}\listxadd{\pyvalues}{#1}}
% \num on a non-numeric body is a FATAL siunitx error. The generator knows
% whether it emitted a number or a word, so it writes \pyval for a number and
% \pyvaltext for a word, and this end only enforces which macro reads which.
\newcommand{\val}[1]{%
  \ifcsdef{pyval@#1}{%
    \ifcsdef{pyvaltext@#1}{%
      \PackageError{pybook}{Value `#1' is not a number}%
        {\string\val\space passes its body to siunitx, which refuses anything
         that is not a number. Use \string\valtext{#1} instead.}%
    }{}%
    \booknum{\csuse{pyval@#1}}%
  }{\textcolor{mred}{\textbf{[\lblValueMissing: #1]}}}%
}
\newcommand{\valtext}[1]{%
  \ifcsdef{pyval@#1}{\csuse{pyval@#1}}%
    {\textcolor{mred}{\textbf{[\lblValueMissing: #1]}}}%
}
\newcommand{\pyvaltext}[2]{%
  \csgdef{pyval@#1}{#2}\csgdef{pyvaltext@#1}{}\listxadd{\pyvalues}{#1}%
}
\makeatother

% figures/values/all.tex is generated by `make numbers` and does nothing but
% \input each value file. A build with no values yet still compiles; every
% \val{} in it prints a visible marker instead of a number.
\IfFileExists{figures/values/all.tex}{\input{figures/values/all}}{%
  \AtBeginDocument{\PackageWarning{pybook}{No computed values found. Run `make numbers`.}}}

% ============================================================
%  PINNED VERSIONS --- single source of truth
%  Change these here; they propagate through both editions and Appendix C.
%  Verified against pypi.org on \pinnedon (lang/<lang>.tex).
%  tools/check_versions.py compares every one of them with PyPI on each build.
%  Python itself is not on PyPI and is checked by hand against python.org.
% ============================================================
\newcommand{\bookedition}{0.1-dev}
\newcommand{\bookyear}{2026}
\newcommand{\pyver}{3.14}              % the minor the book targets
\newcommand{\pypatch}{3.14.7}          % the exact interpreter the code ran on
\newcommand{\uvver}{0.12.13}           % uv
\newcommand{\dotnetver}{10.0.100}      % dotnet
\newcommand{\ruffver}{0.16.7}          % ruff
\newcommand{\pyrightver}{1.1.414}      % pyright
\newcommand{\mypyver}{2.3.1}           % mypy
\newcommand{\pydanticver}{2.13.5}      % pydantic
\newcommand{\pysettingsver}{2.15.0}    % pydantic-settings
\newcommand{\fastapiver}{0.141.1}      % fastapi
\newcommand{\uvicornver}{0.53.0}       % uvicorn
\newcommand{\sqlalchemyver}{2.0.52}    % sqlalchemy
\newcommand{\alembicver}{1.20.0}       % alembic
\newcommand{\aiosqlitever}{0.22.1}     % aiosqlite
\newcommand{\pytestver}{9.1.1}         % pytest
\newcommand{\pytestasyncver}{1.4.0}    % pytest-asyncio
\newcommand{\hypothesisver}{6.168.0}   % hypothesis
\newcommand{\coveragever}{7.16.1}      % coverage
\newcommand{\httpxver}{0.28.1}         % httpx
\newcommand{\structlogver}{26.1.0}     % structlog
\newcommand{\otelver}{1.44.0}          % opentelemetry-sdk
\newcommand{\polarsver}{1.44.2}        % polars
\newcommand{\anthropicver}{1.5.0}      % anthropic
\newcommand{\openaiver}{3.13.0}        % openai
\newcommand{\pipauditver}{2.10.1}      % pip-audit

% ============================================================
%  INDEX
%  Two-column, and its entries include \texttt{} names that do not hyphenate.
%  Ragged right keeps multi-word entries off the margin; it cannot help a
%  single token wider than the column, so keep verbatim index entries under
%  about 29 characters and index longer names by concept instead.
%  \raggedcolumns is the switch imakeidx's multicols actually reads.
% ============================================================
\apptocmd{\theindex}{\raggedcolumns\raggedright}{}{%
  \PackageWarning{pybook}{Could not patch theindex for ragged setting}}
\providecommand*\see[2]{}
\renewcommand*\see[2]{\emph{\lblIndexSee} #1}
\providecommand*\seealso[2]{}
\renewcommand*\seealso[2]{\emph{\lblIndexSeeAlso} #1}

\makeindex
; nothing else differs between them at this level.
%  Every user-visible string lives in lang/en.tex or lang/pl.tex, so a label
%  that appears in one edition and not the other is a missing macro rather
%  than a silent divergence.
%
%  Lineage. The bilingual wiring (\booklang, lang/, body.tex, the generated
%  structure) is the math book's; the chapter machinery (listings, the
%  admonitions, verifybox, \pyfile and \pyregion, the Mermaid pipeline) is the
%  LangChain book's. What is new here is the exercise macro, which ties every
%  exercise to a starter file and a test under code/, because this book is
%  read with an editor open beside it.
% ============================================================

% \booklang must be set by the main file. Default to English rather than
% failing, so a stray % ============================================================
%  preamble.tex --- styling, environments, macros
%
%  Shared by BOTH editions. main-en.tex and main-pl.tex each define \booklang
%  before \input{preamble}; nothing else differs between them at this level.
%  Every user-visible string lives in lang/en.tex or lang/pl.tex, so a label
%  that appears in one edition and not the other is a missing macro rather
%  than a silent divergence.
%
%  Lineage. The bilingual wiring (\booklang, lang/, body.tex, the generated
%  structure) is the math book's; the chapter machinery (listings, the
%  admonitions, verifybox, \pyfile and \pyregion, the Mermaid pipeline) is the
%  LangChain book's. What is new here is the exercise macro, which ties every
%  exercise to a starter file and a test under code/, because this book is
%  read with an editor open beside it.
% ============================================================

% \booklang must be set by the main file. Default to English rather than
% failing, so a stray \input{preamble} still compiles and says what it did.
\providecommand{\booklang}{en}

\usepackage{etoolbox}   % loaded first: the language switch below needs it

% ---------- Page geometry ----------
% ONE format: A4, 12pt, ONE-SIDED. This book is read on a screen with an
% editor open beside it, and never printed as a bound book, so there is no
% recto, no verso, no gutter, no mirrored margin and no blank leaf before a
% chapter -- every one of those is a property of paper that a PDF viewer at
% half a screen's width turns into wasted pixels.
%
% The margins are symmetric and the measure is about seventy-five characters
% of prose, which is where the eye keeps the line return. The number that was
% actually chosen is the MONOSPACE measure: at \footnotesize on this block a
% listing line of 79 characters fits without wrapping, and 79 is the width
% every listing under code/ is held to by tools/check_structure.py --lines.
% A wrapped listing line is not an error TeX reports -- breaklines=true wraps
% it silently and prints an arrow into the middle of what the reader is meant
% to paste into an editor -- so the width is enforced in the source instead.
\usepackage[
  a4paper,
  left=3.1cm, right=3.1cm, top=2.6cm, bottom=3.0cm,
  head=26pt, headsep=0.8cm, footskip=1.4cm
]{geometry}

% ---------- Encoding & fonts ----------
\usepackage[T1]{fontenc}
\usepackage[utf8]{inputenc}
\IfFileExists{lmodern.sty}{\usepackage{lmodern}}{}
% newtxtext takes its TS1 symbols (\textcopyright, \textdegree) from TeX
% Gyre Termes, which Debian packages SEPARATELY (tex-gyre): a machine with
% newtx and without it dies on the copyright page with
% `Font TS1/ntxtlf/m/n/10.95=ts1-qtmr ... not loadable`. There is no
% inert-safe probe for a font metric in pdfTeX -- \IfFileExists searches
% TEXINPUTS and not TFMFONTS, so it reports the metric missing on a machine
% that has it, and \suppressfontnotfounderror is LuaTeX-only -- so nothing
% here guesses. CI's full TeX Live is the reference; locally, install
% tex-gyre. Both facts are recorded in CLAUDE.md under Build traps.
\IfFileExists{newtxtext.sty}{\usepackage{newtxtext}}{}
% amsmath before newtxmath, because newtxmath patches amsmath's internals.
% amssymb only when newtxmath is absent: newtxmath supplies the AMS symbols
% itself and loading both is a fatal clash (\Bbbk already defined) that is
% invisible on a machine without newtx and fatal on a full TeX Live.
\usepackage{amsmath}
\IfFileExists{newtxmath.sty}{\usepackage{newtxmath}}{\usepackage{amssymb}}
\IfFileExists{inconsolata.sty}{\usepackage[varqu,scaled=0.95]{inconsolata}}{}
\IfFileExists{lmodern.sty}{\usepackage{microtype}}{\usepackage[expansion=false]{microtype}}

% upquote is not a refinement. In T1 the input character ' is quoteright, so
% a listings body sets f('x') with U+2019 at both ends of the literal, and
% typed back in that is a SyntaxError. inconsolata's varqu option hides the
% defect on a machine that has inconsolata; upquote makes ' the ASCII
% apostrophe whatever the typewriter font is. In a book whose every listing is
% meant to be pasted into an editor, this line is load-bearing.
\IfFileExists{upquote.sty}{\usepackage{upquote}}{}

% And upquote's twin, for the same reason one character over. In T1 the pair
% -- is a ligature for an en dash, in EVERY family including the typewriter
% one, so \code{uv sync --locked} sets as "uv sync <en dash>locked" and a
% reader who copies it off the page gets an unknown-argument error from a
% command the chapter told them to run. It is invisible in review because it
% looks like a dash, and it bites exactly the flags a packaging chapter is
% made of. A `listings` body is already safe -- listings sets characters
% singly and no ligature forms -- which is what makes the defect so easy to
% miss: the same text is right inside a listing and wrong in the sentence
% above it.
%
% Scoped to tt* on purpose, and measured rather than assumed: prose keeps its
% own dashes, so \dash{} still sets an em dash and a range like 1--2 still
% sets an en dash. Probed against this preamble in all three positions before
% it was believed.
\DisableLigatures[-]{encoding = T1, family = tt*}

% ---------- Language ----------
% Loading babel with a language whose .ldf is absent is a *fatal* error, not
% a warning, so both the package and each language file are probed before
% either is requested. Both editions load the other language too, because the
% glossary rows and the cheat sheet carry words from both.
\IfFileExists{babel.sty}{%
  \IfFileExists{polish.ldf}{%
    \IfFileExists{british.ldf}{%
      \ifdefstring{\booklang}{pl}%
        {\usepackage[british,main=polish]{babel}}%
        {\usepackage[polish,main=british]{babel}}%
    }{\usepackage[polish]{babel}}%
  }{%
    \IfFileExists{british.ldf}{\usepackage[british]{babel}}{}%
  }%
}{}

% babel-polish makes " an active shorthand. listings reads its body verbatim
% and an active character in a verbatim context is a category-code accident
% waiting to happen, so the shorthand is turned off. Polish quotation marks
% come from \enquote, which csquotes sets per language.
\ifdefstring{\booklang}{pl}{\AtBeginDocument{\shorthandoff{"}}}{}

% ---------- Core packages ----------
\usepackage{graphicx}
\usepackage{xcolor}
\usepackage{booktabs}
\usepackage{tabularx}
% Appendix A's cheat sheet runs past a page in every part, and tabularx
% cannot break. longtable can, and works with booktabs.
\usepackage{longtable}
\usepackage{array}
\usepackage{enumitem}
\usepackage{listings}
\usepackage[most]{tcolorbox}
\usepackage{fancyhdr}
\usepackage{titlesec}
\usepackage{caption}
\usepackage{imakeidx}
% csquotes with autostyle follows babel's active language, so the identical
% source \enquote{x} sets 'x' in the English edition and the Polish quotation
% marks in the Polish one.
\IfFileExists{csquotes.sty}{\usepackage[autostyle=true]{csquotes}}{%
  \providecommand{\enquote}[1]{``##1''}}
\usepackage[hidelinks]{hyperref}

% ---------- Palette ----------
% The same family as the two companion volumes, so the three read as a set.
\definecolor{mblue}{HTML}{1F4E79}
\definecolor{mteal}{HTML}{0E7C7B}
\definecolor{mamber}{HTML}{B26A00}
\definecolor{mred}{HTML}{9B2C2C}
\definecolor{mviolet}{HTML}{7B4397}
\definecolor{mgrey}{HTML}{5A5A5A}
\definecolor{mlight}{HTML}{F4F6F8}
\definecolor{mfaint}{HTML}{FBFCFD}
\definecolor{codebg}{HTML}{FAFAFA}
\definecolor{codeframe}{HTML}{D8DEE4}
\definecolor{kw}{HTML}{0B5394}
\definecolor{str}{HTML}{9B2C2C}
\definecolor{cmt}{HTML}{4F7A4F}
\definecolor{typ}{HTML}{0E7C7B}
\definecolor{dec}{HTML}{7B4397}

\hypersetup{
  colorlinks=true,
  linkcolor=mblue,
  urlcolor=mteal,
  citecolor=mblue,
  pdfauthor={Konrad Cinkusz}
}

% \ifpl{Polish text}{English text} -- for the handful of places where a
% string is structural rather than content. Anything longer than a few words
% belongs in lang/ or in the edition's own source, not here.
% \DeclareRobustCommand, not \newcommand: a fragile \ifpl inside a moving
% argument is written to the .toc one level expanded and fails on the next run
% in exactly one of the two editions.
\DeclareRobustCommand{\ifpl}[2]{\ifdefstring{\booklang}{pl}{#1}{#2}}
\makeatletter
\pdfstringdefDisableCommands{%
  \ifdefstring{\booklang}{pl}{\let\ifpl\@firstoftwo}{\let\ifpl\@secondoftwo}%
  % \api{}, \pkg{} and \code{} are allowed in a section title, and a title
  % becomes a PDF bookmark. hyperref drops the \color command from a bookmark
  % string but keeps its ARGUMENT, so a title carrying \api{model\_validate}
  % made a bookmark reading "mtealmodel_validate" and warned "Token not
  % allowed" twice. The colour name is gobbled here instead. Measured on a
  % probe against this preamble before the line was added.
  \def\color#1{}%
}
\makeatother

% ---------- Language strings ----------
\input{lang/\booklang}

% The PDF's subject, keyword list and document language read \lblPdf... from
% the language file input on the line above, so they cannot sit in the colour
% \hypersetup block, which runs before it. Two blocks is the ordering, not an
% oversight; merging them raises Undefined control sequence and still writes
% a PDF, with the three fields empty.
\hypersetup{
  pdfsubject={\lblPdfSubject},
  pdfkeywords={\lblPdfKeywords},
  pdflang={\lblPdfLang}
}

% ---------- Headings ----------
% \raggedright on the title body is load-bearing: at \Huge on this block a
% title that cannot break overflows the margin by 30 pt or more, and several
% chapter titles here run past fifty characters.
\titleformat{\chapter}[display]
  {\normalfont\huge\bfseries\color{mblue}\raggedright}
  {\normalfont\normalsize\color{mgrey}\MakeUppercase{\chaptertitlename}\ \thechapter}
  {8pt}{\Huge\raggedright}
\titlespacing*{\chapter}{0pt}{10pt}{28pt}
\titleformat{\section}{\normalfont\Large\bfseries\color{mblue}\raggedright}{\thesection}{0.8em}{}
\titleformat{\subsection}{\normalfont\large\bfseries\color{mgrey}\raggedright}{\thesubsection}{0.7em}{}

% Sections in the contents, subsections not: fourteen chapters of subsections
% is a contents nobody reads, and the reader navigates by chapter and section.
\setcounter{tocdepth}{1}
\setcounter{secnumdepth}{2}

% One-sided running head: the chapter on the left, the folio on the right.
\pagestyle{fancy}
\fancyhf{}
\fancyhead[L]{\small\nouppercase{\leftmark}}
\fancyhead[R]{\small\thepage}
\renewcommand{\headrulewidth}{0.4pt}
\fancypagestyle{plain}{\fancyhf{}\fancyfoot[R]{\small\thepage}%
  \renewcommand{\headrulewidth}{0pt}}

% Drops the word "Chapter" from the head: the number is still there and the
% longer titles no longer overflow. MUST come after \pagestyle{fancy}, which
% installs its own \chaptermark and silently overwrites anything earlier.
\renewcommand{\chaptermark}[1]{\markboth{\thechapter.\ #1}{}}

% ============================================================
%  CODE LISTINGS
%  ---------------------------------------------------------
%  Every listing in a chapter is a FILE under code/, pulled in with \pyfile or
%  \pyregion, and CI runs it against the pinned versions. An in-source
%  \begin{python} block is allowed for a fragment of three or four lines that
%  exists to show a shape rather than to run; anything longer belongs in a
%  file the reader can open in the editor beside the PDF.
% ============================================================
\lstdefinelanguage{pythonx}{
  language=Python,
  morekeywords={
    async,await,match,case,nonlocal,yield,type,
    TypedDict,Annotated,Protocol,Literal,Generic,Final,ClassVar,Self,
    TypeVar,TypeGuard,TypeIs,Callable,Iterator,Iterable,Awaitable,
    dataclass,field,BaseModel,Field,Enum,StrEnum,NamedTuple,
    override,overload,property,staticmethod,classmethod
  },
  morekeywords=[2]{
    gather,create_task,TaskGroup,to_thread,timeout,wait_for,Semaphore,Queue,
    fixture,parametrize,monkeypatch,raises,mark,
    Depends,FastAPI,APIRouter,HTTPException,
    select,Session,Mapped,mapped_column,relationship,
    model_validate,model_validate_json,model_dump,model_json_schema
  },
  sensitive=true
}

% listings' C# definition predates async, records and pattern matching.
\lstdefinelanguage{csharpx}{
  language=[Sharp]C,
  morekeywords={
    async,await,var,dynamic,record,init,required,nameof,yield,partial,
    global,scoped,when,with,and,or,not,notnull,unmanaged,file,get,set,value,
    where,select,from,orderby,let,join,into,ascending,descending,on,equals,
    Task,ValueTask,CancellationToken,IEnumerable,IAsyncEnumerable,Func,Action
  },
  sensitive=true
}

\lstdefinestyle{bookcode}{
  basicstyle=\ttfamily\footnotesize,
  keywordstyle=\color{kw}\bfseries,
  keywordstyle=[2]\color{typ},
  stringstyle=\color{str},
  % Colour only, no \itshape: inconsolata (zi4) has no italic shape, so an
  % italic comment style is substituted with upright on every machine that
  % has the font, and the log carries "Font shape `T1/zi4/m/it' undefined"
  % on every build. The colour already marks a comment.
  commentstyle=\color{cmt},
  emphstyle=\color{dec}\bfseries,
  numbers=left,
  numberstyle=\ttfamily\tiny\color{mgrey},
  numbersep=8pt,
  backgroundcolor=\color{codebg},
  frame=single,
  rulecolor=\color{codeframe},
  framesep=6pt,
  breaklines=true,
  breakatwhitespace=false,
  postbreak=\mbox{\textcolor{mgrey}{$\hookrightarrow$}\space},
  showstringspaces=false,
  tabsize=4,
  columns=fullflexible,
  keepspaces=true,
  captionpos=b,
  aboveskip=1.2em,
  belowskip=1.2em,
  % Maps the characters most likely to be pasted into a listing by accident.
  %
  % Do NOT add a mapping for U+00A0 (non-breaking space). It looks like the
  % obvious next entry and it makes `listings` abort the run with "Improper
  % alphabetic constant" -- fatal, no PDF, and the message names neither the
  % character nor this line. The ASCII-in-listings convention, checked by
  % parity's C13, is the guard against it instead.
  literate={ł}{{\l}}1 {Ł}{{\L}}1 {ą}{{\k a}}1 {ę}{{\k e}}1
           {ś}{{\'s}}1 {ć}{{\'c}}1 {ż}{{\.z}}1 {ź}{{\'z}}1 {ó}{{\'o}}1 {ń}{{\'n}}1
           {Ą}{{\k A}}1 {Ę}{{\k E}}1 {Ś}{{\'S}}1 {Ć}{{\'C}}1
           {Ż}{{\.Z}}1 {Ź}{{\'Z}}1 {Ó}{{\'O}}1 {Ń}{{\'N}}1
           {—}{{-{}-}}2 {–}{{-}}1 {‘}{{'}}1 {’}{{'}}1
           {“}{{"}}1 {”}{{"}}1 {°}{{\textdegree}}1
}

\lstset{style=bookcode}

\lstnewenvironment{python}[1][]
  {\lstset{language=pythonx,style=bookcode,#1}}
  {}

% C# comparison snippets. The whole book is written for engineers arriving
% from .NET, so listings come in pairs often enough that the C# half has its
% own environment rather than a generic one.
\lstnewenvironment{csharp}[1][]
  {\lstset{language=csharpx,style=bookcode,#1}}
  {}

\lstnewenvironment{shellcmd}[1][]
  {\lstset{language=bash,style=bookcode,numbers=none,keywordstyle=\color{mteal}\bfseries,#1}}
  {}

\lstnewenvironment{tomlcode}[1][]
  {\lstset{style=bookcode,numbers=none,keywordstyle=\color{black},#1}}
  {}

\lstnewenvironment{yamlcode}[1][]
  {\lstset{style=bookcode,numbers=none,keywordstyle=\color{black},#1}}
  {}

\lstnewenvironment{jsoncode}[1][]
  {\lstset{style=bookcode,numbers=none,keywordstyle=\color{black},#1}}
  {}

% Terminal transcripts typed into the source. Allowed for a shell session the
% reader is meant to reproduce; NOT for program output, which is \transcript
% below, because an in-source transcript nobody ran reads exactly like one
% that was, and that is where a remembered number hides.
\lstnewenvironment{console}[1][]
  {\lstset{style=bookcode,numbers=none,basicstyle=\ttfamily\footnotesize,
           keywordstyle=\color{black},backgroundcolor=\color{white},#1}}
  {}

% The path is printed above every file listing, in the repository-relative
% form the reader opens in the editor beside the PDF. That is the whole point
% of a listing being a file, and the front matter promises it. The \nobreak
% keeps the path on the same page as the listing it names.
% ---------------------------------------------------------------------------
%  Appendix B prints the trap catalogue from notes/02-traps.md. One entry is
%  its number -- which a chapter may cite, and which is never reused -- the
%  habit in the reader's own voice, and what Python does instead.
%
%  The body is set RAGGED RIGHT on purpose. An entry is mostly \code{} spans,
%  which are unbreakable runs, and sixty-seven justified paragraphs of them
%  is sixty-seven chances at an overfull hbox in the edition with the longer
%  words. Ragged right removes the class rather than fixing it one box at a
%  time, and a reference list is the one place in this book where it reads
%  better anyway.
%
%  The habit is set with \enquote{} rather than \emph{}, and that is not a
%  style choice. Every habit here carries \code{} spans, and \code{} inside
%  \emph{} asks inconsolata for an italic it does not have: the first build
%  of this appendix raised `Font shape T1/zi4/m/it undefined', which
%  checklog.py is hard on for exactly this reason. It is the trap CLAUDE.md
%  records, met sixty-seven times at once -- and fixed once, because the
%  entries go through a macro. A quotation is also what the habit IS, and
%  \enquote{} is what the Polish edition owes anyway.
% ---------------------------------------------------------------------------
%  Appendix A's cheat sheet: one table per part, C# on the left, Python on
%  the right, and the chapter that teaches the row on the end. The three
%  header strings are arguments so that one environment serves both
%  editions and the structural signature stays identical.
%
%  Columns are ragged right for the reason Appendix B's entries are: a cell
%  is mostly \code{} spans, which are unbreakable runs, and a justified
%  narrow column full of them is an overfull box waiting for the edition
%  with the longer words.
% ---------------------------------------------------------------------------
%  Appendix C's tool matrix: the .NET tool, the Python tool this book uses,
%  the version it was verified against and the chapter that introduces it.
%
%  The version column prints \pydanticver and its siblings and is never
%  typed, which is the whole point of the appendix: it is the one page where
%  a reader sees every pin at once, and a pin that had to be re-typed here
%  would be a second copy of a fact the preamble already owns.
%  tools/check_structure.py --tools fails when a pin macro exists and this
%  table does not print it.
\NewDocumentEnvironment{toolmatrix}{mmmm}{%
  \small
  \begin{longtable}{@{}%
    >{\raggedright\arraybackslash}p{0.27\linewidth}%
    >{\raggedright\arraybackslash}p{0.31\linewidth}%
    >{\raggedright\arraybackslash}p{0.20\linewidth}%
    >{\raggedright\arraybackslash}p{0.08\linewidth}@{}}
  \toprule
  \textbf{#1} & \textbf{#2} & \textbf{#3} & \textbf{#4} \\
  \midrule
  \endfirsthead
  \toprule
  \textbf{#1} & \textbf{#2} & \textbf{#3} & \textbf{#4} \\
  \midrule
  \endhead
  \midrule
  \endfoot
  \bottomrule
  \endlastfoot
}{%
  \end{longtable}%
}

\NewDocumentEnvironment{cheatsheet}{mmm}{%
  \small
  \begin{longtable}{@{}%
    >{\raggedright\arraybackslash}p{0.41\linewidth}%
    >{\raggedright\arraybackslash}p{0.41\linewidth}%
    >{\raggedright\arraybackslash}p{0.08\linewidth}@{}}
  \toprule
  \textbf{#1} & \textbf{#2} & \textbf{#3} \\
  \midrule
  \endfirsthead
  \toprule
  \textbf{#1} & \textbf{#2} & \textbf{#3} \\
  \midrule
  \endhead
  \midrule
  \endfoot
  \bottomrule
  \endlastfoot
}{%
  \end{longtable}%
}

\newcommand{\trapentry}[3]{%
  \par\medskip
  \noindent\textbf{#1.}\enspace\enquote{#2}\par
  \nopagebreak
  {\raggedright\noindent#3\par}%
}

\newcommand{\listingpath}[1]{%
  \par\medskip\noindent
  {\ttfamily\footnotesize\color{mgrey}\detokenize{#1}}%
  \par\nobreak}

% External source file -- the preferred form. The path is relative to the
% repository root, which is also the path the reader opens in the editor.
% Usage: \pyfile{code/ch05/decorators.py}{Caption}{lst:label}
\newcommand{\pyfile}[3]{%
  \listingpath{#1}%
  \lstinputlisting[language=pythonx,style=bookcode,aboveskip=0.3em,
    caption={#2},label={#3}]{#1}%
}

% A named region of an external file, delimited in the source by
% `# --8<-- [start:name]` / `# --8<-- [end:name]`, so the book and the running
% code cannot drift apart. tools/check_structure.py --listings fails when the
% file or the region is missing.
%
% The markers are matched through listings' range prefix and suffix keys,
% with linerange={name-name}. The form this was inherited with -- a linerange
% whose two ends were the whole marker strings -- matched nothing and printed
% the WHOLE FILE, with no warning: measured on a three-region probe file
% before this was rewritten, and the same definition sits unused in the
% sibling book it came from. A region name is a word, never a number, because
% a number is a line number to listings. The line numbers printed are the
% file's own, so a region's listing says where in the file it sits.
% Usage: \pyregion{code/ch05/decorators.py}{timing}{Caption}{lst:label}
\lstset{
  rangebeginprefix=\#\ --8<--\ [start:, rangebeginsuffix=],
  rangeendprefix=\#\ --8<--\ [end:, rangeendsuffix=],
  includerangemarker=false}
\newcommand{\pyregion}[4]{%
  \listingpath{#1}%
  \lstinputlisting[language=pythonx,style=bookcode,aboveskip=0.3em,
    linerange={#2-#2},caption={#3},label={#4}]{#1}%
}

% The C# twin of \pyfile, for the comparison half of a pair.
\newcommand{\csfile}[3]{%
  \listingpath{#1}%
  \lstinputlisting[language=csharpx,style=bookcode,aboveskip=0.3em,
    caption={#2},label={#3}]{#1}%
}

% A program's output, written by a script under code/measure/ and pulled in
% whole. There is no typed form of this and there must not be: \val{} does
% not expand inside a verbatim body, so a transcript typed into a chapter is
% one that cannot quote a computed value, and a transcript nobody ran is
% indistinguishable on the page from one that was.
\newcommand{\transcript}[1]{%
  \par\smallskip
  \IfFileExists{figures/transcripts/#1.txt}{%
    \lstinputlisting[style=bookcode,numbers=none,basicstyle=\ttfamily\footnotesize,
      keywordstyle=\color{black},backgroundcolor=\color{white},
      language={}]{figures/transcripts/#1.txt}%
  }{%
    \begin{tcolorbox}[enhanced, sharp corners, width=\linewidth,
      colback=mlight, colframe=mgrey, boxrule=0.6pt,
      left=8pt, right=8pt, top=6pt, bottom=6pt]
      {\bfseries\color{mgrey}\small \lblTranscriptMissing}\\[2pt]
      {\ttfamily\scriptsize\color{mgrey} figures/transcripts/#1.txt}
    \end{tcolorbox}%
  }%
  \par\smallskip
}

% Inline literals. Underscores inside these must be written \_.
\newcommand{\code}[1]{\texttt{\small #1}}
\newcommand{\pkg}[1]{\texttt{\small\color{mblue}#1}}
\newcommand{\api}[1]{\texttt{\small\color{mteal}#1}}

% ============================================================
%  ADMONITIONS
%  ---------------------------------------------------------
%  A deliberately small, fixed vocabulary. Two visual treatments carry the
%  whole grammar: `tinted` is a block the reader may read in passing -- a
%  tint and a thick left bar, no outline; `outlined` is a block the reader
%  must stop at, with a full rule on white. A third treatment would mean the
%  first two are not carrying their meaning.
% ============================================================
\tcbset{
  mfa tinted/.style={
    enhanced, breakable, sharp corners,
    colback=mlight, boxrule=0pt, leftrule=3pt,
    left=8pt, right=8pt, top=6pt, bottom=6pt,
    before skip=13pt, after skip=13pt,
    fonttitle=\bfseries\small,
    coltitle=black, colbacktitle=mlight,
    attach title to upper={\par\smallskip},
  },
  mfa outlined/.style={
    enhanced, breakable, sharp corners,
    colback=white, boxrule=0.6pt,
    left=8pt, right=8pt, top=6pt, bottom=6pt,
    before skip=13pt, after skip=13pt,
    fonttitle=\bfseries\small,
    coltitle=black, colbacktitle=white,
    attach title to upper={\par\smallskip},
  },
}

% Read in passing.
\newtcolorbox{note}[1][]{mfa tinted, colframe=mblue, title={\lblNote}, #1}
\newtcolorbox{warning}[1][]{mfa tinted, colframe=mred, title={\lblWarning}, #1}
% The translation. The .NET reader is the book's only assumed reader, so this
% box carries the C# side of a comparison and nothing else. Its use is
% counted: a chapter with none is a chapter that has stopped translating.
\newtcolorbox{csbox}[1][]{mfa tinted, colframe=mteal, title={\lblCsharp}, #1}
% Anything that is preview, experimental or likely to move. Python moves more
% slowly than the agent frameworks do, and this box should be rare here.
\newtcolorbox{versionbox}[1][]{mfa tinted, colframe=mamber, title={\lblVersion}, #1}

% Stop and read.
% The misconception, named AFTER the reader has walked into it. Appendix B is
% the catalogue; a trap box in a chapter is one entry of it, in place.
\newtcolorbox{trapbox}[1][]{mfa outlined, colframe=mred, title={\lblTrap}, #1}
% Marks a listing that has not been executed against the pinned versions.
% `make debt` counts these, CI prints the count, and nothing ships to a reader
% with one attached. Removing it means you ran the code.
\newtcolorbox{verifybox}[1][]{mfa outlined, colframe=mgrey, title={\lblVerify}, #1}
\newtcolorbox{exercisebox}[1][]{mfa outlined, colframe=mteal, title={\lblExercise}, #1}
% The guiding project. Points at the stage of trace-assert this section
% builds, under code/src/trace_assert/.
\newtcolorbox{projectbox}[1][]{mfa outlined, colframe=mamber, title={\lblProject}, #1}

% ============================================================
%  EXERCISES
%  ---------------------------------------------------------
%  \begin{exercise}{key}{title} ... \end{exercise}
%
%  This is the mechanism the whole book is read through: the PDF on one half
%  of the screen, an editor on the other. An exercise is therefore not a
%  paragraph but a FILE the reader opens -- a starter under
%  code/exercises/chNN/<key>.py -- and a TEST that decides whether it is done,
%  test_<key>.py beside it. The box prints both paths and the one command
%  that runs the test, and the manifest at the back lists every exercise.
%
%  An ENVIRONMENT, not a macro with a body argument: a macro argument cannot
%  contain verbatim material, so an exercise written as \exercise{key}{title}
%  {body} could never show the reader a code fragment, and the first chapter
%  that tried would have died on `\begin{python}` inside the braces.
%
%  The key is the file stem, with underscores, and it is written with
%  \detokenize so the underscore reaches the page as a character rather than
%  as a subscript. tools/check_structure.py --exercises fails when the starter
%  or the test is missing, or when the key's chapter prefix disagrees with the
%  chapter it sits in.
% ============================================================
\newcounter{exercise}[chapter]
\renewcommand{\theexercise}{\thechapter.\arabic{exercise}}
\makeatletter
\newcommand{\exercisedir}{ch\two@digits{\value{chapter}}}
\newcommand{\listofexercises}{%
  \section*{\lblExerciseManifest}%
  \phantomsection\addcontentsline{toc}{section}{\lblExerciseManifest}%
  % The entries are subsection-level and tocdepth is 1 for the contents,
  % so the depth is raised for this list alone or it prints empty --
  % which it did, on the scaffold's first clean build.
  \begingroup\c@tocdepth=2\relax\@starttoc{exr}\endgroup%
}
\makeatother
\NewDocumentEnvironment{exercise}{mm}{%
  \refstepcounter{exercise}%
  \addcontentsline{exr}{subsection}{\protect\numberline{\theexercise}\texttt{\protect\detokenize{#1}} --- #2}%
  \begin{exercisebox}[title={\lblExercise~\theexercise\ \dash{} #2}]%
}{%
  \par\medskip
  {\small\setlength{\parindent}{0pt}%
  \lblExerciseStarter: \texttt{code/exercises/\exercisedir/\detokenize{#1}.py}\\
  \lblExerciseTest: \texttt{code/exercises/\exercisedir/test\_\detokenize{#1}.py}\\
  \lblExerciseRun: \texttt{cd code \&\& uv run pytest -k \detokenize{#1}}\par}%
  \end{exercisebox}%
}

% ============================================================
%  DIAGRAMS --- Mermaid pipeline
%  ---------------------------------------------------------
%  \mermaidfig{key}{caption}{one-line manifest description}
%
%  Source of truth is figures/mermaid/<lang>/<key>.mmd, committed. `make
%  diagrams` renders each to figures/diagrams/<lang>/<key>.pdf, which is build
%  output and is not committed, so a diagram change reviews as a text diff.
%  The source is per-language because a diagram with English labels in the
%  Polish edition is exactly the half-translation that makes a bilingual book
%  feel machine-made; CI checks that both languages carry the same keys.
%
%  If the rendered PDF is absent the macro draws a placeholder AND typesets
%  the Mermaid source, in the flow rather than as a float, because a source
%  typeset verbatim is often taller than a page and a float cannot break.
% ============================================================
\makeatletter
\newcommand{\listofdiagrams}{%
  \section*{\lblDiagramManifest}%
  \phantomsection\addcontentsline{toc}{section}{\lblDiagramManifest}%
  % The entries are subsection-level and tocdepth is 1 for the contents,
  % so the depth is raised for this list alone or it prints empty --
  % which it did, on the scaffold's first clean build.
  \begingroup\c@tocdepth=2\relax\@starttoc{dgm}\endgroup%
}
\makeatother

\newcommand{\mermaidfig}[3]{%
  % \phantomsection is load-bearing and not decoration. \addcontentsline
  % records hyperref's CURRENT anchor, and here that is whatever preceded
  % the figure -- for a figure placed after a listing it is the listing's
  % line-number anchor, and listings has already advanced the counter past
  % the last line it printed, so the manifest entry linked to a
  % destination that is never created: "name{lstnumber.10.4.46} has been
  % referenced but does not exist, replaced by a fixed one". The two
  % front-matter figures had the same wrong target and were silent about
  % it, because the lines they named happen to be printed. The exercise
  % manifest never had the bug, because \begin{exercise} steps a counter
  % and that is what sets the anchor.
  \phantomsection%
  \addcontentsline{dgm}{subsection}{\protect\numberline{}\texttt{#1.mmd} --- #3}%
  \IfFileExists{figures/diagrams/\booklang/#1.pdf}{%
    \begin{figure}[htbp]
    \centering
    \includegraphics[width=0.95\linewidth,height=0.42\textheight,keepaspectratio]{figures/diagrams/\booklang/#1.pdf}%
    \caption{#2}\label{fig:#1}
    \end{figure}%
  }{%
    \par\medskip
    \begin{tcolorbox}[
      enhanced, breakable, width=\linewidth, sharp corners,
      colback=white, colframe=mgrey, boxrule=0.8pt,
      left=8pt, right=8pt, top=8pt, bottom=8pt]
      {\bfseries\color{mgrey}\small \lblNotRendered}\\[4pt]
      {\ttfamily\scriptsize\color{mgrey} figures/mermaid/\booklang/#1.mmd}\\[6pt]
      \IfFileExists{figures/mermaid/\booklang/#1.mmd}{%
        \lstinputlisting[style=bookcode,numbers=none,basicstyle=\ttfamily\scriptsize,
                         backgroundcolor=\color{mlight},frame=none,
                         keywordstyle=\color{black},language={}]{figures/mermaid/\booklang/#1.mmd}%
      }{%
        {\small\color{mred} \lblSourceMissing: \texttt{figures/mermaid/\booklang/#1.mmd}}%
      }%
    \end{tcolorbox}%
    \captionof{figure}{#2}\label{fig:#1}
    \par\medskip
  }%
}

% ============================================================
%  CHAPTER STUBS
%  ---------------------------------------------------------
%  Prints a visible marker so an unwritten chapter cannot be mistaken for a
%  written one and the page count never flatters the draft. `make debt`
%  counts them, and CI publishes the count on every build.
% ============================================================
\newcommand{\chapterstub}[1]{%
  \begin{tcolorbox}[
    enhanced, breakable, width=\linewidth, sharp corners,
    colback=mlight, colframe=mred, boxrule=1pt,
    borderline={0.6pt}{3pt}{mred, dashed},
    left=12pt, right=12pt, top=12pt, bottom=12pt]
    {\bfseries\color{mred}\large \lblNotWritten}\\[8pt]
    \small\raggedright #1
  \end{tcolorbox}%
}

% ============================================================
%  COMPUTED VALUES
%  ---------------------------------------------------------
%  A number that came out of a measurement is written by a script under
%  code/measure/ into figures/values/<name>.tex as
%
%      \pyval{e01.threads.speedup}{3.71}
%
%  and pulled into the text with \val{e01.threads.speedup}. The script is the
%  source of truth and the book cannot disagree with it, because the book does
%  not contain the digits. `make verify` fails when a committed value no
%  longer matches its script. A value the scripts do not produce prints a
%  visible marker.
%
%  \val also applies the edition's decimal separator, which is how the Polish
%  edition gets its decimal comma without a translator having to remember.
% ============================================================
\newcommand{\bookdecimalmarker}{\ifpl{,}{.}}
\IfFileExists{siunitx.sty}{%
  \usepackage{siunitx}%
  \sisetup{
    group-digits = integer,
    group-minimum-digits = 5,
    group-separator = {\,},
    output-decimal-marker = {\bookdecimalmarker}
  }%
  \newcommand{\booknum}[1]{\num{##1}}%
}{%
  \newcommand{\booknum}[1]{##1}%
}

\newcommand{\pyvalues}{}
\makeatletter
\newcommand{\pyval}[2]{\csgdef{pyval@#1}{#2}\listxadd{\pyvalues}{#1}}
% \num on a non-numeric body is a FATAL siunitx error. The generator knows
% whether it emitted a number or a word, so it writes \pyval for a number and
% \pyvaltext for a word, and this end only enforces which macro reads which.
\newcommand{\val}[1]{%
  \ifcsdef{pyval@#1}{%
    \ifcsdef{pyvaltext@#1}{%
      \PackageError{pybook}{Value `#1' is not a number}%
        {\string\val\space passes its body to siunitx, which refuses anything
         that is not a number. Use \string\valtext{#1} instead.}%
    }{}%
    \booknum{\csuse{pyval@#1}}%
  }{\textcolor{mred}{\textbf{[\lblValueMissing: #1]}}}%
}
\newcommand{\valtext}[1]{%
  \ifcsdef{pyval@#1}{\csuse{pyval@#1}}%
    {\textcolor{mred}{\textbf{[\lblValueMissing: #1]}}}%
}
\newcommand{\pyvaltext}[2]{%
  \csgdef{pyval@#1}{#2}\csgdef{pyvaltext@#1}{}\listxadd{\pyvalues}{#1}%
}
\makeatother

% figures/values/all.tex is generated by `make numbers` and does nothing but
% \input each value file. A build with no values yet still compiles; every
% \val{} in it prints a visible marker instead of a number.
\IfFileExists{figures/values/all.tex}{\input{figures/values/all}}{%
  \AtBeginDocument{\PackageWarning{pybook}{No computed values found. Run `make numbers`.}}}

% ============================================================
%  PINNED VERSIONS --- single source of truth
%  Change these here; they propagate through both editions and Appendix C.
%  Verified against pypi.org on \pinnedon (lang/<lang>.tex).
%  tools/check_versions.py compares every one of them with PyPI on each build.
%  Python itself is not on PyPI and is checked by hand against python.org.
% ============================================================
\newcommand{\bookedition}{0.1-dev}
\newcommand{\bookyear}{2026}
\newcommand{\pyver}{3.14}              % the minor the book targets
\newcommand{\pypatch}{3.14.7}          % the exact interpreter the code ran on
\newcommand{\uvver}{0.12.13}           % uv
\newcommand{\dotnetver}{10.0.100}      % dotnet
\newcommand{\ruffver}{0.16.7}          % ruff
\newcommand{\pyrightver}{1.1.414}      % pyright
\newcommand{\mypyver}{2.3.1}           % mypy
\newcommand{\pydanticver}{2.13.5}      % pydantic
\newcommand{\pysettingsver}{2.15.0}    % pydantic-settings
\newcommand{\fastapiver}{0.141.1}      % fastapi
\newcommand{\uvicornver}{0.53.0}       % uvicorn
\newcommand{\sqlalchemyver}{2.0.52}    % sqlalchemy
\newcommand{\alembicver}{1.20.0}       % alembic
\newcommand{\aiosqlitever}{0.22.1}     % aiosqlite
\newcommand{\pytestver}{9.1.1}         % pytest
\newcommand{\pytestasyncver}{1.4.0}    % pytest-asyncio
\newcommand{\hypothesisver}{6.168.0}   % hypothesis
\newcommand{\coveragever}{7.16.1}      % coverage
\newcommand{\httpxver}{0.28.1}         % httpx
\newcommand{\structlogver}{26.1.0}     % structlog
\newcommand{\otelver}{1.44.0}          % opentelemetry-sdk
\newcommand{\polarsver}{1.44.2}        % polars
\newcommand{\anthropicver}{1.5.0}      % anthropic
\newcommand{\openaiver}{3.13.0}        % openai
\newcommand{\pipauditver}{2.10.1}      % pip-audit

% ============================================================
%  INDEX
%  Two-column, and its entries include \texttt{} names that do not hyphenate.
%  Ragged right keeps multi-word entries off the margin; it cannot help a
%  single token wider than the column, so keep verbatim index entries under
%  about 29 characters and index longer names by concept instead.
%  \raggedcolumns is the switch imakeidx's multicols actually reads.
% ============================================================
\apptocmd{\theindex}{\raggedcolumns\raggedright}{}{%
  \PackageWarning{pybook}{Could not patch theindex for ragged setting}}
\providecommand*\see[2]{}
\renewcommand*\see[2]{\emph{\lblIndexSee} #1}
\providecommand*\seealso[2]{}
\renewcommand*\seealso[2]{\emph{\lblIndexSeeAlso} #1}

\makeindex
 still compiles and says what it did.
\providecommand{\booklang}{en}

\usepackage{etoolbox}   % loaded first: the language switch below needs it

% ---------- Page geometry ----------
% ONE format: A4, 12pt, ONE-SIDED. This book is read on a screen with an
% editor open beside it, and never printed as a bound book, so there is no
% recto, no verso, no gutter, no mirrored margin and no blank leaf before a
% chapter -- every one of those is a property of paper that a PDF viewer at
% half a screen's width turns into wasted pixels.
%
% The margins are symmetric and the measure is about seventy-five characters
% of prose, which is where the eye keeps the line return. The number that was
% actually chosen is the MONOSPACE measure: at \footnotesize on this block a
% listing line of 79 characters fits without wrapping, and 79 is the width
% every listing under code/ is held to by tools/check_structure.py --lines.
% A wrapped listing line is not an error TeX reports -- breaklines=true wraps
% it silently and prints an arrow into the middle of what the reader is meant
% to paste into an editor -- so the width is enforced in the source instead.
\usepackage[
  a4paper,
  left=3.1cm, right=3.1cm, top=2.6cm, bottom=3.0cm,
  head=26pt, headsep=0.8cm, footskip=1.4cm
]{geometry}

% ---------- Encoding & fonts ----------
\usepackage[T1]{fontenc}
\usepackage[utf8]{inputenc}
\IfFileExists{lmodern.sty}{\usepackage{lmodern}}{}
% newtxtext takes its TS1 symbols (\textcopyright, \textdegree) from TeX
% Gyre Termes, which Debian packages SEPARATELY (tex-gyre): a machine with
% newtx and without it dies on the copyright page with
% `Font TS1/ntxtlf/m/n/10.95=ts1-qtmr ... not loadable`. There is no
% inert-safe probe for a font metric in pdfTeX -- \IfFileExists searches
% TEXINPUTS and not TFMFONTS, so it reports the metric missing on a machine
% that has it, and \suppressfontnotfounderror is LuaTeX-only -- so nothing
% here guesses. CI's full TeX Live is the reference; locally, install
% tex-gyre. Both facts are recorded in CLAUDE.md under Build traps.
\IfFileExists{newtxtext.sty}{\usepackage{newtxtext}}{}
% amsmath before newtxmath, because newtxmath patches amsmath's internals.
% amssymb only when newtxmath is absent: newtxmath supplies the AMS symbols
% itself and loading both is a fatal clash (\Bbbk already defined) that is
% invisible on a machine without newtx and fatal on a full TeX Live.
\usepackage{amsmath}
\IfFileExists{newtxmath.sty}{\usepackage{newtxmath}}{\usepackage{amssymb}}
\IfFileExists{inconsolata.sty}{\usepackage[varqu,scaled=0.95]{inconsolata}}{}
\IfFileExists{lmodern.sty}{\usepackage{microtype}}{\usepackage[expansion=false]{microtype}}

% upquote is not a refinement. In T1 the input character ' is quoteright, so
% a listings body sets f('x') with U+2019 at both ends of the literal, and
% typed back in that is a SyntaxError. inconsolata's varqu option hides the
% defect on a machine that has inconsolata; upquote makes ' the ASCII
% apostrophe whatever the typewriter font is. In a book whose every listing is
% meant to be pasted into an editor, this line is load-bearing.
\IfFileExists{upquote.sty}{\usepackage{upquote}}{}

% And upquote's twin, for the same reason one character over. In T1 the pair
% -- is a ligature for an en dash, in EVERY family including the typewriter
% one, so \code{uv sync --locked} sets as "uv sync <en dash>locked" and a
% reader who copies it off the page gets an unknown-argument error from a
% command the chapter told them to run. It is invisible in review because it
% looks like a dash, and it bites exactly the flags a packaging chapter is
% made of. A `listings` body is already safe -- listings sets characters
% singly and no ligature forms -- which is what makes the defect so easy to
% miss: the same text is right inside a listing and wrong in the sentence
% above it.
%
% Scoped to tt* on purpose, and measured rather than assumed: prose keeps its
% own dashes, so \dash{} still sets an em dash and a range like 1--2 still
% sets an en dash. Probed against this preamble in all three positions before
% it was believed.
\DisableLigatures[-]{encoding = T1, family = tt*}

% ---------- Language ----------
% Loading babel with a language whose .ldf is absent is a *fatal* error, not
% a warning, so both the package and each language file are probed before
% either is requested. Both editions load the other language too, because the
% glossary rows and the cheat sheet carry words from both.
\IfFileExists{babel.sty}{%
  \IfFileExists{polish.ldf}{%
    \IfFileExists{british.ldf}{%
      \ifdefstring{\booklang}{pl}%
        {\usepackage[british,main=polish]{babel}}%
        {\usepackage[polish,main=british]{babel}}%
    }{\usepackage[polish]{babel}}%
  }{%
    \IfFileExists{british.ldf}{\usepackage[british]{babel}}{}%
  }%
}{}

% babel-polish makes " an active shorthand. listings reads its body verbatim
% and an active character in a verbatim context is a category-code accident
% waiting to happen, so the shorthand is turned off. Polish quotation marks
% come from \enquote, which csquotes sets per language.
\ifdefstring{\booklang}{pl}{\AtBeginDocument{\shorthandoff{"}}}{}

% ---------- Core packages ----------
\usepackage{graphicx}
\usepackage{xcolor}
\usepackage{booktabs}
\usepackage{tabularx}
% Appendix A's cheat sheet runs past a page in every part, and tabularx
% cannot break. longtable can, and works with booktabs.
\usepackage{longtable}
\usepackage{array}
\usepackage{enumitem}
\usepackage{listings}
\usepackage[most]{tcolorbox}
\usepackage{fancyhdr}
\usepackage{titlesec}
\usepackage{caption}
\usepackage{imakeidx}
% csquotes with autostyle follows babel's active language, so the identical
% source \enquote{x} sets 'x' in the English edition and the Polish quotation
% marks in the Polish one.
\IfFileExists{csquotes.sty}{\usepackage[autostyle=true]{csquotes}}{%
  \providecommand{\enquote}[1]{``##1''}}
\usepackage[hidelinks]{hyperref}

% ---------- Palette ----------
% The same family as the two companion volumes, so the three read as a set.
\definecolor{mblue}{HTML}{1F4E79}
\definecolor{mteal}{HTML}{0E7C7B}
\definecolor{mamber}{HTML}{B26A00}
\definecolor{mred}{HTML}{9B2C2C}
\definecolor{mviolet}{HTML}{7B4397}
\definecolor{mgrey}{HTML}{5A5A5A}
\definecolor{mlight}{HTML}{F4F6F8}
\definecolor{mfaint}{HTML}{FBFCFD}
\definecolor{codebg}{HTML}{FAFAFA}
\definecolor{codeframe}{HTML}{D8DEE4}
\definecolor{kw}{HTML}{0B5394}
\definecolor{str}{HTML}{9B2C2C}
\definecolor{cmt}{HTML}{4F7A4F}
\definecolor{typ}{HTML}{0E7C7B}
\definecolor{dec}{HTML}{7B4397}

\hypersetup{
  colorlinks=true,
  linkcolor=mblue,
  urlcolor=mteal,
  citecolor=mblue,
  pdfauthor={Konrad Cinkusz}
}

% \ifpl{Polish text}{English text} -- for the handful of places where a
% string is structural rather than content. Anything longer than a few words
% belongs in lang/ or in the edition's own source, not here.
% \DeclareRobustCommand, not \newcommand: a fragile \ifpl inside a moving
% argument is written to the .toc one level expanded and fails on the next run
% in exactly one of the two editions.
\DeclareRobustCommand{\ifpl}[2]{\ifdefstring{\booklang}{pl}{#1}{#2}}
\makeatletter
\pdfstringdefDisableCommands{%
  \ifdefstring{\booklang}{pl}{\let\ifpl\@firstoftwo}{\let\ifpl\@secondoftwo}%
  % \api{}, \pkg{} and \code{} are allowed in a section title, and a title
  % becomes a PDF bookmark. hyperref drops the \color command from a bookmark
  % string but keeps its ARGUMENT, so a title carrying \api{model\_validate}
  % made a bookmark reading "mtealmodel_validate" and warned "Token not
  % allowed" twice. The colour name is gobbled here instead. Measured on a
  % probe against this preamble before the line was added.
  \def\color#1{}%
}
\makeatother

% ---------- Language strings ----------
\input{lang/\booklang}

% The PDF's subject, keyword list and document language read \lblPdf... from
% the language file input on the line above, so they cannot sit in the colour
% \hypersetup block, which runs before it. Two blocks is the ordering, not an
% oversight; merging them raises Undefined control sequence and still writes
% a PDF, with the three fields empty.
\hypersetup{
  pdfsubject={\lblPdfSubject},
  pdfkeywords={\lblPdfKeywords},
  pdflang={\lblPdfLang}
}

% ---------- Headings ----------
% \raggedright on the title body is load-bearing: at \Huge on this block a
% title that cannot break overflows the margin by 30 pt or more, and several
% chapter titles here run past fifty characters.
\titleformat{\chapter}[display]
  {\normalfont\huge\bfseries\color{mblue}\raggedright}
  {\normalfont\normalsize\color{mgrey}\MakeUppercase{\chaptertitlename}\ \thechapter}
  {8pt}{\Huge\raggedright}
\titlespacing*{\chapter}{0pt}{10pt}{28pt}
\titleformat{\section}{\normalfont\Large\bfseries\color{mblue}\raggedright}{\thesection}{0.8em}{}
\titleformat{\subsection}{\normalfont\large\bfseries\color{mgrey}\raggedright}{\thesubsection}{0.7em}{}

% Sections in the contents, subsections not: fourteen chapters of subsections
% is a contents nobody reads, and the reader navigates by chapter and section.
\setcounter{tocdepth}{1}
\setcounter{secnumdepth}{2}

% One-sided running head: the chapter on the left, the folio on the right.
\pagestyle{fancy}
\fancyhf{}
\fancyhead[L]{\small\nouppercase{\leftmark}}
\fancyhead[R]{\small\thepage}
\renewcommand{\headrulewidth}{0.4pt}
\fancypagestyle{plain}{\fancyhf{}\fancyfoot[R]{\small\thepage}%
  \renewcommand{\headrulewidth}{0pt}}

% Drops the word "Chapter" from the head: the number is still there and the
% longer titles no longer overflow. MUST come after \pagestyle{fancy}, which
% installs its own \chaptermark and silently overwrites anything earlier.
\renewcommand{\chaptermark}[1]{\markboth{\thechapter.\ #1}{}}

% ============================================================
%  CODE LISTINGS
%  ---------------------------------------------------------
%  Every listing in a chapter is a FILE under code/, pulled in with \pyfile or
%  \pyregion, and CI runs it against the pinned versions. An in-source
%  \begin{python} block is allowed for a fragment of three or four lines that
%  exists to show a shape rather than to run; anything longer belongs in a
%  file the reader can open in the editor beside the PDF.
% ============================================================
\lstdefinelanguage{pythonx}{
  language=Python,
  morekeywords={
    async,await,match,case,nonlocal,yield,type,
    TypedDict,Annotated,Protocol,Literal,Generic,Final,ClassVar,Self,
    TypeVar,TypeGuard,TypeIs,Callable,Iterator,Iterable,Awaitable,
    dataclass,field,BaseModel,Field,Enum,StrEnum,NamedTuple,
    override,overload,property,staticmethod,classmethod
  },
  morekeywords=[2]{
    gather,create_task,TaskGroup,to_thread,timeout,wait_for,Semaphore,Queue,
    fixture,parametrize,monkeypatch,raises,mark,
    Depends,FastAPI,APIRouter,HTTPException,
    select,Session,Mapped,mapped_column,relationship,
    model_validate,model_validate_json,model_dump,model_json_schema
  },
  sensitive=true
}

% listings' C# definition predates async, records and pattern matching.
\lstdefinelanguage{csharpx}{
  language=[Sharp]C,
  morekeywords={
    async,await,var,dynamic,record,init,required,nameof,yield,partial,
    global,scoped,when,with,and,or,not,notnull,unmanaged,file,get,set,value,
    where,select,from,orderby,let,join,into,ascending,descending,on,equals,
    Task,ValueTask,CancellationToken,IEnumerable,IAsyncEnumerable,Func,Action
  },
  sensitive=true
}

\lstdefinestyle{bookcode}{
  basicstyle=\ttfamily\footnotesize,
  keywordstyle=\color{kw}\bfseries,
  keywordstyle=[2]\color{typ},
  stringstyle=\color{str},
  % Colour only, no \itshape: inconsolata (zi4) has no italic shape, so an
  % italic comment style is substituted with upright on every machine that
  % has the font, and the log carries "Font shape `T1/zi4/m/it' undefined"
  % on every build. The colour already marks a comment.
  commentstyle=\color{cmt},
  emphstyle=\color{dec}\bfseries,
  numbers=left,
  numberstyle=\ttfamily\tiny\color{mgrey},
  numbersep=8pt,
  backgroundcolor=\color{codebg},
  frame=single,
  rulecolor=\color{codeframe},
  framesep=6pt,
  breaklines=true,
  breakatwhitespace=false,
  postbreak=\mbox{\textcolor{mgrey}{$\hookrightarrow$}\space},
  showstringspaces=false,
  tabsize=4,
  columns=fullflexible,
  keepspaces=true,
  captionpos=b,
  aboveskip=1.2em,
  belowskip=1.2em,
  % Maps the characters most likely to be pasted into a listing by accident.
  %
  % Do NOT add a mapping for U+00A0 (non-breaking space). It looks like the
  % obvious next entry and it makes `listings` abort the run with "Improper
  % alphabetic constant" -- fatal, no PDF, and the message names neither the
  % character nor this line. The ASCII-in-listings convention, checked by
  % parity's C13, is the guard against it instead.
  literate={ł}{{\l}}1 {Ł}{{\L}}1 {ą}{{\k a}}1 {ę}{{\k e}}1
           {ś}{{\'s}}1 {ć}{{\'c}}1 {ż}{{\.z}}1 {ź}{{\'z}}1 {ó}{{\'o}}1 {ń}{{\'n}}1
           {Ą}{{\k A}}1 {Ę}{{\k E}}1 {Ś}{{\'S}}1 {Ć}{{\'C}}1
           {Ż}{{\.Z}}1 {Ź}{{\'Z}}1 {Ó}{{\'O}}1 {Ń}{{\'N}}1
           {—}{{-{}-}}2 {–}{{-}}1 {‘}{{'}}1 {’}{{'}}1
           {“}{{"}}1 {”}{{"}}1 {°}{{\textdegree}}1
}

\lstset{style=bookcode}

\lstnewenvironment{python}[1][]
  {\lstset{language=pythonx,style=bookcode,#1}}
  {}

% C# comparison snippets. The whole book is written for engineers arriving
% from .NET, so listings come in pairs often enough that the C# half has its
% own environment rather than a generic one.
\lstnewenvironment{csharp}[1][]
  {\lstset{language=csharpx,style=bookcode,#1}}
  {}

\lstnewenvironment{shellcmd}[1][]
  {\lstset{language=bash,style=bookcode,numbers=none,keywordstyle=\color{mteal}\bfseries,#1}}
  {}

\lstnewenvironment{tomlcode}[1][]
  {\lstset{style=bookcode,numbers=none,keywordstyle=\color{black},#1}}
  {}

\lstnewenvironment{yamlcode}[1][]
  {\lstset{style=bookcode,numbers=none,keywordstyle=\color{black},#1}}
  {}

\lstnewenvironment{jsoncode}[1][]
  {\lstset{style=bookcode,numbers=none,keywordstyle=\color{black},#1}}
  {}

% Terminal transcripts typed into the source. Allowed for a shell session the
% reader is meant to reproduce; NOT for program output, which is \transcript
% below, because an in-source transcript nobody ran reads exactly like one
% that was, and that is where a remembered number hides.
\lstnewenvironment{console}[1][]
  {\lstset{style=bookcode,numbers=none,basicstyle=\ttfamily\footnotesize,
           keywordstyle=\color{black},backgroundcolor=\color{white},#1}}
  {}

% The path is printed above every file listing, in the repository-relative
% form the reader opens in the editor beside the PDF. That is the whole point
% of a listing being a file, and the front matter promises it. The \nobreak
% keeps the path on the same page as the listing it names.
% ---------------------------------------------------------------------------
%  Appendix B prints the trap catalogue from notes/02-traps.md. One entry is
%  its number -- which a chapter may cite, and which is never reused -- the
%  habit in the reader's own voice, and what Python does instead.
%
%  The body is set RAGGED RIGHT on purpose. An entry is mostly \code{} spans,
%  which are unbreakable runs, and sixty-seven justified paragraphs of them
%  is sixty-seven chances at an overfull hbox in the edition with the longer
%  words. Ragged right removes the class rather than fixing it one box at a
%  time, and a reference list is the one place in this book where it reads
%  better anyway.
%
%  The habit is set with \enquote{} rather than \emph{}, and that is not a
%  style choice. Every habit here carries \code{} spans, and \code{} inside
%  \emph{} asks inconsolata for an italic it does not have: the first build
%  of this appendix raised `Font shape T1/zi4/m/it undefined', which
%  checklog.py is hard on for exactly this reason. It is the trap CLAUDE.md
%  records, met sixty-seven times at once -- and fixed once, because the
%  entries go through a macro. A quotation is also what the habit IS, and
%  \enquote{} is what the Polish edition owes anyway.
% ---------------------------------------------------------------------------
%  Appendix A's cheat sheet: one table per part, C# on the left, Python on
%  the right, and the chapter that teaches the row on the end. The three
%  header strings are arguments so that one environment serves both
%  editions and the structural signature stays identical.
%
%  Columns are ragged right for the reason Appendix B's entries are: a cell
%  is mostly \code{} spans, which are unbreakable runs, and a justified
%  narrow column full of them is an overfull box waiting for the edition
%  with the longer words.
% ---------------------------------------------------------------------------
%  Appendix C's tool matrix: the .NET tool, the Python tool this book uses,
%  the version it was verified against and the chapter that introduces it.
%
%  The version column prints \pydanticver and its siblings and is never
%  typed, which is the whole point of the appendix: it is the one page where
%  a reader sees every pin at once, and a pin that had to be re-typed here
%  would be a second copy of a fact the preamble already owns.
%  tools/check_structure.py --tools fails when a pin macro exists and this
%  table does not print it.
\NewDocumentEnvironment{toolmatrix}{mmmm}{%
  \small
  \begin{longtable}{@{}%
    >{\raggedright\arraybackslash}p{0.27\linewidth}%
    >{\raggedright\arraybackslash}p{0.31\linewidth}%
    >{\raggedright\arraybackslash}p{0.20\linewidth}%
    >{\raggedright\arraybackslash}p{0.08\linewidth}@{}}
  \toprule
  \textbf{#1} & \textbf{#2} & \textbf{#3} & \textbf{#4} \\
  \midrule
  \endfirsthead
  \toprule
  \textbf{#1} & \textbf{#2} & \textbf{#3} & \textbf{#4} \\
  \midrule
  \endhead
  \midrule
  \endfoot
  \bottomrule
  \endlastfoot
}{%
  \end{longtable}%
}

\NewDocumentEnvironment{cheatsheet}{mmm}{%
  \small
  \begin{longtable}{@{}%
    >{\raggedright\arraybackslash}p{0.41\linewidth}%
    >{\raggedright\arraybackslash}p{0.41\linewidth}%
    >{\raggedright\arraybackslash}p{0.08\linewidth}@{}}
  \toprule
  \textbf{#1} & \textbf{#2} & \textbf{#3} \\
  \midrule
  \endfirsthead
  \toprule
  \textbf{#1} & \textbf{#2} & \textbf{#3} \\
  \midrule
  \endhead
  \midrule
  \endfoot
  \bottomrule
  \endlastfoot
}{%
  \end{longtable}%
}

\newcommand{\trapentry}[3]{%
  \par\medskip
  \noindent\textbf{#1.}\enspace\enquote{#2}\par
  \nopagebreak
  {\raggedright\noindent#3\par}%
}

\newcommand{\listingpath}[1]{%
  \par\medskip\noindent
  {\ttfamily\footnotesize\color{mgrey}\detokenize{#1}}%
  \par\nobreak}

% External source file -- the preferred form. The path is relative to the
% repository root, which is also the path the reader opens in the editor.
% Usage: \pyfile{code/ch05/decorators.py}{Caption}{lst:label}
\newcommand{\pyfile}[3]{%
  \listingpath{#1}%
  \lstinputlisting[language=pythonx,style=bookcode,aboveskip=0.3em,
    caption={#2},label={#3}]{#1}%
}

% A named region of an external file, delimited in the source by
% `# --8<-- [start:name]` / `# --8<-- [end:name]`, so the book and the running
% code cannot drift apart. tools/check_structure.py --listings fails when the
% file or the region is missing.
%
% The markers are matched through listings' range prefix and suffix keys,
% with linerange={name-name}. The form this was inherited with -- a linerange
% whose two ends were the whole marker strings -- matched nothing and printed
% the WHOLE FILE, with no warning: measured on a three-region probe file
% before this was rewritten, and the same definition sits unused in the
% sibling book it came from. A region name is a word, never a number, because
% a number is a line number to listings. The line numbers printed are the
% file's own, so a region's listing says where in the file it sits.
% Usage: \pyregion{code/ch05/decorators.py}{timing}{Caption}{lst:label}
\lstset{
  rangebeginprefix=\#\ --8<--\ [start:, rangebeginsuffix=],
  rangeendprefix=\#\ --8<--\ [end:, rangeendsuffix=],
  includerangemarker=false}
\newcommand{\pyregion}[4]{%
  \listingpath{#1}%
  \lstinputlisting[language=pythonx,style=bookcode,aboveskip=0.3em,
    linerange={#2-#2},caption={#3},label={#4}]{#1}%
}

% The C# twin of \pyfile, for the comparison half of a pair.
\newcommand{\csfile}[3]{%
  \listingpath{#1}%
  \lstinputlisting[language=csharpx,style=bookcode,aboveskip=0.3em,
    caption={#2},label={#3}]{#1}%
}

% A program's output, written by a script under code/measure/ and pulled in
% whole. There is no typed form of this and there must not be: \val{} does
% not expand inside a verbatim body, so a transcript typed into a chapter is
% one that cannot quote a computed value, and a transcript nobody ran is
% indistinguishable on the page from one that was.
\newcommand{\transcript}[1]{%
  \par\smallskip
  \IfFileExists{figures/transcripts/#1.txt}{%
    \lstinputlisting[style=bookcode,numbers=none,basicstyle=\ttfamily\footnotesize,
      keywordstyle=\color{black},backgroundcolor=\color{white},
      language={}]{figures/transcripts/#1.txt}%
  }{%
    \begin{tcolorbox}[enhanced, sharp corners, width=\linewidth,
      colback=mlight, colframe=mgrey, boxrule=0.6pt,
      left=8pt, right=8pt, top=6pt, bottom=6pt]
      {\bfseries\color{mgrey}\small \lblTranscriptMissing}\\[2pt]
      {\ttfamily\scriptsize\color{mgrey} figures/transcripts/#1.txt}
    \end{tcolorbox}%
  }%
  \par\smallskip
}

% Inline literals. Underscores inside these must be written \_.
\newcommand{\code}[1]{\texttt{\small #1}}
\newcommand{\pkg}[1]{\texttt{\small\color{mblue}#1}}
\newcommand{\api}[1]{\texttt{\small\color{mteal}#1}}

% ============================================================
%  ADMONITIONS
%  ---------------------------------------------------------
%  A deliberately small, fixed vocabulary. Two visual treatments carry the
%  whole grammar: `tinted` is a block the reader may read in passing -- a
%  tint and a thick left bar, no outline; `outlined` is a block the reader
%  must stop at, with a full rule on white. A third treatment would mean the
%  first two are not carrying their meaning.
% ============================================================
\tcbset{
  mfa tinted/.style={
    enhanced, breakable, sharp corners,
    colback=mlight, boxrule=0pt, leftrule=3pt,
    left=8pt, right=8pt, top=6pt, bottom=6pt,
    before skip=13pt, after skip=13pt,
    fonttitle=\bfseries\small,
    coltitle=black, colbacktitle=mlight,
    attach title to upper={\par\smallskip},
  },
  mfa outlined/.style={
    enhanced, breakable, sharp corners,
    colback=white, boxrule=0.6pt,
    left=8pt, right=8pt, top=6pt, bottom=6pt,
    before skip=13pt, after skip=13pt,
    fonttitle=\bfseries\small,
    coltitle=black, colbacktitle=white,
    attach title to upper={\par\smallskip},
  },
}

% Read in passing.
\newtcolorbox{note}[1][]{mfa tinted, colframe=mblue, title={\lblNote}, #1}
\newtcolorbox{warning}[1][]{mfa tinted, colframe=mred, title={\lblWarning}, #1}
% The translation. The .NET reader is the book's only assumed reader, so this
% box carries the C# side of a comparison and nothing else. Its use is
% counted: a chapter with none is a chapter that has stopped translating.
\newtcolorbox{csbox}[1][]{mfa tinted, colframe=mteal, title={\lblCsharp}, #1}
% Anything that is preview, experimental or likely to move. Python moves more
% slowly than the agent frameworks do, and this box should be rare here.
\newtcolorbox{versionbox}[1][]{mfa tinted, colframe=mamber, title={\lblVersion}, #1}

% Stop and read.
% The misconception, named AFTER the reader has walked into it. Appendix B is
% the catalogue; a trap box in a chapter is one entry of it, in place.
\newtcolorbox{trapbox}[1][]{mfa outlined, colframe=mred, title={\lblTrap}, #1}
% Marks a listing that has not been executed against the pinned versions.
% `make debt` counts these, CI prints the count, and nothing ships to a reader
% with one attached. Removing it means you ran the code.
\newtcolorbox{verifybox}[1][]{mfa outlined, colframe=mgrey, title={\lblVerify}, #1}
\newtcolorbox{exercisebox}[1][]{mfa outlined, colframe=mteal, title={\lblExercise}, #1}
% The guiding project. Points at the stage of trace-assert this section
% builds, under code/src/trace_assert/.
\newtcolorbox{projectbox}[1][]{mfa outlined, colframe=mamber, title={\lblProject}, #1}

% ============================================================
%  EXERCISES
%  ---------------------------------------------------------
%  \begin{exercise}{key}{title} ... \end{exercise}
%
%  This is the mechanism the whole book is read through: the PDF on one half
%  of the screen, an editor on the other. An exercise is therefore not a
%  paragraph but a FILE the reader opens -- a starter under
%  code/exercises/chNN/<key>.py -- and a TEST that decides whether it is done,
%  test_<key>.py beside it. The box prints both paths and the one command
%  that runs the test, and the manifest at the back lists every exercise.
%
%  An ENVIRONMENT, not a macro with a body argument: a macro argument cannot
%  contain verbatim material, so an exercise written as \exercise{key}{title}
%  {body} could never show the reader a code fragment, and the first chapter
%  that tried would have died on `\begin{python}` inside the braces.
%
%  The key is the file stem, with underscores, and it is written with
%  \detokenize so the underscore reaches the page as a character rather than
%  as a subscript. tools/check_structure.py --exercises fails when the starter
%  or the test is missing, or when the key's chapter prefix disagrees with the
%  chapter it sits in.
% ============================================================
\newcounter{exercise}[chapter]
\renewcommand{\theexercise}{\thechapter.\arabic{exercise}}
\makeatletter
\newcommand{\exercisedir}{ch\two@digits{\value{chapter}}}
\newcommand{\listofexercises}{%
  \section*{\lblExerciseManifest}%
  \phantomsection\addcontentsline{toc}{section}{\lblExerciseManifest}%
  % The entries are subsection-level and tocdepth is 1 for the contents,
  % so the depth is raised for this list alone or it prints empty --
  % which it did, on the scaffold's first clean build.
  \begingroup\c@tocdepth=2\relax\@starttoc{exr}\endgroup%
}
\makeatother
\NewDocumentEnvironment{exercise}{mm}{%
  \refstepcounter{exercise}%
  \addcontentsline{exr}{subsection}{\protect\numberline{\theexercise}\texttt{\protect\detokenize{#1}} --- #2}%
  \begin{exercisebox}[title={\lblExercise~\theexercise\ \dash{} #2}]%
}{%
  \par\medskip
  {\small\setlength{\parindent}{0pt}%
  \lblExerciseStarter: \texttt{code/exercises/\exercisedir/\detokenize{#1}.py}\\
  \lblExerciseTest: \texttt{code/exercises/\exercisedir/test\_\detokenize{#1}.py}\\
  \lblExerciseRun: \texttt{cd code \&\& uv run pytest -k \detokenize{#1}}\par}%
  \end{exercisebox}%
}

% ============================================================
%  DIAGRAMS --- Mermaid pipeline
%  ---------------------------------------------------------
%  \mermaidfig{key}{caption}{one-line manifest description}
%
%  Source of truth is figures/mermaid/<lang>/<key>.mmd, committed. `make
%  diagrams` renders each to figures/diagrams/<lang>/<key>.pdf, which is build
%  output and is not committed, so a diagram change reviews as a text diff.
%  The source is per-language because a diagram with English labels in the
%  Polish edition is exactly the half-translation that makes a bilingual book
%  feel machine-made; CI checks that both languages carry the same keys.
%
%  If the rendered PDF is absent the macro draws a placeholder AND typesets
%  the Mermaid source, in the flow rather than as a float, because a source
%  typeset verbatim is often taller than a page and a float cannot break.
% ============================================================
\makeatletter
\newcommand{\listofdiagrams}{%
  \section*{\lblDiagramManifest}%
  \phantomsection\addcontentsline{toc}{section}{\lblDiagramManifest}%
  % The entries are subsection-level and tocdepth is 1 for the contents,
  % so the depth is raised for this list alone or it prints empty --
  % which it did, on the scaffold's first clean build.
  \begingroup\c@tocdepth=2\relax\@starttoc{dgm}\endgroup%
}
\makeatother

\newcommand{\mermaidfig}[3]{%
  % \phantomsection is load-bearing and not decoration. \addcontentsline
  % records hyperref's CURRENT anchor, and here that is whatever preceded
  % the figure -- for a figure placed after a listing it is the listing's
  % line-number anchor, and listings has already advanced the counter past
  % the last line it printed, so the manifest entry linked to a
  % destination that is never created: "name{lstnumber.10.4.46} has been
  % referenced but does not exist, replaced by a fixed one". The two
  % front-matter figures had the same wrong target and were silent about
  % it, because the lines they named happen to be printed. The exercise
  % manifest never had the bug, because \begin{exercise} steps a counter
  % and that is what sets the anchor.
  \phantomsection%
  \addcontentsline{dgm}{subsection}{\protect\numberline{}\texttt{#1.mmd} --- #3}%
  \IfFileExists{figures/diagrams/\booklang/#1.pdf}{%
    \begin{figure}[htbp]
    \centering
    \includegraphics[width=0.95\linewidth,height=0.42\textheight,keepaspectratio]{figures/diagrams/\booklang/#1.pdf}%
    \caption{#2}\label{fig:#1}
    \end{figure}%
  }{%
    \par\medskip
    \begin{tcolorbox}[
      enhanced, breakable, width=\linewidth, sharp corners,
      colback=white, colframe=mgrey, boxrule=0.8pt,
      left=8pt, right=8pt, top=8pt, bottom=8pt]
      {\bfseries\color{mgrey}\small \lblNotRendered}\\[4pt]
      {\ttfamily\scriptsize\color{mgrey} figures/mermaid/\booklang/#1.mmd}\\[6pt]
      \IfFileExists{figures/mermaid/\booklang/#1.mmd}{%
        \lstinputlisting[style=bookcode,numbers=none,basicstyle=\ttfamily\scriptsize,
                         backgroundcolor=\color{mlight},frame=none,
                         keywordstyle=\color{black},language={}]{figures/mermaid/\booklang/#1.mmd}%
      }{%
        {\small\color{mred} \lblSourceMissing: \texttt{figures/mermaid/\booklang/#1.mmd}}%
      }%
    \end{tcolorbox}%
    \captionof{figure}{#2}\label{fig:#1}
    \par\medskip
  }%
}

% ============================================================
%  CHAPTER STUBS
%  ---------------------------------------------------------
%  Prints a visible marker so an unwritten chapter cannot be mistaken for a
%  written one and the page count never flatters the draft. `make debt`
%  counts them, and CI publishes the count on every build.
% ============================================================
\newcommand{\chapterstub}[1]{%
  \begin{tcolorbox}[
    enhanced, breakable, width=\linewidth, sharp corners,
    colback=mlight, colframe=mred, boxrule=1pt,
    borderline={0.6pt}{3pt}{mred, dashed},
    left=12pt, right=12pt, top=12pt, bottom=12pt]
    {\bfseries\color{mred}\large \lblNotWritten}\\[8pt]
    \small\raggedright #1
  \end{tcolorbox}%
}

% ============================================================
%  COMPUTED VALUES
%  ---------------------------------------------------------
%  A number that came out of a measurement is written by a script under
%  code/measure/ into figures/values/<name>.tex as
%
%      \pyval{e01.threads.speedup}{3.71}
%
%  and pulled into the text with \val{e01.threads.speedup}. The script is the
%  source of truth and the book cannot disagree with it, because the book does
%  not contain the digits. `make verify` fails when a committed value no
%  longer matches its script. A value the scripts do not produce prints a
%  visible marker.
%
%  \val also applies the edition's decimal separator, which is how the Polish
%  edition gets its decimal comma without a translator having to remember.
% ============================================================
\newcommand{\bookdecimalmarker}{\ifpl{,}{.}}
\IfFileExists{siunitx.sty}{%
  \usepackage{siunitx}%
  \sisetup{
    group-digits = integer,
    group-minimum-digits = 5,
    group-separator = {\,},
    output-decimal-marker = {\bookdecimalmarker}
  }%
  \newcommand{\booknum}[1]{\num{##1}}%
}{%
  \newcommand{\booknum}[1]{##1}%
}

\newcommand{\pyvalues}{}
\makeatletter
\newcommand{\pyval}[2]{\csgdef{pyval@#1}{#2}\listxadd{\pyvalues}{#1}}
% \num on a non-numeric body is a FATAL siunitx error. The generator knows
% whether it emitted a number or a word, so it writes \pyval for a number and
% \pyvaltext for a word, and this end only enforces which macro reads which.
\newcommand{\val}[1]{%
  \ifcsdef{pyval@#1}{%
    \ifcsdef{pyvaltext@#1}{%
      \PackageError{pybook}{Value `#1' is not a number}%
        {\string\val\space passes its body to siunitx, which refuses anything
         that is not a number. Use \string\valtext{#1} instead.}%
    }{}%
    \booknum{\csuse{pyval@#1}}%
  }{\textcolor{mred}{\textbf{[\lblValueMissing: #1]}}}%
}
\newcommand{\valtext}[1]{%
  \ifcsdef{pyval@#1}{\csuse{pyval@#1}}%
    {\textcolor{mred}{\textbf{[\lblValueMissing: #1]}}}%
}
\newcommand{\pyvaltext}[2]{%
  \csgdef{pyval@#1}{#2}\csgdef{pyvaltext@#1}{}\listxadd{\pyvalues}{#1}%
}
\makeatother

% figures/values/all.tex is generated by `make numbers` and does nothing but
% \input each value file. A build with no values yet still compiles; every
% \val{} in it prints a visible marker instead of a number.
\IfFileExists{figures/values/all.tex}{\input{figures/values/all}}{%
  \AtBeginDocument{\PackageWarning{pybook}{No computed values found. Run `make numbers`.}}}

% ============================================================
%  PINNED VERSIONS --- single source of truth
%  Change these here; they propagate through both editions and Appendix C.
%  Verified against pypi.org on \pinnedon (lang/<lang>.tex).
%  tools/check_versions.py compares every one of them with PyPI on each build.
%  Python itself is not on PyPI and is checked by hand against python.org.
% ============================================================
\newcommand{\bookedition}{0.1-dev}
\newcommand{\bookyear}{2026}
\newcommand{\pyver}{3.14}              % the minor the book targets
\newcommand{\pypatch}{3.14.7}          % the exact interpreter the code ran on
\newcommand{\uvver}{0.12.13}           % uv
\newcommand{\dotnetver}{10.0.100}      % dotnet
\newcommand{\ruffver}{0.16.7}          % ruff
\newcommand{\pyrightver}{1.1.414}      % pyright
\newcommand{\mypyver}{2.3.1}           % mypy
\newcommand{\pydanticver}{2.13.5}      % pydantic
\newcommand{\pysettingsver}{2.15.0}    % pydantic-settings
\newcommand{\fastapiver}{0.141.1}      % fastapi
\newcommand{\uvicornver}{0.53.0}       % uvicorn
\newcommand{\sqlalchemyver}{2.0.52}    % sqlalchemy
\newcommand{\alembicver}{1.20.0}       % alembic
\newcommand{\aiosqlitever}{0.22.1}     % aiosqlite
\newcommand{\pytestver}{9.1.1}         % pytest
\newcommand{\pytestasyncver}{1.4.0}    % pytest-asyncio
\newcommand{\hypothesisver}{6.168.0}   % hypothesis
\newcommand{\coveragever}{7.16.1}      % coverage
\newcommand{\httpxver}{0.28.1}         % httpx
\newcommand{\structlogver}{26.1.0}     % structlog
\newcommand{\otelver}{1.44.0}          % opentelemetry-sdk
\newcommand{\polarsver}{1.44.2}        % polars
\newcommand{\anthropicver}{1.5.0}      % anthropic
\newcommand{\openaiver}{3.13.0}        % openai
\newcommand{\pipauditver}{2.10.1}      % pip-audit

% ============================================================
%  INDEX
%  Two-column, and its entries include \texttt{} names that do not hyphenate.
%  Ragged right keeps multi-word entries off the margin; it cannot help a
%  single token wider than the column, so keep verbatim index entries under
%  about 29 characters and index longer names by concept instead.
%  \raggedcolumns is the switch imakeidx's multicols actually reads.
% ============================================================
\apptocmd{\theindex}{\raggedcolumns\raggedright}{}{%
  \PackageWarning{pybook}{Could not patch theindex for ragged setting}}
\providecommand*\see[2]{}
\renewcommand*\see[2]{\emph{\lblIndexSee} #1}
\providecommand*\seealso[2]{}
\renewcommand*\seealso[2]{\emph{\lblIndexSeeAlso} #1}

\makeindex
; nothing else differs between them at this level.
%  Every user-visible string lives in lang/en.tex or lang/pl.tex, so a label
%  that appears in one edition and not the other is a missing macro rather
%  than a silent divergence.
%
%  Lineage. The bilingual wiring (\booklang, lang/, body.tex, the generated
%  structure) is the math book's; the chapter machinery (listings, the
%  admonitions, verifybox, \pyfile and \pyregion, the Mermaid pipeline) is the
%  LangChain book's. What is new here is the exercise macro, which ties every
%  exercise to a starter file and a test under code/, because this book is
%  read with an editor open beside it.
% ============================================================

% \booklang must be set by the main file. Default to English rather than
% failing, so a stray % ============================================================
%  preamble.tex --- styling, environments, macros
%
%  Shared by BOTH editions. main-en.tex and main-pl.tex each define \booklang
%  before % ============================================================
%  preamble.tex --- styling, environments, macros
%
%  Shared by BOTH editions. main-en.tex and main-pl.tex each define \booklang
%  before \input{preamble}; nothing else differs between them at this level.
%  Every user-visible string lives in lang/en.tex or lang/pl.tex, so a label
%  that appears in one edition and not the other is a missing macro rather
%  than a silent divergence.
%
%  Lineage. The bilingual wiring (\booklang, lang/, body.tex, the generated
%  structure) is the math book's; the chapter machinery (listings, the
%  admonitions, verifybox, \pyfile and \pyregion, the Mermaid pipeline) is the
%  LangChain book's. What is new here is the exercise macro, which ties every
%  exercise to a starter file and a test under code/, because this book is
%  read with an editor open beside it.
% ============================================================

% \booklang must be set by the main file. Default to English rather than
% failing, so a stray \input{preamble} still compiles and says what it did.
\providecommand{\booklang}{en}

\usepackage{etoolbox}   % loaded first: the language switch below needs it

% ---------- Page geometry ----------
% ONE format: A4, 12pt, ONE-SIDED. This book is read on a screen with an
% editor open beside it, and never printed as a bound book, so there is no
% recto, no verso, no gutter, no mirrored margin and no blank leaf before a
% chapter -- every one of those is a property of paper that a PDF viewer at
% half a screen's width turns into wasted pixels.
%
% The margins are symmetric and the measure is about seventy-five characters
% of prose, which is where the eye keeps the line return. The number that was
% actually chosen is the MONOSPACE measure: at \footnotesize on this block a
% listing line of 79 characters fits without wrapping, and 79 is the width
% every listing under code/ is held to by tools/check_structure.py --lines.
% A wrapped listing line is not an error TeX reports -- breaklines=true wraps
% it silently and prints an arrow into the middle of what the reader is meant
% to paste into an editor -- so the width is enforced in the source instead.
\usepackage[
  a4paper,
  left=3.1cm, right=3.1cm, top=2.6cm, bottom=3.0cm,
  head=26pt, headsep=0.8cm, footskip=1.4cm
]{geometry}

% ---------- Encoding & fonts ----------
\usepackage[T1]{fontenc}
\usepackage[utf8]{inputenc}
\IfFileExists{lmodern.sty}{\usepackage{lmodern}}{}
% newtxtext takes its TS1 symbols (\textcopyright, \textdegree) from TeX
% Gyre Termes, which Debian packages SEPARATELY (tex-gyre): a machine with
% newtx and without it dies on the copyright page with
% `Font TS1/ntxtlf/m/n/10.95=ts1-qtmr ... not loadable`. There is no
% inert-safe probe for a font metric in pdfTeX -- \IfFileExists searches
% TEXINPUTS and not TFMFONTS, so it reports the metric missing on a machine
% that has it, and \suppressfontnotfounderror is LuaTeX-only -- so nothing
% here guesses. CI's full TeX Live is the reference; locally, install
% tex-gyre. Both facts are recorded in CLAUDE.md under Build traps.
\IfFileExists{newtxtext.sty}{\usepackage{newtxtext}}{}
% amsmath before newtxmath, because newtxmath patches amsmath's internals.
% amssymb only when newtxmath is absent: newtxmath supplies the AMS symbols
% itself and loading both is a fatal clash (\Bbbk already defined) that is
% invisible on a machine without newtx and fatal on a full TeX Live.
\usepackage{amsmath}
\IfFileExists{newtxmath.sty}{\usepackage{newtxmath}}{\usepackage{amssymb}}
\IfFileExists{inconsolata.sty}{\usepackage[varqu,scaled=0.95]{inconsolata}}{}
\IfFileExists{lmodern.sty}{\usepackage{microtype}}{\usepackage[expansion=false]{microtype}}

% upquote is not a refinement. In T1 the input character ' is quoteright, so
% a listings body sets f('x') with U+2019 at both ends of the literal, and
% typed back in that is a SyntaxError. inconsolata's varqu option hides the
% defect on a machine that has inconsolata; upquote makes ' the ASCII
% apostrophe whatever the typewriter font is. In a book whose every listing is
% meant to be pasted into an editor, this line is load-bearing.
\IfFileExists{upquote.sty}{\usepackage{upquote}}{}

% And upquote's twin, for the same reason one character over. In T1 the pair
% -- is a ligature for an en dash, in EVERY family including the typewriter
% one, so \code{uv sync --locked} sets as "uv sync <en dash>locked" and a
% reader who copies it off the page gets an unknown-argument error from a
% command the chapter told them to run. It is invisible in review because it
% looks like a dash, and it bites exactly the flags a packaging chapter is
% made of. A `listings` body is already safe -- listings sets characters
% singly and no ligature forms -- which is what makes the defect so easy to
% miss: the same text is right inside a listing and wrong in the sentence
% above it.
%
% Scoped to tt* on purpose, and measured rather than assumed: prose keeps its
% own dashes, so \dash{} still sets an em dash and a range like 1--2 still
% sets an en dash. Probed against this preamble in all three positions before
% it was believed.
\DisableLigatures[-]{encoding = T1, family = tt*}

% ---------- Language ----------
% Loading babel with a language whose .ldf is absent is a *fatal* error, not
% a warning, so both the package and each language file are probed before
% either is requested. Both editions load the other language too, because the
% glossary rows and the cheat sheet carry words from both.
\IfFileExists{babel.sty}{%
  \IfFileExists{polish.ldf}{%
    \IfFileExists{british.ldf}{%
      \ifdefstring{\booklang}{pl}%
        {\usepackage[british,main=polish]{babel}}%
        {\usepackage[polish,main=british]{babel}}%
    }{\usepackage[polish]{babel}}%
  }{%
    \IfFileExists{british.ldf}{\usepackage[british]{babel}}{}%
  }%
}{}

% babel-polish makes " an active shorthand. listings reads its body verbatim
% and an active character in a verbatim context is a category-code accident
% waiting to happen, so the shorthand is turned off. Polish quotation marks
% come from \enquote, which csquotes sets per language.
\ifdefstring{\booklang}{pl}{\AtBeginDocument{\shorthandoff{"}}}{}

% ---------- Core packages ----------
\usepackage{graphicx}
\usepackage{xcolor}
\usepackage{booktabs}
\usepackage{tabularx}
% Appendix A's cheat sheet runs past a page in every part, and tabularx
% cannot break. longtable can, and works with booktabs.
\usepackage{longtable}
\usepackage{array}
\usepackage{enumitem}
\usepackage{listings}
\usepackage[most]{tcolorbox}
\usepackage{fancyhdr}
\usepackage{titlesec}
\usepackage{caption}
\usepackage{imakeidx}
% csquotes with autostyle follows babel's active language, so the identical
% source \enquote{x} sets 'x' in the English edition and the Polish quotation
% marks in the Polish one.
\IfFileExists{csquotes.sty}{\usepackage[autostyle=true]{csquotes}}{%
  \providecommand{\enquote}[1]{``##1''}}
\usepackage[hidelinks]{hyperref}

% ---------- Palette ----------
% The same family as the two companion volumes, so the three read as a set.
\definecolor{mblue}{HTML}{1F4E79}
\definecolor{mteal}{HTML}{0E7C7B}
\definecolor{mamber}{HTML}{B26A00}
\definecolor{mred}{HTML}{9B2C2C}
\definecolor{mviolet}{HTML}{7B4397}
\definecolor{mgrey}{HTML}{5A5A5A}
\definecolor{mlight}{HTML}{F4F6F8}
\definecolor{mfaint}{HTML}{FBFCFD}
\definecolor{codebg}{HTML}{FAFAFA}
\definecolor{codeframe}{HTML}{D8DEE4}
\definecolor{kw}{HTML}{0B5394}
\definecolor{str}{HTML}{9B2C2C}
\definecolor{cmt}{HTML}{4F7A4F}
\definecolor{typ}{HTML}{0E7C7B}
\definecolor{dec}{HTML}{7B4397}

\hypersetup{
  colorlinks=true,
  linkcolor=mblue,
  urlcolor=mteal,
  citecolor=mblue,
  pdfauthor={Konrad Cinkusz}
}

% \ifpl{Polish text}{English text} -- for the handful of places where a
% string is structural rather than content. Anything longer than a few words
% belongs in lang/ or in the edition's own source, not here.
% \DeclareRobustCommand, not \newcommand: a fragile \ifpl inside a moving
% argument is written to the .toc one level expanded and fails on the next run
% in exactly one of the two editions.
\DeclareRobustCommand{\ifpl}[2]{\ifdefstring{\booklang}{pl}{#1}{#2}}
\makeatletter
\pdfstringdefDisableCommands{%
  \ifdefstring{\booklang}{pl}{\let\ifpl\@firstoftwo}{\let\ifpl\@secondoftwo}%
  % \api{}, \pkg{} and \code{} are allowed in a section title, and a title
  % becomes a PDF bookmark. hyperref drops the \color command from a bookmark
  % string but keeps its ARGUMENT, so a title carrying \api{model\_validate}
  % made a bookmark reading "mtealmodel_validate" and warned "Token not
  % allowed" twice. The colour name is gobbled here instead. Measured on a
  % probe against this preamble before the line was added.
  \def\color#1{}%
}
\makeatother

% ---------- Language strings ----------
\input{lang/\booklang}

% The PDF's subject, keyword list and document language read \lblPdf... from
% the language file input on the line above, so they cannot sit in the colour
% \hypersetup block, which runs before it. Two blocks is the ordering, not an
% oversight; merging them raises Undefined control sequence and still writes
% a PDF, with the three fields empty.
\hypersetup{
  pdfsubject={\lblPdfSubject},
  pdfkeywords={\lblPdfKeywords},
  pdflang={\lblPdfLang}
}

% ---------- Headings ----------
% \raggedright on the title body is load-bearing: at \Huge on this block a
% title that cannot break overflows the margin by 30 pt or more, and several
% chapter titles here run past fifty characters.
\titleformat{\chapter}[display]
  {\normalfont\huge\bfseries\color{mblue}\raggedright}
  {\normalfont\normalsize\color{mgrey}\MakeUppercase{\chaptertitlename}\ \thechapter}
  {8pt}{\Huge\raggedright}
\titlespacing*{\chapter}{0pt}{10pt}{28pt}
\titleformat{\section}{\normalfont\Large\bfseries\color{mblue}\raggedright}{\thesection}{0.8em}{}
\titleformat{\subsection}{\normalfont\large\bfseries\color{mgrey}\raggedright}{\thesubsection}{0.7em}{}

% Sections in the contents, subsections not: fourteen chapters of subsections
% is a contents nobody reads, and the reader navigates by chapter and section.
\setcounter{tocdepth}{1}
\setcounter{secnumdepth}{2}

% One-sided running head: the chapter on the left, the folio on the right.
\pagestyle{fancy}
\fancyhf{}
\fancyhead[L]{\small\nouppercase{\leftmark}}
\fancyhead[R]{\small\thepage}
\renewcommand{\headrulewidth}{0.4pt}
\fancypagestyle{plain}{\fancyhf{}\fancyfoot[R]{\small\thepage}%
  \renewcommand{\headrulewidth}{0pt}}

% Drops the word "Chapter" from the head: the number is still there and the
% longer titles no longer overflow. MUST come after \pagestyle{fancy}, which
% installs its own \chaptermark and silently overwrites anything earlier.
\renewcommand{\chaptermark}[1]{\markboth{\thechapter.\ #1}{}}

% ============================================================
%  CODE LISTINGS
%  ---------------------------------------------------------
%  Every listing in a chapter is a FILE under code/, pulled in with \pyfile or
%  \pyregion, and CI runs it against the pinned versions. An in-source
%  \begin{python} block is allowed for a fragment of three or four lines that
%  exists to show a shape rather than to run; anything longer belongs in a
%  file the reader can open in the editor beside the PDF.
% ============================================================
\lstdefinelanguage{pythonx}{
  language=Python,
  morekeywords={
    async,await,match,case,nonlocal,yield,type,
    TypedDict,Annotated,Protocol,Literal,Generic,Final,ClassVar,Self,
    TypeVar,TypeGuard,TypeIs,Callable,Iterator,Iterable,Awaitable,
    dataclass,field,BaseModel,Field,Enum,StrEnum,NamedTuple,
    override,overload,property,staticmethod,classmethod
  },
  morekeywords=[2]{
    gather,create_task,TaskGroup,to_thread,timeout,wait_for,Semaphore,Queue,
    fixture,parametrize,monkeypatch,raises,mark,
    Depends,FastAPI,APIRouter,HTTPException,
    select,Session,Mapped,mapped_column,relationship,
    model_validate,model_validate_json,model_dump,model_json_schema
  },
  sensitive=true
}

% listings' C# definition predates async, records and pattern matching.
\lstdefinelanguage{csharpx}{
  language=[Sharp]C,
  morekeywords={
    async,await,var,dynamic,record,init,required,nameof,yield,partial,
    global,scoped,when,with,and,or,not,notnull,unmanaged,file,get,set,value,
    where,select,from,orderby,let,join,into,ascending,descending,on,equals,
    Task,ValueTask,CancellationToken,IEnumerable,IAsyncEnumerable,Func,Action
  },
  sensitive=true
}

\lstdefinestyle{bookcode}{
  basicstyle=\ttfamily\footnotesize,
  keywordstyle=\color{kw}\bfseries,
  keywordstyle=[2]\color{typ},
  stringstyle=\color{str},
  % Colour only, no \itshape: inconsolata (zi4) has no italic shape, so an
  % italic comment style is substituted with upright on every machine that
  % has the font, and the log carries "Font shape `T1/zi4/m/it' undefined"
  % on every build. The colour already marks a comment.
  commentstyle=\color{cmt},
  emphstyle=\color{dec}\bfseries,
  numbers=left,
  numberstyle=\ttfamily\tiny\color{mgrey},
  numbersep=8pt,
  backgroundcolor=\color{codebg},
  frame=single,
  rulecolor=\color{codeframe},
  framesep=6pt,
  breaklines=true,
  breakatwhitespace=false,
  postbreak=\mbox{\textcolor{mgrey}{$\hookrightarrow$}\space},
  showstringspaces=false,
  tabsize=4,
  columns=fullflexible,
  keepspaces=true,
  captionpos=b,
  aboveskip=1.2em,
  belowskip=1.2em,
  % Maps the characters most likely to be pasted into a listing by accident.
  %
  % Do NOT add a mapping for U+00A0 (non-breaking space). It looks like the
  % obvious next entry and it makes `listings` abort the run with "Improper
  % alphabetic constant" -- fatal, no PDF, and the message names neither the
  % character nor this line. The ASCII-in-listings convention, checked by
  % parity's C13, is the guard against it instead.
  literate={ł}{{\l}}1 {Ł}{{\L}}1 {ą}{{\k a}}1 {ę}{{\k e}}1
           {ś}{{\'s}}1 {ć}{{\'c}}1 {ż}{{\.z}}1 {ź}{{\'z}}1 {ó}{{\'o}}1 {ń}{{\'n}}1
           {Ą}{{\k A}}1 {Ę}{{\k E}}1 {Ś}{{\'S}}1 {Ć}{{\'C}}1
           {Ż}{{\.Z}}1 {Ź}{{\'Z}}1 {Ó}{{\'O}}1 {Ń}{{\'N}}1
           {—}{{-{}-}}2 {–}{{-}}1 {‘}{{'}}1 {’}{{'}}1
           {“}{{"}}1 {”}{{"}}1 {°}{{\textdegree}}1
}

\lstset{style=bookcode}

\lstnewenvironment{python}[1][]
  {\lstset{language=pythonx,style=bookcode,#1}}
  {}

% C# comparison snippets. The whole book is written for engineers arriving
% from .NET, so listings come in pairs often enough that the C# half has its
% own environment rather than a generic one.
\lstnewenvironment{csharp}[1][]
  {\lstset{language=csharpx,style=bookcode,#1}}
  {}

\lstnewenvironment{shellcmd}[1][]
  {\lstset{language=bash,style=bookcode,numbers=none,keywordstyle=\color{mteal}\bfseries,#1}}
  {}

\lstnewenvironment{tomlcode}[1][]
  {\lstset{style=bookcode,numbers=none,keywordstyle=\color{black},#1}}
  {}

\lstnewenvironment{yamlcode}[1][]
  {\lstset{style=bookcode,numbers=none,keywordstyle=\color{black},#1}}
  {}

\lstnewenvironment{jsoncode}[1][]
  {\lstset{style=bookcode,numbers=none,keywordstyle=\color{black},#1}}
  {}

% Terminal transcripts typed into the source. Allowed for a shell session the
% reader is meant to reproduce; NOT for program output, which is \transcript
% below, because an in-source transcript nobody ran reads exactly like one
% that was, and that is where a remembered number hides.
\lstnewenvironment{console}[1][]
  {\lstset{style=bookcode,numbers=none,basicstyle=\ttfamily\footnotesize,
           keywordstyle=\color{black},backgroundcolor=\color{white},#1}}
  {}

% The path is printed above every file listing, in the repository-relative
% form the reader opens in the editor beside the PDF. That is the whole point
% of a listing being a file, and the front matter promises it. The \nobreak
% keeps the path on the same page as the listing it names.
% ---------------------------------------------------------------------------
%  Appendix B prints the trap catalogue from notes/02-traps.md. One entry is
%  its number -- which a chapter may cite, and which is never reused -- the
%  habit in the reader's own voice, and what Python does instead.
%
%  The body is set RAGGED RIGHT on purpose. An entry is mostly \code{} spans,
%  which are unbreakable runs, and sixty-seven justified paragraphs of them
%  is sixty-seven chances at an overfull hbox in the edition with the longer
%  words. Ragged right removes the class rather than fixing it one box at a
%  time, and a reference list is the one place in this book where it reads
%  better anyway.
%
%  The habit is set with \enquote{} rather than \emph{}, and that is not a
%  style choice. Every habit here carries \code{} spans, and \code{} inside
%  \emph{} asks inconsolata for an italic it does not have: the first build
%  of this appendix raised `Font shape T1/zi4/m/it undefined', which
%  checklog.py is hard on for exactly this reason. It is the trap CLAUDE.md
%  records, met sixty-seven times at once -- and fixed once, because the
%  entries go through a macro. A quotation is also what the habit IS, and
%  \enquote{} is what the Polish edition owes anyway.
% ---------------------------------------------------------------------------
%  Appendix A's cheat sheet: one table per part, C# on the left, Python on
%  the right, and the chapter that teaches the row on the end. The three
%  header strings are arguments so that one environment serves both
%  editions and the structural signature stays identical.
%
%  Columns are ragged right for the reason Appendix B's entries are: a cell
%  is mostly \code{} spans, which are unbreakable runs, and a justified
%  narrow column full of them is an overfull box waiting for the edition
%  with the longer words.
% ---------------------------------------------------------------------------
%  Appendix C's tool matrix: the .NET tool, the Python tool this book uses,
%  the version it was verified against and the chapter that introduces it.
%
%  The version column prints \pydanticver and its siblings and is never
%  typed, which is the whole point of the appendix: it is the one page where
%  a reader sees every pin at once, and a pin that had to be re-typed here
%  would be a second copy of a fact the preamble already owns.
%  tools/check_structure.py --tools fails when a pin macro exists and this
%  table does not print it.
\NewDocumentEnvironment{toolmatrix}{mmmm}{%
  \small
  \begin{longtable}{@{}%
    >{\raggedright\arraybackslash}p{0.27\linewidth}%
    >{\raggedright\arraybackslash}p{0.31\linewidth}%
    >{\raggedright\arraybackslash}p{0.20\linewidth}%
    >{\raggedright\arraybackslash}p{0.08\linewidth}@{}}
  \toprule
  \textbf{#1} & \textbf{#2} & \textbf{#3} & \textbf{#4} \\
  \midrule
  \endfirsthead
  \toprule
  \textbf{#1} & \textbf{#2} & \textbf{#3} & \textbf{#4} \\
  \midrule
  \endhead
  \midrule
  \endfoot
  \bottomrule
  \endlastfoot
}{%
  \end{longtable}%
}

\NewDocumentEnvironment{cheatsheet}{mmm}{%
  \small
  \begin{longtable}{@{}%
    >{\raggedright\arraybackslash}p{0.41\linewidth}%
    >{\raggedright\arraybackslash}p{0.41\linewidth}%
    >{\raggedright\arraybackslash}p{0.08\linewidth}@{}}
  \toprule
  \textbf{#1} & \textbf{#2} & \textbf{#3} \\
  \midrule
  \endfirsthead
  \toprule
  \textbf{#1} & \textbf{#2} & \textbf{#3} \\
  \midrule
  \endhead
  \midrule
  \endfoot
  \bottomrule
  \endlastfoot
}{%
  \end{longtable}%
}

\newcommand{\trapentry}[3]{%
  \par\medskip
  \noindent\textbf{#1.}\enspace\enquote{#2}\par
  \nopagebreak
  {\raggedright\noindent#3\par}%
}

\newcommand{\listingpath}[1]{%
  \par\medskip\noindent
  {\ttfamily\footnotesize\color{mgrey}\detokenize{#1}}%
  \par\nobreak}

% External source file -- the preferred form. The path is relative to the
% repository root, which is also the path the reader opens in the editor.
% Usage: \pyfile{code/ch05/decorators.py}{Caption}{lst:label}
\newcommand{\pyfile}[3]{%
  \listingpath{#1}%
  \lstinputlisting[language=pythonx,style=bookcode,aboveskip=0.3em,
    caption={#2},label={#3}]{#1}%
}

% A named region of an external file, delimited in the source by
% `# --8<-- [start:name]` / `# --8<-- [end:name]`, so the book and the running
% code cannot drift apart. tools/check_structure.py --listings fails when the
% file or the region is missing.
%
% The markers are matched through listings' range prefix and suffix keys,
% with linerange={name-name}. The form this was inherited with -- a linerange
% whose two ends were the whole marker strings -- matched nothing and printed
% the WHOLE FILE, with no warning: measured on a three-region probe file
% before this was rewritten, and the same definition sits unused in the
% sibling book it came from. A region name is a word, never a number, because
% a number is a line number to listings. The line numbers printed are the
% file's own, so a region's listing says where in the file it sits.
% Usage: \pyregion{code/ch05/decorators.py}{timing}{Caption}{lst:label}
\lstset{
  rangebeginprefix=\#\ --8<--\ [start:, rangebeginsuffix=],
  rangeendprefix=\#\ --8<--\ [end:, rangeendsuffix=],
  includerangemarker=false}
\newcommand{\pyregion}[4]{%
  \listingpath{#1}%
  \lstinputlisting[language=pythonx,style=bookcode,aboveskip=0.3em,
    linerange={#2-#2},caption={#3},label={#4}]{#1}%
}

% The C# twin of \pyfile, for the comparison half of a pair.
\newcommand{\csfile}[3]{%
  \listingpath{#1}%
  \lstinputlisting[language=csharpx,style=bookcode,aboveskip=0.3em,
    caption={#2},label={#3}]{#1}%
}

% A program's output, written by a script under code/measure/ and pulled in
% whole. There is no typed form of this and there must not be: \val{} does
% not expand inside a verbatim body, so a transcript typed into a chapter is
% one that cannot quote a computed value, and a transcript nobody ran is
% indistinguishable on the page from one that was.
\newcommand{\transcript}[1]{%
  \par\smallskip
  \IfFileExists{figures/transcripts/#1.txt}{%
    \lstinputlisting[style=bookcode,numbers=none,basicstyle=\ttfamily\footnotesize,
      keywordstyle=\color{black},backgroundcolor=\color{white},
      language={}]{figures/transcripts/#1.txt}%
  }{%
    \begin{tcolorbox}[enhanced, sharp corners, width=\linewidth,
      colback=mlight, colframe=mgrey, boxrule=0.6pt,
      left=8pt, right=8pt, top=6pt, bottom=6pt]
      {\bfseries\color{mgrey}\small \lblTranscriptMissing}\\[2pt]
      {\ttfamily\scriptsize\color{mgrey} figures/transcripts/#1.txt}
    \end{tcolorbox}%
  }%
  \par\smallskip
}

% Inline literals. Underscores inside these must be written \_.
\newcommand{\code}[1]{\texttt{\small #1}}
\newcommand{\pkg}[1]{\texttt{\small\color{mblue}#1}}
\newcommand{\api}[1]{\texttt{\small\color{mteal}#1}}

% ============================================================
%  ADMONITIONS
%  ---------------------------------------------------------
%  A deliberately small, fixed vocabulary. Two visual treatments carry the
%  whole grammar: `tinted` is a block the reader may read in passing -- a
%  tint and a thick left bar, no outline; `outlined` is a block the reader
%  must stop at, with a full rule on white. A third treatment would mean the
%  first two are not carrying their meaning.
% ============================================================
\tcbset{
  mfa tinted/.style={
    enhanced, breakable, sharp corners,
    colback=mlight, boxrule=0pt, leftrule=3pt,
    left=8pt, right=8pt, top=6pt, bottom=6pt,
    before skip=13pt, after skip=13pt,
    fonttitle=\bfseries\small,
    coltitle=black, colbacktitle=mlight,
    attach title to upper={\par\smallskip},
  },
  mfa outlined/.style={
    enhanced, breakable, sharp corners,
    colback=white, boxrule=0.6pt,
    left=8pt, right=8pt, top=6pt, bottom=6pt,
    before skip=13pt, after skip=13pt,
    fonttitle=\bfseries\small,
    coltitle=black, colbacktitle=white,
    attach title to upper={\par\smallskip},
  },
}

% Read in passing.
\newtcolorbox{note}[1][]{mfa tinted, colframe=mblue, title={\lblNote}, #1}
\newtcolorbox{warning}[1][]{mfa tinted, colframe=mred, title={\lblWarning}, #1}
% The translation. The .NET reader is the book's only assumed reader, so this
% box carries the C# side of a comparison and nothing else. Its use is
% counted: a chapter with none is a chapter that has stopped translating.
\newtcolorbox{csbox}[1][]{mfa tinted, colframe=mteal, title={\lblCsharp}, #1}
% Anything that is preview, experimental or likely to move. Python moves more
% slowly than the agent frameworks do, and this box should be rare here.
\newtcolorbox{versionbox}[1][]{mfa tinted, colframe=mamber, title={\lblVersion}, #1}

% Stop and read.
% The misconception, named AFTER the reader has walked into it. Appendix B is
% the catalogue; a trap box in a chapter is one entry of it, in place.
\newtcolorbox{trapbox}[1][]{mfa outlined, colframe=mred, title={\lblTrap}, #1}
% Marks a listing that has not been executed against the pinned versions.
% `make debt` counts these, CI prints the count, and nothing ships to a reader
% with one attached. Removing it means you ran the code.
\newtcolorbox{verifybox}[1][]{mfa outlined, colframe=mgrey, title={\lblVerify}, #1}
\newtcolorbox{exercisebox}[1][]{mfa outlined, colframe=mteal, title={\lblExercise}, #1}
% The guiding project. Points at the stage of trace-assert this section
% builds, under code/src/trace_assert/.
\newtcolorbox{projectbox}[1][]{mfa outlined, colframe=mamber, title={\lblProject}, #1}

% ============================================================
%  EXERCISES
%  ---------------------------------------------------------
%  \begin{exercise}{key}{title} ... \end{exercise}
%
%  This is the mechanism the whole book is read through: the PDF on one half
%  of the screen, an editor on the other. An exercise is therefore not a
%  paragraph but a FILE the reader opens -- a starter under
%  code/exercises/chNN/<key>.py -- and a TEST that decides whether it is done,
%  test_<key>.py beside it. The box prints both paths and the one command
%  that runs the test, and the manifest at the back lists every exercise.
%
%  An ENVIRONMENT, not a macro with a body argument: a macro argument cannot
%  contain verbatim material, so an exercise written as \exercise{key}{title}
%  {body} could never show the reader a code fragment, and the first chapter
%  that tried would have died on `\begin{python}` inside the braces.
%
%  The key is the file stem, with underscores, and it is written with
%  \detokenize so the underscore reaches the page as a character rather than
%  as a subscript. tools/check_structure.py --exercises fails when the starter
%  or the test is missing, or when the key's chapter prefix disagrees with the
%  chapter it sits in.
% ============================================================
\newcounter{exercise}[chapter]
\renewcommand{\theexercise}{\thechapter.\arabic{exercise}}
\makeatletter
\newcommand{\exercisedir}{ch\two@digits{\value{chapter}}}
\newcommand{\listofexercises}{%
  \section*{\lblExerciseManifest}%
  \phantomsection\addcontentsline{toc}{section}{\lblExerciseManifest}%
  % The entries are subsection-level and tocdepth is 1 for the contents,
  % so the depth is raised for this list alone or it prints empty --
  % which it did, on the scaffold's first clean build.
  \begingroup\c@tocdepth=2\relax\@starttoc{exr}\endgroup%
}
\makeatother
\NewDocumentEnvironment{exercise}{mm}{%
  \refstepcounter{exercise}%
  \addcontentsline{exr}{subsection}{\protect\numberline{\theexercise}\texttt{\protect\detokenize{#1}} --- #2}%
  \begin{exercisebox}[title={\lblExercise~\theexercise\ \dash{} #2}]%
}{%
  \par\medskip
  {\small\setlength{\parindent}{0pt}%
  \lblExerciseStarter: \texttt{code/exercises/\exercisedir/\detokenize{#1}.py}\\
  \lblExerciseTest: \texttt{code/exercises/\exercisedir/test\_\detokenize{#1}.py}\\
  \lblExerciseRun: \texttt{cd code \&\& uv run pytest -k \detokenize{#1}}\par}%
  \end{exercisebox}%
}

% ============================================================
%  DIAGRAMS --- Mermaid pipeline
%  ---------------------------------------------------------
%  \mermaidfig{key}{caption}{one-line manifest description}
%
%  Source of truth is figures/mermaid/<lang>/<key>.mmd, committed. `make
%  diagrams` renders each to figures/diagrams/<lang>/<key>.pdf, which is build
%  output and is not committed, so a diagram change reviews as a text diff.
%  The source is per-language because a diagram with English labels in the
%  Polish edition is exactly the half-translation that makes a bilingual book
%  feel machine-made; CI checks that both languages carry the same keys.
%
%  If the rendered PDF is absent the macro draws a placeholder AND typesets
%  the Mermaid source, in the flow rather than as a float, because a source
%  typeset verbatim is often taller than a page and a float cannot break.
% ============================================================
\makeatletter
\newcommand{\listofdiagrams}{%
  \section*{\lblDiagramManifest}%
  \phantomsection\addcontentsline{toc}{section}{\lblDiagramManifest}%
  % The entries are subsection-level and tocdepth is 1 for the contents,
  % so the depth is raised for this list alone or it prints empty --
  % which it did, on the scaffold's first clean build.
  \begingroup\c@tocdepth=2\relax\@starttoc{dgm}\endgroup%
}
\makeatother

\newcommand{\mermaidfig}[3]{%
  % \phantomsection is load-bearing and not decoration. \addcontentsline
  % records hyperref's CURRENT anchor, and here that is whatever preceded
  % the figure -- for a figure placed after a listing it is the listing's
  % line-number anchor, and listings has already advanced the counter past
  % the last line it printed, so the manifest entry linked to a
  % destination that is never created: "name{lstnumber.10.4.46} has been
  % referenced but does not exist, replaced by a fixed one". The two
  % front-matter figures had the same wrong target and were silent about
  % it, because the lines they named happen to be printed. The exercise
  % manifest never had the bug, because \begin{exercise} steps a counter
  % and that is what sets the anchor.
  \phantomsection%
  \addcontentsline{dgm}{subsection}{\protect\numberline{}\texttt{#1.mmd} --- #3}%
  \IfFileExists{figures/diagrams/\booklang/#1.pdf}{%
    \begin{figure}[htbp]
    \centering
    \includegraphics[width=0.95\linewidth,height=0.42\textheight,keepaspectratio]{figures/diagrams/\booklang/#1.pdf}%
    \caption{#2}\label{fig:#1}
    \end{figure}%
  }{%
    \par\medskip
    \begin{tcolorbox}[
      enhanced, breakable, width=\linewidth, sharp corners,
      colback=white, colframe=mgrey, boxrule=0.8pt,
      left=8pt, right=8pt, top=8pt, bottom=8pt]
      {\bfseries\color{mgrey}\small \lblNotRendered}\\[4pt]
      {\ttfamily\scriptsize\color{mgrey} figures/mermaid/\booklang/#1.mmd}\\[6pt]
      \IfFileExists{figures/mermaid/\booklang/#1.mmd}{%
        \lstinputlisting[style=bookcode,numbers=none,basicstyle=\ttfamily\scriptsize,
                         backgroundcolor=\color{mlight},frame=none,
                         keywordstyle=\color{black},language={}]{figures/mermaid/\booklang/#1.mmd}%
      }{%
        {\small\color{mred} \lblSourceMissing: \texttt{figures/mermaid/\booklang/#1.mmd}}%
      }%
    \end{tcolorbox}%
    \captionof{figure}{#2}\label{fig:#1}
    \par\medskip
  }%
}

% ============================================================
%  CHAPTER STUBS
%  ---------------------------------------------------------
%  Prints a visible marker so an unwritten chapter cannot be mistaken for a
%  written one and the page count never flatters the draft. `make debt`
%  counts them, and CI publishes the count on every build.
% ============================================================
\newcommand{\chapterstub}[1]{%
  \begin{tcolorbox}[
    enhanced, breakable, width=\linewidth, sharp corners,
    colback=mlight, colframe=mred, boxrule=1pt,
    borderline={0.6pt}{3pt}{mred, dashed},
    left=12pt, right=12pt, top=12pt, bottom=12pt]
    {\bfseries\color{mred}\large \lblNotWritten}\\[8pt]
    \small\raggedright #1
  \end{tcolorbox}%
}

% ============================================================
%  COMPUTED VALUES
%  ---------------------------------------------------------
%  A number that came out of a measurement is written by a script under
%  code/measure/ into figures/values/<name>.tex as
%
%      \pyval{e01.threads.speedup}{3.71}
%
%  and pulled into the text with \val{e01.threads.speedup}. The script is the
%  source of truth and the book cannot disagree with it, because the book does
%  not contain the digits. `make verify` fails when a committed value no
%  longer matches its script. A value the scripts do not produce prints a
%  visible marker.
%
%  \val also applies the edition's decimal separator, which is how the Polish
%  edition gets its decimal comma without a translator having to remember.
% ============================================================
\newcommand{\bookdecimalmarker}{\ifpl{,}{.}}
\IfFileExists{siunitx.sty}{%
  \usepackage{siunitx}%
  \sisetup{
    group-digits = integer,
    group-minimum-digits = 5,
    group-separator = {\,},
    output-decimal-marker = {\bookdecimalmarker}
  }%
  \newcommand{\booknum}[1]{\num{##1}}%
}{%
  \newcommand{\booknum}[1]{##1}%
}

\newcommand{\pyvalues}{}
\makeatletter
\newcommand{\pyval}[2]{\csgdef{pyval@#1}{#2}\listxadd{\pyvalues}{#1}}
% \num on a non-numeric body is a FATAL siunitx error. The generator knows
% whether it emitted a number or a word, so it writes \pyval for a number and
% \pyvaltext for a word, and this end only enforces which macro reads which.
\newcommand{\val}[1]{%
  \ifcsdef{pyval@#1}{%
    \ifcsdef{pyvaltext@#1}{%
      \PackageError{pybook}{Value `#1' is not a number}%
        {\string\val\space passes its body to siunitx, which refuses anything
         that is not a number. Use \string\valtext{#1} instead.}%
    }{}%
    \booknum{\csuse{pyval@#1}}%
  }{\textcolor{mred}{\textbf{[\lblValueMissing: #1]}}}%
}
\newcommand{\valtext}[1]{%
  \ifcsdef{pyval@#1}{\csuse{pyval@#1}}%
    {\textcolor{mred}{\textbf{[\lblValueMissing: #1]}}}%
}
\newcommand{\pyvaltext}[2]{%
  \csgdef{pyval@#1}{#2}\csgdef{pyvaltext@#1}{}\listxadd{\pyvalues}{#1}%
}
\makeatother

% figures/values/all.tex is generated by `make numbers` and does nothing but
% \input each value file. A build with no values yet still compiles; every
% \val{} in it prints a visible marker instead of a number.
\IfFileExists{figures/values/all.tex}{\input{figures/values/all}}{%
  \AtBeginDocument{\PackageWarning{pybook}{No computed values found. Run `make numbers`.}}}

% ============================================================
%  PINNED VERSIONS --- single source of truth
%  Change these here; they propagate through both editions and Appendix C.
%  Verified against pypi.org on \pinnedon (lang/<lang>.tex).
%  tools/check_versions.py compares every one of them with PyPI on each build.
%  Python itself is not on PyPI and is checked by hand against python.org.
% ============================================================
\newcommand{\bookedition}{0.1-dev}
\newcommand{\bookyear}{2026}
\newcommand{\pyver}{3.14}              % the minor the book targets
\newcommand{\pypatch}{3.14.7}          % the exact interpreter the code ran on
\newcommand{\uvver}{0.12.13}           % uv
\newcommand{\dotnetver}{10.0.100}      % dotnet
\newcommand{\ruffver}{0.16.7}          % ruff
\newcommand{\pyrightver}{1.1.414}      % pyright
\newcommand{\mypyver}{2.3.1}           % mypy
\newcommand{\pydanticver}{2.13.5}      % pydantic
\newcommand{\pysettingsver}{2.15.0}    % pydantic-settings
\newcommand{\fastapiver}{0.141.1}      % fastapi
\newcommand{\uvicornver}{0.53.0}       % uvicorn
\newcommand{\sqlalchemyver}{2.0.52}    % sqlalchemy
\newcommand{\alembicver}{1.20.0}       % alembic
\newcommand{\aiosqlitever}{0.22.1}     % aiosqlite
\newcommand{\pytestver}{9.1.1}         % pytest
\newcommand{\pytestasyncver}{1.4.0}    % pytest-asyncio
\newcommand{\hypothesisver}{6.168.0}   % hypothesis
\newcommand{\coveragever}{7.16.1}      % coverage
\newcommand{\httpxver}{0.28.1}         % httpx
\newcommand{\structlogver}{26.1.0}     % structlog
\newcommand{\otelver}{1.44.0}          % opentelemetry-sdk
\newcommand{\polarsver}{1.44.2}        % polars
\newcommand{\anthropicver}{1.5.0}      % anthropic
\newcommand{\openaiver}{3.13.0}        % openai
\newcommand{\pipauditver}{2.10.1}      % pip-audit

% ============================================================
%  INDEX
%  Two-column, and its entries include \texttt{} names that do not hyphenate.
%  Ragged right keeps multi-word entries off the margin; it cannot help a
%  single token wider than the column, so keep verbatim index entries under
%  about 29 characters and index longer names by concept instead.
%  \raggedcolumns is the switch imakeidx's multicols actually reads.
% ============================================================
\apptocmd{\theindex}{\raggedcolumns\raggedright}{}{%
  \PackageWarning{pybook}{Could not patch theindex for ragged setting}}
\providecommand*\see[2]{}
\renewcommand*\see[2]{\emph{\lblIndexSee} #1}
\providecommand*\seealso[2]{}
\renewcommand*\seealso[2]{\emph{\lblIndexSeeAlso} #1}

\makeindex
; nothing else differs between them at this level.
%  Every user-visible string lives in lang/en.tex or lang/pl.tex, so a label
%  that appears in one edition and not the other is a missing macro rather
%  than a silent divergence.
%
%  Lineage. The bilingual wiring (\booklang, lang/, body.tex, the generated
%  structure) is the math book's; the chapter machinery (listings, the
%  admonitions, verifybox, \pyfile and \pyregion, the Mermaid pipeline) is the
%  LangChain book's. What is new here is the exercise macro, which ties every
%  exercise to a starter file and a test under code/, because this book is
%  read with an editor open beside it.
% ============================================================

% \booklang must be set by the main file. Default to English rather than
% failing, so a stray % ============================================================
%  preamble.tex --- styling, environments, macros
%
%  Shared by BOTH editions. main-en.tex and main-pl.tex each define \booklang
%  before \input{preamble}; nothing else differs between them at this level.
%  Every user-visible string lives in lang/en.tex or lang/pl.tex, so a label
%  that appears in one edition and not the other is a missing macro rather
%  than a silent divergence.
%
%  Lineage. The bilingual wiring (\booklang, lang/, body.tex, the generated
%  structure) is the math book's; the chapter machinery (listings, the
%  admonitions, verifybox, \pyfile and \pyregion, the Mermaid pipeline) is the
%  LangChain book's. What is new here is the exercise macro, which ties every
%  exercise to a starter file and a test under code/, because this book is
%  read with an editor open beside it.
% ============================================================

% \booklang must be set by the main file. Default to English rather than
% failing, so a stray \input{preamble} still compiles and says what it did.
\providecommand{\booklang}{en}

\usepackage{etoolbox}   % loaded first: the language switch below needs it

% ---------- Page geometry ----------
% ONE format: A4, 12pt, ONE-SIDED. This book is read on a screen with an
% editor open beside it, and never printed as a bound book, so there is no
% recto, no verso, no gutter, no mirrored margin and no blank leaf before a
% chapter -- every one of those is a property of paper that a PDF viewer at
% half a screen's width turns into wasted pixels.
%
% The margins are symmetric and the measure is about seventy-five characters
% of prose, which is where the eye keeps the line return. The number that was
% actually chosen is the MONOSPACE measure: at \footnotesize on this block a
% listing line of 79 characters fits without wrapping, and 79 is the width
% every listing under code/ is held to by tools/check_structure.py --lines.
% A wrapped listing line is not an error TeX reports -- breaklines=true wraps
% it silently and prints an arrow into the middle of what the reader is meant
% to paste into an editor -- so the width is enforced in the source instead.
\usepackage[
  a4paper,
  left=3.1cm, right=3.1cm, top=2.6cm, bottom=3.0cm,
  head=26pt, headsep=0.8cm, footskip=1.4cm
]{geometry}

% ---------- Encoding & fonts ----------
\usepackage[T1]{fontenc}
\usepackage[utf8]{inputenc}
\IfFileExists{lmodern.sty}{\usepackage{lmodern}}{}
% newtxtext takes its TS1 symbols (\textcopyright, \textdegree) from TeX
% Gyre Termes, which Debian packages SEPARATELY (tex-gyre): a machine with
% newtx and without it dies on the copyright page with
% `Font TS1/ntxtlf/m/n/10.95=ts1-qtmr ... not loadable`. There is no
% inert-safe probe for a font metric in pdfTeX -- \IfFileExists searches
% TEXINPUTS and not TFMFONTS, so it reports the metric missing on a machine
% that has it, and \suppressfontnotfounderror is LuaTeX-only -- so nothing
% here guesses. CI's full TeX Live is the reference; locally, install
% tex-gyre. Both facts are recorded in CLAUDE.md under Build traps.
\IfFileExists{newtxtext.sty}{\usepackage{newtxtext}}{}
% amsmath before newtxmath, because newtxmath patches amsmath's internals.
% amssymb only when newtxmath is absent: newtxmath supplies the AMS symbols
% itself and loading both is a fatal clash (\Bbbk already defined) that is
% invisible on a machine without newtx and fatal on a full TeX Live.
\usepackage{amsmath}
\IfFileExists{newtxmath.sty}{\usepackage{newtxmath}}{\usepackage{amssymb}}
\IfFileExists{inconsolata.sty}{\usepackage[varqu,scaled=0.95]{inconsolata}}{}
\IfFileExists{lmodern.sty}{\usepackage{microtype}}{\usepackage[expansion=false]{microtype}}

% upquote is not a refinement. In T1 the input character ' is quoteright, so
% a listings body sets f('x') with U+2019 at both ends of the literal, and
% typed back in that is a SyntaxError. inconsolata's varqu option hides the
% defect on a machine that has inconsolata; upquote makes ' the ASCII
% apostrophe whatever the typewriter font is. In a book whose every listing is
% meant to be pasted into an editor, this line is load-bearing.
\IfFileExists{upquote.sty}{\usepackage{upquote}}{}

% And upquote's twin, for the same reason one character over. In T1 the pair
% -- is a ligature for an en dash, in EVERY family including the typewriter
% one, so \code{uv sync --locked} sets as "uv sync <en dash>locked" and a
% reader who copies it off the page gets an unknown-argument error from a
% command the chapter told them to run. It is invisible in review because it
% looks like a dash, and it bites exactly the flags a packaging chapter is
% made of. A `listings` body is already safe -- listings sets characters
% singly and no ligature forms -- which is what makes the defect so easy to
% miss: the same text is right inside a listing and wrong in the sentence
% above it.
%
% Scoped to tt* on purpose, and measured rather than assumed: prose keeps its
% own dashes, so \dash{} still sets an em dash and a range like 1--2 still
% sets an en dash. Probed against this preamble in all three positions before
% it was believed.
\DisableLigatures[-]{encoding = T1, family = tt*}

% ---------- Language ----------
% Loading babel with a language whose .ldf is absent is a *fatal* error, not
% a warning, so both the package and each language file are probed before
% either is requested. Both editions load the other language too, because the
% glossary rows and the cheat sheet carry words from both.
\IfFileExists{babel.sty}{%
  \IfFileExists{polish.ldf}{%
    \IfFileExists{british.ldf}{%
      \ifdefstring{\booklang}{pl}%
        {\usepackage[british,main=polish]{babel}}%
        {\usepackage[polish,main=british]{babel}}%
    }{\usepackage[polish]{babel}}%
  }{%
    \IfFileExists{british.ldf}{\usepackage[british]{babel}}{}%
  }%
}{}

% babel-polish makes " an active shorthand. listings reads its body verbatim
% and an active character in a verbatim context is a category-code accident
% waiting to happen, so the shorthand is turned off. Polish quotation marks
% come from \enquote, which csquotes sets per language.
\ifdefstring{\booklang}{pl}{\AtBeginDocument{\shorthandoff{"}}}{}

% ---------- Core packages ----------
\usepackage{graphicx}
\usepackage{xcolor}
\usepackage{booktabs}
\usepackage{tabularx}
% Appendix A's cheat sheet runs past a page in every part, and tabularx
% cannot break. longtable can, and works with booktabs.
\usepackage{longtable}
\usepackage{array}
\usepackage{enumitem}
\usepackage{listings}
\usepackage[most]{tcolorbox}
\usepackage{fancyhdr}
\usepackage{titlesec}
\usepackage{caption}
\usepackage{imakeidx}
% csquotes with autostyle follows babel's active language, so the identical
% source \enquote{x} sets 'x' in the English edition and the Polish quotation
% marks in the Polish one.
\IfFileExists{csquotes.sty}{\usepackage[autostyle=true]{csquotes}}{%
  \providecommand{\enquote}[1]{``##1''}}
\usepackage[hidelinks]{hyperref}

% ---------- Palette ----------
% The same family as the two companion volumes, so the three read as a set.
\definecolor{mblue}{HTML}{1F4E79}
\definecolor{mteal}{HTML}{0E7C7B}
\definecolor{mamber}{HTML}{B26A00}
\definecolor{mred}{HTML}{9B2C2C}
\definecolor{mviolet}{HTML}{7B4397}
\definecolor{mgrey}{HTML}{5A5A5A}
\definecolor{mlight}{HTML}{F4F6F8}
\definecolor{mfaint}{HTML}{FBFCFD}
\definecolor{codebg}{HTML}{FAFAFA}
\definecolor{codeframe}{HTML}{D8DEE4}
\definecolor{kw}{HTML}{0B5394}
\definecolor{str}{HTML}{9B2C2C}
\definecolor{cmt}{HTML}{4F7A4F}
\definecolor{typ}{HTML}{0E7C7B}
\definecolor{dec}{HTML}{7B4397}

\hypersetup{
  colorlinks=true,
  linkcolor=mblue,
  urlcolor=mteal,
  citecolor=mblue,
  pdfauthor={Konrad Cinkusz}
}

% \ifpl{Polish text}{English text} -- for the handful of places where a
% string is structural rather than content. Anything longer than a few words
% belongs in lang/ or in the edition's own source, not here.
% \DeclareRobustCommand, not \newcommand: a fragile \ifpl inside a moving
% argument is written to the .toc one level expanded and fails on the next run
% in exactly one of the two editions.
\DeclareRobustCommand{\ifpl}[2]{\ifdefstring{\booklang}{pl}{#1}{#2}}
\makeatletter
\pdfstringdefDisableCommands{%
  \ifdefstring{\booklang}{pl}{\let\ifpl\@firstoftwo}{\let\ifpl\@secondoftwo}%
  % \api{}, \pkg{} and \code{} are allowed in a section title, and a title
  % becomes a PDF bookmark. hyperref drops the \color command from a bookmark
  % string but keeps its ARGUMENT, so a title carrying \api{model\_validate}
  % made a bookmark reading "mtealmodel_validate" and warned "Token not
  % allowed" twice. The colour name is gobbled here instead. Measured on a
  % probe against this preamble before the line was added.
  \def\color#1{}%
}
\makeatother

% ---------- Language strings ----------
\input{lang/\booklang}

% The PDF's subject, keyword list and document language read \lblPdf... from
% the language file input on the line above, so they cannot sit in the colour
% \hypersetup block, which runs before it. Two blocks is the ordering, not an
% oversight; merging them raises Undefined control sequence and still writes
% a PDF, with the three fields empty.
\hypersetup{
  pdfsubject={\lblPdfSubject},
  pdfkeywords={\lblPdfKeywords},
  pdflang={\lblPdfLang}
}

% ---------- Headings ----------
% \raggedright on the title body is load-bearing: at \Huge on this block a
% title that cannot break overflows the margin by 30 pt or more, and several
% chapter titles here run past fifty characters.
\titleformat{\chapter}[display]
  {\normalfont\huge\bfseries\color{mblue}\raggedright}
  {\normalfont\normalsize\color{mgrey}\MakeUppercase{\chaptertitlename}\ \thechapter}
  {8pt}{\Huge\raggedright}
\titlespacing*{\chapter}{0pt}{10pt}{28pt}
\titleformat{\section}{\normalfont\Large\bfseries\color{mblue}\raggedright}{\thesection}{0.8em}{}
\titleformat{\subsection}{\normalfont\large\bfseries\color{mgrey}\raggedright}{\thesubsection}{0.7em}{}

% Sections in the contents, subsections not: fourteen chapters of subsections
% is a contents nobody reads, and the reader navigates by chapter and section.
\setcounter{tocdepth}{1}
\setcounter{secnumdepth}{2}

% One-sided running head: the chapter on the left, the folio on the right.
\pagestyle{fancy}
\fancyhf{}
\fancyhead[L]{\small\nouppercase{\leftmark}}
\fancyhead[R]{\small\thepage}
\renewcommand{\headrulewidth}{0.4pt}
\fancypagestyle{plain}{\fancyhf{}\fancyfoot[R]{\small\thepage}%
  \renewcommand{\headrulewidth}{0pt}}

% Drops the word "Chapter" from the head: the number is still there and the
% longer titles no longer overflow. MUST come after \pagestyle{fancy}, which
% installs its own \chaptermark and silently overwrites anything earlier.
\renewcommand{\chaptermark}[1]{\markboth{\thechapter.\ #1}{}}

% ============================================================
%  CODE LISTINGS
%  ---------------------------------------------------------
%  Every listing in a chapter is a FILE under code/, pulled in with \pyfile or
%  \pyregion, and CI runs it against the pinned versions. An in-source
%  \begin{python} block is allowed for a fragment of three or four lines that
%  exists to show a shape rather than to run; anything longer belongs in a
%  file the reader can open in the editor beside the PDF.
% ============================================================
\lstdefinelanguage{pythonx}{
  language=Python,
  morekeywords={
    async,await,match,case,nonlocal,yield,type,
    TypedDict,Annotated,Protocol,Literal,Generic,Final,ClassVar,Self,
    TypeVar,TypeGuard,TypeIs,Callable,Iterator,Iterable,Awaitable,
    dataclass,field,BaseModel,Field,Enum,StrEnum,NamedTuple,
    override,overload,property,staticmethod,classmethod
  },
  morekeywords=[2]{
    gather,create_task,TaskGroup,to_thread,timeout,wait_for,Semaphore,Queue,
    fixture,parametrize,monkeypatch,raises,mark,
    Depends,FastAPI,APIRouter,HTTPException,
    select,Session,Mapped,mapped_column,relationship,
    model_validate,model_validate_json,model_dump,model_json_schema
  },
  sensitive=true
}

% listings' C# definition predates async, records and pattern matching.
\lstdefinelanguage{csharpx}{
  language=[Sharp]C,
  morekeywords={
    async,await,var,dynamic,record,init,required,nameof,yield,partial,
    global,scoped,when,with,and,or,not,notnull,unmanaged,file,get,set,value,
    where,select,from,orderby,let,join,into,ascending,descending,on,equals,
    Task,ValueTask,CancellationToken,IEnumerable,IAsyncEnumerable,Func,Action
  },
  sensitive=true
}

\lstdefinestyle{bookcode}{
  basicstyle=\ttfamily\footnotesize,
  keywordstyle=\color{kw}\bfseries,
  keywordstyle=[2]\color{typ},
  stringstyle=\color{str},
  % Colour only, no \itshape: inconsolata (zi4) has no italic shape, so an
  % italic comment style is substituted with upright on every machine that
  % has the font, and the log carries "Font shape `T1/zi4/m/it' undefined"
  % on every build. The colour already marks a comment.
  commentstyle=\color{cmt},
  emphstyle=\color{dec}\bfseries,
  numbers=left,
  numberstyle=\ttfamily\tiny\color{mgrey},
  numbersep=8pt,
  backgroundcolor=\color{codebg},
  frame=single,
  rulecolor=\color{codeframe},
  framesep=6pt,
  breaklines=true,
  breakatwhitespace=false,
  postbreak=\mbox{\textcolor{mgrey}{$\hookrightarrow$}\space},
  showstringspaces=false,
  tabsize=4,
  columns=fullflexible,
  keepspaces=true,
  captionpos=b,
  aboveskip=1.2em,
  belowskip=1.2em,
  % Maps the characters most likely to be pasted into a listing by accident.
  %
  % Do NOT add a mapping for U+00A0 (non-breaking space). It looks like the
  % obvious next entry and it makes `listings` abort the run with "Improper
  % alphabetic constant" -- fatal, no PDF, and the message names neither the
  % character nor this line. The ASCII-in-listings convention, checked by
  % parity's C13, is the guard against it instead.
  literate={ł}{{\l}}1 {Ł}{{\L}}1 {ą}{{\k a}}1 {ę}{{\k e}}1
           {ś}{{\'s}}1 {ć}{{\'c}}1 {ż}{{\.z}}1 {ź}{{\'z}}1 {ó}{{\'o}}1 {ń}{{\'n}}1
           {Ą}{{\k A}}1 {Ę}{{\k E}}1 {Ś}{{\'S}}1 {Ć}{{\'C}}1
           {Ż}{{\.Z}}1 {Ź}{{\'Z}}1 {Ó}{{\'O}}1 {Ń}{{\'N}}1
           {—}{{-{}-}}2 {–}{{-}}1 {‘}{{'}}1 {’}{{'}}1
           {“}{{"}}1 {”}{{"}}1 {°}{{\textdegree}}1
}

\lstset{style=bookcode}

\lstnewenvironment{python}[1][]
  {\lstset{language=pythonx,style=bookcode,#1}}
  {}

% C# comparison snippets. The whole book is written for engineers arriving
% from .NET, so listings come in pairs often enough that the C# half has its
% own environment rather than a generic one.
\lstnewenvironment{csharp}[1][]
  {\lstset{language=csharpx,style=bookcode,#1}}
  {}

\lstnewenvironment{shellcmd}[1][]
  {\lstset{language=bash,style=bookcode,numbers=none,keywordstyle=\color{mteal}\bfseries,#1}}
  {}

\lstnewenvironment{tomlcode}[1][]
  {\lstset{style=bookcode,numbers=none,keywordstyle=\color{black},#1}}
  {}

\lstnewenvironment{yamlcode}[1][]
  {\lstset{style=bookcode,numbers=none,keywordstyle=\color{black},#1}}
  {}

\lstnewenvironment{jsoncode}[1][]
  {\lstset{style=bookcode,numbers=none,keywordstyle=\color{black},#1}}
  {}

% Terminal transcripts typed into the source. Allowed for a shell session the
% reader is meant to reproduce; NOT for program output, which is \transcript
% below, because an in-source transcript nobody ran reads exactly like one
% that was, and that is where a remembered number hides.
\lstnewenvironment{console}[1][]
  {\lstset{style=bookcode,numbers=none,basicstyle=\ttfamily\footnotesize,
           keywordstyle=\color{black},backgroundcolor=\color{white},#1}}
  {}

% The path is printed above every file listing, in the repository-relative
% form the reader opens in the editor beside the PDF. That is the whole point
% of a listing being a file, and the front matter promises it. The \nobreak
% keeps the path on the same page as the listing it names.
% ---------------------------------------------------------------------------
%  Appendix B prints the trap catalogue from notes/02-traps.md. One entry is
%  its number -- which a chapter may cite, and which is never reused -- the
%  habit in the reader's own voice, and what Python does instead.
%
%  The body is set RAGGED RIGHT on purpose. An entry is mostly \code{} spans,
%  which are unbreakable runs, and sixty-seven justified paragraphs of them
%  is sixty-seven chances at an overfull hbox in the edition with the longer
%  words. Ragged right removes the class rather than fixing it one box at a
%  time, and a reference list is the one place in this book where it reads
%  better anyway.
%
%  The habit is set with \enquote{} rather than \emph{}, and that is not a
%  style choice. Every habit here carries \code{} spans, and \code{} inside
%  \emph{} asks inconsolata for an italic it does not have: the first build
%  of this appendix raised `Font shape T1/zi4/m/it undefined', which
%  checklog.py is hard on for exactly this reason. It is the trap CLAUDE.md
%  records, met sixty-seven times at once -- and fixed once, because the
%  entries go through a macro. A quotation is also what the habit IS, and
%  \enquote{} is what the Polish edition owes anyway.
% ---------------------------------------------------------------------------
%  Appendix A's cheat sheet: one table per part, C# on the left, Python on
%  the right, and the chapter that teaches the row on the end. The three
%  header strings are arguments so that one environment serves both
%  editions and the structural signature stays identical.
%
%  Columns are ragged right for the reason Appendix B's entries are: a cell
%  is mostly \code{} spans, which are unbreakable runs, and a justified
%  narrow column full of them is an overfull box waiting for the edition
%  with the longer words.
% ---------------------------------------------------------------------------
%  Appendix C's tool matrix: the .NET tool, the Python tool this book uses,
%  the version it was verified against and the chapter that introduces it.
%
%  The version column prints \pydanticver and its siblings and is never
%  typed, which is the whole point of the appendix: it is the one page where
%  a reader sees every pin at once, and a pin that had to be re-typed here
%  would be a second copy of a fact the preamble already owns.
%  tools/check_structure.py --tools fails when a pin macro exists and this
%  table does not print it.
\NewDocumentEnvironment{toolmatrix}{mmmm}{%
  \small
  \begin{longtable}{@{}%
    >{\raggedright\arraybackslash}p{0.27\linewidth}%
    >{\raggedright\arraybackslash}p{0.31\linewidth}%
    >{\raggedright\arraybackslash}p{0.20\linewidth}%
    >{\raggedright\arraybackslash}p{0.08\linewidth}@{}}
  \toprule
  \textbf{#1} & \textbf{#2} & \textbf{#3} & \textbf{#4} \\
  \midrule
  \endfirsthead
  \toprule
  \textbf{#1} & \textbf{#2} & \textbf{#3} & \textbf{#4} \\
  \midrule
  \endhead
  \midrule
  \endfoot
  \bottomrule
  \endlastfoot
}{%
  \end{longtable}%
}

\NewDocumentEnvironment{cheatsheet}{mmm}{%
  \small
  \begin{longtable}{@{}%
    >{\raggedright\arraybackslash}p{0.41\linewidth}%
    >{\raggedright\arraybackslash}p{0.41\linewidth}%
    >{\raggedright\arraybackslash}p{0.08\linewidth}@{}}
  \toprule
  \textbf{#1} & \textbf{#2} & \textbf{#3} \\
  \midrule
  \endfirsthead
  \toprule
  \textbf{#1} & \textbf{#2} & \textbf{#3} \\
  \midrule
  \endhead
  \midrule
  \endfoot
  \bottomrule
  \endlastfoot
}{%
  \end{longtable}%
}

\newcommand{\trapentry}[3]{%
  \par\medskip
  \noindent\textbf{#1.}\enspace\enquote{#2}\par
  \nopagebreak
  {\raggedright\noindent#3\par}%
}

\newcommand{\listingpath}[1]{%
  \par\medskip\noindent
  {\ttfamily\footnotesize\color{mgrey}\detokenize{#1}}%
  \par\nobreak}

% External source file -- the preferred form. The path is relative to the
% repository root, which is also the path the reader opens in the editor.
% Usage: \pyfile{code/ch05/decorators.py}{Caption}{lst:label}
\newcommand{\pyfile}[3]{%
  \listingpath{#1}%
  \lstinputlisting[language=pythonx,style=bookcode,aboveskip=0.3em,
    caption={#2},label={#3}]{#1}%
}

% A named region of an external file, delimited in the source by
% `# --8<-- [start:name]` / `# --8<-- [end:name]`, so the book and the running
% code cannot drift apart. tools/check_structure.py --listings fails when the
% file or the region is missing.
%
% The markers are matched through listings' range prefix and suffix keys,
% with linerange={name-name}. The form this was inherited with -- a linerange
% whose two ends were the whole marker strings -- matched nothing and printed
% the WHOLE FILE, with no warning: measured on a three-region probe file
% before this was rewritten, and the same definition sits unused in the
% sibling book it came from. A region name is a word, never a number, because
% a number is a line number to listings. The line numbers printed are the
% file's own, so a region's listing says where in the file it sits.
% Usage: \pyregion{code/ch05/decorators.py}{timing}{Caption}{lst:label}
\lstset{
  rangebeginprefix=\#\ --8<--\ [start:, rangebeginsuffix=],
  rangeendprefix=\#\ --8<--\ [end:, rangeendsuffix=],
  includerangemarker=false}
\newcommand{\pyregion}[4]{%
  \listingpath{#1}%
  \lstinputlisting[language=pythonx,style=bookcode,aboveskip=0.3em,
    linerange={#2-#2},caption={#3},label={#4}]{#1}%
}

% The C# twin of \pyfile, for the comparison half of a pair.
\newcommand{\csfile}[3]{%
  \listingpath{#1}%
  \lstinputlisting[language=csharpx,style=bookcode,aboveskip=0.3em,
    caption={#2},label={#3}]{#1}%
}

% A program's output, written by a script under code/measure/ and pulled in
% whole. There is no typed form of this and there must not be: \val{} does
% not expand inside a verbatim body, so a transcript typed into a chapter is
% one that cannot quote a computed value, and a transcript nobody ran is
% indistinguishable on the page from one that was.
\newcommand{\transcript}[1]{%
  \par\smallskip
  \IfFileExists{figures/transcripts/#1.txt}{%
    \lstinputlisting[style=bookcode,numbers=none,basicstyle=\ttfamily\footnotesize,
      keywordstyle=\color{black},backgroundcolor=\color{white},
      language={}]{figures/transcripts/#1.txt}%
  }{%
    \begin{tcolorbox}[enhanced, sharp corners, width=\linewidth,
      colback=mlight, colframe=mgrey, boxrule=0.6pt,
      left=8pt, right=8pt, top=6pt, bottom=6pt]
      {\bfseries\color{mgrey}\small \lblTranscriptMissing}\\[2pt]
      {\ttfamily\scriptsize\color{mgrey} figures/transcripts/#1.txt}
    \end{tcolorbox}%
  }%
  \par\smallskip
}

% Inline literals. Underscores inside these must be written \_.
\newcommand{\code}[1]{\texttt{\small #1}}
\newcommand{\pkg}[1]{\texttt{\small\color{mblue}#1}}
\newcommand{\api}[1]{\texttt{\small\color{mteal}#1}}

% ============================================================
%  ADMONITIONS
%  ---------------------------------------------------------
%  A deliberately small, fixed vocabulary. Two visual treatments carry the
%  whole grammar: `tinted` is a block the reader may read in passing -- a
%  tint and a thick left bar, no outline; `outlined` is a block the reader
%  must stop at, with a full rule on white. A third treatment would mean the
%  first two are not carrying their meaning.
% ============================================================
\tcbset{
  mfa tinted/.style={
    enhanced, breakable, sharp corners,
    colback=mlight, boxrule=0pt, leftrule=3pt,
    left=8pt, right=8pt, top=6pt, bottom=6pt,
    before skip=13pt, after skip=13pt,
    fonttitle=\bfseries\small,
    coltitle=black, colbacktitle=mlight,
    attach title to upper={\par\smallskip},
  },
  mfa outlined/.style={
    enhanced, breakable, sharp corners,
    colback=white, boxrule=0.6pt,
    left=8pt, right=8pt, top=6pt, bottom=6pt,
    before skip=13pt, after skip=13pt,
    fonttitle=\bfseries\small,
    coltitle=black, colbacktitle=white,
    attach title to upper={\par\smallskip},
  },
}

% Read in passing.
\newtcolorbox{note}[1][]{mfa tinted, colframe=mblue, title={\lblNote}, #1}
\newtcolorbox{warning}[1][]{mfa tinted, colframe=mred, title={\lblWarning}, #1}
% The translation. The .NET reader is the book's only assumed reader, so this
% box carries the C# side of a comparison and nothing else. Its use is
% counted: a chapter with none is a chapter that has stopped translating.
\newtcolorbox{csbox}[1][]{mfa tinted, colframe=mteal, title={\lblCsharp}, #1}
% Anything that is preview, experimental or likely to move. Python moves more
% slowly than the agent frameworks do, and this box should be rare here.
\newtcolorbox{versionbox}[1][]{mfa tinted, colframe=mamber, title={\lblVersion}, #1}

% Stop and read.
% The misconception, named AFTER the reader has walked into it. Appendix B is
% the catalogue; a trap box in a chapter is one entry of it, in place.
\newtcolorbox{trapbox}[1][]{mfa outlined, colframe=mred, title={\lblTrap}, #1}
% Marks a listing that has not been executed against the pinned versions.
% `make debt` counts these, CI prints the count, and nothing ships to a reader
% with one attached. Removing it means you ran the code.
\newtcolorbox{verifybox}[1][]{mfa outlined, colframe=mgrey, title={\lblVerify}, #1}
\newtcolorbox{exercisebox}[1][]{mfa outlined, colframe=mteal, title={\lblExercise}, #1}
% The guiding project. Points at the stage of trace-assert this section
% builds, under code/src/trace_assert/.
\newtcolorbox{projectbox}[1][]{mfa outlined, colframe=mamber, title={\lblProject}, #1}

% ============================================================
%  EXERCISES
%  ---------------------------------------------------------
%  \begin{exercise}{key}{title} ... \end{exercise}
%
%  This is the mechanism the whole book is read through: the PDF on one half
%  of the screen, an editor on the other. An exercise is therefore not a
%  paragraph but a FILE the reader opens -- a starter under
%  code/exercises/chNN/<key>.py -- and a TEST that decides whether it is done,
%  test_<key>.py beside it. The box prints both paths and the one command
%  that runs the test, and the manifest at the back lists every exercise.
%
%  An ENVIRONMENT, not a macro with a body argument: a macro argument cannot
%  contain verbatim material, so an exercise written as \exercise{key}{title}
%  {body} could never show the reader a code fragment, and the first chapter
%  that tried would have died on `\begin{python}` inside the braces.
%
%  The key is the file stem, with underscores, and it is written with
%  \detokenize so the underscore reaches the page as a character rather than
%  as a subscript. tools/check_structure.py --exercises fails when the starter
%  or the test is missing, or when the key's chapter prefix disagrees with the
%  chapter it sits in.
% ============================================================
\newcounter{exercise}[chapter]
\renewcommand{\theexercise}{\thechapter.\arabic{exercise}}
\makeatletter
\newcommand{\exercisedir}{ch\two@digits{\value{chapter}}}
\newcommand{\listofexercises}{%
  \section*{\lblExerciseManifest}%
  \phantomsection\addcontentsline{toc}{section}{\lblExerciseManifest}%
  % The entries are subsection-level and tocdepth is 1 for the contents,
  % so the depth is raised for this list alone or it prints empty --
  % which it did, on the scaffold's first clean build.
  \begingroup\c@tocdepth=2\relax\@starttoc{exr}\endgroup%
}
\makeatother
\NewDocumentEnvironment{exercise}{mm}{%
  \refstepcounter{exercise}%
  \addcontentsline{exr}{subsection}{\protect\numberline{\theexercise}\texttt{\protect\detokenize{#1}} --- #2}%
  \begin{exercisebox}[title={\lblExercise~\theexercise\ \dash{} #2}]%
}{%
  \par\medskip
  {\small\setlength{\parindent}{0pt}%
  \lblExerciseStarter: \texttt{code/exercises/\exercisedir/\detokenize{#1}.py}\\
  \lblExerciseTest: \texttt{code/exercises/\exercisedir/test\_\detokenize{#1}.py}\\
  \lblExerciseRun: \texttt{cd code \&\& uv run pytest -k \detokenize{#1}}\par}%
  \end{exercisebox}%
}

% ============================================================
%  DIAGRAMS --- Mermaid pipeline
%  ---------------------------------------------------------
%  \mermaidfig{key}{caption}{one-line manifest description}
%
%  Source of truth is figures/mermaid/<lang>/<key>.mmd, committed. `make
%  diagrams` renders each to figures/diagrams/<lang>/<key>.pdf, which is build
%  output and is not committed, so a diagram change reviews as a text diff.
%  The source is per-language because a diagram with English labels in the
%  Polish edition is exactly the half-translation that makes a bilingual book
%  feel machine-made; CI checks that both languages carry the same keys.
%
%  If the rendered PDF is absent the macro draws a placeholder AND typesets
%  the Mermaid source, in the flow rather than as a float, because a source
%  typeset verbatim is often taller than a page and a float cannot break.
% ============================================================
\makeatletter
\newcommand{\listofdiagrams}{%
  \section*{\lblDiagramManifest}%
  \phantomsection\addcontentsline{toc}{section}{\lblDiagramManifest}%
  % The entries are subsection-level and tocdepth is 1 for the contents,
  % so the depth is raised for this list alone or it prints empty --
  % which it did, on the scaffold's first clean build.
  \begingroup\c@tocdepth=2\relax\@starttoc{dgm}\endgroup%
}
\makeatother

\newcommand{\mermaidfig}[3]{%
  % \phantomsection is load-bearing and not decoration. \addcontentsline
  % records hyperref's CURRENT anchor, and here that is whatever preceded
  % the figure -- for a figure placed after a listing it is the listing's
  % line-number anchor, and listings has already advanced the counter past
  % the last line it printed, so the manifest entry linked to a
  % destination that is never created: "name{lstnumber.10.4.46} has been
  % referenced but does not exist, replaced by a fixed one". The two
  % front-matter figures had the same wrong target and were silent about
  % it, because the lines they named happen to be printed. The exercise
  % manifest never had the bug, because \begin{exercise} steps a counter
  % and that is what sets the anchor.
  \phantomsection%
  \addcontentsline{dgm}{subsection}{\protect\numberline{}\texttt{#1.mmd} --- #3}%
  \IfFileExists{figures/diagrams/\booklang/#1.pdf}{%
    \begin{figure}[htbp]
    \centering
    \includegraphics[width=0.95\linewidth,height=0.42\textheight,keepaspectratio]{figures/diagrams/\booklang/#1.pdf}%
    \caption{#2}\label{fig:#1}
    \end{figure}%
  }{%
    \par\medskip
    \begin{tcolorbox}[
      enhanced, breakable, width=\linewidth, sharp corners,
      colback=white, colframe=mgrey, boxrule=0.8pt,
      left=8pt, right=8pt, top=8pt, bottom=8pt]
      {\bfseries\color{mgrey}\small \lblNotRendered}\\[4pt]
      {\ttfamily\scriptsize\color{mgrey} figures/mermaid/\booklang/#1.mmd}\\[6pt]
      \IfFileExists{figures/mermaid/\booklang/#1.mmd}{%
        \lstinputlisting[style=bookcode,numbers=none,basicstyle=\ttfamily\scriptsize,
                         backgroundcolor=\color{mlight},frame=none,
                         keywordstyle=\color{black},language={}]{figures/mermaid/\booklang/#1.mmd}%
      }{%
        {\small\color{mred} \lblSourceMissing: \texttt{figures/mermaid/\booklang/#1.mmd}}%
      }%
    \end{tcolorbox}%
    \captionof{figure}{#2}\label{fig:#1}
    \par\medskip
  }%
}

% ============================================================
%  CHAPTER STUBS
%  ---------------------------------------------------------
%  Prints a visible marker so an unwritten chapter cannot be mistaken for a
%  written one and the page count never flatters the draft. `make debt`
%  counts them, and CI publishes the count on every build.
% ============================================================
\newcommand{\chapterstub}[1]{%
  \begin{tcolorbox}[
    enhanced, breakable, width=\linewidth, sharp corners,
    colback=mlight, colframe=mred, boxrule=1pt,
    borderline={0.6pt}{3pt}{mred, dashed},
    left=12pt, right=12pt, top=12pt, bottom=12pt]
    {\bfseries\color{mred}\large \lblNotWritten}\\[8pt]
    \small\raggedright #1
  \end{tcolorbox}%
}

% ============================================================
%  COMPUTED VALUES
%  ---------------------------------------------------------
%  A number that came out of a measurement is written by a script under
%  code/measure/ into figures/values/<name>.tex as
%
%      \pyval{e01.threads.speedup}{3.71}
%
%  and pulled into the text with \val{e01.threads.speedup}. The script is the
%  source of truth and the book cannot disagree with it, because the book does
%  not contain the digits. `make verify` fails when a committed value no
%  longer matches its script. A value the scripts do not produce prints a
%  visible marker.
%
%  \val also applies the edition's decimal separator, which is how the Polish
%  edition gets its decimal comma without a translator having to remember.
% ============================================================
\newcommand{\bookdecimalmarker}{\ifpl{,}{.}}
\IfFileExists{siunitx.sty}{%
  \usepackage{siunitx}%
  \sisetup{
    group-digits = integer,
    group-minimum-digits = 5,
    group-separator = {\,},
    output-decimal-marker = {\bookdecimalmarker}
  }%
  \newcommand{\booknum}[1]{\num{##1}}%
}{%
  \newcommand{\booknum}[1]{##1}%
}

\newcommand{\pyvalues}{}
\makeatletter
\newcommand{\pyval}[2]{\csgdef{pyval@#1}{#2}\listxadd{\pyvalues}{#1}}
% \num on a non-numeric body is a FATAL siunitx error. The generator knows
% whether it emitted a number or a word, so it writes \pyval for a number and
% \pyvaltext for a word, and this end only enforces which macro reads which.
\newcommand{\val}[1]{%
  \ifcsdef{pyval@#1}{%
    \ifcsdef{pyvaltext@#1}{%
      \PackageError{pybook}{Value `#1' is not a number}%
        {\string\val\space passes its body to siunitx, which refuses anything
         that is not a number. Use \string\valtext{#1} instead.}%
    }{}%
    \booknum{\csuse{pyval@#1}}%
  }{\textcolor{mred}{\textbf{[\lblValueMissing: #1]}}}%
}
\newcommand{\valtext}[1]{%
  \ifcsdef{pyval@#1}{\csuse{pyval@#1}}%
    {\textcolor{mred}{\textbf{[\lblValueMissing: #1]}}}%
}
\newcommand{\pyvaltext}[2]{%
  \csgdef{pyval@#1}{#2}\csgdef{pyvaltext@#1}{}\listxadd{\pyvalues}{#1}%
}
\makeatother

% figures/values/all.tex is generated by `make numbers` and does nothing but
% \input each value file. A build with no values yet still compiles; every
% \val{} in it prints a visible marker instead of a number.
\IfFileExists{figures/values/all.tex}{\input{figures/values/all}}{%
  \AtBeginDocument{\PackageWarning{pybook}{No computed values found. Run `make numbers`.}}}

% ============================================================
%  PINNED VERSIONS --- single source of truth
%  Change these here; they propagate through both editions and Appendix C.
%  Verified against pypi.org on \pinnedon (lang/<lang>.tex).
%  tools/check_versions.py compares every one of them with PyPI on each build.
%  Python itself is not on PyPI and is checked by hand against python.org.
% ============================================================
\newcommand{\bookedition}{0.1-dev}
\newcommand{\bookyear}{2026}
\newcommand{\pyver}{3.14}              % the minor the book targets
\newcommand{\pypatch}{3.14.7}          % the exact interpreter the code ran on
\newcommand{\uvver}{0.12.13}           % uv
\newcommand{\dotnetver}{10.0.100}      % dotnet
\newcommand{\ruffver}{0.16.7}          % ruff
\newcommand{\pyrightver}{1.1.414}      % pyright
\newcommand{\mypyver}{2.3.1}           % mypy
\newcommand{\pydanticver}{2.13.5}      % pydantic
\newcommand{\pysettingsver}{2.15.0}    % pydantic-settings
\newcommand{\fastapiver}{0.141.1}      % fastapi
\newcommand{\uvicornver}{0.53.0}       % uvicorn
\newcommand{\sqlalchemyver}{2.0.52}    % sqlalchemy
\newcommand{\alembicver}{1.20.0}       % alembic
\newcommand{\aiosqlitever}{0.22.1}     % aiosqlite
\newcommand{\pytestver}{9.1.1}         % pytest
\newcommand{\pytestasyncver}{1.4.0}    % pytest-asyncio
\newcommand{\hypothesisver}{6.168.0}   % hypothesis
\newcommand{\coveragever}{7.16.1}      % coverage
\newcommand{\httpxver}{0.28.1}         % httpx
\newcommand{\structlogver}{26.1.0}     % structlog
\newcommand{\otelver}{1.44.0}          % opentelemetry-sdk
\newcommand{\polarsver}{1.44.2}        % polars
\newcommand{\anthropicver}{1.5.0}      % anthropic
\newcommand{\openaiver}{3.13.0}        % openai
\newcommand{\pipauditver}{2.10.1}      % pip-audit

% ============================================================
%  INDEX
%  Two-column, and its entries include \texttt{} names that do not hyphenate.
%  Ragged right keeps multi-word entries off the margin; it cannot help a
%  single token wider than the column, so keep verbatim index entries under
%  about 29 characters and index longer names by concept instead.
%  \raggedcolumns is the switch imakeidx's multicols actually reads.
% ============================================================
\apptocmd{\theindex}{\raggedcolumns\raggedright}{}{%
  \PackageWarning{pybook}{Could not patch theindex for ragged setting}}
\providecommand*\see[2]{}
\renewcommand*\see[2]{\emph{\lblIndexSee} #1}
\providecommand*\seealso[2]{}
\renewcommand*\seealso[2]{\emph{\lblIndexSeeAlso} #1}

\makeindex
 still compiles and says what it did.
\providecommand{\booklang}{en}

\usepackage{etoolbox}   % loaded first: the language switch below needs it

% ---------- Page geometry ----------
% ONE format: A4, 12pt, ONE-SIDED. This book is read on a screen with an
% editor open beside it, and never printed as a bound book, so there is no
% recto, no verso, no gutter, no mirrored margin and no blank leaf before a
% chapter -- every one of those is a property of paper that a PDF viewer at
% half a screen's width turns into wasted pixels.
%
% The margins are symmetric and the measure is about seventy-five characters
% of prose, which is where the eye keeps the line return. The number that was
% actually chosen is the MONOSPACE measure: at \footnotesize on this block a
% listing line of 79 characters fits without wrapping, and 79 is the width
% every listing under code/ is held to by tools/check_structure.py --lines.
% A wrapped listing line is not an error TeX reports -- breaklines=true wraps
% it silently and prints an arrow into the middle of what the reader is meant
% to paste into an editor -- so the width is enforced in the source instead.
\usepackage[
  a4paper,
  left=3.1cm, right=3.1cm, top=2.6cm, bottom=3.0cm,
  head=26pt, headsep=0.8cm, footskip=1.4cm
]{geometry}

% ---------- Encoding & fonts ----------
\usepackage[T1]{fontenc}
\usepackage[utf8]{inputenc}
\IfFileExists{lmodern.sty}{\usepackage{lmodern}}{}
% newtxtext takes its TS1 symbols (\textcopyright, \textdegree) from TeX
% Gyre Termes, which Debian packages SEPARATELY (tex-gyre): a machine with
% newtx and without it dies on the copyright page with
% `Font TS1/ntxtlf/m/n/10.95=ts1-qtmr ... not loadable`. There is no
% inert-safe probe for a font metric in pdfTeX -- \IfFileExists searches
% TEXINPUTS and not TFMFONTS, so it reports the metric missing on a machine
% that has it, and \suppressfontnotfounderror is LuaTeX-only -- so nothing
% here guesses. CI's full TeX Live is the reference; locally, install
% tex-gyre. Both facts are recorded in CLAUDE.md under Build traps.
\IfFileExists{newtxtext.sty}{\usepackage{newtxtext}}{}
% amsmath before newtxmath, because newtxmath patches amsmath's internals.
% amssymb only when newtxmath is absent: newtxmath supplies the AMS symbols
% itself and loading both is a fatal clash (\Bbbk already defined) that is
% invisible on a machine without newtx and fatal on a full TeX Live.
\usepackage{amsmath}
\IfFileExists{newtxmath.sty}{\usepackage{newtxmath}}{\usepackage{amssymb}}
\IfFileExists{inconsolata.sty}{\usepackage[varqu,scaled=0.95]{inconsolata}}{}
\IfFileExists{lmodern.sty}{\usepackage{microtype}}{\usepackage[expansion=false]{microtype}}

% upquote is not a refinement. In T1 the input character ' is quoteright, so
% a listings body sets f('x') with U+2019 at both ends of the literal, and
% typed back in that is a SyntaxError. inconsolata's varqu option hides the
% defect on a machine that has inconsolata; upquote makes ' the ASCII
% apostrophe whatever the typewriter font is. In a book whose every listing is
% meant to be pasted into an editor, this line is load-bearing.
\IfFileExists{upquote.sty}{\usepackage{upquote}}{}

% And upquote's twin, for the same reason one character over. In T1 the pair
% -- is a ligature for an en dash, in EVERY family including the typewriter
% one, so \code{uv sync --locked} sets as "uv sync <en dash>locked" and a
% reader who copies it off the page gets an unknown-argument error from a
% command the chapter told them to run. It is invisible in review because it
% looks like a dash, and it bites exactly the flags a packaging chapter is
% made of. A `listings` body is already safe -- listings sets characters
% singly and no ligature forms -- which is what makes the defect so easy to
% miss: the same text is right inside a listing and wrong in the sentence
% above it.
%
% Scoped to tt* on purpose, and measured rather than assumed: prose keeps its
% own dashes, so \dash{} still sets an em dash and a range like 1--2 still
% sets an en dash. Probed against this preamble in all three positions before
% it was believed.
\DisableLigatures[-]{encoding = T1, family = tt*}

% ---------- Language ----------
% Loading babel with a language whose .ldf is absent is a *fatal* error, not
% a warning, so both the package and each language file are probed before
% either is requested. Both editions load the other language too, because the
% glossary rows and the cheat sheet carry words from both.
\IfFileExists{babel.sty}{%
  \IfFileExists{polish.ldf}{%
    \IfFileExists{british.ldf}{%
      \ifdefstring{\booklang}{pl}%
        {\usepackage[british,main=polish]{babel}}%
        {\usepackage[polish,main=british]{babel}}%
    }{\usepackage[polish]{babel}}%
  }{%
    \IfFileExists{british.ldf}{\usepackage[british]{babel}}{}%
  }%
}{}

% babel-polish makes " an active shorthand. listings reads its body verbatim
% and an active character in a verbatim context is a category-code accident
% waiting to happen, so the shorthand is turned off. Polish quotation marks
% come from \enquote, which csquotes sets per language.
\ifdefstring{\booklang}{pl}{\AtBeginDocument{\shorthandoff{"}}}{}

% ---------- Core packages ----------
\usepackage{graphicx}
\usepackage{xcolor}
\usepackage{booktabs}
\usepackage{tabularx}
% Appendix A's cheat sheet runs past a page in every part, and tabularx
% cannot break. longtable can, and works with booktabs.
\usepackage{longtable}
\usepackage{array}
\usepackage{enumitem}
\usepackage{listings}
\usepackage[most]{tcolorbox}
\usepackage{fancyhdr}
\usepackage{titlesec}
\usepackage{caption}
\usepackage{imakeidx}
% csquotes with autostyle follows babel's active language, so the identical
% source \enquote{x} sets 'x' in the English edition and the Polish quotation
% marks in the Polish one.
\IfFileExists{csquotes.sty}{\usepackage[autostyle=true]{csquotes}}{%
  \providecommand{\enquote}[1]{``##1''}}
\usepackage[hidelinks]{hyperref}

% ---------- Palette ----------
% The same family as the two companion volumes, so the three read as a set.
\definecolor{mblue}{HTML}{1F4E79}
\definecolor{mteal}{HTML}{0E7C7B}
\definecolor{mamber}{HTML}{B26A00}
\definecolor{mred}{HTML}{9B2C2C}
\definecolor{mviolet}{HTML}{7B4397}
\definecolor{mgrey}{HTML}{5A5A5A}
\definecolor{mlight}{HTML}{F4F6F8}
\definecolor{mfaint}{HTML}{FBFCFD}
\definecolor{codebg}{HTML}{FAFAFA}
\definecolor{codeframe}{HTML}{D8DEE4}
\definecolor{kw}{HTML}{0B5394}
\definecolor{str}{HTML}{9B2C2C}
\definecolor{cmt}{HTML}{4F7A4F}
\definecolor{typ}{HTML}{0E7C7B}
\definecolor{dec}{HTML}{7B4397}

\hypersetup{
  colorlinks=true,
  linkcolor=mblue,
  urlcolor=mteal,
  citecolor=mblue,
  pdfauthor={Konrad Cinkusz}
}

% \ifpl{Polish text}{English text} -- for the handful of places where a
% string is structural rather than content. Anything longer than a few words
% belongs in lang/ or in the edition's own source, not here.
% \DeclareRobustCommand, not \newcommand: a fragile \ifpl inside a moving
% argument is written to the .toc one level expanded and fails on the next run
% in exactly one of the two editions.
\DeclareRobustCommand{\ifpl}[2]{\ifdefstring{\booklang}{pl}{#1}{#2}}
\makeatletter
\pdfstringdefDisableCommands{%
  \ifdefstring{\booklang}{pl}{\let\ifpl\@firstoftwo}{\let\ifpl\@secondoftwo}%
  % \api{}, \pkg{} and \code{} are allowed in a section title, and a title
  % becomes a PDF bookmark. hyperref drops the \color command from a bookmark
  % string but keeps its ARGUMENT, so a title carrying \api{model\_validate}
  % made a bookmark reading "mtealmodel_validate" and warned "Token not
  % allowed" twice. The colour name is gobbled here instead. Measured on a
  % probe against this preamble before the line was added.
  \def\color#1{}%
}
\makeatother

% ---------- Language strings ----------
\input{lang/\booklang}

% The PDF's subject, keyword list and document language read \lblPdf... from
% the language file input on the line above, so they cannot sit in the colour
% \hypersetup block, which runs before it. Two blocks is the ordering, not an
% oversight; merging them raises Undefined control sequence and still writes
% a PDF, with the three fields empty.
\hypersetup{
  pdfsubject={\lblPdfSubject},
  pdfkeywords={\lblPdfKeywords},
  pdflang={\lblPdfLang}
}

% ---------- Headings ----------
% \raggedright on the title body is load-bearing: at \Huge on this block a
% title that cannot break overflows the margin by 30 pt or more, and several
% chapter titles here run past fifty characters.
\titleformat{\chapter}[display]
  {\normalfont\huge\bfseries\color{mblue}\raggedright}
  {\normalfont\normalsize\color{mgrey}\MakeUppercase{\chaptertitlename}\ \thechapter}
  {8pt}{\Huge\raggedright}
\titlespacing*{\chapter}{0pt}{10pt}{28pt}
\titleformat{\section}{\normalfont\Large\bfseries\color{mblue}\raggedright}{\thesection}{0.8em}{}
\titleformat{\subsection}{\normalfont\large\bfseries\color{mgrey}\raggedright}{\thesubsection}{0.7em}{}

% Sections in the contents, subsections not: fourteen chapters of subsections
% is a contents nobody reads, and the reader navigates by chapter and section.
\setcounter{tocdepth}{1}
\setcounter{secnumdepth}{2}

% One-sided running head: the chapter on the left, the folio on the right.
\pagestyle{fancy}
\fancyhf{}
\fancyhead[L]{\small\nouppercase{\leftmark}}
\fancyhead[R]{\small\thepage}
\renewcommand{\headrulewidth}{0.4pt}
\fancypagestyle{plain}{\fancyhf{}\fancyfoot[R]{\small\thepage}%
  \renewcommand{\headrulewidth}{0pt}}

% Drops the word "Chapter" from the head: the number is still there and the
% longer titles no longer overflow. MUST come after \pagestyle{fancy}, which
% installs its own \chaptermark and silently overwrites anything earlier.
\renewcommand{\chaptermark}[1]{\markboth{\thechapter.\ #1}{}}

% ============================================================
%  CODE LISTINGS
%  ---------------------------------------------------------
%  Every listing in a chapter is a FILE under code/, pulled in with \pyfile or
%  \pyregion, and CI runs it against the pinned versions. An in-source
%  \begin{python} block is allowed for a fragment of three or four lines that
%  exists to show a shape rather than to run; anything longer belongs in a
%  file the reader can open in the editor beside the PDF.
% ============================================================
\lstdefinelanguage{pythonx}{
  language=Python,
  morekeywords={
    async,await,match,case,nonlocal,yield,type,
    TypedDict,Annotated,Protocol,Literal,Generic,Final,ClassVar,Self,
    TypeVar,TypeGuard,TypeIs,Callable,Iterator,Iterable,Awaitable,
    dataclass,field,BaseModel,Field,Enum,StrEnum,NamedTuple,
    override,overload,property,staticmethod,classmethod
  },
  morekeywords=[2]{
    gather,create_task,TaskGroup,to_thread,timeout,wait_for,Semaphore,Queue,
    fixture,parametrize,monkeypatch,raises,mark,
    Depends,FastAPI,APIRouter,HTTPException,
    select,Session,Mapped,mapped_column,relationship,
    model_validate,model_validate_json,model_dump,model_json_schema
  },
  sensitive=true
}

% listings' C# definition predates async, records and pattern matching.
\lstdefinelanguage{csharpx}{
  language=[Sharp]C,
  morekeywords={
    async,await,var,dynamic,record,init,required,nameof,yield,partial,
    global,scoped,when,with,and,or,not,notnull,unmanaged,file,get,set,value,
    where,select,from,orderby,let,join,into,ascending,descending,on,equals,
    Task,ValueTask,CancellationToken,IEnumerable,IAsyncEnumerable,Func,Action
  },
  sensitive=true
}

\lstdefinestyle{bookcode}{
  basicstyle=\ttfamily\footnotesize,
  keywordstyle=\color{kw}\bfseries,
  keywordstyle=[2]\color{typ},
  stringstyle=\color{str},
  % Colour only, no \itshape: inconsolata (zi4) has no italic shape, so an
  % italic comment style is substituted with upright on every machine that
  % has the font, and the log carries "Font shape `T1/zi4/m/it' undefined"
  % on every build. The colour already marks a comment.
  commentstyle=\color{cmt},
  emphstyle=\color{dec}\bfseries,
  numbers=left,
  numberstyle=\ttfamily\tiny\color{mgrey},
  numbersep=8pt,
  backgroundcolor=\color{codebg},
  frame=single,
  rulecolor=\color{codeframe},
  framesep=6pt,
  breaklines=true,
  breakatwhitespace=false,
  postbreak=\mbox{\textcolor{mgrey}{$\hookrightarrow$}\space},
  showstringspaces=false,
  tabsize=4,
  columns=fullflexible,
  keepspaces=true,
  captionpos=b,
  aboveskip=1.2em,
  belowskip=1.2em,
  % Maps the characters most likely to be pasted into a listing by accident.
  %
  % Do NOT add a mapping for U+00A0 (non-breaking space). It looks like the
  % obvious next entry and it makes `listings` abort the run with "Improper
  % alphabetic constant" -- fatal, no PDF, and the message names neither the
  % character nor this line. The ASCII-in-listings convention, checked by
  % parity's C13, is the guard against it instead.
  literate={ł}{{\l}}1 {Ł}{{\L}}1 {ą}{{\k a}}1 {ę}{{\k e}}1
           {ś}{{\'s}}1 {ć}{{\'c}}1 {ż}{{\.z}}1 {ź}{{\'z}}1 {ó}{{\'o}}1 {ń}{{\'n}}1
           {Ą}{{\k A}}1 {Ę}{{\k E}}1 {Ś}{{\'S}}1 {Ć}{{\'C}}1
           {Ż}{{\.Z}}1 {Ź}{{\'Z}}1 {Ó}{{\'O}}1 {Ń}{{\'N}}1
           {—}{{-{}-}}2 {–}{{-}}1 {‘}{{'}}1 {’}{{'}}1
           {“}{{"}}1 {”}{{"}}1 {°}{{\textdegree}}1
}

\lstset{style=bookcode}

\lstnewenvironment{python}[1][]
  {\lstset{language=pythonx,style=bookcode,#1}}
  {}

% C# comparison snippets. The whole book is written for engineers arriving
% from .NET, so listings come in pairs often enough that the C# half has its
% own environment rather than a generic one.
\lstnewenvironment{csharp}[1][]
  {\lstset{language=csharpx,style=bookcode,#1}}
  {}

\lstnewenvironment{shellcmd}[1][]
  {\lstset{language=bash,style=bookcode,numbers=none,keywordstyle=\color{mteal}\bfseries,#1}}
  {}

\lstnewenvironment{tomlcode}[1][]
  {\lstset{style=bookcode,numbers=none,keywordstyle=\color{black},#1}}
  {}

\lstnewenvironment{yamlcode}[1][]
  {\lstset{style=bookcode,numbers=none,keywordstyle=\color{black},#1}}
  {}

\lstnewenvironment{jsoncode}[1][]
  {\lstset{style=bookcode,numbers=none,keywordstyle=\color{black},#1}}
  {}

% Terminal transcripts typed into the source. Allowed for a shell session the
% reader is meant to reproduce; NOT for program output, which is \transcript
% below, because an in-source transcript nobody ran reads exactly like one
% that was, and that is where a remembered number hides.
\lstnewenvironment{console}[1][]
  {\lstset{style=bookcode,numbers=none,basicstyle=\ttfamily\footnotesize,
           keywordstyle=\color{black},backgroundcolor=\color{white},#1}}
  {}

% The path is printed above every file listing, in the repository-relative
% form the reader opens in the editor beside the PDF. That is the whole point
% of a listing being a file, and the front matter promises it. The \nobreak
% keeps the path on the same page as the listing it names.
% ---------------------------------------------------------------------------
%  Appendix B prints the trap catalogue from notes/02-traps.md. One entry is
%  its number -- which a chapter may cite, and which is never reused -- the
%  habit in the reader's own voice, and what Python does instead.
%
%  The body is set RAGGED RIGHT on purpose. An entry is mostly \code{} spans,
%  which are unbreakable runs, and sixty-seven justified paragraphs of them
%  is sixty-seven chances at an overfull hbox in the edition with the longer
%  words. Ragged right removes the class rather than fixing it one box at a
%  time, and a reference list is the one place in this book where it reads
%  better anyway.
%
%  The habit is set with \enquote{} rather than \emph{}, and that is not a
%  style choice. Every habit here carries \code{} spans, and \code{} inside
%  \emph{} asks inconsolata for an italic it does not have: the first build
%  of this appendix raised `Font shape T1/zi4/m/it undefined', which
%  checklog.py is hard on for exactly this reason. It is the trap CLAUDE.md
%  records, met sixty-seven times at once -- and fixed once, because the
%  entries go through a macro. A quotation is also what the habit IS, and
%  \enquote{} is what the Polish edition owes anyway.
% ---------------------------------------------------------------------------
%  Appendix A's cheat sheet: one table per part, C# on the left, Python on
%  the right, and the chapter that teaches the row on the end. The three
%  header strings are arguments so that one environment serves both
%  editions and the structural signature stays identical.
%
%  Columns are ragged right for the reason Appendix B's entries are: a cell
%  is mostly \code{} spans, which are unbreakable runs, and a justified
%  narrow column full of them is an overfull box waiting for the edition
%  with the longer words.
% ---------------------------------------------------------------------------
%  Appendix C's tool matrix: the .NET tool, the Python tool this book uses,
%  the version it was verified against and the chapter that introduces it.
%
%  The version column prints \pydanticver and its siblings and is never
%  typed, which is the whole point of the appendix: it is the one page where
%  a reader sees every pin at once, and a pin that had to be re-typed here
%  would be a second copy of a fact the preamble already owns.
%  tools/check_structure.py --tools fails when a pin macro exists and this
%  table does not print it.
\NewDocumentEnvironment{toolmatrix}{mmmm}{%
  \small
  \begin{longtable}{@{}%
    >{\raggedright\arraybackslash}p{0.27\linewidth}%
    >{\raggedright\arraybackslash}p{0.31\linewidth}%
    >{\raggedright\arraybackslash}p{0.20\linewidth}%
    >{\raggedright\arraybackslash}p{0.08\linewidth}@{}}
  \toprule
  \textbf{#1} & \textbf{#2} & \textbf{#3} & \textbf{#4} \\
  \midrule
  \endfirsthead
  \toprule
  \textbf{#1} & \textbf{#2} & \textbf{#3} & \textbf{#4} \\
  \midrule
  \endhead
  \midrule
  \endfoot
  \bottomrule
  \endlastfoot
}{%
  \end{longtable}%
}

\NewDocumentEnvironment{cheatsheet}{mmm}{%
  \small
  \begin{longtable}{@{}%
    >{\raggedright\arraybackslash}p{0.41\linewidth}%
    >{\raggedright\arraybackslash}p{0.41\linewidth}%
    >{\raggedright\arraybackslash}p{0.08\linewidth}@{}}
  \toprule
  \textbf{#1} & \textbf{#2} & \textbf{#3} \\
  \midrule
  \endfirsthead
  \toprule
  \textbf{#1} & \textbf{#2} & \textbf{#3} \\
  \midrule
  \endhead
  \midrule
  \endfoot
  \bottomrule
  \endlastfoot
}{%
  \end{longtable}%
}

\newcommand{\trapentry}[3]{%
  \par\medskip
  \noindent\textbf{#1.}\enspace\enquote{#2}\par
  \nopagebreak
  {\raggedright\noindent#3\par}%
}

\newcommand{\listingpath}[1]{%
  \par\medskip\noindent
  {\ttfamily\footnotesize\color{mgrey}\detokenize{#1}}%
  \par\nobreak}

% External source file -- the preferred form. The path is relative to the
% repository root, which is also the path the reader opens in the editor.
% Usage: \pyfile{code/ch05/decorators.py}{Caption}{lst:label}
\newcommand{\pyfile}[3]{%
  \listingpath{#1}%
  \lstinputlisting[language=pythonx,style=bookcode,aboveskip=0.3em,
    caption={#2},label={#3}]{#1}%
}

% A named region of an external file, delimited in the source by
% `# --8<-- [start:name]` / `# --8<-- [end:name]`, so the book and the running
% code cannot drift apart. tools/check_structure.py --listings fails when the
% file or the region is missing.
%
% The markers are matched through listings' range prefix and suffix keys,
% with linerange={name-name}. The form this was inherited with -- a linerange
% whose two ends were the whole marker strings -- matched nothing and printed
% the WHOLE FILE, with no warning: measured on a three-region probe file
% before this was rewritten, and the same definition sits unused in the
% sibling book it came from. A region name is a word, never a number, because
% a number is a line number to listings. The line numbers printed are the
% file's own, so a region's listing says where in the file it sits.
% Usage: \pyregion{code/ch05/decorators.py}{timing}{Caption}{lst:label}
\lstset{
  rangebeginprefix=\#\ --8<--\ [start:, rangebeginsuffix=],
  rangeendprefix=\#\ --8<--\ [end:, rangeendsuffix=],
  includerangemarker=false}
\newcommand{\pyregion}[4]{%
  \listingpath{#1}%
  \lstinputlisting[language=pythonx,style=bookcode,aboveskip=0.3em,
    linerange={#2-#2},caption={#3},label={#4}]{#1}%
}

% The C# twin of \pyfile, for the comparison half of a pair.
\newcommand{\csfile}[3]{%
  \listingpath{#1}%
  \lstinputlisting[language=csharpx,style=bookcode,aboveskip=0.3em,
    caption={#2},label={#3}]{#1}%
}

% A program's output, written by a script under code/measure/ and pulled in
% whole. There is no typed form of this and there must not be: \val{} does
% not expand inside a verbatim body, so a transcript typed into a chapter is
% one that cannot quote a computed value, and a transcript nobody ran is
% indistinguishable on the page from one that was.
\newcommand{\transcript}[1]{%
  \par\smallskip
  \IfFileExists{figures/transcripts/#1.txt}{%
    \lstinputlisting[style=bookcode,numbers=none,basicstyle=\ttfamily\footnotesize,
      keywordstyle=\color{black},backgroundcolor=\color{white},
      language={}]{figures/transcripts/#1.txt}%
  }{%
    \begin{tcolorbox}[enhanced, sharp corners, width=\linewidth,
      colback=mlight, colframe=mgrey, boxrule=0.6pt,
      left=8pt, right=8pt, top=6pt, bottom=6pt]
      {\bfseries\color{mgrey}\small \lblTranscriptMissing}\\[2pt]
      {\ttfamily\scriptsize\color{mgrey} figures/transcripts/#1.txt}
    \end{tcolorbox}%
  }%
  \par\smallskip
}

% Inline literals. Underscores inside these must be written \_.
\newcommand{\code}[1]{\texttt{\small #1}}
\newcommand{\pkg}[1]{\texttt{\small\color{mblue}#1}}
\newcommand{\api}[1]{\texttt{\small\color{mteal}#1}}

% ============================================================
%  ADMONITIONS
%  ---------------------------------------------------------
%  A deliberately small, fixed vocabulary. Two visual treatments carry the
%  whole grammar: `tinted` is a block the reader may read in passing -- a
%  tint and a thick left bar, no outline; `outlined` is a block the reader
%  must stop at, with a full rule on white. A third treatment would mean the
%  first two are not carrying their meaning.
% ============================================================
\tcbset{
  mfa tinted/.style={
    enhanced, breakable, sharp corners,
    colback=mlight, boxrule=0pt, leftrule=3pt,
    left=8pt, right=8pt, top=6pt, bottom=6pt,
    before skip=13pt, after skip=13pt,
    fonttitle=\bfseries\small,
    coltitle=black, colbacktitle=mlight,
    attach title to upper={\par\smallskip},
  },
  mfa outlined/.style={
    enhanced, breakable, sharp corners,
    colback=white, boxrule=0.6pt,
    left=8pt, right=8pt, top=6pt, bottom=6pt,
    before skip=13pt, after skip=13pt,
    fonttitle=\bfseries\small,
    coltitle=black, colbacktitle=white,
    attach title to upper={\par\smallskip},
  },
}

% Read in passing.
\newtcolorbox{note}[1][]{mfa tinted, colframe=mblue, title={\lblNote}, #1}
\newtcolorbox{warning}[1][]{mfa tinted, colframe=mred, title={\lblWarning}, #1}
% The translation. The .NET reader is the book's only assumed reader, so this
% box carries the C# side of a comparison and nothing else. Its use is
% counted: a chapter with none is a chapter that has stopped translating.
\newtcolorbox{csbox}[1][]{mfa tinted, colframe=mteal, title={\lblCsharp}, #1}
% Anything that is preview, experimental or likely to move. Python moves more
% slowly than the agent frameworks do, and this box should be rare here.
\newtcolorbox{versionbox}[1][]{mfa tinted, colframe=mamber, title={\lblVersion}, #1}

% Stop and read.
% The misconception, named AFTER the reader has walked into it. Appendix B is
% the catalogue; a trap box in a chapter is one entry of it, in place.
\newtcolorbox{trapbox}[1][]{mfa outlined, colframe=mred, title={\lblTrap}, #1}
% Marks a listing that has not been executed against the pinned versions.
% `make debt` counts these, CI prints the count, and nothing ships to a reader
% with one attached. Removing it means you ran the code.
\newtcolorbox{verifybox}[1][]{mfa outlined, colframe=mgrey, title={\lblVerify}, #1}
\newtcolorbox{exercisebox}[1][]{mfa outlined, colframe=mteal, title={\lblExercise}, #1}
% The guiding project. Points at the stage of trace-assert this section
% builds, under code/src/trace_assert/.
\newtcolorbox{projectbox}[1][]{mfa outlined, colframe=mamber, title={\lblProject}, #1}

% ============================================================
%  EXERCISES
%  ---------------------------------------------------------
%  \begin{exercise}{key}{title} ... \end{exercise}
%
%  This is the mechanism the whole book is read through: the PDF on one half
%  of the screen, an editor on the other. An exercise is therefore not a
%  paragraph but a FILE the reader opens -- a starter under
%  code/exercises/chNN/<key>.py -- and a TEST that decides whether it is done,
%  test_<key>.py beside it. The box prints both paths and the one command
%  that runs the test, and the manifest at the back lists every exercise.
%
%  An ENVIRONMENT, not a macro with a body argument: a macro argument cannot
%  contain verbatim material, so an exercise written as \exercise{key}{title}
%  {body} could never show the reader a code fragment, and the first chapter
%  that tried would have died on `\begin{python}` inside the braces.
%
%  The key is the file stem, with underscores, and it is written with
%  \detokenize so the underscore reaches the page as a character rather than
%  as a subscript. tools/check_structure.py --exercises fails when the starter
%  or the test is missing, or when the key's chapter prefix disagrees with the
%  chapter it sits in.
% ============================================================
\newcounter{exercise}[chapter]
\renewcommand{\theexercise}{\thechapter.\arabic{exercise}}
\makeatletter
\newcommand{\exercisedir}{ch\two@digits{\value{chapter}}}
\newcommand{\listofexercises}{%
  \section*{\lblExerciseManifest}%
  \phantomsection\addcontentsline{toc}{section}{\lblExerciseManifest}%
  % The entries are subsection-level and tocdepth is 1 for the contents,
  % so the depth is raised for this list alone or it prints empty --
  % which it did, on the scaffold's first clean build.
  \begingroup\c@tocdepth=2\relax\@starttoc{exr}\endgroup%
}
\makeatother
\NewDocumentEnvironment{exercise}{mm}{%
  \refstepcounter{exercise}%
  \addcontentsline{exr}{subsection}{\protect\numberline{\theexercise}\texttt{\protect\detokenize{#1}} --- #2}%
  \begin{exercisebox}[title={\lblExercise~\theexercise\ \dash{} #2}]%
}{%
  \par\medskip
  {\small\setlength{\parindent}{0pt}%
  \lblExerciseStarter: \texttt{code/exercises/\exercisedir/\detokenize{#1}.py}\\
  \lblExerciseTest: \texttt{code/exercises/\exercisedir/test\_\detokenize{#1}.py}\\
  \lblExerciseRun: \texttt{cd code \&\& uv run pytest -k \detokenize{#1}}\par}%
  \end{exercisebox}%
}

% ============================================================
%  DIAGRAMS --- Mermaid pipeline
%  ---------------------------------------------------------
%  \mermaidfig{key}{caption}{one-line manifest description}
%
%  Source of truth is figures/mermaid/<lang>/<key>.mmd, committed. `make
%  diagrams` renders each to figures/diagrams/<lang>/<key>.pdf, which is build
%  output and is not committed, so a diagram change reviews as a text diff.
%  The source is per-language because a diagram with English labels in the
%  Polish edition is exactly the half-translation that makes a bilingual book
%  feel machine-made; CI checks that both languages carry the same keys.
%
%  If the rendered PDF is absent the macro draws a placeholder AND typesets
%  the Mermaid source, in the flow rather than as a float, because a source
%  typeset verbatim is often taller than a page and a float cannot break.
% ============================================================
\makeatletter
\newcommand{\listofdiagrams}{%
  \section*{\lblDiagramManifest}%
  \phantomsection\addcontentsline{toc}{section}{\lblDiagramManifest}%
  % The entries are subsection-level and tocdepth is 1 for the contents,
  % so the depth is raised for this list alone or it prints empty --
  % which it did, on the scaffold's first clean build.
  \begingroup\c@tocdepth=2\relax\@starttoc{dgm}\endgroup%
}
\makeatother

\newcommand{\mermaidfig}[3]{%
  % \phantomsection is load-bearing and not decoration. \addcontentsline
  % records hyperref's CURRENT anchor, and here that is whatever preceded
  % the figure -- for a figure placed after a listing it is the listing's
  % line-number anchor, and listings has already advanced the counter past
  % the last line it printed, so the manifest entry linked to a
  % destination that is never created: "name{lstnumber.10.4.46} has been
  % referenced but does not exist, replaced by a fixed one". The two
  % front-matter figures had the same wrong target and were silent about
  % it, because the lines they named happen to be printed. The exercise
  % manifest never had the bug, because \begin{exercise} steps a counter
  % and that is what sets the anchor.
  \phantomsection%
  \addcontentsline{dgm}{subsection}{\protect\numberline{}\texttt{#1.mmd} --- #3}%
  \IfFileExists{figures/diagrams/\booklang/#1.pdf}{%
    \begin{figure}[htbp]
    \centering
    \includegraphics[width=0.95\linewidth,height=0.42\textheight,keepaspectratio]{figures/diagrams/\booklang/#1.pdf}%
    \caption{#2}\label{fig:#1}
    \end{figure}%
  }{%
    \par\medskip
    \begin{tcolorbox}[
      enhanced, breakable, width=\linewidth, sharp corners,
      colback=white, colframe=mgrey, boxrule=0.8pt,
      left=8pt, right=8pt, top=8pt, bottom=8pt]
      {\bfseries\color{mgrey}\small \lblNotRendered}\\[4pt]
      {\ttfamily\scriptsize\color{mgrey} figures/mermaid/\booklang/#1.mmd}\\[6pt]
      \IfFileExists{figures/mermaid/\booklang/#1.mmd}{%
        \lstinputlisting[style=bookcode,numbers=none,basicstyle=\ttfamily\scriptsize,
                         backgroundcolor=\color{mlight},frame=none,
                         keywordstyle=\color{black},language={}]{figures/mermaid/\booklang/#1.mmd}%
      }{%
        {\small\color{mred} \lblSourceMissing: \texttt{figures/mermaid/\booklang/#1.mmd}}%
      }%
    \end{tcolorbox}%
    \captionof{figure}{#2}\label{fig:#1}
    \par\medskip
  }%
}

% ============================================================
%  CHAPTER STUBS
%  ---------------------------------------------------------
%  Prints a visible marker so an unwritten chapter cannot be mistaken for a
%  written one and the page count never flatters the draft. `make debt`
%  counts them, and CI publishes the count on every build.
% ============================================================
\newcommand{\chapterstub}[1]{%
  \begin{tcolorbox}[
    enhanced, breakable, width=\linewidth, sharp corners,
    colback=mlight, colframe=mred, boxrule=1pt,
    borderline={0.6pt}{3pt}{mred, dashed},
    left=12pt, right=12pt, top=12pt, bottom=12pt]
    {\bfseries\color{mred}\large \lblNotWritten}\\[8pt]
    \small\raggedright #1
  \end{tcolorbox}%
}

% ============================================================
%  COMPUTED VALUES
%  ---------------------------------------------------------
%  A number that came out of a measurement is written by a script under
%  code/measure/ into figures/values/<name>.tex as
%
%      \pyval{e01.threads.speedup}{3.71}
%
%  and pulled into the text with \val{e01.threads.speedup}. The script is the
%  source of truth and the book cannot disagree with it, because the book does
%  not contain the digits. `make verify` fails when a committed value no
%  longer matches its script. A value the scripts do not produce prints a
%  visible marker.
%
%  \val also applies the edition's decimal separator, which is how the Polish
%  edition gets its decimal comma without a translator having to remember.
% ============================================================
\newcommand{\bookdecimalmarker}{\ifpl{,}{.}}
\IfFileExists{siunitx.sty}{%
  \usepackage{siunitx}%
  \sisetup{
    group-digits = integer,
    group-minimum-digits = 5,
    group-separator = {\,},
    output-decimal-marker = {\bookdecimalmarker}
  }%
  \newcommand{\booknum}[1]{\num{##1}}%
}{%
  \newcommand{\booknum}[1]{##1}%
}

\newcommand{\pyvalues}{}
\makeatletter
\newcommand{\pyval}[2]{\csgdef{pyval@#1}{#2}\listxadd{\pyvalues}{#1}}
% \num on a non-numeric body is a FATAL siunitx error. The generator knows
% whether it emitted a number or a word, so it writes \pyval for a number and
% \pyvaltext for a word, and this end only enforces which macro reads which.
\newcommand{\val}[1]{%
  \ifcsdef{pyval@#1}{%
    \ifcsdef{pyvaltext@#1}{%
      \PackageError{pybook}{Value `#1' is not a number}%
        {\string\val\space passes its body to siunitx, which refuses anything
         that is not a number. Use \string\valtext{#1} instead.}%
    }{}%
    \booknum{\csuse{pyval@#1}}%
  }{\textcolor{mred}{\textbf{[\lblValueMissing: #1]}}}%
}
\newcommand{\valtext}[1]{%
  \ifcsdef{pyval@#1}{\csuse{pyval@#1}}%
    {\textcolor{mred}{\textbf{[\lblValueMissing: #1]}}}%
}
\newcommand{\pyvaltext}[2]{%
  \csgdef{pyval@#1}{#2}\csgdef{pyvaltext@#1}{}\listxadd{\pyvalues}{#1}%
}
\makeatother

% figures/values/all.tex is generated by `make numbers` and does nothing but
% \input each value file. A build with no values yet still compiles; every
% \val{} in it prints a visible marker instead of a number.
\IfFileExists{figures/values/all.tex}{\input{figures/values/all}}{%
  \AtBeginDocument{\PackageWarning{pybook}{No computed values found. Run `make numbers`.}}}

% ============================================================
%  PINNED VERSIONS --- single source of truth
%  Change these here; they propagate through both editions and Appendix C.
%  Verified against pypi.org on \pinnedon (lang/<lang>.tex).
%  tools/check_versions.py compares every one of them with PyPI on each build.
%  Python itself is not on PyPI and is checked by hand against python.org.
% ============================================================
\newcommand{\bookedition}{0.1-dev}
\newcommand{\bookyear}{2026}
\newcommand{\pyver}{3.14}              % the minor the book targets
\newcommand{\pypatch}{3.14.7}          % the exact interpreter the code ran on
\newcommand{\uvver}{0.12.13}           % uv
\newcommand{\dotnetver}{10.0.100}      % dotnet
\newcommand{\ruffver}{0.16.7}          % ruff
\newcommand{\pyrightver}{1.1.414}      % pyright
\newcommand{\mypyver}{2.3.1}           % mypy
\newcommand{\pydanticver}{2.13.5}      % pydantic
\newcommand{\pysettingsver}{2.15.0}    % pydantic-settings
\newcommand{\fastapiver}{0.141.1}      % fastapi
\newcommand{\uvicornver}{0.53.0}       % uvicorn
\newcommand{\sqlalchemyver}{2.0.52}    % sqlalchemy
\newcommand{\alembicver}{1.20.0}       % alembic
\newcommand{\aiosqlitever}{0.22.1}     % aiosqlite
\newcommand{\pytestver}{9.1.1}         % pytest
\newcommand{\pytestasyncver}{1.4.0}    % pytest-asyncio
\newcommand{\hypothesisver}{6.168.0}   % hypothesis
\newcommand{\coveragever}{7.16.1}      % coverage
\newcommand{\httpxver}{0.28.1}         % httpx
\newcommand{\structlogver}{26.1.0}     % structlog
\newcommand{\otelver}{1.44.0}          % opentelemetry-sdk
\newcommand{\polarsver}{1.44.2}        % polars
\newcommand{\anthropicver}{1.5.0}      % anthropic
\newcommand{\openaiver}{3.13.0}        % openai
\newcommand{\pipauditver}{2.10.1}      % pip-audit

% ============================================================
%  INDEX
%  Two-column, and its entries include \texttt{} names that do not hyphenate.
%  Ragged right keeps multi-word entries off the margin; it cannot help a
%  single token wider than the column, so keep verbatim index entries under
%  about 29 characters and index longer names by concept instead.
%  \raggedcolumns is the switch imakeidx's multicols actually reads.
% ============================================================
\apptocmd{\theindex}{\raggedcolumns\raggedright}{}{%
  \PackageWarning{pybook}{Could not patch theindex for ragged setting}}
\providecommand*\see[2]{}
\renewcommand*\see[2]{\emph{\lblIndexSee} #1}
\providecommand*\seealso[2]{}
\renewcommand*\seealso[2]{\emph{\lblIndexSeeAlso} #1}

\makeindex
 still compiles and says what it did.
\providecommand{\booklang}{en}

\usepackage{etoolbox}   % loaded first: the language switch below needs it

% ---------- Page geometry ----------
% ONE format: A4, 12pt, ONE-SIDED. This book is read on a screen with an
% editor open beside it, and never printed as a bound book, so there is no
% recto, no verso, no gutter, no mirrored margin and no blank leaf before a
% chapter -- every one of those is a property of paper that a PDF viewer at
% half a screen's width turns into wasted pixels.
%
% The margins are symmetric and the measure is about seventy-five characters
% of prose, which is where the eye keeps the line return. The number that was
% actually chosen is the MONOSPACE measure: at \footnotesize on this block a
% listing line of 79 characters fits without wrapping, and 79 is the width
% every listing under code/ is held to by tools/check_structure.py --lines.
% A wrapped listing line is not an error TeX reports -- breaklines=true wraps
% it silently and prints an arrow into the middle of what the reader is meant
% to paste into an editor -- so the width is enforced in the source instead.
\usepackage[
  a4paper,
  left=3.1cm, right=3.1cm, top=2.6cm, bottom=3.0cm,
  head=26pt, headsep=0.8cm, footskip=1.4cm
]{geometry}

% ---------- Encoding & fonts ----------
\usepackage[T1]{fontenc}
\usepackage[utf8]{inputenc}
\IfFileExists{lmodern.sty}{\usepackage{lmodern}}{}
% newtxtext takes its TS1 symbols (\textcopyright, \textdegree) from TeX
% Gyre Termes, which Debian packages SEPARATELY (tex-gyre): a machine with
% newtx and without it dies on the copyright page with
% `Font TS1/ntxtlf/m/n/10.95=ts1-qtmr ... not loadable`. There is no
% inert-safe probe for a font metric in pdfTeX -- \IfFileExists searches
% TEXINPUTS and not TFMFONTS, so it reports the metric missing on a machine
% that has it, and \suppressfontnotfounderror is LuaTeX-only -- so nothing
% here guesses. CI's full TeX Live is the reference; locally, install
% tex-gyre. Both facts are recorded in CLAUDE.md under Build traps.
\IfFileExists{newtxtext.sty}{\usepackage{newtxtext}}{}
% amsmath before newtxmath, because newtxmath patches amsmath's internals.
% amssymb only when newtxmath is absent: newtxmath supplies the AMS symbols
% itself and loading both is a fatal clash (\Bbbk already defined) that is
% invisible on a machine without newtx and fatal on a full TeX Live.
\usepackage{amsmath}
\IfFileExists{newtxmath.sty}{\usepackage{newtxmath}}{\usepackage{amssymb}}
\IfFileExists{inconsolata.sty}{\usepackage[varqu,scaled=0.95]{inconsolata}}{}
\IfFileExists{lmodern.sty}{\usepackage{microtype}}{\usepackage[expansion=false]{microtype}}

% upquote is not a refinement. In T1 the input character ' is quoteright, so
% a listings body sets f('x') with U+2019 at both ends of the literal, and
% typed back in that is a SyntaxError. inconsolata's varqu option hides the
% defect on a machine that has inconsolata; upquote makes ' the ASCII
% apostrophe whatever the typewriter font is. In a book whose every listing is
% meant to be pasted into an editor, this line is load-bearing.
\IfFileExists{upquote.sty}{\usepackage{upquote}}{}

% And upquote's twin, for the same reason one character over. In T1 the pair
% -- is a ligature for an en dash, in EVERY family including the typewriter
% one, so \code{uv sync --locked} sets as "uv sync <en dash>locked" and a
% reader who copies it off the page gets an unknown-argument error from a
% command the chapter told them to run. It is invisible in review because it
% looks like a dash, and it bites exactly the flags a packaging chapter is
% made of. A `listings` body is already safe -- listings sets characters
% singly and no ligature forms -- which is what makes the defect so easy to
% miss: the same text is right inside a listing and wrong in the sentence
% above it.
%
% Scoped to tt* on purpose, and measured rather than assumed: prose keeps its
% own dashes, so \dash{} still sets an em dash and a range like 1--2 still
% sets an en dash. Probed against this preamble in all three positions before
% it was believed.
\DisableLigatures[-]{encoding = T1, family = tt*}

% ---------- Language ----------
% Loading babel with a language whose .ldf is absent is a *fatal* error, not
% a warning, so both the package and each language file are probed before
% either is requested. Both editions load the other language too, because the
% glossary rows and the cheat sheet carry words from both.
\IfFileExists{babel.sty}{%
  \IfFileExists{polish.ldf}{%
    \IfFileExists{british.ldf}{%
      \ifdefstring{\booklang}{pl}%
        {\usepackage[british,main=polish]{babel}}%
        {\usepackage[polish,main=british]{babel}}%
    }{\usepackage[polish]{babel}}%
  }{%
    \IfFileExists{british.ldf}{\usepackage[british]{babel}}{}%
  }%
}{}

% babel-polish makes " an active shorthand. listings reads its body verbatim
% and an active character in a verbatim context is a category-code accident
% waiting to happen, so the shorthand is turned off. Polish quotation marks
% come from \enquote, which csquotes sets per language.
\ifdefstring{\booklang}{pl}{\AtBeginDocument{\shorthandoff{"}}}{}

% ---------- Core packages ----------
\usepackage{graphicx}
\usepackage{xcolor}
\usepackage{booktabs}
\usepackage{tabularx}
% Appendix A's cheat sheet runs past a page in every part, and tabularx
% cannot break. longtable can, and works with booktabs.
\usepackage{longtable}
\usepackage{array}
\usepackage{enumitem}
\usepackage{listings}
\usepackage[most]{tcolorbox}
\usepackage{fancyhdr}
\usepackage{titlesec}
\usepackage{caption}
\usepackage{imakeidx}
% csquotes with autostyle follows babel's active language, so the identical
% source \enquote{x} sets 'x' in the English edition and the Polish quotation
% marks in the Polish one.
\IfFileExists{csquotes.sty}{\usepackage[autostyle=true]{csquotes}}{%
  \providecommand{\enquote}[1]{``##1''}}
\usepackage[hidelinks]{hyperref}

% ---------- Palette ----------
% The same family as the two companion volumes, so the three read as a set.
\definecolor{mblue}{HTML}{1F4E79}
\definecolor{mteal}{HTML}{0E7C7B}
\definecolor{mamber}{HTML}{B26A00}
\definecolor{mred}{HTML}{9B2C2C}
\definecolor{mviolet}{HTML}{7B4397}
\definecolor{mgrey}{HTML}{5A5A5A}
\definecolor{mlight}{HTML}{F4F6F8}
\definecolor{mfaint}{HTML}{FBFCFD}
\definecolor{codebg}{HTML}{FAFAFA}
\definecolor{codeframe}{HTML}{D8DEE4}
\definecolor{kw}{HTML}{0B5394}
\definecolor{str}{HTML}{9B2C2C}
\definecolor{cmt}{HTML}{4F7A4F}
\definecolor{typ}{HTML}{0E7C7B}
\definecolor{dec}{HTML}{7B4397}

\hypersetup{
  colorlinks=true,
  linkcolor=mblue,
  urlcolor=mteal,
  citecolor=mblue,
  pdfauthor={Konrad Cinkusz}
}

% \ifpl{Polish text}{English text} -- for the handful of places where a
% string is structural rather than content. Anything longer than a few words
% belongs in lang/ or in the edition's own source, not here.
% \DeclareRobustCommand, not \newcommand: a fragile \ifpl inside a moving
% argument is written to the .toc one level expanded and fails on the next run
% in exactly one of the two editions.
\DeclareRobustCommand{\ifpl}[2]{\ifdefstring{\booklang}{pl}{#1}{#2}}
\makeatletter
\pdfstringdefDisableCommands{%
  \ifdefstring{\booklang}{pl}{\let\ifpl\@firstoftwo}{\let\ifpl\@secondoftwo}%
  % \api{}, \pkg{} and \code{} are allowed in a section title, and a title
  % becomes a PDF bookmark. hyperref drops the \color command from a bookmark
  % string but keeps its ARGUMENT, so a title carrying \api{model\_validate}
  % made a bookmark reading "mtealmodel_validate" and warned "Token not
  % allowed" twice. The colour name is gobbled here instead. Measured on a
  % probe against this preamble before the line was added.
  \def\color#1{}%
}
\makeatother

% ---------- Language strings ----------
\input{lang/\booklang}

% The PDF's subject, keyword list and document language read \lblPdf... from
% the language file input on the line above, so they cannot sit in the colour
% \hypersetup block, which runs before it. Two blocks is the ordering, not an
% oversight; merging them raises Undefined control sequence and still writes
% a PDF, with the three fields empty.
\hypersetup{
  pdfsubject={\lblPdfSubject},
  pdfkeywords={\lblPdfKeywords},
  pdflang={\lblPdfLang}
}

% ---------- Headings ----------
% \raggedright on the title body is load-bearing: at \Huge on this block a
% title that cannot break overflows the margin by 30 pt or more, and several
% chapter titles here run past fifty characters.
\titleformat{\chapter}[display]
  {\normalfont\huge\bfseries\color{mblue}\raggedright}
  {\normalfont\normalsize\color{mgrey}\MakeUppercase{\chaptertitlename}\ \thechapter}
  {8pt}{\Huge\raggedright}
\titlespacing*{\chapter}{0pt}{10pt}{28pt}
\titleformat{\section}{\normalfont\Large\bfseries\color{mblue}\raggedright}{\thesection}{0.8em}{}
\titleformat{\subsection}{\normalfont\large\bfseries\color{mgrey}\raggedright}{\thesubsection}{0.7em}{}

% Sections in the contents, subsections not: fourteen chapters of subsections
% is a contents nobody reads, and the reader navigates by chapter and section.
\setcounter{tocdepth}{1}
\setcounter{secnumdepth}{2}

% One-sided running head: the chapter on the left, the folio on the right.
\pagestyle{fancy}
\fancyhf{}
\fancyhead[L]{\small\nouppercase{\leftmark}}
\fancyhead[R]{\small\thepage}
\renewcommand{\headrulewidth}{0.4pt}
\fancypagestyle{plain}{\fancyhf{}\fancyfoot[R]{\small\thepage}%
  \renewcommand{\headrulewidth}{0pt}}

% Drops the word "Chapter" from the head: the number is still there and the
% longer titles no longer overflow. MUST come after \pagestyle{fancy}, which
% installs its own \chaptermark and silently overwrites anything earlier.
\renewcommand{\chaptermark}[1]{\markboth{\thechapter.\ #1}{}}

% ============================================================
%  CODE LISTINGS
%  ---------------------------------------------------------
%  Every listing in a chapter is a FILE under code/, pulled in with \pyfile or
%  \pyregion, and CI runs it against the pinned versions. An in-source
%  \begin{python} block is allowed for a fragment of three or four lines that
%  exists to show a shape rather than to run; anything longer belongs in a
%  file the reader can open in the editor beside the PDF.
% ============================================================
\lstdefinelanguage{pythonx}{
  language=Python,
  morekeywords={
    async,await,match,case,nonlocal,yield,type,
    TypedDict,Annotated,Protocol,Literal,Generic,Final,ClassVar,Self,
    TypeVar,TypeGuard,TypeIs,Callable,Iterator,Iterable,Awaitable,
    dataclass,field,BaseModel,Field,Enum,StrEnum,NamedTuple,
    override,overload,property,staticmethod,classmethod
  },
  morekeywords=[2]{
    gather,create_task,TaskGroup,to_thread,timeout,wait_for,Semaphore,Queue,
    fixture,parametrize,monkeypatch,raises,mark,
    Depends,FastAPI,APIRouter,HTTPException,
    select,Session,Mapped,mapped_column,relationship,
    model_validate,model_validate_json,model_dump,model_json_schema
  },
  sensitive=true
}

% listings' C# definition predates async, records and pattern matching.
\lstdefinelanguage{csharpx}{
  language=[Sharp]C,
  morekeywords={
    async,await,var,dynamic,record,init,required,nameof,yield,partial,
    global,scoped,when,with,and,or,not,notnull,unmanaged,file,get,set,value,
    where,select,from,orderby,let,join,into,ascending,descending,on,equals,
    Task,ValueTask,CancellationToken,IEnumerable,IAsyncEnumerable,Func,Action
  },
  sensitive=true
}

\lstdefinestyle{bookcode}{
  basicstyle=\ttfamily\footnotesize,
  keywordstyle=\color{kw}\bfseries,
  keywordstyle=[2]\color{typ},
  stringstyle=\color{str},
  % Colour only, no \itshape: inconsolata (zi4) has no italic shape, so an
  % italic comment style is substituted with upright on every machine that
  % has the font, and the log carries "Font shape `T1/zi4/m/it' undefined"
  % on every build. The colour already marks a comment.
  commentstyle=\color{cmt},
  emphstyle=\color{dec}\bfseries,
  numbers=left,
  numberstyle=\ttfamily\tiny\color{mgrey},
  numbersep=8pt,
  backgroundcolor=\color{codebg},
  frame=single,
  rulecolor=\color{codeframe},
  framesep=6pt,
  breaklines=true,
  breakatwhitespace=false,
  postbreak=\mbox{\textcolor{mgrey}{$\hookrightarrow$}\space},
  showstringspaces=false,
  tabsize=4,
  columns=fullflexible,
  keepspaces=true,
  captionpos=b,
  aboveskip=1.2em,
  belowskip=1.2em,
  % Maps the characters most likely to be pasted into a listing by accident.
  %
  % Do NOT add a mapping for U+00A0 (non-breaking space). It looks like the
  % obvious next entry and it makes `listings` abort the run with "Improper
  % alphabetic constant" -- fatal, no PDF, and the message names neither the
  % character nor this line. The ASCII-in-listings convention, checked by
  % parity's C13, is the guard against it instead.
  literate={ł}{{\l}}1 {Ł}{{\L}}1 {ą}{{\k a}}1 {ę}{{\k e}}1
           {ś}{{\'s}}1 {ć}{{\'c}}1 {ż}{{\.z}}1 {ź}{{\'z}}1 {ó}{{\'o}}1 {ń}{{\'n}}1
           {Ą}{{\k A}}1 {Ę}{{\k E}}1 {Ś}{{\'S}}1 {Ć}{{\'C}}1
           {Ż}{{\.Z}}1 {Ź}{{\'Z}}1 {Ó}{{\'O}}1 {Ń}{{\'N}}1
           {—}{{-{}-}}2 {–}{{-}}1 {‘}{{'}}1 {’}{{'}}1
           {“}{{"}}1 {”}{{"}}1 {°}{{\textdegree}}1
}

\lstset{style=bookcode}

\lstnewenvironment{python}[1][]
  {\lstset{language=pythonx,style=bookcode,#1}}
  {}

% C# comparison snippets. The whole book is written for engineers arriving
% from .NET, so listings come in pairs often enough that the C# half has its
% own environment rather than a generic one.
\lstnewenvironment{csharp}[1][]
  {\lstset{language=csharpx,style=bookcode,#1}}
  {}

\lstnewenvironment{shellcmd}[1][]
  {\lstset{language=bash,style=bookcode,numbers=none,keywordstyle=\color{mteal}\bfseries,#1}}
  {}

\lstnewenvironment{tomlcode}[1][]
  {\lstset{style=bookcode,numbers=none,keywordstyle=\color{black},#1}}
  {}

\lstnewenvironment{yamlcode}[1][]
  {\lstset{style=bookcode,numbers=none,keywordstyle=\color{black},#1}}
  {}

\lstnewenvironment{jsoncode}[1][]
  {\lstset{style=bookcode,numbers=none,keywordstyle=\color{black},#1}}
  {}

% Terminal transcripts typed into the source. Allowed for a shell session the
% reader is meant to reproduce; NOT for program output, which is \transcript
% below, because an in-source transcript nobody ran reads exactly like one
% that was, and that is where a remembered number hides.
\lstnewenvironment{console}[1][]
  {\lstset{style=bookcode,numbers=none,basicstyle=\ttfamily\footnotesize,
           keywordstyle=\color{black},backgroundcolor=\color{white},#1}}
  {}

% The path is printed above every file listing, in the repository-relative
% form the reader opens in the editor beside the PDF. That is the whole point
% of a listing being a file, and the front matter promises it. The \nobreak
% keeps the path on the same page as the listing it names.
% ---------------------------------------------------------------------------
%  Appendix B prints the trap catalogue from notes/02-traps.md. One entry is
%  its number -- which a chapter may cite, and which is never reused -- the
%  habit in the reader's own voice, and what Python does instead.
%
%  The body is set RAGGED RIGHT on purpose. An entry is mostly \code{} spans,
%  which are unbreakable runs, and sixty-seven justified paragraphs of them
%  is sixty-seven chances at an overfull hbox in the edition with the longer
%  words. Ragged right removes the class rather than fixing it one box at a
%  time, and a reference list is the one place in this book where it reads
%  better anyway.
%
%  The habit is set with \enquote{} rather than \emph{}, and that is not a
%  style choice. Every habit here carries \code{} spans, and \code{} inside
%  \emph{} asks inconsolata for an italic it does not have: the first build
%  of this appendix raised `Font shape T1/zi4/m/it undefined', which
%  checklog.py is hard on for exactly this reason. It is the trap CLAUDE.md
%  records, met sixty-seven times at once -- and fixed once, because the
%  entries go through a macro. A quotation is also what the habit IS, and
%  \enquote{} is what the Polish edition owes anyway.
% ---------------------------------------------------------------------------
%  Appendix A's cheat sheet: one table per part, C# on the left, Python on
%  the right, and the chapter that teaches the row on the end. The three
%  header strings are arguments so that one environment serves both
%  editions and the structural signature stays identical.
%
%  Columns are ragged right for the reason Appendix B's entries are: a cell
%  is mostly \code{} spans, which are unbreakable runs, and a justified
%  narrow column full of them is an overfull box waiting for the edition
%  with the longer words.
% ---------------------------------------------------------------------------
%  Appendix C's tool matrix: the .NET tool, the Python tool this book uses,
%  the version it was verified against and the chapter that introduces it.
%
%  The version column prints \pydanticver and its siblings and is never
%  typed, which is the whole point of the appendix: it is the one page where
%  a reader sees every pin at once, and a pin that had to be re-typed here
%  would be a second copy of a fact the preamble already owns.
%  tools/check_structure.py --tools fails when a pin macro exists and this
%  table does not print it.
\NewDocumentEnvironment{toolmatrix}{mmmm}{%
  \small
  \begin{longtable}{@{}%
    >{\raggedright\arraybackslash}p{0.27\linewidth}%
    >{\raggedright\arraybackslash}p{0.31\linewidth}%
    >{\raggedright\arraybackslash}p{0.20\linewidth}%
    >{\raggedright\arraybackslash}p{0.08\linewidth}@{}}
  \toprule
  \textbf{#1} & \textbf{#2} & \textbf{#3} & \textbf{#4} \\
  \midrule
  \endfirsthead
  \toprule
  \textbf{#1} & \textbf{#2} & \textbf{#3} & \textbf{#4} \\
  \midrule
  \endhead
  \midrule
  \endfoot
  \bottomrule
  \endlastfoot
}{%
  \end{longtable}%
}

\NewDocumentEnvironment{cheatsheet}{mmm}{%
  \small
  \begin{longtable}{@{}%
    >{\raggedright\arraybackslash}p{0.41\linewidth}%
    >{\raggedright\arraybackslash}p{0.41\linewidth}%
    >{\raggedright\arraybackslash}p{0.08\linewidth}@{}}
  \toprule
  \textbf{#1} & \textbf{#2} & \textbf{#3} \\
  \midrule
  \endfirsthead
  \toprule
  \textbf{#1} & \textbf{#2} & \textbf{#3} \\
  \midrule
  \endhead
  \midrule
  \endfoot
  \bottomrule
  \endlastfoot
}{%
  \end{longtable}%
}

\newcommand{\trapentry}[3]{%
  \par\medskip
  \noindent\textbf{#1.}\enspace\enquote{#2}\par
  \nopagebreak
  {\raggedright\noindent#3\par}%
}

\newcommand{\listingpath}[1]{%
  \par\medskip\noindent
  {\ttfamily\footnotesize\color{mgrey}\detokenize{#1}}%
  \par\nobreak}

% External source file -- the preferred form. The path is relative to the
% repository root, which is also the path the reader opens in the editor.
% Usage: \pyfile{code/ch05/decorators.py}{Caption}{lst:label}
\newcommand{\pyfile}[3]{%
  \listingpath{#1}%
  \lstinputlisting[language=pythonx,style=bookcode,aboveskip=0.3em,
    caption={#2},label={#3}]{#1}%
}

% A named region of an external file, delimited in the source by
% `# --8<-- [start:name]` / `# --8<-- [end:name]`, so the book and the running
% code cannot drift apart. tools/check_structure.py --listings fails when the
% file or the region is missing.
%
% The markers are matched through listings' range prefix and suffix keys,
% with linerange={name-name}. The form this was inherited with -- a linerange
% whose two ends were the whole marker strings -- matched nothing and printed
% the WHOLE FILE, with no warning: measured on a three-region probe file
% before this was rewritten, and the same definition sits unused in the
% sibling book it came from. A region name is a word, never a number, because
% a number is a line number to listings. The line numbers printed are the
% file's own, so a region's listing says where in the file it sits.
% Usage: \pyregion{code/ch05/decorators.py}{timing}{Caption}{lst:label}
\lstset{
  rangebeginprefix=\#\ --8<--\ [start:, rangebeginsuffix=],
  rangeendprefix=\#\ --8<--\ [end:, rangeendsuffix=],
  includerangemarker=false}
\newcommand{\pyregion}[4]{%
  \listingpath{#1}%
  \lstinputlisting[language=pythonx,style=bookcode,aboveskip=0.3em,
    linerange={#2-#2},caption={#3},label={#4}]{#1}%
}

% The C# twin of \pyfile, for the comparison half of a pair.
\newcommand{\csfile}[3]{%
  \listingpath{#1}%
  \lstinputlisting[language=csharpx,style=bookcode,aboveskip=0.3em,
    caption={#2},label={#3}]{#1}%
}

% A program's output, written by a script under code/measure/ and pulled in
% whole. There is no typed form of this and there must not be: \val{} does
% not expand inside a verbatim body, so a transcript typed into a chapter is
% one that cannot quote a computed value, and a transcript nobody ran is
% indistinguishable on the page from one that was.
\newcommand{\transcript}[1]{%
  \par\smallskip
  \IfFileExists{figures/transcripts/#1.txt}{%
    \lstinputlisting[style=bookcode,numbers=none,basicstyle=\ttfamily\footnotesize,
      keywordstyle=\color{black},backgroundcolor=\color{white},
      language={}]{figures/transcripts/#1.txt}%
  }{%
    \begin{tcolorbox}[enhanced, sharp corners, width=\linewidth,
      colback=mlight, colframe=mgrey, boxrule=0.6pt,
      left=8pt, right=8pt, top=6pt, bottom=6pt]
      {\bfseries\color{mgrey}\small \lblTranscriptMissing}\\[2pt]
      {\ttfamily\scriptsize\color{mgrey} figures/transcripts/#1.txt}
    \end{tcolorbox}%
  }%
  \par\smallskip
}

% Inline literals. Underscores inside these must be written \_.
\newcommand{\code}[1]{\texttt{\small #1}}
\newcommand{\pkg}[1]{\texttt{\small\color{mblue}#1}}
\newcommand{\api}[1]{\texttt{\small\color{mteal}#1}}

% ============================================================
%  ADMONITIONS
%  ---------------------------------------------------------
%  A deliberately small, fixed vocabulary. Two visual treatments carry the
%  whole grammar: `tinted` is a block the reader may read in passing -- a
%  tint and a thick left bar, no outline; `outlined` is a block the reader
%  must stop at, with a full rule on white. A third treatment would mean the
%  first two are not carrying their meaning.
% ============================================================
\tcbset{
  mfa tinted/.style={
    enhanced, breakable, sharp corners,
    colback=mlight, boxrule=0pt, leftrule=3pt,
    left=8pt, right=8pt, top=6pt, bottom=6pt,
    before skip=13pt, after skip=13pt,
    fonttitle=\bfseries\small,
    coltitle=black, colbacktitle=mlight,
    attach title to upper={\par\smallskip},
  },
  mfa outlined/.style={
    enhanced, breakable, sharp corners,
    colback=white, boxrule=0.6pt,
    left=8pt, right=8pt, top=6pt, bottom=6pt,
    before skip=13pt, after skip=13pt,
    fonttitle=\bfseries\small,
    coltitle=black, colbacktitle=white,
    attach title to upper={\par\smallskip},
  },
}

% Read in passing.
\newtcolorbox{note}[1][]{mfa tinted, colframe=mblue, title={\lblNote}, #1}
\newtcolorbox{warning}[1][]{mfa tinted, colframe=mred, title={\lblWarning}, #1}
% The translation. The .NET reader is the book's only assumed reader, so this
% box carries the C# side of a comparison and nothing else. Its use is
% counted: a chapter with none is a chapter that has stopped translating.
\newtcolorbox{csbox}[1][]{mfa tinted, colframe=mteal, title={\lblCsharp}, #1}
% Anything that is preview, experimental or likely to move. Python moves more
% slowly than the agent frameworks do, and this box should be rare here.
\newtcolorbox{versionbox}[1][]{mfa tinted, colframe=mamber, title={\lblVersion}, #1}

% Stop and read.
% The misconception, named AFTER the reader has walked into it. Appendix B is
% the catalogue; a trap box in a chapter is one entry of it, in place.
\newtcolorbox{trapbox}[1][]{mfa outlined, colframe=mred, title={\lblTrap}, #1}
% Marks a listing that has not been executed against the pinned versions.
% `make debt` counts these, CI prints the count, and nothing ships to a reader
% with one attached. Removing it means you ran the code.
\newtcolorbox{verifybox}[1][]{mfa outlined, colframe=mgrey, title={\lblVerify}, #1}
\newtcolorbox{exercisebox}[1][]{mfa outlined, colframe=mteal, title={\lblExercise}, #1}
% The guiding project. Points at the stage of trace-assert this section
% builds, under code/src/trace_assert/.
\newtcolorbox{projectbox}[1][]{mfa outlined, colframe=mamber, title={\lblProject}, #1}

% ============================================================
%  EXERCISES
%  ---------------------------------------------------------
%  \begin{exercise}{key}{title} ... \end{exercise}
%
%  This is the mechanism the whole book is read through: the PDF on one half
%  of the screen, an editor on the other. An exercise is therefore not a
%  paragraph but a FILE the reader opens -- a starter under
%  code/exercises/chNN/<key>.py -- and a TEST that decides whether it is done,
%  test_<key>.py beside it. The box prints both paths and the one command
%  that runs the test, and the manifest at the back lists every exercise.
%
%  An ENVIRONMENT, not a macro with a body argument: a macro argument cannot
%  contain verbatim material, so an exercise written as \exercise{key}{title}
%  {body} could never show the reader a code fragment, and the first chapter
%  that tried would have died on `\begin{python}` inside the braces.
%
%  The key is the file stem, with underscores, and it is written with
%  \detokenize so the underscore reaches the page as a character rather than
%  as a subscript. tools/check_structure.py --exercises fails when the starter
%  or the test is missing, or when the key's chapter prefix disagrees with the
%  chapter it sits in.
% ============================================================
\newcounter{exercise}[chapter]
\renewcommand{\theexercise}{\thechapter.\arabic{exercise}}
\makeatletter
\newcommand{\exercisedir}{ch\two@digits{\value{chapter}}}
\newcommand{\listofexercises}{%
  \section*{\lblExerciseManifest}%
  \phantomsection\addcontentsline{toc}{section}{\lblExerciseManifest}%
  % The entries are subsection-level and tocdepth is 1 for the contents,
  % so the depth is raised for this list alone or it prints empty --
  % which it did, on the scaffold's first clean build.
  \begingroup\c@tocdepth=2\relax\@starttoc{exr}\endgroup%
}
\makeatother
\NewDocumentEnvironment{exercise}{mm}{%
  \refstepcounter{exercise}%
  \addcontentsline{exr}{subsection}{\protect\numberline{\theexercise}\texttt{\protect\detokenize{#1}} --- #2}%
  \begin{exercisebox}[title={\lblExercise~\theexercise\ \dash{} #2}]%
}{%
  \par\medskip
  {\small\setlength{\parindent}{0pt}%
  \lblExerciseStarter: \texttt{code/exercises/\exercisedir/\detokenize{#1}.py}\\
  \lblExerciseTest: \texttt{code/exercises/\exercisedir/test\_\detokenize{#1}.py}\\
  \lblExerciseRun: \texttt{cd code \&\& uv run pytest -k \detokenize{#1}}\par}%
  \end{exercisebox}%
}

% ============================================================
%  DIAGRAMS --- Mermaid pipeline
%  ---------------------------------------------------------
%  \mermaidfig{key}{caption}{one-line manifest description}
%
%  Source of truth is figures/mermaid/<lang>/<key>.mmd, committed. `make
%  diagrams` renders each to figures/diagrams/<lang>/<key>.pdf, which is build
%  output and is not committed, so a diagram change reviews as a text diff.
%  The source is per-language because a diagram with English labels in the
%  Polish edition is exactly the half-translation that makes a bilingual book
%  feel machine-made; CI checks that both languages carry the same keys.
%
%  If the rendered PDF is absent the macro draws a placeholder AND typesets
%  the Mermaid source, in the flow rather than as a float, because a source
%  typeset verbatim is often taller than a page and a float cannot break.
% ============================================================
\makeatletter
\newcommand{\listofdiagrams}{%
  \section*{\lblDiagramManifest}%
  \phantomsection\addcontentsline{toc}{section}{\lblDiagramManifest}%
  % The entries are subsection-level and tocdepth is 1 for the contents,
  % so the depth is raised for this list alone or it prints empty --
  % which it did, on the scaffold's first clean build.
  \begingroup\c@tocdepth=2\relax\@starttoc{dgm}\endgroup%
}
\makeatother

\newcommand{\mermaidfig}[3]{%
  % \phantomsection is load-bearing and not decoration. \addcontentsline
  % records hyperref's CURRENT anchor, and here that is whatever preceded
  % the figure -- for a figure placed after a listing it is the listing's
  % line-number anchor, and listings has already advanced the counter past
  % the last line it printed, so the manifest entry linked to a
  % destination that is never created: "name{lstnumber.10.4.46} has been
  % referenced but does not exist, replaced by a fixed one". The two
  % front-matter figures had the same wrong target and were silent about
  % it, because the lines they named happen to be printed. The exercise
  % manifest never had the bug, because \begin{exercise} steps a counter
  % and that is what sets the anchor.
  \phantomsection%
  \addcontentsline{dgm}{subsection}{\protect\numberline{}\texttt{#1.mmd} --- #3}%
  \IfFileExists{figures/diagrams/\booklang/#1.pdf}{%
    \begin{figure}[htbp]
    \centering
    \includegraphics[width=0.95\linewidth,height=0.42\textheight,keepaspectratio]{figures/diagrams/\booklang/#1.pdf}%
    \caption{#2}\label{fig:#1}
    \end{figure}%
  }{%
    \par\medskip
    \begin{tcolorbox}[
      enhanced, breakable, width=\linewidth, sharp corners,
      colback=white, colframe=mgrey, boxrule=0.8pt,
      left=8pt, right=8pt, top=8pt, bottom=8pt]
      {\bfseries\color{mgrey}\small \lblNotRendered}\\[4pt]
      {\ttfamily\scriptsize\color{mgrey} figures/mermaid/\booklang/#1.mmd}\\[6pt]
      \IfFileExists{figures/mermaid/\booklang/#1.mmd}{%
        \lstinputlisting[style=bookcode,numbers=none,basicstyle=\ttfamily\scriptsize,
                         backgroundcolor=\color{mlight},frame=none,
                         keywordstyle=\color{black},language={}]{figures/mermaid/\booklang/#1.mmd}%
      }{%
        {\small\color{mred} \lblSourceMissing: \texttt{figures/mermaid/\booklang/#1.mmd}}%
      }%
    \end{tcolorbox}%
    \captionof{figure}{#2}\label{fig:#1}
    \par\medskip
  }%
}

% ============================================================
%  CHAPTER STUBS
%  ---------------------------------------------------------
%  Prints a visible marker so an unwritten chapter cannot be mistaken for a
%  written one and the page count never flatters the draft. `make debt`
%  counts them, and CI publishes the count on every build.
% ============================================================
\newcommand{\chapterstub}[1]{%
  \begin{tcolorbox}[
    enhanced, breakable, width=\linewidth, sharp corners,
    colback=mlight, colframe=mred, boxrule=1pt,
    borderline={0.6pt}{3pt}{mred, dashed},
    left=12pt, right=12pt, top=12pt, bottom=12pt]
    {\bfseries\color{mred}\large \lblNotWritten}\\[8pt]
    \small\raggedright #1
  \end{tcolorbox}%
}

% ============================================================
%  COMPUTED VALUES
%  ---------------------------------------------------------
%  A number that came out of a measurement is written by a script under
%  code/measure/ into figures/values/<name>.tex as
%
%      \pyval{e01.threads.speedup}{3.71}
%
%  and pulled into the text with \val{e01.threads.speedup}. The script is the
%  source of truth and the book cannot disagree with it, because the book does
%  not contain the digits. `make verify` fails when a committed value no
%  longer matches its script. A value the scripts do not produce prints a
%  visible marker.
%
%  \val also applies the edition's decimal separator, which is how the Polish
%  edition gets its decimal comma without a translator having to remember.
% ============================================================
\newcommand{\bookdecimalmarker}{\ifpl{,}{.}}
\IfFileExists{siunitx.sty}{%
  \usepackage{siunitx}%
  \sisetup{
    group-digits = integer,
    group-minimum-digits = 5,
    group-separator = {\,},
    output-decimal-marker = {\bookdecimalmarker}
  }%
  \newcommand{\booknum}[1]{\num{##1}}%
}{%
  \newcommand{\booknum}[1]{##1}%
}

\newcommand{\pyvalues}{}
\makeatletter
\newcommand{\pyval}[2]{\csgdef{pyval@#1}{#2}\listxadd{\pyvalues}{#1}}
% \num on a non-numeric body is a FATAL siunitx error. The generator knows
% whether it emitted a number or a word, so it writes \pyval for a number and
% \pyvaltext for a word, and this end only enforces which macro reads which.
\newcommand{\val}[1]{%
  \ifcsdef{pyval@#1}{%
    \ifcsdef{pyvaltext@#1}{%
      \PackageError{pybook}{Value `#1' is not a number}%
        {\string\val\space passes its body to siunitx, which refuses anything
         that is not a number. Use \string\valtext{#1} instead.}%
    }{}%
    \booknum{\csuse{pyval@#1}}%
  }{\textcolor{mred}{\textbf{[\lblValueMissing: #1]}}}%
}
\newcommand{\valtext}[1]{%
  \ifcsdef{pyval@#1}{\csuse{pyval@#1}}%
    {\textcolor{mred}{\textbf{[\lblValueMissing: #1]}}}%
}
\newcommand{\pyvaltext}[2]{%
  \csgdef{pyval@#1}{#2}\csgdef{pyvaltext@#1}{}\listxadd{\pyvalues}{#1}%
}
\makeatother

% figures/values/all.tex is generated by `make numbers` and does nothing but
% \input each value file. A build with no values yet still compiles; every
% \val{} in it prints a visible marker instead of a number.
\IfFileExists{figures/values/all.tex}{\input{figures/values/all}}{%
  \AtBeginDocument{\PackageWarning{pybook}{No computed values found. Run `make numbers`.}}}

% ============================================================
%  PINNED VERSIONS --- single source of truth
%  Change these here; they propagate through both editions and Appendix C.
%  Verified against pypi.org on \pinnedon (lang/<lang>.tex).
%  tools/check_versions.py compares every one of them with PyPI on each build.
%  Python itself is not on PyPI and is checked by hand against python.org.
% ============================================================
\newcommand{\bookedition}{0.1-dev}
\newcommand{\bookyear}{2026}
\newcommand{\pyver}{3.14}              % the minor the book targets
\newcommand{\pypatch}{3.14.7}          % the exact interpreter the code ran on
\newcommand{\uvver}{0.12.13}           % uv
\newcommand{\dotnetver}{10.0.100}      % dotnet
\newcommand{\ruffver}{0.16.7}          % ruff
\newcommand{\pyrightver}{1.1.414}      % pyright
\newcommand{\mypyver}{2.3.1}           % mypy
\newcommand{\pydanticver}{2.13.5}      % pydantic
\newcommand{\pysettingsver}{2.15.0}    % pydantic-settings
\newcommand{\fastapiver}{0.141.1}      % fastapi
\newcommand{\uvicornver}{0.53.0}       % uvicorn
\newcommand{\sqlalchemyver}{2.0.52}    % sqlalchemy
\newcommand{\alembicver}{1.20.0}       % alembic
\newcommand{\aiosqlitever}{0.22.1}     % aiosqlite
\newcommand{\pytestver}{9.1.1}         % pytest
\newcommand{\pytestasyncver}{1.4.0}    % pytest-asyncio
\newcommand{\hypothesisver}{6.168.0}   % hypothesis
\newcommand{\coveragever}{7.16.1}      % coverage
\newcommand{\httpxver}{0.28.1}         % httpx
\newcommand{\structlogver}{26.1.0}     % structlog
\newcommand{\otelver}{1.44.0}          % opentelemetry-sdk
\newcommand{\polarsver}{1.44.2}        % polars
\newcommand{\anthropicver}{1.5.0}      % anthropic
\newcommand{\openaiver}{3.13.0}        % openai
\newcommand{\pipauditver}{2.10.1}      % pip-audit

% ============================================================
%  INDEX
%  Two-column, and its entries include \texttt{} names that do not hyphenate.
%  Ragged right keeps multi-word entries off the margin; it cannot help a
%  single token wider than the column, so keep verbatim index entries under
%  about 29 characters and index longer names by concept instead.
%  \raggedcolumns is the switch imakeidx's multicols actually reads.
% ============================================================
\apptocmd{\theindex}{\raggedcolumns\raggedright}{}{%
  \PackageWarning{pybook}{Could not patch theindex for ragged setting}}
\providecommand*\see[2]{}
\renewcommand*\see[2]{\emph{\lblIndexSee} #1}
\providecommand*\seealso[2]{}
\renewcommand*\seealso[2]{\emph{\lblIndexSeeAlso} #1}

\makeindex
 still compiles and says what it did.
\providecommand{\booklang}{en}

\usepackage{etoolbox}   % loaded first: the language switch below needs it

% ---------- Page geometry ----------
% ONE format: A4, 12pt, ONE-SIDED. This book is read on a screen with an
% editor open beside it, and never printed as a bound book, so there is no
% recto, no verso, no gutter, no mirrored margin and no blank leaf before a
% chapter -- every one of those is a property of paper that a PDF viewer at
% half a screen's width turns into wasted pixels.
%
% The margins are symmetric and the measure is about seventy-five characters
% of prose, which is where the eye keeps the line return. The number that was
% actually chosen is the MONOSPACE measure: at \footnotesize on this block a
% listing line of 79 characters fits without wrapping, and 79 is the width
% every listing under code/ is held to by tools/check_structure.py --lines.
% A wrapped listing line is not an error TeX reports -- breaklines=true wraps
% it silently and prints an arrow into the middle of what the reader is meant
% to paste into an editor -- so the width is enforced in the source instead.
\usepackage[
  a4paper,
  left=3.1cm, right=3.1cm, top=2.6cm, bottom=3.0cm,
  head=26pt, headsep=0.8cm, footskip=1.4cm
]{geometry}

% ---------- Encoding & fonts ----------
\usepackage[T1]{fontenc}
\usepackage[utf8]{inputenc}
\IfFileExists{lmodern.sty}{\usepackage{lmodern}}{}
% newtxtext takes its TS1 symbols (\textcopyright, \textdegree) from TeX
% Gyre Termes, which Debian packages SEPARATELY (tex-gyre): a machine with
% newtx and without it dies on the copyright page with
% `Font TS1/ntxtlf/m/n/10.95=ts1-qtmr ... not loadable`. There is no
% inert-safe probe for a font metric in pdfTeX -- \IfFileExists searches
% TEXINPUTS and not TFMFONTS, so it reports the metric missing on a machine
% that has it, and \suppressfontnotfounderror is LuaTeX-only -- so nothing
% here guesses. CI's full TeX Live is the reference; locally, install
% tex-gyre. Both facts are recorded in CLAUDE.md under Build traps.
\IfFileExists{newtxtext.sty}{\usepackage{newtxtext}}{}
% amsmath before newtxmath, because newtxmath patches amsmath's internals.
% amssymb only when newtxmath is absent: newtxmath supplies the AMS symbols
% itself and loading both is a fatal clash (\Bbbk already defined) that is
% invisible on a machine without newtx and fatal on a full TeX Live.
\usepackage{amsmath}
\IfFileExists{newtxmath.sty}{\usepackage{newtxmath}}{\usepackage{amssymb}}
\IfFileExists{inconsolata.sty}{\usepackage[varqu,scaled=0.95]{inconsolata}}{}
\IfFileExists{lmodern.sty}{\usepackage{microtype}}{\usepackage[expansion=false]{microtype}}

% upquote is not a refinement. In T1 the input character ' is quoteright, so
% a listings body sets f('x') with U+2019 at both ends of the literal, and
% typed back in that is a SyntaxError. inconsolata's varqu option hides the
% defect on a machine that has inconsolata; upquote makes ' the ASCII
% apostrophe whatever the typewriter font is. In a book whose every listing is
% meant to be pasted into an editor, this line is load-bearing.
\IfFileExists{upquote.sty}{\usepackage{upquote}}{}

% And upquote's twin, for the same reason one character over. In T1 the pair
% -- is a ligature for an en dash, in EVERY family including the typewriter
% one, so \code{uv sync --locked} sets as "uv sync <en dash>locked" and a
% reader who copies it off the page gets an unknown-argument error from a
% command the chapter told them to run. It is invisible in review because it
% looks like a dash, and it bites exactly the flags a packaging chapter is
% made of. A `listings` body is already safe -- listings sets characters
% singly and no ligature forms -- which is what makes the defect so easy to
% miss: the same text is right inside a listing and wrong in the sentence
% above it.
%
% Scoped to tt* on purpose, and measured rather than assumed: prose keeps its
% own dashes, so \dash{} still sets an em dash and a range like 1--2 still
% sets an en dash. Probed against this preamble in all three positions before
% it was believed.
\DisableLigatures[-]{encoding = T1, family = tt*}

% ---------- Language ----------
% Loading babel with a language whose .ldf is absent is a *fatal* error, not
% a warning, so both the package and each language file are probed before
% either is requested. Both editions load the other language too, because the
% glossary rows and the cheat sheet carry words from both.
\IfFileExists{babel.sty}{%
  \IfFileExists{polish.ldf}{%
    \IfFileExists{british.ldf}{%
      \ifdefstring{\booklang}{pl}%
        {\usepackage[british,main=polish]{babel}}%
        {\usepackage[polish,main=british]{babel}}%
    }{\usepackage[polish]{babel}}%
  }{%
    \IfFileExists{british.ldf}{\usepackage[british]{babel}}{}%
  }%
}{}

% babel-polish makes " an active shorthand. listings reads its body verbatim
% and an active character in a verbatim context is a category-code accident
% waiting to happen, so the shorthand is turned off. Polish quotation marks
% come from \enquote, which csquotes sets per language.
\ifdefstring{\booklang}{pl}{\AtBeginDocument{\shorthandoff{"}}}{}

% ---------- Core packages ----------
\usepackage{graphicx}
\usepackage{xcolor}
\usepackage{booktabs}
\usepackage{tabularx}
% Appendix A's cheat sheet runs past a page in every part, and tabularx
% cannot break. longtable can, and works with booktabs.
\usepackage{longtable}
\usepackage{array}
\usepackage{enumitem}
\usepackage{listings}
\usepackage[most]{tcolorbox}
\usepackage{fancyhdr}
\usepackage{titlesec}
\usepackage{caption}
\usepackage{imakeidx}
% csquotes with autostyle follows babel's active language, so the identical
% source \enquote{x} sets 'x' in the English edition and the Polish quotation
% marks in the Polish one.
\IfFileExists{csquotes.sty}{\usepackage[autostyle=true]{csquotes}}{%
  \providecommand{\enquote}[1]{``##1''}}
\usepackage[hidelinks]{hyperref}

% ---------- Palette ----------
% The same family as the two companion volumes, so the three read as a set.
\definecolor{mblue}{HTML}{1F4E79}
\definecolor{mteal}{HTML}{0E7C7B}
\definecolor{mamber}{HTML}{B26A00}
\definecolor{mred}{HTML}{9B2C2C}
\definecolor{mviolet}{HTML}{7B4397}
\definecolor{mgrey}{HTML}{5A5A5A}
\definecolor{mlight}{HTML}{F4F6F8}
\definecolor{mfaint}{HTML}{FBFCFD}
\definecolor{codebg}{HTML}{FAFAFA}
\definecolor{codeframe}{HTML}{D8DEE4}
\definecolor{kw}{HTML}{0B5394}
\definecolor{str}{HTML}{9B2C2C}
\definecolor{cmt}{HTML}{4F7A4F}
\definecolor{typ}{HTML}{0E7C7B}
\definecolor{dec}{HTML}{7B4397}

\hypersetup{
  colorlinks=true,
  linkcolor=mblue,
  urlcolor=mteal,
  citecolor=mblue,
  pdfauthor={Konrad Cinkusz}
}

% \ifpl{Polish text}{English text} -- for the handful of places where a
% string is structural rather than content. Anything longer than a few words
% belongs in lang/ or in the edition's own source, not here.
% \DeclareRobustCommand, not \newcommand: a fragile \ifpl inside a moving
% argument is written to the .toc one level expanded and fails on the next run
% in exactly one of the two editions.
\DeclareRobustCommand{\ifpl}[2]{\ifdefstring{\booklang}{pl}{#1}{#2}}
\makeatletter
\pdfstringdefDisableCommands{%
  \ifdefstring{\booklang}{pl}{\let\ifpl\@firstoftwo}{\let\ifpl\@secondoftwo}%
  % \api{}, \pkg{} and \code{} are allowed in a section title, and a title
  % becomes a PDF bookmark. hyperref drops the \color command from a bookmark
  % string but keeps its ARGUMENT, so a title carrying \api{model\_validate}
  % made a bookmark reading "mtealmodel_validate" and warned "Token not
  % allowed" twice. The colour name is gobbled here instead. Measured on a
  % probe against this preamble before the line was added.
  \def\color#1{}%
}
\makeatother

% ---------- Language strings ----------
\input{lang/\booklang}

% The PDF's subject, keyword list and document language read \lblPdf... from
% the language file input on the line above, so they cannot sit in the colour
% \hypersetup block, which runs before it. Two blocks is the ordering, not an
% oversight; merging them raises Undefined control sequence and still writes
% a PDF, with the three fields empty.
\hypersetup{
  pdfsubject={\lblPdfSubject},
  pdfkeywords={\lblPdfKeywords},
  pdflang={\lblPdfLang}
}

% ---------- Headings ----------
% \raggedright on the title body is load-bearing: at \Huge on this block a
% title that cannot break overflows the margin by 30 pt or more, and several
% chapter titles here run past fifty characters.
\titleformat{\chapter}[display]
  {\normalfont\huge\bfseries\color{mblue}\raggedright}
  {\normalfont\normalsize\color{mgrey}\MakeUppercase{\chaptertitlename}\ \thechapter}
  {8pt}{\Huge\raggedright}
\titlespacing*{\chapter}{0pt}{10pt}{28pt}
\titleformat{\section}{\normalfont\Large\bfseries\color{mblue}\raggedright}{\thesection}{0.8em}{}
\titleformat{\subsection}{\normalfont\large\bfseries\color{mgrey}\raggedright}{\thesubsection}{0.7em}{}

% Sections in the contents, subsections not: fourteen chapters of subsections
% is a contents nobody reads, and the reader navigates by chapter and section.
\setcounter{tocdepth}{1}
\setcounter{secnumdepth}{2}

% One-sided running head: the chapter on the left, the folio on the right.
\pagestyle{fancy}
\fancyhf{}
\fancyhead[L]{\small\nouppercase{\leftmark}}
\fancyhead[R]{\small\thepage}
\renewcommand{\headrulewidth}{0.4pt}
\fancypagestyle{plain}{\fancyhf{}\fancyfoot[R]{\small\thepage}%
  \renewcommand{\headrulewidth}{0pt}}

% Drops the word "Chapter" from the head: the number is still there and the
% longer titles no longer overflow. MUST come after \pagestyle{fancy}, which
% installs its own \chaptermark and silently overwrites anything earlier.
\renewcommand{\chaptermark}[1]{\markboth{\thechapter.\ #1}{}}

% ============================================================
%  CODE LISTINGS
%  ---------------------------------------------------------
%  Every listing in a chapter is a FILE under code/, pulled in with \pyfile or
%  \pyregion, and CI runs it against the pinned versions. An in-source
%  \begin{python} block is allowed for a fragment of three or four lines that
%  exists to show a shape rather than to run; anything longer belongs in a
%  file the reader can open in the editor beside the PDF.
% ============================================================
\lstdefinelanguage{pythonx}{
  language=Python,
  morekeywords={
    async,await,match,case,nonlocal,yield,type,
    TypedDict,Annotated,Protocol,Literal,Generic,Final,ClassVar,Self,
    TypeVar,TypeGuard,TypeIs,Callable,Iterator,Iterable,Awaitable,
    dataclass,field,BaseModel,Field,Enum,StrEnum,NamedTuple,
    override,overload,property,staticmethod,classmethod
  },
  morekeywords=[2]{
    gather,create_task,TaskGroup,to_thread,timeout,wait_for,Semaphore,Queue,
    fixture,parametrize,monkeypatch,raises,mark,
    Depends,FastAPI,APIRouter,HTTPException,
    select,Session,Mapped,mapped_column,relationship,
    model_validate,model_validate_json,model_dump,model_json_schema
  },
  sensitive=true
}

% listings' C# definition predates async, records and pattern matching.
\lstdefinelanguage{csharpx}{
  language=[Sharp]C,
  morekeywords={
    async,await,var,dynamic,record,init,required,nameof,yield,partial,
    global,scoped,when,with,and,or,not,notnull,unmanaged,file,get,set,value,
    where,select,from,orderby,let,join,into,ascending,descending,on,equals,
    Task,ValueTask,CancellationToken,IEnumerable,IAsyncEnumerable,Func,Action
  },
  sensitive=true
}

\lstdefinestyle{bookcode}{
  basicstyle=\ttfamily\footnotesize,
  keywordstyle=\color{kw}\bfseries,
  keywordstyle=[2]\color{typ},
  stringstyle=\color{str},
  % Colour only, no \itshape: inconsolata (zi4) has no italic shape, so an
  % italic comment style is substituted with upright on every machine that
  % has the font, and the log carries "Font shape `T1/zi4/m/it' undefined"
  % on every build. The colour already marks a comment.
  commentstyle=\color{cmt},
  emphstyle=\color{dec}\bfseries,
  numbers=left,
  numberstyle=\ttfamily\tiny\color{mgrey},
  numbersep=8pt,
  backgroundcolor=\color{codebg},
  frame=single,
  rulecolor=\color{codeframe},
  framesep=6pt,
  breaklines=true,
  breakatwhitespace=false,
  postbreak=\mbox{\textcolor{mgrey}{$\hookrightarrow$}\space},
  showstringspaces=false,
  tabsize=4,
  columns=fullflexible,
  keepspaces=true,
  captionpos=b,
  aboveskip=1.2em,
  belowskip=1.2em,
  % Maps the characters most likely to be pasted into a listing by accident.
  %
  % Do NOT add a mapping for U+00A0 (non-breaking space). It looks like the
  % obvious next entry and it makes `listings` abort the run with "Improper
  % alphabetic constant" -- fatal, no PDF, and the message names neither the
  % character nor this line. The ASCII-in-listings convention, checked by
  % parity's C13, is the guard against it instead.
  literate={ł}{{\l}}1 {Ł}{{\L}}1 {ą}{{\k a}}1 {ę}{{\k e}}1
           {ś}{{\'s}}1 {ć}{{\'c}}1 {ż}{{\.z}}1 {ź}{{\'z}}1 {ó}{{\'o}}1 {ń}{{\'n}}1
           {Ą}{{\k A}}1 {Ę}{{\k E}}1 {Ś}{{\'S}}1 {Ć}{{\'C}}1
           {Ż}{{\.Z}}1 {Ź}{{\'Z}}1 {Ó}{{\'O}}1 {Ń}{{\'N}}1
           {—}{{-{}-}}2 {–}{{-}}1 {‘}{{'}}1 {’}{{'}}1
           {“}{{"}}1 {”}{{"}}1 {°}{{\textdegree}}1
}

\lstset{style=bookcode}

\lstnewenvironment{python}[1][]
  {\lstset{language=pythonx,style=bookcode,#1}}
  {}

% C# comparison snippets. The whole book is written for engineers arriving
% from .NET, so listings come in pairs often enough that the C# half has its
% own environment rather than a generic one.
\lstnewenvironment{csharp}[1][]
  {\lstset{language=csharpx,style=bookcode,#1}}
  {}

\lstnewenvironment{shellcmd}[1][]
  {\lstset{language=bash,style=bookcode,numbers=none,keywordstyle=\color{mteal}\bfseries,#1}}
  {}

\lstnewenvironment{tomlcode}[1][]
  {\lstset{style=bookcode,numbers=none,keywordstyle=\color{black},#1}}
  {}

\lstnewenvironment{yamlcode}[1][]
  {\lstset{style=bookcode,numbers=none,keywordstyle=\color{black},#1}}
  {}

\lstnewenvironment{jsoncode}[1][]
  {\lstset{style=bookcode,numbers=none,keywordstyle=\color{black},#1}}
  {}

% Terminal transcripts typed into the source. Allowed for a shell session the
% reader is meant to reproduce; NOT for program output, which is \transcript
% below, because an in-source transcript nobody ran reads exactly like one
% that was, and that is where a remembered number hides.
\lstnewenvironment{console}[1][]
  {\lstset{style=bookcode,numbers=none,basicstyle=\ttfamily\footnotesize,
           keywordstyle=\color{black},backgroundcolor=\color{white},#1}}
  {}

% The path is printed above every file listing, in the repository-relative
% form the reader opens in the editor beside the PDF. That is the whole point
% of a listing being a file, and the front matter promises it. The \nobreak
% keeps the path on the same page as the listing it names.
% ---------------------------------------------------------------------------
%  Appendix B prints the trap catalogue from notes/02-traps.md. One entry is
%  its number -- which a chapter may cite, and which is never reused -- the
%  habit in the reader's own voice, and what Python does instead.
%
%  The body is set RAGGED RIGHT on purpose. An entry is mostly \code{} spans,
%  which are unbreakable runs, and sixty-seven justified paragraphs of them
%  is sixty-seven chances at an overfull hbox in the edition with the longer
%  words. Ragged right removes the class rather than fixing it one box at a
%  time, and a reference list is the one place in this book where it reads
%  better anyway.
%
%  The habit is set with \enquote{} rather than \emph{}, and that is not a
%  style choice. Every habit here carries \code{} spans, and \code{} inside
%  \emph{} asks inconsolata for an italic it does not have: the first build
%  of this appendix raised `Font shape T1/zi4/m/it undefined', which
%  checklog.py is hard on for exactly this reason. It is the trap CLAUDE.md
%  records, met sixty-seven times at once -- and fixed once, because the
%  entries go through a macro. A quotation is also what the habit IS, and
%  \enquote{} is what the Polish edition owes anyway.
% ---------------------------------------------------------------------------
%  Appendix A's cheat sheet: one table per part, C# on the left, Python on
%  the right, and the chapter that teaches the row on the end. The three
%  header strings are arguments so that one environment serves both
%  editions and the structural signature stays identical.
%
%  Columns are ragged right for the reason Appendix B's entries are: a cell
%  is mostly \code{} spans, which are unbreakable runs, and a justified
%  narrow column full of them is an overfull box waiting for the edition
%  with the longer words.
% ---------------------------------------------------------------------------
%  Appendix C's tool matrix: the .NET tool, the Python tool this book uses,
%  the version it was verified against and the chapter that introduces it.
%
%  The version column prints \pydanticver and its siblings and is never
%  typed, which is the whole point of the appendix: it is the one page where
%  a reader sees every pin at once, and a pin that had to be re-typed here
%  would be a second copy of a fact the preamble already owns.
%  tools/check_structure.py --tools fails when a pin macro exists and this
%  table does not print it.
\NewDocumentEnvironment{toolmatrix}{mmmm}{%
  \small
  \begin{longtable}{@{}%
    >{\raggedright\arraybackslash}p{0.27\linewidth}%
    >{\raggedright\arraybackslash}p{0.31\linewidth}%
    >{\raggedright\arraybackslash}p{0.20\linewidth}%
    >{\raggedright\arraybackslash}p{0.08\linewidth}@{}}
  \toprule
  \textbf{#1} & \textbf{#2} & \textbf{#3} & \textbf{#4} \\
  \midrule
  \endfirsthead
  \toprule
  \textbf{#1} & \textbf{#2} & \textbf{#3} & \textbf{#4} \\
  \midrule
  \endhead
  \midrule
  \endfoot
  \bottomrule
  \endlastfoot
}{%
  \end{longtable}%
}

\NewDocumentEnvironment{cheatsheet}{mmm}{%
  \small
  \begin{longtable}{@{}%
    >{\raggedright\arraybackslash}p{0.41\linewidth}%
    >{\raggedright\arraybackslash}p{0.41\linewidth}%
    >{\raggedright\arraybackslash}p{0.08\linewidth}@{}}
  \toprule
  \textbf{#1} & \textbf{#2} & \textbf{#3} \\
  \midrule
  \endfirsthead
  \toprule
  \textbf{#1} & \textbf{#2} & \textbf{#3} \\
  \midrule
  \endhead
  \midrule
  \endfoot
  \bottomrule
  \endlastfoot
}{%
  \end{longtable}%
}

\newcommand{\trapentry}[3]{%
  \par\medskip
  \noindent\textbf{#1.}\enspace\enquote{#2}\par
  \nopagebreak
  {\raggedright\noindent#3\par}%
}

\newcommand{\listingpath}[1]{%
  \par\medskip\noindent
  {\ttfamily\footnotesize\color{mgrey}\detokenize{#1}}%
  \par\nobreak}

% External source file -- the preferred form. The path is relative to the
% repository root, which is also the path the reader opens in the editor.
% Usage: \pyfile{code/ch05/decorators.py}{Caption}{lst:label}
\newcommand{\pyfile}[3]{%
  \listingpath{#1}%
  \lstinputlisting[language=pythonx,style=bookcode,aboveskip=0.3em,
    caption={#2},label={#3}]{#1}%
}

% A named region of an external file, delimited in the source by
% `# --8<-- [start:name]` / `# --8<-- [end:name]`, so the book and the running
% code cannot drift apart. tools/check_structure.py --listings fails when the
% file or the region is missing.
%
% The markers are matched through listings' range prefix and suffix keys,
% with linerange={name-name}. The form this was inherited with -- a linerange
% whose two ends were the whole marker strings -- matched nothing and printed
% the WHOLE FILE, with no warning: measured on a three-region probe file
% before this was rewritten, and the same definition sits unused in the
% sibling book it came from. A region name is a word, never a number, because
% a number is a line number to listings. The line numbers printed are the
% file's own, so a region's listing says where in the file it sits.
% Usage: \pyregion{code/ch05/decorators.py}{timing}{Caption}{lst:label}
\lstset{
  rangebeginprefix=\#\ --8<--\ [start:, rangebeginsuffix=],
  rangeendprefix=\#\ --8<--\ [end:, rangeendsuffix=],
  includerangemarker=false}
\newcommand{\pyregion}[4]{%
  \listingpath{#1}%
  \lstinputlisting[language=pythonx,style=bookcode,aboveskip=0.3em,
    linerange={#2-#2},caption={#3},label={#4}]{#1}%
}

% The C# twin of \pyfile, for the comparison half of a pair.
\newcommand{\csfile}[3]{%
  \listingpath{#1}%
  \lstinputlisting[language=csharpx,style=bookcode,aboveskip=0.3em,
    caption={#2},label={#3}]{#1}%
}

% A program's output, written by a script under code/measure/ and pulled in
% whole. There is no typed form of this and there must not be: \val{} does
% not expand inside a verbatim body, so a transcript typed into a chapter is
% one that cannot quote a computed value, and a transcript nobody ran is
% indistinguishable on the page from one that was.
\newcommand{\transcript}[1]{%
  \par\smallskip
  \IfFileExists{figures/transcripts/#1.txt}{%
    \lstinputlisting[style=bookcode,numbers=none,basicstyle=\ttfamily\footnotesize,
      keywordstyle=\color{black},backgroundcolor=\color{white},
      language={}]{figures/transcripts/#1.txt}%
  }{%
    \begin{tcolorbox}[enhanced, sharp corners, width=\linewidth,
      colback=mlight, colframe=mgrey, boxrule=0.6pt,
      left=8pt, right=8pt, top=6pt, bottom=6pt]
      {\bfseries\color{mgrey}\small \lblTranscriptMissing}\\[2pt]
      {\ttfamily\scriptsize\color{mgrey} figures/transcripts/#1.txt}
    \end{tcolorbox}%
  }%
  \par\smallskip
}

% Inline literals. Underscores inside these must be written \_.
\newcommand{\code}[1]{\texttt{\small #1}}
\newcommand{\pkg}[1]{\texttt{\small\color{mblue}#1}}
\newcommand{\api}[1]{\texttt{\small\color{mteal}#1}}

% ============================================================
%  ADMONITIONS
%  ---------------------------------------------------------
%  A deliberately small, fixed vocabulary. Two visual treatments carry the
%  whole grammar: `tinted` is a block the reader may read in passing -- a
%  tint and a thick left bar, no outline; `outlined` is a block the reader
%  must stop at, with a full rule on white. A third treatment would mean the
%  first two are not carrying their meaning.
% ============================================================
\tcbset{
  mfa tinted/.style={
    enhanced, breakable, sharp corners,
    colback=mlight, boxrule=0pt, leftrule=3pt,
    left=8pt, right=8pt, top=6pt, bottom=6pt,
    before skip=13pt, after skip=13pt,
    fonttitle=\bfseries\small,
    coltitle=black, colbacktitle=mlight,
    attach title to upper={\par\smallskip},
  },
  mfa outlined/.style={
    enhanced, breakable, sharp corners,
    colback=white, boxrule=0.6pt,
    left=8pt, right=8pt, top=6pt, bottom=6pt,
    before skip=13pt, after skip=13pt,
    fonttitle=\bfseries\small,
    coltitle=black, colbacktitle=white,
    attach title to upper={\par\smallskip},
  },
}

% Read in passing.
\newtcolorbox{note}[1][]{mfa tinted, colframe=mblue, title={\lblNote}, #1}
\newtcolorbox{warning}[1][]{mfa tinted, colframe=mred, title={\lblWarning}, #1}
% The translation. The .NET reader is the book's only assumed reader, so this
% box carries the C# side of a comparison and nothing else. Its use is
% counted: a chapter with none is a chapter that has stopped translating.
\newtcolorbox{csbox}[1][]{mfa tinted, colframe=mteal, title={\lblCsharp}, #1}
% Anything that is preview, experimental or likely to move. Python moves more
% slowly than the agent frameworks do, and this box should be rare here.
\newtcolorbox{versionbox}[1][]{mfa tinted, colframe=mamber, title={\lblVersion}, #1}

% Stop and read.
% The misconception, named AFTER the reader has walked into it. Appendix B is
% the catalogue; a trap box in a chapter is one entry of it, in place.
\newtcolorbox{trapbox}[1][]{mfa outlined, colframe=mred, title={\lblTrap}, #1}
% Marks a listing that has not been executed against the pinned versions.
% `make debt` counts these, CI prints the count, and nothing ships to a reader
% with one attached. Removing it means you ran the code.
\newtcolorbox{verifybox}[1][]{mfa outlined, colframe=mgrey, title={\lblVerify}, #1}
\newtcolorbox{exercisebox}[1][]{mfa outlined, colframe=mteal, title={\lblExercise}, #1}
% The guiding project. Points at the stage of trace-assert this section
% builds, under code/src/trace_assert/.
\newtcolorbox{projectbox}[1][]{mfa outlined, colframe=mamber, title={\lblProject}, #1}

% ============================================================
%  EXERCISES
%  ---------------------------------------------------------
%  \begin{exercise}{key}{title} ... \end{exercise}
%
%  This is the mechanism the whole book is read through: the PDF on one half
%  of the screen, an editor on the other. An exercise is therefore not a
%  paragraph but a FILE the reader opens -- a starter under
%  code/exercises/chNN/<key>.py -- and a TEST that decides whether it is done,
%  test_<key>.py beside it. The box prints both paths and the one command
%  that runs the test, and the manifest at the back lists every exercise.
%
%  An ENVIRONMENT, not a macro with a body argument: a macro argument cannot
%  contain verbatim material, so an exercise written as \exercise{key}{title}
%  {body} could never show the reader a code fragment, and the first chapter
%  that tried would have died on `\begin{python}` inside the braces.
%
%  The key is the file stem, with underscores, and it is written with
%  \detokenize so the underscore reaches the page as a character rather than
%  as a subscript. tools/check_structure.py --exercises fails when the starter
%  or the test is missing, or when the key's chapter prefix disagrees with the
%  chapter it sits in.
% ============================================================
\newcounter{exercise}[chapter]
\renewcommand{\theexercise}{\thechapter.\arabic{exercise}}
\makeatletter
\newcommand{\exercisedir}{ch\two@digits{\value{chapter}}}
\newcommand{\listofexercises}{%
  \section*{\lblExerciseManifest}%
  \phantomsection\addcontentsline{toc}{section}{\lblExerciseManifest}%
  % The entries are subsection-level and tocdepth is 1 for the contents,
  % so the depth is raised for this list alone or it prints empty --
  % which it did, on the scaffold's first clean build.
  \begingroup\c@tocdepth=2\relax\@starttoc{exr}\endgroup%
}
\makeatother
\NewDocumentEnvironment{exercise}{mm}{%
  \refstepcounter{exercise}%
  \addcontentsline{exr}{subsection}{\protect\numberline{\theexercise}\texttt{\protect\detokenize{#1}} --- #2}%
  \begin{exercisebox}[title={\lblExercise~\theexercise\ \dash{} #2}]%
}{%
  \par\medskip
  {\small\setlength{\parindent}{0pt}%
  \lblExerciseStarter: \texttt{code/exercises/\exercisedir/\detokenize{#1}.py}\\
  \lblExerciseTest: \texttt{code/exercises/\exercisedir/test\_\detokenize{#1}.py}\\
  \lblExerciseRun: \texttt{cd code \&\& uv run pytest -k \detokenize{#1}}\par}%
  \end{exercisebox}%
}

% ============================================================
%  DIAGRAMS --- Mermaid pipeline
%  ---------------------------------------------------------
%  \mermaidfig{key}{caption}{one-line manifest description}
%
%  Source of truth is figures/mermaid/<lang>/<key>.mmd, committed. `make
%  diagrams` renders each to figures/diagrams/<lang>/<key>.pdf, which is build
%  output and is not committed, so a diagram change reviews as a text diff.
%  The source is per-language because a diagram with English labels in the
%  Polish edition is exactly the half-translation that makes a bilingual book
%  feel machine-made; CI checks that both languages carry the same keys.
%
%  If the rendered PDF is absent the macro draws a placeholder AND typesets
%  the Mermaid source, in the flow rather than as a float, because a source
%  typeset verbatim is often taller than a page and a float cannot break.
% ============================================================
\makeatletter
\newcommand{\listofdiagrams}{%
  \section*{\lblDiagramManifest}%
  \phantomsection\addcontentsline{toc}{section}{\lblDiagramManifest}%
  % The entries are subsection-level and tocdepth is 1 for the contents,
  % so the depth is raised for this list alone or it prints empty --
  % which it did, on the scaffold's first clean build.
  \begingroup\c@tocdepth=2\relax\@starttoc{dgm}\endgroup%
}
\makeatother

\newcommand{\mermaidfig}[3]{%
  % \phantomsection is load-bearing and not decoration. \addcontentsline
  % records hyperref's CURRENT anchor, and here that is whatever preceded
  % the figure -- for a figure placed after a listing it is the listing's
  % line-number anchor, and listings has already advanced the counter past
  % the last line it printed, so the manifest entry linked to a
  % destination that is never created: "name{lstnumber.10.4.46} has been
  % referenced but does not exist, replaced by a fixed one". The two
  % front-matter figures had the same wrong target and were silent about
  % it, because the lines they named happen to be printed. The exercise
  % manifest never had the bug, because \begin{exercise} steps a counter
  % and that is what sets the anchor.
  \phantomsection%
  \addcontentsline{dgm}{subsection}{\protect\numberline{}\texttt{#1.mmd} --- #3}%
  \IfFileExists{figures/diagrams/\booklang/#1.pdf}{%
    \begin{figure}[htbp]
    \centering
    \includegraphics[width=0.95\linewidth,height=0.42\textheight,keepaspectratio]{figures/diagrams/\booklang/#1.pdf}%
    \caption{#2}\label{fig:#1}
    \end{figure}%
  }{%
    \par\medskip
    \begin{tcolorbox}[
      enhanced, breakable, width=\linewidth, sharp corners,
      colback=white, colframe=mgrey, boxrule=0.8pt,
      left=8pt, right=8pt, top=8pt, bottom=8pt]
      {\bfseries\color{mgrey}\small \lblNotRendered}\\[4pt]
      {\ttfamily\scriptsize\color{mgrey} figures/mermaid/\booklang/#1.mmd}\\[6pt]
      \IfFileExists{figures/mermaid/\booklang/#1.mmd}{%
        \lstinputlisting[style=bookcode,numbers=none,basicstyle=\ttfamily\scriptsize,
                         backgroundcolor=\color{mlight},frame=none,
                         keywordstyle=\color{black},language={}]{figures/mermaid/\booklang/#1.mmd}%
      }{%
        {\small\color{mred} \lblSourceMissing: \texttt{figures/mermaid/\booklang/#1.mmd}}%
      }%
    \end{tcolorbox}%
    \captionof{figure}{#2}\label{fig:#1}
    \par\medskip
  }%
}

% ============================================================
%  CHAPTER STUBS
%  ---------------------------------------------------------
%  Prints a visible marker so an unwritten chapter cannot be mistaken for a
%  written one and the page count never flatters the draft. `make debt`
%  counts them, and CI publishes the count on every build.
% ============================================================
\newcommand{\chapterstub}[1]{%
  \begin{tcolorbox}[
    enhanced, breakable, width=\linewidth, sharp corners,
    colback=mlight, colframe=mred, boxrule=1pt,
    borderline={0.6pt}{3pt}{mred, dashed},
    left=12pt, right=12pt, top=12pt, bottom=12pt]
    {\bfseries\color{mred}\large \lblNotWritten}\\[8pt]
    \small\raggedright #1
  \end{tcolorbox}%
}

% ============================================================
%  COMPUTED VALUES
%  ---------------------------------------------------------
%  A number that came out of a measurement is written by a script under
%  code/measure/ into figures/values/<name>.tex as
%
%      \pyval{e01.threads.speedup}{3.71}
%
%  and pulled into the text with \val{e01.threads.speedup}. The script is the
%  source of truth and the book cannot disagree with it, because the book does
%  not contain the digits. `make verify` fails when a committed value no
%  longer matches its script. A value the scripts do not produce prints a
%  visible marker.
%
%  \val also applies the edition's decimal separator, which is how the Polish
%  edition gets its decimal comma without a translator having to remember.
% ============================================================
\newcommand{\bookdecimalmarker}{\ifpl{,}{.}}
\IfFileExists{siunitx.sty}{%
  \usepackage{siunitx}%
  \sisetup{
    group-digits = integer,
    group-minimum-digits = 5,
    group-separator = {\,},
    output-decimal-marker = {\bookdecimalmarker}
  }%
  \newcommand{\booknum}[1]{\num{##1}}%
}{%
  \newcommand{\booknum}[1]{##1}%
}

\newcommand{\pyvalues}{}
\makeatletter
\newcommand{\pyval}[2]{\csgdef{pyval@#1}{#2}\listxadd{\pyvalues}{#1}}
% \num on a non-numeric body is a FATAL siunitx error. The generator knows
% whether it emitted a number or a word, so it writes \pyval for a number and
% \pyvaltext for a word, and this end only enforces which macro reads which.
\newcommand{\val}[1]{%
  \ifcsdef{pyval@#1}{%
    \ifcsdef{pyvaltext@#1}{%
      \PackageError{pybook}{Value `#1' is not a number}%
        {\string\val\space passes its body to siunitx, which refuses anything
         that is not a number. Use \string\valtext{#1} instead.}%
    }{}%
    \booknum{\csuse{pyval@#1}}%
  }{\textcolor{mred}{\textbf{[\lblValueMissing: #1]}}}%
}
\newcommand{\valtext}[1]{%
  \ifcsdef{pyval@#1}{\csuse{pyval@#1}}%
    {\textcolor{mred}{\textbf{[\lblValueMissing: #1]}}}%
}
\newcommand{\pyvaltext}[2]{%
  \csgdef{pyval@#1}{#2}\csgdef{pyvaltext@#1}{}\listxadd{\pyvalues}{#1}%
}
\makeatother

% figures/values/all.tex is generated by `make numbers` and does nothing but
% \input each value file. A build with no values yet still compiles; every
% \val{} in it prints a visible marker instead of a number.
\IfFileExists{figures/values/all.tex}{\input{figures/values/all}}{%
  \AtBeginDocument{\PackageWarning{pybook}{No computed values found. Run `make numbers`.}}}

% ============================================================
%  PINNED VERSIONS --- single source of truth
%  Change these here; they propagate through both editions and Appendix C.
%  Verified against pypi.org on \pinnedon (lang/<lang>.tex).
%  tools/check_versions.py compares every one of them with PyPI on each build.
%  Python itself is not on PyPI and is checked by hand against python.org.
% ============================================================
\newcommand{\bookedition}{0.1-dev}
\newcommand{\bookyear}{2026}
\newcommand{\pyver}{3.14}              % the minor the book targets
\newcommand{\pypatch}{3.14.7}          % the exact interpreter the code ran on
\newcommand{\uvver}{0.12.13}           % uv
\newcommand{\dotnetver}{10.0.100}      % dotnet
\newcommand{\ruffver}{0.16.7}          % ruff
\newcommand{\pyrightver}{1.1.414}      % pyright
\newcommand{\mypyver}{2.3.1}           % mypy
\newcommand{\pydanticver}{2.13.5}      % pydantic
\newcommand{\pysettingsver}{2.15.0}    % pydantic-settings
\newcommand{\fastapiver}{0.141.1}      % fastapi
\newcommand{\uvicornver}{0.53.0}       % uvicorn
\newcommand{\sqlalchemyver}{2.0.52}    % sqlalchemy
\newcommand{\alembicver}{1.20.0}       % alembic
\newcommand{\aiosqlitever}{0.22.1}     % aiosqlite
\newcommand{\pytestver}{9.1.1}         % pytest
\newcommand{\pytestasyncver}{1.4.0}    % pytest-asyncio
\newcommand{\hypothesisver}{6.168.0}   % hypothesis
\newcommand{\coveragever}{7.16.1}      % coverage
\newcommand{\httpxver}{0.28.1}         % httpx
\newcommand{\structlogver}{26.1.0}     % structlog
\newcommand{\otelver}{1.44.0}          % opentelemetry-sdk
\newcommand{\polarsver}{1.44.2}        % polars
\newcommand{\anthropicver}{1.5.0}      % anthropic
\newcommand{\openaiver}{3.13.0}        % openai
\newcommand{\pipauditver}{2.10.1}      % pip-audit

% ============================================================
%  INDEX
%  Two-column, and its entries include \texttt{} names that do not hyphenate.
%  Ragged right keeps multi-word entries off the margin; it cannot help a
%  single token wider than the column, so keep verbatim index entries under
%  about 29 characters and index longer names by concept instead.
%  \raggedcolumns is the switch imakeidx's multicols actually reads.
% ============================================================
\apptocmd{\theindex}{\raggedcolumns\raggedright}{}{%
  \PackageWarning{pybook}{Could not patch theindex for ragged setting}}
\providecommand*\see[2]{}
\renewcommand*\see[2]{\emph{\lblIndexSee} #1}
\providecommand*\seealso[2]{}
\renewcommand*\seealso[2]{\emph{\lblIndexSeeAlso} #1}

\makeindex
